%%%%%%%%%%%%%%%%%%%%%%%%%%%%%%%%%%%%%%%%%%%%%%%%%%%%%%%%%%%%%%%%%%%%%%
% Lớp 11 | Quan hệ vuông góc trong không gian | Khoảng cách
% Dạng 1H8?5-3: Khoảng cách từ một điểm đến một mặt phẳng
% Số câu: 145  (sinh tự động bởi tach_cau.py; từ nay thầy sửa trực tiếp file này)
%%%%%%%%%%%%%%%%%%%%%%%%%%%%%%%%%%%%%%%%%%%%%%%%%%%%%%%%%%%%%%%%%%%%%%

%%%%%%%%%%%%%%%% TRẮC NGHIỆM 4 PHƯƠNG ÁN %%%%%%%%%%%%%%%%

% nguồn: Ok_Toan11_TracNghiem_Lan2
\begin{ex}%[1H8N5-3]%[ma:0cd9b49f]
	[ĐỀ THAM KHẢO BGD\&ĐT NĂM  2020-2021]
	\immini[thm]{Cho hình chóp tứ giác  đều $S.ABCD$ có độ tài cạnh đáy bằng $2$ và độ dài cạnh bên bằng $3$ (tham khảo hình bên). Khoảng cách từ $S$ đến mặt phẳng $(ABCD)$ bằng
	\choice 
		{\True  $\sqrt{7}$}
		{$1$} 
		 {$7$}
		{$\sqrt{11}$}} 
		 {
		\begin{tikzpicture}[scale=0.9,font=\footnotesize,line join=round,line cap=round,>=stealth]
				\def\a{4} 
				\def\h{4} 
				\path (0:0) coordinate  (A)
					++(0:\a)  coordinate (D)
					++(-130:\a/2) coordinate  (C)
					($(A)+(C)-(D)$) coordinate  (B)
					(intersection of A--C and B--D) coordinate  (O)
					($(O)+(90:\h)$) coordinate  (S);
				\draw[dashed]  (B)--(A)--(D)
					(A)--(S) 
					 (A)--(C)
					 (B)--(D)
					(S)--(O); 
				 \draw[] (B)--(C)--(D)
					 (B)--(S)
					(C)--(S) 
					 (D)--(S);
				\foreach \x/\g in {A/135,B/-135,C/-45,D/45,S/90,O/-90} 
				\fill[] (\x) circle  (1pt)
					 ($(\g:3mm)+(\x)$) node {$\x$};
				\draw  pic[draw,angle radius=6pt]{right angle=D--O--S};
			\end{tikzpicture}
		} 
	 \loigiai{
		Gọi $O$ là tâm  của đáy thì $\mathrm{d}(S,(ABCD))=SO$. \\
		Ta có $OA=\dfrac{AC}{2}=\dfrac{2\sqrt{2}}{2}=\sqrt{2}$ và $SA=3$ nên $SO=\sqrt{SA^2-OA^2}=\sqrt{3^2-2}=\sqrt{7}$. 
	} 
\end{ex}

% nguồn: Ok_TracNghiemLop11_Lan1
\begin{ex}%[1H8N5-3]%[ma:35e90e5a]
	Cho hình chóp $S.ABCD$ có đáy  là hình chữ nhật với $AB=a$, $AD=3a$. Biết $SA=2a$ và $SA$ vuông góc với mặt phẳng $(ABCD)$. Khoảng cách từ điểm $C$ đến đường thẳng $AD$ bằng
	\choice 
		{\True  $a$}
		 {$2a$}
		{$3a$} 
		{$a\sqrt{10}$} 
	 \loigiai{
		\begin{center}
			\begin{tikzpicture}[line join=round,line cap=round]
					\coordinate (A)  at (0,0);
					\coordinate (B) at (-2,-3); 
					\coordinate  (D) at (7,0);
					\coordinate  (C) at ($(B)+(D)-(A)$);
					\coordinate (S)  at ($(A)+(0,5)$);
					\draw (S)--(B) 
						 (S)--(C)
						(S)--(D) 
						(B)--(C)--(D); 
					\draw[dashed](S)--(A) 
						 (A)--(B)
						 (A)--(D);
					 \pic[draw,angle radius=3mm] {right angle = S--A--D} pic[draw,angle radius=3mm] {right angle = S--A--B};
					\foreach \i/\g in  {S/90,A/-90,B/-90,C/-90,D/0}{\draw[fill=white](\i) circle (1pt)
						($(\i)+(\g:3mm)$) node[scale=1]{$\i$};} 
				\end{tikzpicture}
		\end{center}
		 Ta có $CD\perp AD$ nên $\mathrm{d}(C,AD)=CD=a$.
	 }
\end{ex}

% nguồn: Ok_Toan11_TracNghiem_Lan2
\begin{ex}%[1H8H5-3]%[ma:4cff2b84]
	[Mã  104-2022]
	\immini[thm]{Cho hình lập phương $ABCD.A'B'C'D'$ có cạnh bằng $3$ (tham khảo hình bên). Khoảng  cách từ điểm $B$ đến mặt phẳng $\left(ACC'A'\right)$ bằng
	 \choice
		 {$3$}
		 {$3\sqrt{2}$}
		 {\True $\dfrac{3\sqrt{2}}{2}$}
		 {$\dfrac{3}{2}$}}
		{ 
		\begin{tikzpicture}[scale=0.9,font=\footnotesize,line join=round,line cap=round,>=stealth]
				\def\a{3} 
				 \def\b{2}
				\def\h{3} 
				\path (0:0) coordinate  (A)
					++(0:\a)  coordinate (D)
					 ++(-130:\b) coordinate (C)
					($(A)+(C)-(D)$) coordinate  (B)
					($(A)+(90:\h)$)  coordinate (A')
					 ($(B)+(90:\h)$) coordinate (B')
					 ($(C)+(90:\h)$) coordinate (C')
					($(D)+(90:\h)$) coordinate  (D')
					 ($(B)!1/2!(D)$) coordinate (O);
				\draw[dashed] (B)--(A)--(D)--(B) 
					(C)--(A)--(A'); 
				\draw[]  (C)--(C')
					 (D)--(D')
					(B)--(B') 
					(C)--(C')--(A') 
					(B)--(C)--(D) 
					 (A')--(B')--(C')--(D')--cycle;
				\foreach \x/\g in {A/180,B/180,C/0,D/0,A'/180,B'/180,C'/0,D'/0,O/-90} 
				\fill[] (\x) circle (1pt) 
					($(\g:4mm)+(\x)$)  node {$\x$};
			\end{tikzpicture}
		} 
	\loigiai{ 
		Gọi $O$ là giao điểm của $AC$ và $BD$.\\ 
		Ta  có $BO\perp AC$, $BO\perp AA'$ (vì $AA'\perp(ABCD)$.\\
		Nên  $BO\perp\left(ACC'A'\right)$
		 $\Rightarrow \mathrm{d}\left(B;\left(ACCC'A'\right)\right)=BO$.\\
		Xét  tam giác vuông $ABD$ có $BD=\sqrt{AB^2+AD^2}=\sqrt{3^2+3^2}=\sqrt{18}=3\sqrt{2}$.\\
		 Khi đó $BO=\dfrac{1}{2} BD=\dfrac{3\sqrt{2}}{2}$.\\
		Vậy  khoảng cách từ điểm $B$ đến mặt phẳng $\left(ACC'A'\right)$ bằng $\dfrac{3\sqrt{2}}{2}$.
	} 
\end{ex}

% nguồn: Ok_Toan11_TracNghiem_Lan2
\begin{ex}%[1H8H5-3]%[ma:fd2456b1]
	Cho hình lăng trụ đứng $ABC.A'B'C'$ có mặt đáy  $ABC$ là tam giác vuông tại $B$ có $AB=a$, $AC=a\sqrt{3}$, $A'B=2a$. Gọi $M$ là trung điểm của $AC$. Khoảng cách từ $M$ đến mặt phẳng $(A'BC)$ là
	 \choice
		{\True  $\dfrac{a\sqrt{3}}{4}$}
		 {$\dfrac{a\sqrt{3}}{2}$}
		 {$\dfrac{3a}{2}$}
		{$\dfrac{3a}{4}$} 
	 \loigiai{
		\begin{center}
			\begin{tikzpicture}[scale=1,font=\footnotesize,line join=round,line cap=round,>=stealth]
					 \coordinate (A) at (0,0);
					\coordinate (B) at (1,-1.5); 
					\coordinate (C)  at (4,0);
					 \coordinate (A') at ($(A)+(0,3)$);
					\coordinate (B')  at ($(B)+(0,3)$);
					\coordinate (C') at  ($(C)+(0,3)$);
					\coordinate  (M) at ($(A)!0.5!(C)$);
					\coordinate  (H) at ($(A')!0.45!(B)$);
					\draw  (H)--(A)--(B)--(C)--(C')--(B')--(A')--(A)
						 (B')--(B)
						(B)--(A')--(C'); 
					 \draw[dashed,thin](A)--(C)--(A');
					\pic[draw,thin,angle  radius=6pt] {right angle = A'--H--A};
					\foreach \i/\g in {A'/90,A/-90,B/-90,C/-90,C'/0,B'/180,H/0,M/90}{\draw[fill](\i)  circle (1pt)
						 ($(\i)+(\g:3mm)$) node[scale=1]{$\i$};}
				\end{tikzpicture}
		\end{center}
		 Ta có $\mathrm{d}(M,(A'BC)) = \dfrac{1}{2}\mathrm{d}(A,(A'BC))$.\\
		Kẻ $AH \perp  A'B$.\\
		Ta có  $A'A \perp (ABC) \Rightarrow A'A \perp BC$.
		Mà $AB \perp BC  \Rightarrow BC \perp (A'AB)$.\\
		Có $\begin{cases} AH \perp A'B \\ BC \perp AH
	\end{cases} \Rightarrow AH \perp (A'BC)$.\\
	$\Rightarrow  \mathrm{d}(A,(A'BC)) = AH$.\\
	Ta có $AA' =  \sqrt{A'B^2-AB^2} = \sqrt{(2a)^2-a^2} = a\sqrt{3}$.\\
	$S_{A'AB} = \dfrac{1}{2}\cdot AH  \cdot A'B = \dfrac{1}{2}AA' \cdot AB \Rightarrow AH = \dfrac{AA' \cdot AB}{A'B} = \dfrac{a\sqrt{3} \cdot a}{2a} = \dfrac{a\sqrt{3}}{2}$.\\
	Vậy $\mathrm{d}(M,(A'BC)) = \dfrac{1}{2}\mathrm{d}(A,(A'BC))  = \dfrac{1}{2}AH = \dfrac{1}{2}\dfrac{a\sqrt{3}}{2} = \dfrac{a\sqrt{3}}{4}$.
} 
\end{ex}

% nguồn: Ok_Toan11_TracNghiem_Lan2
\begin{ex}%[1H8H5-3]%[ma:f4c15a97]
	Cho lăng trụ đứng $ABCD.A'B'C'D$ có đáy là hình thoi cạnh $a$, $\widehat{BAC}=60^\circ$. Khoảng cách từ điểm $C$ đến mặt  phẳng $(ABA'B')$ bằng
	\choice 
		 {$2a$}
		 {\True $\dfrac{a\sqrt{3}}{2}$}
		{$a\sqrt{3}$} 
		 {$a$}
	 \loigiai{
		\begin{center}
			\begin{tikzpicture}[scale=0.75,font=\footnotesize,line join=round,line cap=round,>=stealth]
					 \coordinate (A) at (0,0);
					\coordinate (B) at  (-1,-1.5);
					\coordinate (D)  at (4,0);
					\coordinate  (C) at ($(B)+(D)-(A)$);
					\coordinate (A') at  ($(A)+(0,3)$);
					\coordinate (B')  at ($(B)+(0,3)$);
					\coordinate (C') at ($(C)+(0,3)$); 
					\coordinate  (D') at ($(D)+(0,3)$);
					 \coordinate (H) at ($(A)!0.5!(B)$);
					\draw  (B)--(C)--(D)--(D')--(C')--(B')--(A')--(D')
						(B')--(B) 
						(C')--(C); 
					 \draw[dashed,thin](A')--(A)--(B)
						(A)--(D) 
						 (A)--(C)
						(B)--(D) 
						(C)--(H); 
					\pic[draw,thin,angle radius=2mm] {right  angle = C--H--B};
					\foreach \i/\g in {A'/90,A/-90,B/-90,C/-90,D/0,C'/0,D'/30,B'/180,H/90}{\draw[fill](\i) circle  (1pt)
						($(\i)+(\g:3mm)$) node[scale=1]{$\i$};} 
				\end{tikzpicture}
		\end{center}
		Ta có $\widehat {BAC} = 60^\circ \Rightarrow \triangle ABC$ đều.\\ 
		Gọi $H$ là trung điểm của  $AB \Rightarrow CH \perp AB \Rightarrow CH \perp (ABA'B')$.\\
		Ta  có $\mathrm{d}(C,(ABA'B')) = CH = \dfrac{a\sqrt{3}}{2}$.
	} 
\end{ex}

% nguồn: Ok_Toan11_TracNghiem_Lan2
\begin{ex}%[1H8H5-3]%[ma:cf29cccc]
	Cho hình lăng trụ tam giác đều  $ABC.A'B'C'$ có cạnh đáy bằng $2a$. Khoảng cách từ $B$ đến mặt phẳng $(ACC'A')$ bằng
	\choice 
		{$2a$} 
		{$2\sqrt{2}a$} 
		 {$\sqrt{2}a$}
		{\True  $\sqrt{3}a$}
	 \loigiai{
		\begin{center}
			\begin{tikzpicture}[scale=1,font=\footnotesize,line join=round,line cap=round,>=stealth]
					\coordinate  (A) at (0,0);
					\coordinate (B) at  (1,-1.5);
					 \coordinate (C) at (4,0);
					\coordinate (A')  at ($(A)+(0,3)$);
					\coordinate (B') at  ($(B)+(0,3)$);
					\coordinate (C')  at ($(C)+(0,3)$);
					\coordinate (H) at ($(A)!0.5!(C)$); 
					 \draw (A)--(B)--(C)--(C')--(B')--(A')--(A)
						(B')--(B) 
						(A')--(C'); 
					 \draw[dashed,thin](A)--(C)
						(B)--(H); 
					 \pic[draw,thin,angle radius=6pt] {right angle = B--H--A};
					 \foreach \i/\g in {A'/90,A/-90,B/-90,C/-90,C'/0,B'/180,H/60}{\draw[fill](\i) circle (1pt)
						 ($(\i)+(\g:3mm)$) node[scale=1]{$\i$};}
				\end{tikzpicture}
		\end{center}
		Kẻ $BH \perp AC$, mà $BH \perp AA'$ nên $BH \perp (ACC'A')$.\\ 
		$\Rightarrow \mathrm{d}(B,(ACC'A')) = BH = \dfrac{2a\sqrt{3}}{2} =  a\sqrt{3}$.
	 }
\end{ex}

% nguồn: Ok_Toan11_TracNghiem_Lan2
\begin{ex}%[1H8H5-3]%[ma:41686dc8]
	Cho hình lăng trụ đứng $ABC.A'B'C'$ có đáy $ABC$ là tam giác vuông tại $B$, $AB=a$, $AA'=2a$. Khoảng cách từ điểm $C'$ đến mặt  phẳng $(A'BC)$ bằng
	\choice 
		{\True  $\dfrac{2a\sqrt{5}}{5}$}
		{$2a\sqrt{5}$} 
		 {$\dfrac{a\sqrt{5}}{5}$}
		{$\dfrac{3a\sqrt{5}}{5}$} 
	\loigiai{ 
		\begin{center}
			\begin{tikzpicture}[scale=1,font=\footnotesize,line join=round,line cap=round,>=stealth]
					\coordinate (A) at (0,0); 
					\coordinate (B) at  (1,-1.5);
					\coordinate (C) at (4,0); 
					\coordinate (A') at ($(A)+(0,3)$); 
					\coordinate (B') at  ($(B)+(0,3)$);
					\coordinate  (C') at ($(C)+(0,3)$);
					\coordinate (O) at  ($(A)!0.5!(C')$);
					\coordinate  (H) at ($(A')!0.45!(B)$);
					\draw  (H)--(A)--(B)--(C)--(C')--(B')--(A')--(A)
						(B')--(B) 
						(B)--(A')--(C'); 
					 \draw[dashed,thin](A)--(C)--(A')
						(A)--(C'); 
					\pic[draw,thin,angle radius=6pt] {right  angle = A'--H--A};
					\foreach \i/\g in  {A'/90,A/-90,B/-90,C/-90,C'/0,B'/180,H/0,O/0}{\draw[fill](\i) circle (1pt)
						 ($(\i)+(\g:3mm)$) node[scale=1]{$\i$};}
				\end{tikzpicture}
		\end{center}
		Vì $ABC.A'B'C'$ là lăng trụ đứng nên $A'C'CA$ là hình chữ  nhật.\\
		Gọi  $O=AC' \cap A'C$, khi đó $AO=C'O$.\\
		Mà $AC' \cap (A'BC) = O$ nên khoảng cách từ điểm $C'$ đến mặt phẳng $(A'BC)$ bằng  khoảng cách từ điểm $A$ đến mặt phẳng $(A'BC)$.\\
		Ta có 
		$\heva{&AA' \perp BC \\  &AB \perp BC}\Rightarrow BC \perp (A'AB)$.\\
		Từ $A$ hạ đường cao $AH$ xuống $A'B$.\\ 
		Khi đó ta có $AH \perp A'B$ mà $BC \perp AH$.\\ 
		$\Rightarrow AH \perp  (A'BC)$ nên khoảng cách từ điểm $A$ đến mặt phẳng $(A'BC)$ bằng $AH$.\\
		Xét $\Delta A'AB$ vuông tại $A$, đường cao  $AH$ có
		$$\dfrac{1}{AH^2} = \dfrac{1}{AB^2} + \dfrac{1}{A'A^2} = \dfrac{1}{a^2} +  \dfrac{1}{4a^2} = \dfrac{5}{4a^2}
		\Rightarrow AH = \dfrac{2a}{\sqrt{5}} =  \dfrac{2a\sqrt{5}}{5}.$$
	 }
\end{ex}

% nguồn: Ok_Toan11_TracNghiem_Lan2
\begin{ex}%[1H8H5-3]%[ma:34a6d0a0]
	\immini[thm]{Cho hình chóp $S.ABC$ có đáy là tam giác vuông cân tại  $B$, $AB=2a$ và $SA$ vuông góc với mặt phẳng đáy. Khoảng cách từ $C$ đến mặt phẳng $(SAB)$ bằng
	\choice 
		 {$\sqrt{2} a$}
		{\True $2a$} 
		 {$a$}
		{$2\sqrt{2}  a$}}
		{ 
		\begin{tikzpicture}[scale=0.9,font=\footnotesize,line join=round,line cap=round,>=stealth]
				\def\a{4} 
				\def\h{3} 
				\path (0:0) coordinate (A) 
					 ++(0:\a) coordinate (C)
					(-40:4*\a/5)  coordinate (B)
					 ($(A)+(90:\h)$) coordinate (S)
					($(S)!1/3!(B)$) coordinate (H); 
				 \draw[] (A)--(B)--(C)
					 (A)--(S)
					(B)--(S) 
					 (C)--(S);
				\draw[dashed]  (A)--(C);
				\foreach \x /\goc in  {A/180,C/0,B/-135,S/90}
				\fill[] (\x) circle  (1pt)
					($(\x)+(\goc:3mm)$) node  {$\x$};
				 \draw pic[draw,angle radius=2mm]{right angle=C--A--S};
				\draw pic[draw,angle radius=2mm]{right angle=C--B--A}; 
			\end{tikzpicture}
		 }
	 \loigiai{
		Ta có $SA\perp(ABC)\Rightarrow  SA\perp CB$.\\
		 Ta có $\heva{&CB\perp AB\\
		&CB\perp SA} \Rightarrow CB\perp(SAB)$.\\ 
		Do đó $\mathrm{d}(C,(SAB))=CB=AB=2a$. 
	} 
\end{ex}

% nguồn: Ok_Toan11_TracNghiem_Lan2
\begin{ex}%[1H8H5-3]%[ma:c7c6a862]
	[MĐ 103 BGD\&ĐT  NĂM 2017-2018]
	Cho hình chóp $S. ABCD$  có đáy là hình vuông cạnh $\sqrt{3} a$, $SA$ vuông góc với mặt phẳng đáy và $SA=a$. Khoảng cách từ $A$ đến mặt phẳng $(SBC)$ bằng
	 \choice
		 {$\dfrac{\sqrt{5} a}{3}$}
		{\True $\dfrac{\sqrt{3} a}{2}$} 
		 {$\dfrac{\sqrt{6} a}{6}$}
		{$\dfrac{\sqrt{3} a}{3}$} 
	\loigiai{ 
	\immini{Ta có  $\heva{&BC\perp AB\\
		&BC\perp  SA} \Rightarrow BC\perp(SAB)$.\\
		$\Rightarrow\heva{&(SAB)  \perp(SBC)\\
		 &(SAB) \cap(SBC)=SB.}$\\
		Trong mặt phẳng $(SAB)$, kẻ $AH\perp SB\Rightarrow AH=\mathrm{d}(A,(SBC))$.\\ 
		Áp dụng hệ thức lượng trong tam giác vuông  $SAB$ có\\ $\dfrac{1}{AH^2}=\dfrac{1}{SA^2}+\dfrac{1}{AB^2}=\dfrac{1}{a^2}+\dfrac{1}{3a^2}=\dfrac{4}{3a^2}$.\\
		 Suy ra $AH=\dfrac{\sqrt{3}a}{2}$.}
		 {
		\begin{tikzpicture}[scale=0.9,font=\footnotesize,line join=round,line cap=round,>=stealth]
				\def\a{4} 
				\def\h{4} 
				\path (0:0) coordinate (A) 
					++(0:\a) coordinate (D) 
					++(-130:\a/2)  coordinate (C)
					 ($(A)+(C)-(D)$) coordinate (B)
					 ($(A)+(90:\h)$) coordinate (S)
					($(S)!2/3!(B)$) coordinate  (H)
					(intersection of A--C and B--D)  coordinate (O);
				\draw[dashed] (B)--(A)--(D) 
					(H)--(A)--(S); 
				\draw[] (B)--  (C)--(D)
					 (B)--(S)
					 (C)--(S)
					 (D)--(S);
				\foreach \x/\g  in {A/135,B/-135,C/-45,D/45,S/90,H/135}
				\fill[]  (\x) circle (1pt)
					($(\g:3mm)+(\x)$) node  {$\x$};
				\draw pic[draw,angle radius=6pt]{right angle=D--A--S}; 
				 \draw pic[draw,angle radius=6pt]{right angle=B--H--A};
			\end{tikzpicture}
		 }
	} 
\end{ex}

% nguồn: Ok_Toan11_TracNghiem_Lan2
\begin{ex}%[1H8H5-3]%[ma:16a409fc]
	Cho hình chóp $S.ABCD$ có đáy $ABCD$ là hình vuông cạnh $a$, $SA \perp (ABCD)$ và $SA=2a$. Khoảng cách từ điểm $A$ đến mặt phẳng  $(SBD)$ bằng
	\choice 
		{$\dfrac{a}{3}$} 
		{\True  $\dfrac{2a}{3}$}
		 {$\dfrac{a\sqrt{2}}{3}$}
		 {$\dfrac{4a}{9}$}
	\loigiai{ 
		\begin{center}
			\begin{tikzpicture}[scale=1,font=\footnotesize,line join=round,line cap=round,>=stealth]
					 \coordinate (A) at (0,0);
					\coordinate  (B) at (-2.2,-1.5);
					\coordinate  (D) at (4.8,0);
					 \coordinate (C) at (2.6,-1.5);
					\coordinate  (S) at (0,3.5);
					\coordinate (O) at  ($(A)!0.5!(C)$);
					\coordinate (H) at ($(S)!0.45!(O)$); 
					\draw[dashed] (O)  -- (S) -- (A) -- (B) -- (D)
						 (C) -- (A) -- (D)
						(A) -- (H); 
					\draw  (S)--(B)--(C)--(D)--(S)--(C);
					\pic[draw,thin,angle radius=6pt] {right angle  = S--A--D}
						pic[draw,thin,angle  radius=6pt] {right angle = A--H--O};
					 \foreach \i/\g in {S/90,A/160,B/-90,C/-90,D/0,O/-90,H/170}{\draw[fill](\i) circle (1pt)
						 ($(\i)+(\g:3mm)$) node[scale=1]{$\i$};}
				\end{tikzpicture}
		\end{center}
		Gọi  $O$ là tâm hình vuông $ABCD$, dựng $AH \perp SO$.\\
		Khi đó,  $\mathrm{d}(A,(SBD)) = AH$.\\
		 Trong tam giác $SAO$ vuông tại $O$ có $AH$ là chiều cao nên
		$$\dfrac{1}{AH^2} =  \dfrac{1}{SA^2} + \dfrac{1}{AO^2} = \dfrac{1}{(2a)^2} + \dfrac{1}{\left(\dfrac{a\sqrt{2}}{2}\right)^2} = \dfrac{1}{4a^2} + \dfrac{2}{a^2} = \dfrac{9}{4a^2} \Rightarrow AH = \dfrac{2a}{3}.$$
	 }
\end{ex}

% nguồn: Ok_Toan11_TracNghiem_Lan2
\begin{ex}%[1H8H5-3]%[ma:107f1a74]
	Cho hình lập phương $ABCD.A'B'C'D'$ có cạnh bằng $a$. Khoảng cách từ điểm $A$ tới mặt phẳng $(A'BD)$ bằng 
	\choice 
		{$\dfrac{a\sqrt{3}}{2}$} 
		{$\dfrac{a\sqrt{2}}{3}$} 
		 {$\dfrac{a\sqrt{6}}{3}$}
		{\True $\dfrac{a\sqrt{3}}{3}$} 
	\loigiai{ 
		\begin{center}
			\begin{tikzpicture}[scale=1,font=\footnotesize,line join=round,line cap=round,>=stealth]
					\coordinate  (A) at (0,0);
					 \coordinate (B) at (-1,-1.5);
					\coordinate (D) at (3,0); 
					 \coordinate (C) at ($(B)+(D)-(A)$);
					\coordinate (A') at ($(A)+(0,3)$); 
					\coordinate (B')  at ($(B)+(0,3)$);
					\coordinate (C') at  ($(C)+(0,3)$);
					\coordinate (D') at  ($(D)+(0,3)$);
					\coordinate (O) at  ($(A)!0.5!(C)$);
					\coordinate (H) at  ($(A')!0.45!(O)$);
					\draw (B)--(C)--(D)--(D')--(C')--(B')--(A')--(D') 
						(B')--(B) 
						 (C')--(C);
					\draw[dashed,thin](A')--(A)--(B)--(A') 
						 (A)--(D)
						(A)--(O)--(A') 
						(B)--(D)--(A') 
						 (A)--(H);
					\pic[draw,thin,angle radius=6pt] {right angle  = A'--H--A} pic[draw,thin,angle radius=6pt] {right angle = A--O--B};
					\foreach \i/\g in {A'/90,A/-90,B/-90,C/-90,D/0,C'/0,D'/30,B'/180,O/-90,H/0}{\draw[fill](\i) circle  (1pt)
						 ($(\i)+(\g:3mm)$) node[scale=1]{$\i$};}
				\end{tikzpicture}
		\end{center}
		Gọi $O$ là tâm của hình vuông $ABCD$, $H$ là hình chiếu vuông góc của $A$ trên $A'O$ thì  $AH \perp A'O$.\\
		Ta có $\heva{ &BD \perp AO \\ &BD \perp AA'} \Rightarrow BD \perp (AA'O) \Rightarrow BD  \perp AH$.\\
		Vậy $AH \perp (A'BD)  \Rightarrow \mathrm{d}(A,(A'BD))=AH$.\\
		Xét tam giác $AA'O$ vuông tại $A$ có đường cao $AH$, ta  có
		$$\dfrac{1}{AH^2} = \dfrac{1}{AA'^2}+\dfrac{1}{AO^2} = \dfrac{1}{a^2}+\dfrac{1}{\left(\dfrac{a\sqrt{2}}{2}\right)^2} = \dfrac{1}{a^2}+\dfrac{2}{a^2}=\dfrac{3}{a^2} \Rightarrow AH  = \dfrac{a\sqrt{3}}{3}.$$
	 }
\end{ex}

% nguồn: Ok_Toan11_TracNghiem_Lan2
\begin{ex}%[1H8H5-3]%[ma:1ab0b3c1]
	\immini[thm]{Cho hình lập phương $ABCD.A'B'C'D'$ có  cạnh bên bằng $2a$ (tham khảo hình vẽ). Khoảng cách từ $C$ đến mặt phẳng $\left(BDD'B'\right)$ bằng
	\choice 
		 {$2\sqrt{2} a$}
		{$2\sqrt{3}  a$}
		{\True  $\sqrt{2} a$}
		{$\sqrt{3}  a$}}
		 {
		\begin{tikzpicture}[scale=0.9,font=\footnotesize,line join=round,line cap=round,>=stealth]
				 \def\a{3}
				\def\b{2} 
				\def\h{3} 
				\path (0:0) coordinate  (A)
					 ++(0:\a) coordinate (D)
					++(-130:\b)  coordinate (C)
					($(A)+(C)-(D)$) coordinate  (B)
					($(A)+(90:\h)$) coordinate (A') 
					($(B)+(90:\h)$)  coordinate (B')
					($(C)+(90:\h)$)  coordinate (C')
					($(D)+(90:\h)$) coordinate (D') 
					 ($(B)!1/2!(D)$) coordinate (O);
				 \draw[dashed] (B)--(A)--(D)--(B)
					 (C)--(A)--(A');
				\draw[] (C)--(C') 
					(D)--(D') 
					 (B)--(B')--(D')
					(C)--(C') 
					(B)--(C)--(D) 
					(A')--(B')--(C')--(D')--cycle; 
				 \foreach \x/\g in {A/180,B/180,C/0,D/0,A'/180,B'/180,C'/0,D'/0,O/-90}
				\fill[] (\x) circle  (1pt)
					 ($(\g:4mm)+(\x)$) node {$\x$};
			\end{tikzpicture}
		} 
	\loigiai{ 
		Gọi $O=AC\cap BD$, khi đó ta có $CO\perp BD$ (do tứ giác  $ABCD$ là hình vuông).\\
		Lại có  $CO\perp DD'$ (do $DD'\perp (ABCD)$ và $CO\subset (ABCD)$).\\
		Suy ra $CO\perp  (BDD'B')$, do đó $CO=\mathrm{d}(C, (BDD'B'))$.\\
		 Hình lập phương $ABCD.A'B'C'D'$ có cạnh bằng $2a$ nên $AC=2a \sqrt{2}$.\\
		Suy ra  $CO=\dfrac{1}{2} AC=a \sqrt{2}$.\\
		Vậy khoảng cách từ  $C$ đến mặt phẳng $(BDD'B')$ bằng $a \sqrt{2}$.
	} 
\end{ex}

% nguồn: Ok_Toan11_TracNghiem_Lan2
\begin{ex}%[1H8H5-3]%[ma:59635c38]
	[MĐ 102 BGD\&ĐT  NĂM 2017-2018]
	Cho hình chóp $S. ABC$ có đáy là tam  giác vuông đỉnh $B$, $AB=a$, $SA$ vuông góc với mặt phẳng đáy và $SA=a$. Khoảng cách từ điểm $A$ đến mặt phẳng $(SBC)$ bằng
	 \choice
		{$\dfrac{a}{2}$} 
		{$a$} 
		 {$\dfrac{a \sqrt{6}}{3}$}
		 {\True $\dfrac{a \sqrt{2}}{2}$}
	\loigiai{ 
	\immini{Kẻ $AH\perp SB$ trong  mặt phẳng $(SAB)$.\\
		Ta  có $\heva{&BC\perp AB\\
		&BC\perp SA} \Rightarrow  BC\perp(SAB) \Rightarrow BC\perp AH$.\\
		Ta có $\heva{&AH\perp  BC\\
		&AH\perp  SB} \Rightarrow AH\perp(SBC)$.\\
		Khi đó  $\mathrm{d}(A,(SBC))=AH=\dfrac{1}{2} SB=\dfrac{a \sqrt{2}}{2}$.}
		{ 
		\begin{tikzpicture}[scale=0.9,font=\footnotesize,line join=round,line cap=round,>=stealth]
				\def\a{4} 
				 \def\h{3}
				\path (0:0) coordinate (A) 
					++(0:\a) coordinate (C) 
					(-40:4*\a/5) coordinate  (B)
					 ($(A)+(90:\h)$) coordinate (S)
					($(S)!1/3!(B)$) coordinate (H); 
				\draw[]  (A)--(B)--(C)
					(A)--(H) 
					 (A)--(S)
					 (B)--(S)
					(C)--(S); 
				\draw[dashed] (A)--(C); 
				\foreach \x /\goc in  {A/180,C/0,B/-135,S/90,H/45}
				\fill[] (\x)  circle (1pt)
					($(\x)+(\goc:3mm)$) node {$\x$}; 
				 \draw pic[draw,angle radius=2mm]{right angle=C--A--S};
				 \draw pic[draw,angle radius=2mm]{right angle=C--B--A};
				\draw pic[draw,angle radius=2mm]{right  angle=A--H--B};
			\end{tikzpicture}
		 }
	} 
\end{ex}

% nguồn: Ok_Toan11_TracNghiem_Lan2
\begin{ex}%[1H8H5-3]%[ma:3e8d1289]
	 [Mã 103-2022]
	\immini[thm]{Cho  hình lập phương $ABCD.A'B'C'D'$ có cạnh bằng $3$ (tham khảo hình vẽ). Khoảng cách từ $B$ đến mặt phẳng $\left(ACC'A'\right)$ bằng
	 \choice
		{\True  $\dfrac{3\sqrt{2}}{2}$}
		{$\dfrac{3}{2}$} 
		{$3\sqrt{2}$} 
		 {$3$}}
		{ 
		\begin{tikzpicture}[scale=0.7,font=\footnotesize,line join=round,line cap=round,>=stealth]
				 \def\a{3}
				 \def\b{2}
				\def\h{3} 
				\path (0:0) coordinate (A) 
					++(0:\a) coordinate (D) 
					++(-130:\b)  coordinate (C)
					($(A)+(C)-(D)$)  coordinate (B)
					($(A)+(90:\h)$) coordinate (A') 
					($(B)+(90:\h)$)  coordinate (B')
					($(C)+(90:\h)$)  coordinate (C')
					($(D)+(90:\h)$) coordinate (D') 
					($(B)!1/2!(D)$) coordinate (I); 
				\draw[dashed] (B)--(A)--(D)--(B) 
					 (C)--(A)--(A');
				 \draw[] (C)--(C')
					(D)--(D') 
					(B)--(B') 
					 (C)--(C')--(A')
					(B)--(C)--(D) 
					(A')--(B')--(C')--(D')--cycle; 
				 \foreach \x/\g in {A/180,B/180,C/0,D/0,A'/180,B'/180,C'/0,D'/0,I/-90}
				 \fill[] (\x) circle (1pt)
					($(\g:4mm)+(\x)$) node  {$\x$};
			\end{tikzpicture}
		 }
	\loigiai{ 
		 Gọi $I=AC\cap BD$.\\
		 Ta có $BI\perp\left(ACC'A'\right) \Rightarrow \mathrm{d}\left(B;\left(ACC'A'\right)\right)=BI=\dfrac{1}{2} BD=\dfrac{3\sqrt{2}}{2}$.
	 }
\end{ex}

% nguồn: Ok_Toan11_TracNghiem_Lan2
\begin{ex}%[1H8H5-3]%[ma:ec4fd261]
	Cho hình chóp $S.ABC$ có đáy $ABC$ là tam giác vuông tại $B$, $AB=a\sqrt{2}$, cạnh bên $SA=a$ và vuông góc với mặt  phẳng đáy. Khoảng cách từ $A$ đến $(SBC)$ bằng
	\choice 
		{$\dfrac{a\sqrt{6}}{2}$} 
		{\True  $\dfrac{a\sqrt{6}}{3}$}
		{$\dfrac{a\sqrt{2}}{3}$} 
		{$\dfrac{a\sqrt{6}}{3}$} 
	 \loigiai{
		\begin{center}
			\begin{tikzpicture}[scale=1,font=\footnotesize,line join=round,line cap=round,>=stealth]
					\coordinate (A) at (0,0); 
					\coordinate (C) at (4,0); 
					 \coordinate (B) at (1.5,-2);
					\coordinate (S)  at ($(A)+(0,3)$);
					 \coordinate (H) at ($(S)!0.4!(B)$);
					\coordinate  (H) at ($(S)!0.4!(B)$);
					\draw  (S)--(B)--(A)--(S)
						 (S)--(C)
						(C)--(B)--(A)--(H); 
					\draw[dashed,thin] (A)--(C); 
					\pic[draw,thin,angle  radius=6pt] {right angle = S--A--C} pic[draw,thin,angle radius=6pt] {right angle = S--A--B} pic[draw,thin,angle radius=6pt] {right angle = A--B--C} pic[draw,thin,angle radius=6pt] {right angle = A--H--B};
					\foreach \i/\g in  {S/90,A/180,C/0,B/-90,H/60}{\draw[fill](\i) circle (1pt)
						 ($(\i)+(\g:3mm)$) node[scale=1]{$\i$};}
				\end{tikzpicture}
		\end{center}
		Trong mặt phẳng $(SAB)$, kẻ $AH \perp SB$.\\ 
		 Ta có
		$\heva{ &BC \perp AB \\ &BC \perp  SA} \Rightarrow BC \perp (SAB)$.\\
		Khi đó, $AH \perp (SBC)  \Rightarrow \mathrm{d}(A,(SBC)) = AH$.\\
		Ta có $\dfrac{1}{AH^2} = \dfrac{1}{SA^2}  + \dfrac{1}{AB^2} = \dfrac{1}{a^2} + \dfrac{1}{(a\sqrt{2})^2} = \dfrac{1}{a^2} + \dfrac{1}{2a^2} = \dfrac{3}{2a^2} \Rightarrow AH = \sqrt{\dfrac{2a^2}{3}} = \dfrac{a\sqrt{6}}{3}$.
	 }
\end{ex}

% nguồn: Ok_Toan11_TracNghiem_Lan2
\begin{ex}%[1H8H5-3]%[ma:47c21072]
	Cho hình chóp $S.ABC$ có đáy là tam giác đều  cạnh $a$. Cạnh bên $SA = a\sqrt{3}$ và vuông góc với mặt phẳng đáy $(ABC)$. Tính khoảng cách $\mathrm{d}$ từ $A$ đến mặt phẳng $(SBC)$.
	\choice 
		{\True $\mathrm{d}=\dfrac{a\sqrt{15}}{5}$} 
		 {$\mathrm{d}=a$}
		 {$\mathrm{d}=\dfrac{a\sqrt{15}}{5}$}
		 {$\mathrm{d}=\dfrac{a\sqrt{3}}{2}$}
	\loigiai{ 
		\begin{center}
			\begin{tikzpicture}[scale=1,font=\footnotesize,line join=round,line cap=round,>=stealth]
					\coordinate (A)  at (0,0);
					 \coordinate (B) at (4,0);
					\coordinate  (C) at (1.5,-2);
					\coordinate (S)  at ($(A)+(0,3)$);
					\coordinate (H) at  ($(B)!0.5!(C)$);
					\coordinate  (K) at ($(S)!0.4!(H)$);
					 \draw (H)--(S)--(C)--(A)--(S)
						(S)--(B) 
						 (B)--(C);
					 \draw[dashed,thin] (H)--(A)--(B)
						(A)--(K); 
					\pic[draw,thin,angle radius=6pt] {right angle = S--A--B} pic[draw,thin,angle  radius=6pt] {right angle = S--A--C} pic[draw,thin,angle radius=6pt] {right angle = A--K--H};
					\foreach \i/\g in {S/90,A/180,B/0,C/-90,K/60,H/0}{\draw[fill](\i)  circle (1pt)
						($(\i)+(\g:3mm)$) node[scale=1]{$\i$};} 
				\end{tikzpicture}
		\end{center}
		Vẽ $AH \perp  BC$ tại $H \Rightarrow BC \perp (SAH)$.\\
		Vẽ $AK \perp SH$ tại $K$ mà $AK \perp BC \Rightarrow AK \perp (SBC)$ tại  $K$.\\
		Do  đó $AK = \mathrm{d}(A,(SBC))$.
		$H$ là trung điểm $BC$  nên $AH = \dfrac{a\sqrt{3}}{2}$.\\
		Vậy $AK  = \dfrac{SA \cdot AH}{\sqrt{SA^2+AH^2}} = \dfrac{a\sqrt{3} \cdot \dfrac{a\sqrt{3}}{2}}{\sqrt{(a\sqrt{3})^2 + \left(\dfrac{a\sqrt{3}}{2}\right)^2}} = \dfrac{a\sqrt{15}}{5}$.
	} 
\end{ex}

% nguồn: Ok_Toan11_TracNghiem_Lan2
\begin{ex}%[1H8H5-3]%[ma:c05eb83e]
	Cho hình chóp $S.ABCD$ có đáy là  hình vuông cạnh $a$. Đường thẳng $SA$ vuông góc với mặt phẳng đáy và $SA=a$. Gọi $M$ là trung điểm của $CD$. Khoảng cách từ $M$ đến mặt phẳng $(SAB)$ bằng
	 \choice
		 {$\dfrac{a\sqrt{2}}{2}$}
		{\True $a$} 
		{$a\sqrt{2}$} 
		 {$2a$}
	\loigiai{ 
		\begin{center}
			\begin{tikzpicture}[scale=0.75,font=\footnotesize,line join=round,line cap=round,>=stealth]
					\coordinate (A)  at (0,0);
					\coordinate  (B) at (-2.2,-0.8);
					\coordinate (D) at (3.8,0); 
					 \coordinate (C) at (1.6,-0.8);
					\coordinate (S) at (0,3.5); 
					\coordinate (H) at  ($(A)!0.5!(B)$);
					\coordinate (M) at ($(C)!0.5!(D)$); 
					\draw[dashed] (S) -- (A) --  (D)
						(A) -- (B) 
						(H) -- (M); 
					\draw (S)--(B)--(C)--(D)--(S)--(C); 
					 \draw (0.25,0) -- (0.25,0.25) -- (0,0.25);
					 \foreach \i/\g in {S/90,A/160,B/-90,C/-90,D/0,H/120,M/-90}{\draw[fill](\i) circle (1pt)
						 ($(\i)+(\g:3mm)$) node[scale=1]{$\i$};}
				\end{tikzpicture}
		\end{center}
		Gọi $H$  là trung điểm $AB \Rightarrow MH \perp AB$.\\
		Ta có $SA  \perp (ABCD) \Rightarrow SA \perp MH$
		$\Rightarrow MH \perp (SAB) \Rightarrow \mathrm{d}(M,(SAB)) = MH = a$. 
	 }
\end{ex}

% nguồn: Ok_Toan11_TracNghiem_Lan2
\begin{ex}%[1H8H5-3]%[ma:31855d50]
	[Mã 102 BGD\&ĐT NĂM 2020-2021-ĐỢT 1] 
	\immini[thm]{Cho hình chóp $S. ABC$ có đáy  là tam giác vuông cân tại $C$, $AC=3a$ và $SA$ vuông góc với mặt phẳng đáy. Khoảng cách từ $B$ đến mặt phẳng ($SAC$) bằng
	 \choice
		{$\dfrac{3}{2} a$} 
		 {$\dfrac{3\sqrt{2}}{2} a$}
		{\True $3a$} 
		{$3\sqrt{2}  a$}}
		 {
		\begin{tikzpicture}[scale=0.9,font=\footnotesize,line join=round,line cap=round,>=stealth]
				 \def\a{4}
				\def\h{3} 
				 \path (0:0) coordinate (A)
					++(0:\a)  coordinate (B)
					++(-150:4*\a/5)  coordinate (C)
					 ($(A)+(90:\h)$) coordinate (S)
					($(A)!0.5!(B)$) coordinate  (M);
				\draw[]  (A)--(C)--(B)
					(A)--(S) 
					(B)--(S) 
					 (C)--(S);
				\draw[dashed] (A)--(B); 
				\foreach \x /\goc  in {A/180,B/0,C/-135,S/90}
				 \fill[] (\x) circle (1pt)
					($(\x)+(\goc:3mm)$) node {$\x$}; 
				\draw pic[draw,angle  radius=2mm]{right angle=B--A--S};
				 \draw pic[draw,angle radius=2mm]{right angle=B--C--A};
			\end{tikzpicture}
		 }
	 \loigiai{
		Ta có $\triangle ABC$ vuông cân tại  $C$ nên $BC\perp AC$.\\
		Ta  có $\heva{& BC\perp AC\\ & BC\perp SA \,(SA\perp(ABC)} \Rightarrow BC\perp(SAC)\Rightarrow \mathrm{d}(B,(SAC))=BC=3a$.
	 }
\end{ex}

% nguồn: Ok_Toan11_TracNghiem_Lan2
\begin{ex}%[1H8H5-3]%[ma:89757643]
	Cho hình chóp $S.ABC$ có $SA$ vuông góc với đáy, $SA=a$. Tam giác $ABC$ vuông cân tại $A$, $BC=a\sqrt{2}$. Gọi $M$ là trung điểm của $AB$. Khoảng cách từ $M$  đến mặt phẳng $(SBC)$ bằng
	\choice 
		 {$\dfrac{a\sqrt{3}}{3}$}
		 {\True $\dfrac{a\sqrt{3}}{6}$}
		 {$\dfrac{a\sqrt{3}}{4}$}
		{$a\sqrt{3}$} 
	\loigiai{ 
		\begin{center}
			\begin{tikzpicture}[scale=1,font=\footnotesize,line join=round,line cap=round,>=stealth]
					\coordinate  (A) at (0,0);
					\coordinate (C) at (4,0); 
					\coordinate (B) at (1.5,-2); 
					\coordinate (S) at  ($(A)+(0,3)$);
					\coordinate  (I) at ($(B)!0.5!(C)$);
					 \coordinate (H) at ($(S)!0.4!(I)$);
					\coordinate  (M) at ($(A)!0.5!(B)$);
					\draw  (I)--(S)--(B)--(A)--(S)
						(S)--(C) 
						 (B)--(C);
					\draw[dashed,thin] (I)--(A)--(C) 
						 (A)--(H);
					\pic[draw,thin,angle radius=6pt] {right angle =  S--A--C} pic[draw,thin,angle radius=6pt] {right angle = S--A--B} pic[draw,thin,angle radius=6pt] {right angle = A--H--I};
					\foreach \i/\g  in {S/90,A/180,C/0,B/-90,H/60,I/0,M/180}{\draw[fill](\i) circle (1pt)
						($(\i)+(\g:3mm)$)  node[scale=1]{$\i$};}
				\end{tikzpicture}
		\end{center}
		Gọi $I$ là trung điểm của $BC$. Vì tam giác $ABC$ vuông cân tại $A$ nên  $AI \perp BC$.\\
		Theo giả thiết  $SA \perp (ABC) \Rightarrow SA \perp BC$. Do đó $BC \perp (SAI)$.\\
		Trong  mặt phẳng $(SAI)$, kẻ $AH \perp SI$. Mà $BC \perp (SAI) \Rightarrow BC \perp AH$.\\
		Từ đó suy ra $AH \perp (SBC) \Rightarrow \mathrm{d}(A,(SBC))  = AH$.\\
		 Ta có $AI = \dfrac{BC}{2} = \dfrac{a\sqrt{2}}{2}$.\\
		$AH = \dfrac{AI \cdot  AS}{\sqrt{AI^2+AS^2}} = \dfrac{\dfrac{a\sqrt{2}}{2} \cdot a\sqrt{3}}{\sqrt{\left(\dfrac{a\sqrt{2}}{2}\right)^2+(a\sqrt{3})^2}} = \dfrac{a\sqrt{21}}{7}$.\\
		Vì $M$ là trung điểm của $BC$ nên $\mathrm{d}(M,(SBC)) = \dfrac{1}{2}\mathrm{d}(A,(SBC)) =  \dfrac{1}{2}AH = \dfrac{a\sqrt{21}}{14}$.
	} 
\end{ex}

% nguồn: Ok_Toan11_TracNghiem_Lan2
\begin{ex}%[1H8H5-3]%[ma:5f83b47c]
	Cho hình lập phương $ABCD.A'B'C'D'$ có cạnh bằng $a$. Khoảng cách từ $M$ đến mặt  phẳng $(A'BD)$ bằng
	\choice 
		{$\dfrac{a\sqrt{3}}{2}$} 
		{$\dfrac{a\sqrt{2}}{3}$} 
		 {$\dfrac{a\sqrt{6}}{3}$}
		{\True  $\dfrac{a\sqrt{3}}{3}$}
	 \loigiai{
		\begin{center}
			\begin{tikzpicture}[scale=1,font=\footnotesize,line join=round,line cap=round,>=stealth]
					\coordinate (A)  at (0,0);
					\coordinate (B) at (-1,-1.5); 
					\coordinate (D) at  (3,0);
					\coordinate (C) at  ($(B)+(D)-(A)$);
					\coordinate (A')  at ($(A)+(0,3)$);
					\coordinate (B') at ($(B)+(0,3)$); 
					 \coordinate (C') at ($(C)+(0,3)$);
					\coordinate (D')  at ($(D)+(0,3)$);
					 \coordinate (O) at ($(A)!0.5!(C)$);
					\coordinate (H) at  ($(A')!0.45!(O)$);
					\draw  (B)--(C)--(D)--(D')--(C')--(B')--(A')--(D')
						 (B')--(B)
						 (C')--(C);
					 \draw[dashed,thin](A')--(A)--(B)--(A')
						(A)--(D) 
						(A)--(O)--(A') 
						 (B)--(D)
						 (A)--(H);
					\pic[draw,thin,angle radius=6pt] {right angle =  A'--H--A} pic[draw,thin,angle radius=6pt] {right angle = A--O--B};
					\foreach \i/\g  in {A'/90,A/-90,B/-90,C/-90,D/0,C'/0,D'/30,B'/180,O/-90,H/0}{\draw[fill](\i) circle (1pt)
						($(\i)+(\g:3mm)$) node[scale=1]{$\i$};} 
				\end{tikzpicture}
		\end{center}
		Có  $OA = \dfrac{a\sqrt{2}}{2}$.\\
		Khi đó $\dfrac{1}{AH^2} = \dfrac{1}{A'A^2}+\dfrac{1}{OA^2} = \dfrac{1}{a^2}+\dfrac{2}{a^2} = \dfrac{3}{a^2} \Rightarrow  AH = \dfrac{a\sqrt{3}}{3}$.\\
		 Vậy $\mathrm{d}(A,(A'BD)) = \dfrac{a\sqrt{3}}{3}$.
	} 
\end{ex}

% nguồn: Ok_Toan11_TracNghiem_Lan2
\begin{ex}%[1H8H5-3]%[ma:ba550879]
	[MĐ  104 BGD\&ĐT NĂM 2020-2021-ĐỢT 2]
	\immini[thm]{Cho hình lập phương $ABCD.A'B'C'D'$ có cạnh $2a$. Khoảng cách từ $A$ đến mặt phẳng $\left(BDD'B'\right)$ bằng 
	 \choice
		{$2\sqrt{2}  a$}
		{$2\sqrt{3}  a$}
		 {\True $\sqrt{2} a$}
		{$\sqrt{3}  a$}}
		 {
		\begin{tikzpicture}[scale=0.9,font=\footnotesize,line join=round,line cap=round,>=stealth]
				\def\a{3} 
				 \def\b{2}
				 \def\h{3}
				 \path (0:0) coordinate (A)
					++(0:\a) coordinate  (D)
					++(-130:\b) coordinate  (C)
					 ($(A)+(C)-(D)$) coordinate (B)
					 ($(A)+(90:\h)$) coordinate (A')
					($(B)+(90:\h)$) coordinate  (B')
					($(C)+(90:\h)$) coordinate  (C')
					($(D)+(90:\h)$)  coordinate (D')
					($(B)!1/2!(D)$) coordinate  (O);
				 \draw[dashed] (B)--(A)--(D)--(B)
					 (C)--(A)--(A');
				 \draw[] (C)--(C')
					 (D)--(D')
					(B)--(B')--(D') 
					 (C)--(C')
					(B)--(C)--(D) 
					(A')--(B')--(C')--(D')--cycle; 
				\foreach  \x/\g in {A/180,B/180,C/0,D/0,A'/180,B'/180,C'/0,D'/0,O/-90}
				 \fill[] (\x) circle (1pt)
					($(\g:4mm)+(\x)$)  node {$\x$};
			\end{tikzpicture}
		} 
	 \loigiai{
		Gọi  $O$ là giao điểm của $AC$ và $BD$.\\
		Ta có $AC$ cắt $BD$ tại $O$  hay $AO\perp BD$.\\
		Lại  có $ABCD.A'B'C'D'$ là hình lập phương cạnh $2a$ nên $BB'\perp(ABCD) \Rightarrow AO\perp BB'$. $(2)$\\
		Từ $(1)$ và $(2)$ suy ra $AO\perp\left(BDD'B'\right)$.\\ 
		Vậy $\mathrm{d}\left(A,\left(BDD'B'\right)\right)=AO=\dfrac{1}{2}  AC=\dfrac{1}{2} 2\sqrt{2} a=\sqrt{2} a$.
	 }
\end{ex}

% nguồn: Ok_Toan11_TracNghiem_Lan2
\begin{ex}%[1H8H5-3]%[ma:10e7162d]
	Cho hình chóp tứ giác đều $S.ABCD$ có cạnh đáy bằng $a$ và chiều cao bằng $a\sqrt{2}$. Tính khoảng cách $\mathrm{d}$ từ tâm $O$ của đáy $ABCD$ đến một  mặt bên theo $a$.
	 \choice
		{\True  $\mathrm{d} = \dfrac{a\sqrt{2}}{3}$}
		{$\mathrm{d}  = \dfrac{2a\sqrt{5}}{3}$}
		{$\mathrm{d} = \dfrac{a\sqrt{3}}{2}$} 
		{$\mathrm{d} =  \dfrac{a\sqrt{5}}{2}$}
	\loigiai{ 
		\begin{center}
			\begin{tikzpicture}[scale=1,font=\footnotesize,line join=round,line cap=round,>=stealth]
					\coordinate  (A) at (0,0);
					 \coordinate (D) at (-2,-1);
					\coordinate (B) at  (4,0);
					\coordinate (C)  at ($(B)+(D)-(A)$);
					\coordinate (O)  at ($(A)!0.5!(C)$);
					\coordinate  (S) at ($(O)+(0,4)$);
					\coordinate (I) at  ($(C)!0.5!(B)$);
					\coordinate (M) at ($(S)!0.7!(I)$); 
					\draw (S)--(B) 
						 (S)--(C)
						(S)--(D) 
						 (B)--(C)--(D)
						(S)--(I); 
					\draw[dashed,thin](S)--(O)--(M) 
						(O)--(I) 
						 (S)--(A)--(C)
						(A)--(B) 
						(A)--(D)--(B); 
					 \foreach \i/\g in {S/90,A/150,D/-90,C/-90,B/0,O/-90,I/-30,M/30}{\draw[fill](\i) circle (1pt)
						($(\i)+(\g:3mm)$)  node[scale=1]{$\i$};}
				\end{tikzpicture}
		\end{center}
		Gọi $I$ là hình chiếu  của $O$ lên $CD$, $M$ là hình chiếu của $O$ lên $SI$.\\
		 Suy ra đoạn $OM$ là khoảng cách từ $O$ đến mặt phẳng $(SCD)$.\\
		Vậy $$\mathrm{d} = OM = \dfrac{OI \cdot OS}{\sqrt{OI^2+OS^2}} = \dfrac{\dfrac{a}{2} \cdot a\sqrt{2}}{\sqrt{\left(\dfrac{a}{2}\right)^2 + (a\sqrt{2})^2}} = \dfrac{\dfrac{a^2\sqrt{2}}{2}}{\sqrt{\dfrac{a^2}{4}+2a^2}} = \dfrac{\dfrac{a^2\sqrt{2}}{2}}{\sqrt{\dfrac{9a^2}{4}}} = \dfrac{\dfrac{a^2\sqrt{2}}{2}}{\dfrac{3a}{2}}  = \dfrac{a\sqrt{2}}{3}.$$
	 }
\end{ex}

% nguồn: Ok_Toan11_TracNghiem_Lan2
\begin{ex}%[1H8H5-3]%[ma:a0d0b7f6]
	Cho hình chóp tứ giác đều $S.ABCD$ có cạnh đáy bằng $2a$. Gọi $O$ là giao điểm của $AC$ và $BD$. Biết $SO =  a$, khoảng cách từ $O$ đến mặt phẳng $(SBC)$ bằng
	 \choice
		 {$\dfrac{a\sqrt{5}}{5}$}
		 {$\dfrac{a\sqrt{3}}{2}$}
		 {$\dfrac{a}{2}$}
		 {\True $\dfrac{a\sqrt{2}}{2}$}
	\loigiai{ 
		\begin{center}
			\begin{tikzpicture}[scale=1,font=\footnotesize,line join=round,line cap=round,>=stealth]
					\coordinate  (A) at (0,0);
					\coordinate (D) at (-2,-1); 
					 \coordinate (B) at (4,0);
					 \coordinate (C) at ($(B)+(D)-(A)$);
					\coordinate (O)  at ($(A)!0.5!(C)$);
					\coordinate (S) at ($(O)+(0,4)$); 
					\coordinate (I)  at ($(C)!0.5!(B)$);
					\coordinate  (H) at ($(S)!0.7!(I)$);
					\draw (S)--(B) 
						 (S)--(C)
						 (S)--(D)
						(B)--(C)--(D) 
						(S)--(I); 
					 \draw[dashed,thin](S)--(O)--(H)
						 (O)--(I)
						(S)--(A)--(C) 
						 (A)--(B)
						 (A)--(D)--(B);
					\foreach \i/\g in {S/90,A/150,D/-90,C/-90,B/0,O/-90,I/-30,H/30}{\draw[fill](\i) circle  (1pt)
						 ($(\i)+(\g:3mm)$) node[scale=1]{$\i$};}
				\end{tikzpicture}
		\end{center}
		Gọi  $I$ là trung điểm $BC$ và $H$ là hình chiếu của $O$ trên $SI$.\\
		Khi đó $\heva{ &BC \perp OI \\  &BC \perp SI } \Rightarrow BC \perp (SOI) \Rightarrow BC \perp OH$.\\
		Nên $OH \perp (SBC) \Rightarrow \mathrm{d}(O;(SBC))  = OH = \dfrac{a\sqrt{2}}{2}$.
	 }
\end{ex}

% nguồn: Ok_Toan11_TracNghiem_Lan2
\begin{ex}%[1H8H5-3]%[ma:a6871275]
	Cho hình  chóp $S.ABCD$ có đáy là hình vuông tâm $O$ cạnh $a$ và $SA \perp (ABCD)$, $SA=a\sqrt{3}$. Khoảng cách từ $O$ đến mặt phẳng $(SCD)$ là
	\choice 
		 {$\dfrac{a\sqrt{3}}{3}$}
		{\True $\dfrac{a\sqrt{3}}{4}$} 
		 {$\dfrac{3a}{2}$}
		 {$\dfrac{3a}{4}$}
	\loigiai{ 
		\begin{center}
			\begin{tikzpicture}[scale=0.75,font=\footnotesize,line join=round,line cap=round,>=stealth]
					\coordinate (A)  at (0,0);
					 \coordinate (B) at (-2.2,-1.8);
					\coordinate (D) at  (3.8,0);
					\coordinate (C) at  (1.6,-1.8);
					 \coordinate (S) at (0,3.5);
					\coordinate  (H) at ($(S)!0.45!(D)$);
					\coordinate (O)  at ($(A)!0.5!(C)$);
					\draw[dashed]  (S) -- (A) -- (B)
						(H)  -- (A) -- (D)
						 (A)--(C)
						(B)--(D); 
					\draw (S)--(B)--(C)--(D)--(S)--(C); 
					\pic[draw,thin,angle radius=6pt] {right angle = S--A--D} pic[draw,thin,angle radius=6pt] {right  angle = A--H--D};
					 \foreach \i/\g in {S/90,A/150,B/-90,C/-90,D/0,H/60,O/-100}{\draw[fill](\i) circle (1pt)
						 ($(\i)+(\g:3mm)$) node[scale=1]{$\i$};}
				\end{tikzpicture}
		\end{center}
		Kẻ đường cao $OH$ của tam giác  $SAD$.\\
		Ta có $AC \cap  (SCD) = C$.\\
		Mà $O$ là trung điểm của  $AC$,
		nên $\mathrm{d}(O,(SCD))  = \dfrac{1}{2}\mathrm{d}(A,(SCD))$.\\
		Suy ra $\mathrm{d}(O,(SCD)) = \dfrac{1}{2}\mathrm{d}(A,(SCD))$.\\ 
		Ta có $AH \perp SD$, $AH \perp CD$ (vì $CD \perp (SAD)$) và $SD \cap  CD = D$.\\
		Nên suy ra $\mathrm{d}(A,(SCD)) =  AH$.\\
		Xét tam  giác $SAD$ vuông tại $A$ có đường cao $AH$
		$$\dfrac{1}{AH^2} = \dfrac{1}{SA^2}  + \dfrac{1}{AD^2} = \dfrac{1}{(a\sqrt{3})^2} + \dfrac{1}{a^2} = \dfrac{1}{3a^2} + \dfrac{1}{a^2} = \dfrac{4}{3a^2} \Rightarrow AH = \dfrac{a\sqrt{3}}{2}.$$
		Nên $\mathrm{d}(O,(SCD)) =  \dfrac{1}{2}\mathrm{d}(A,(SCD)) = \dfrac{a\sqrt{3}}{4}$.
	 }
\end{ex}

% nguồn: Ok_Toan11_TracNghiem_Lan2
\begin{ex}%[1H8H5-3]%[ma:3869c99c]
	Cho hình chóp $S.ABCD$ có đáy là hình vuông cạnh $a$, $SA \perp (ABCD)$. Khoảng cách từ $C$  đến mặt phẳng $(SAB)$ bằng
	 \choice
		{\True $a$} 
		{$2a$} 
		{$a\sqrt{2}$} 
		 {$\dfrac{a}{2}$}
	 \loigiai{
		\begin{center}
			\begin{tikzpicture}[scale=0.75,font=\footnotesize,line join=round,line cap=round,>=stealth]
					 \coordinate (A) at (0,0);
					\coordinate (B) at (-2.2,-0.8); 
					\coordinate (D) at  (3.8,0);
					\coordinate (C)  at (1.6,-0.8);
					\coordinate (S) at (0,3.5); 
					\draw[dashed] (S) -- (A) -- (B) 
						 (A) -- (D);
					 \draw (S)--(B)--(C)--(D)--(S)--(C);
					 \draw (0.25,0) -- (0.25,0.25) -- (0,0.25);
					\foreach \i/\g in {S/90,A/-90,B/-90,C/-90,D/0}{\draw[fill](\i) circle  (1pt)
						($(\i)+(\g:3mm)$)  node[scale=1]{$\i$};}
				\end{tikzpicture}
		\end{center}
		Vì $BC  \perp AB$ và $BC \perp SA$, suy ra $BC \perp (SAB)$.\\
		 Vậy $\mathrm{d}(C,(SAB)) = BC = a$.
	 }
\end{ex}

% nguồn: Ok_Toan11_TracNghiem_Lan2
\begin{ex}%[1H8H5-3]%[ma:9fadbfef]
	Cho hình chóp $S.ABC$ có đáy là tam giác vuông cân tại $B$, $AB=a\sqrt{2}$ và $SA$ vuông góc với mặt  phẳng đáy. Khoảng cách từ $B$ đến mặt phẳng $(SAC)$ bằng
	\choice 
		{$3\sqrt{2}a$} 
		{\True  $a$}
		{$\dfrac{3}{2}a$} 
		{$3a$} 
	\loigiai{ 
		\begin{center}
			\begin{tikzpicture}[scale=1,font=\footnotesize,line join=round,line cap=round,>=stealth]
					\coordinate (A) at  (0,0);
					\coordinate (C)  at (4,0);
					\coordinate (B) at  (1.5,-2);
					\coordinate (S) at ($(A)+(0,3)$); 
					\coordinate (H) at ($(S)!0.4!(B)$); 
					\draw (S)--(B)--(A)--(S) 
						(S)--(C) 
						(C)--(B)--(A); 
					\draw[dashed,thin]  (A)--(C);
					\pic[draw,thin,angle  radius=6pt] {right angle = S--A--C} pic[draw,thin,angle radius=6pt] {right angle = S--A--B} pic[draw,thin,angle radius=6pt] {right angle = A--B--C};
					\foreach \i/\g in {S/90,A/180,C/0,B/-90}{\draw[fill](\i)  circle (1pt)
						($(\i)+(\g:3mm)$) node[scale=1]{$\i$};} 
				\end{tikzpicture}
		\end{center}
		Vì $SA \perp (ABC)$ nên $(ABC)  \perp (SAC)$.\\
		 Hạ $BH \perp AC$, khi đó $BH \perp (SAC)$, suy ra $\mathrm{d}(B,(SAC))=BH$.\\
		Vì tam giác $ABC$ vuông cân tại $B$, $AB=a\sqrt{2}$  nên $AC=2a$,.\\
		Suy ra $BH=\dfrac{AB  \cdot BC}{AC}=\dfrac{a\sqrt{2} \cdot a\sqrt{2}}{2a}=\dfrac{2a^2}{2a}=a$.\\
		Vậy  $\mathrm{d}(B,(SAC))=a$.
	 }
\end{ex}

% nguồn: Ok_Toan11_TracNghiem_Lan2
\begin{ex}%[1H8H5-3]%[ma:4d0a70b9]
	Cho lăng trụ đều $ABC.A'B'C'$, biết $AB=AA'=a$. Khoảng cách từ điểm $A$ đến mặt phẳng  $(BCC'B')$ bằng
	\choice 
		{$a\sqrt{3}$} 
		{$\dfrac{a\sqrt{3}}{3}$} 
		{\True  $\dfrac{a\sqrt{3}}{2}$}
		{$a$} 
	\loigiai{ 
		\begin{center}
			\begin{tikzpicture}[scale=1,font=\footnotesize,line join=round,line cap=round,>=stealth]
					\coordinate (A) at  (0,0);
					\coordinate (B)  at (1,-1.5);
					\coordinate (C) at (4,0); 
					\coordinate (A') at ($(A)+(0,3)$); 
					\coordinate  (B') at ($(B)+(0,3)$);
					\coordinate (C')  at ($(C)+(0,3)$);
					\coordinate  (M) at ($(B)!0.5!(C)$);
					\draw  (A)--(B)--(C)--(C')--(B')--(A')--(A)
						(B')--(B) 
						 (A')--(C');
					 \draw[dashed,thin](A)--(C)
						 (A)--(M);
					 \pic[draw,thin,angle radius=6pt] {right angle = B--M--A};
					\foreach \i/\g in {A'/90,A/-90,B/-90,C/-90,C'/0,B'/180,M/-90}{\draw[fill](\i) circle (1pt) 
						($(\i)+(\g:3mm)$) node[scale=1]{$\i$};} 
				\end{tikzpicture}
		\end{center}
		Gọi $M$ là trung điểm $BC \Rightarrow AM = \dfrac{a\sqrt{3}}{2}$.\\ 
		Ta có $AM \perp BC$.\\ 
		Mặt  khác $AM \perp BB'$.\\
		Suy ra  $AM \perp (BCC'B') \Rightarrow \mathrm{d}(A,(BCC'B')) = AM = \dfrac{a\sqrt{3}}{2}$.
	} 
\end{ex}

% nguồn: Ok_Toan11_TracNghiem_Lan2
\begin{ex}%[1H8H5-3]%[ma:ee355fec]
	Cho  hình chóp $S.ABCD$ có đáy là hình vuông cạnh $a$. Đường thẳng $SA$ vuông góc với mặt phẳng đáy, $SA=a$. Gọi $M$ là trung điểm của $CD$. Khoảng cách từ $M$ đến $(SAB)$ nhận giá trị nào trong các giá trị sau?
	 \choice
		{$\dfrac{a\sqrt{2}}{2}$} 
		 {$2a$}
		 {$a\sqrt{2}$}
		{\True  $a$}
	\loigiai{ 
		\begin{center}
			\begin{tikzpicture}[scale=0.75,font=\footnotesize,line join=round,line cap=round,>=stealth]
					\coordinate (A)  at (0,0);
					\coordinate (B) at  (-2.2,-0.8);
					\coordinate  (D) at (3.8,0);
					\coordinate (C)  at (1.6,-0.8);
					 \coordinate (S) at (0,3.5);
					\coordinate  (M) at ($(C)!0.5!(D)$);
					\draw[dashed] (S) -- (A) -- (B) 
						 (A) -- (D);
					\draw  (S)--(B)--(C)--(D)--(S)--(C);
					\draw (0.25,0)  -- (0.25,0.25) -- (0,0.25);
					\foreach \i/\g  in {S/90,A/-90,B/-90,C/-90,D/0,M/-30}{\draw[fill](\i) circle (1pt)
						($(\i)+(\g:3mm)$)  node[scale=1]{$\i$};}
				\end{tikzpicture}
		\end{center}
		Ta có $CD \parallel AB$, mà $AB  \subset (SAB)$ nên $CD\parallel (SAB)$.\\
		Từ đó suy ra $\mathrm{d}(M;(SAB))  = \mathrm{d}(D;(SAB))$.\\
		Ta có $AD \perp AB$, $AD \perp SA$  suy ra $AD \perp (SAB)$.\\
		Suy ra $\mathrm{d}(D;(SAB)) =  AD = a$.\\
		Vậy $\mathrm{d}(M;(SAB)) = a$. 
	} 
\end{ex}

% nguồn: Ok_Toan11_TracNghiem_Lan2
\begin{ex}%[1H8H5-3]%[ma:b6fa4d81]
	Cho hình chóp $S.ABCD$ có đáy là hình  vuông cạnh $a$, $SA \perp (ABCD)$. Tính khoảng cách từ điểm $B$ đến mặt phẳng $(SAC)$.
	 \choice
		{$\dfrac{a}{2}$} 
		{\True  $\dfrac{a\sqrt{2}}{2}$}
		 {$\dfrac{a\sqrt{2}}{3}$}
		 {$\dfrac{a\sqrt{2}}{4}$}
	\loigiai{ 
		\begin{center}
			\begin{tikzpicture}[scale=1,font=\footnotesize,line join=round,line cap=round,>=stealth]
					\coordinate (A) at (0,0); 
					\coordinate  (B) at (-2.2,-1.5);
					\coordinate (D) at  (3.8,0);
					\coordinate (C) at (1.6,-1.5); 
					\coordinate (S)  at (0,3.5);
					\coordinate  (O) at ($(A)!0.5!(C)$);
					\draw[dashed] (S) -- (A) -- (B) --  (D)
						 (C) -- (A) -- (D);
					\draw  (S)--(B)--(C)--(D)--(S)--(C);
					\draw (0.25,0) -- (0.25,0.25)  -- (0,0.25);
					\foreach \i/\g in {S/90,A/160,B/-90,C/-90,D/0,O/90}{\draw[fill](\i)  circle (1pt)
						 ($(\i)+(\g:3mm)$) node[scale=1]{$\i$};}
				\end{tikzpicture}
		\end{center}
		Gọi $AC  \cap BD = \{O\}$.\\
		Vì $SA \perp  (ABCD) \Rightarrow SA \perp BO$.\\
		Ta có 
		$\heva{&BO \perp SA, BO \perp AC \\ &SA \subset (SAC),  AC \subset (SAC) \\ &SA \cap AC = \{A\} } \Rightarrow BO \perp (SAC)$.\\
		$\Rightarrow  \mathrm{d}(B,(SAC)) = BO = \dfrac{1}{2}BD = \dfrac{1}{2}\sqrt{AB^2+AD^2} = \dfrac{1}{2}\sqrt{a^2+a^2} = \dfrac{a\sqrt{2}}{2}$.
	} 
\end{ex}

% nguồn: Ok_Toan11_TracNghiem_Lan2
\begin{ex}%[1H8H5-3]%[ma:41e1c86b]
	Cho hình chóp $S.ABCD$ có đáy là hình vuông cạnh $2a$, $SA \perp  (ABCD)$, góc giữa đường thẳng $SC$ và mặt phẳng $(ABCD)$ bằng $60^\circ$. Khoảng cách từ $B$ đến mặt phẳng $(SCD)$ bằng
	\choice 
		{$\dfrac{a\sqrt{6}}{3}$} 
		{$\dfrac{a\sqrt{6}}{4}$} 
		 {$\dfrac{a\sqrt{6}}{2}$}
		 {\True $\dfrac{2a\sqrt{6}}{3}$}
	 \loigiai{
		\begin{center}
			\begin{tikzpicture}[scale=0.75,font=\footnotesize,line join=round,line cap=round,>=stealth]
					\coordinate (A)  at (0,0);
					\coordinate (B) at  (-2.2,-0.8);
					\coordinate (D) at (3.8,0); 
					\coordinate  (C) at (1.6,-0.8);
					\coordinate (S) at  (0,3.5);
					\coordinate (H) at ($(S)!0.45!(D)$); 
					 \draw[dashed] (S) -- (A) -- (B)
						(H)--(A) --  (D);
					\draw  (S)--(B)--(C)--(D)--(S)--(C);
					 \pic[draw,thin,angle radius=2mm] {right angle = S--A--D} pic[draw,thin,angle radius=2mm] {right angle = A--H--D};
					\foreach \i/\g in {S/90,A/-90,B/-90,C/-90,D/0,H/60}{\draw[fill](\i) circle (1pt) 
						($(\i)+(\g:3mm)$) node[scale=1]{$\i$};} 
				\end{tikzpicture}
		\end{center}
		Ta có $SA \perp  (ABCD)$ nên $(SC,(ABCD)) = \widehat {SCA} = 60^\circ$.\\
		Khi đó  $AC = 2a\sqrt{2} \Rightarrow SA = AC \cdot \tan 60^\circ = 2a\sqrt{2} \cdot \sqrt{3} = 2a\sqrt{6}$.\\
		Mà $AB//CD  \Rightarrow AB//(SCD) \Rightarrow \mathrm{d}(B,(SCD)) = \mathrm{d}(A,(SCD))$.\\
		Kẻ $AH \perp SD  \Rightarrow AH \perp (SCD)$.\\
		Khi đó  $\mathrm{d}(B,(SCD)) = \mathrm{d}(A,(SCD)) = AH = \dfrac{SA \cdot AD}{\sqrt{SA^2+AD^2}} = \dfrac{2a\sqrt{6}}{3}$.
	 }
\end{ex}

% nguồn: Ok_Toan11_TracNghiem_Lan2
\begin{ex}%[1H8H5-3]%[ma:5a8e17ef]
	[Mã 104 BGD\&ĐT NĂM 2020-2021-ĐỢT  1]
	\immini[thm]{Cho hình chóp $S. ABC$ có đáy là tam giác vuông cân tại $B$, $AB=4a$ và $SA$ vuông góc với  mặt phẳng đáy. Khoảng cách từ $C$ đến mặt phẳng $(SAB)$ bằng
	 \choice
		 {\True $4a$}
		 {$4\sqrt{2} a$}
		{$2\sqrt{2} a$} 
		 {$2a$}}
		{ 
		\begin{tikzpicture}[scale=0.9,font=\footnotesize,line join=round,line cap=round,>=stealth]
				 \def\a{4}
				\def\h{3} 
				\path  (0:0) coordinate (A)
					++(0:\a) coordinate  (C)
					++(-150:4*\a/5) coordinate  (B)
					($(A)+(90:\h)$)  coordinate (S)
					($(A)!0.5!(B)$) coordinate  (M);
				\draw[]  (A)--(B)--(C)
					 (A)--(S)
					(B)--(S) 
					(C)--(S); 
				 \draw[dashed] (A)--(C);
				\foreach \x  /\goc in {A/180,C/0,B/-135,S/90}
				\fill[]  (\x) circle (1pt)
					 ($(\x)+(\goc:3mm)$) node {$\x$};
				\draw pic[draw,angle  radius=2mm]{right angle=C--A--S};
				\draw pic[draw,angle radius=2mm]{right  angle=C--B--A};
			\end{tikzpicture}
		} 
	 \loigiai{Ta có $\heva{&BC\perp AB \\
		&BC\perp SA \,(\text{do}\,SA\perp(ABC))  \\
		&AB\subset(SAB) \\ 
		&SA\subset(SAB) \\ 
		&AB\cap  SA=A} \Rightarrow BC\perp(SAB)$ tại $B$.\\
		 Suy ra $\mathrm{d}(C,(SAB))=CB$.\\
		Xét $\triangle ABC$ vuông cân tại  $B$ có $BC=AB=4a$.\\
		Vậy $\mathrm{d}(C,(SAB))=4a$. 
	} 
\end{ex}

% nguồn: Ok_Toan11_TracNghiem_Lan2
\begin{ex}%[1H8H5-3]%[ma:a66ec506]
	\immini[thm]{Cho hình lập phương $ABCD.A'B'C'D'$ có cạnh bằng $a$ (tham  khảo hình vẽ). Khoảng cách từ $C$ đến mặt phẳng $\left(BDB'D'\right)$ bằng
	 \choice
		{$\sqrt{3} a$} 
		 {\True $\dfrac{\sqrt{2}}{2} a$}
		{$\dfrac{3}{2} a$} 
		{$\sqrt{2}  a$}}
		 {
		\begin{tikzpicture}[scale=0.9,font=\footnotesize,line join=round,line cap=round,>=stealth]
				\def\a{3} 
				 \def\b{2}
				 \def\h{3}
				\path (0:0)  coordinate (A)
					 ++(0:\a) coordinate (D)
					++(-130:\b) coordinate (C) 
					 ($(A)+(C)-(D)$) coordinate (B)
					($(A)+(90:\h)$) coordinate  (A')
					 ($(B)+(90:\h)$) coordinate (B')
					($(C)+(90:\h)$) coordinate (C') 
					($(D)+(90:\h)$) coordinate  (D')
					($(B)!1/2!(D)$) coordinate  (O);
				\draw[dashed] (B)--(A)--(D)--(B) 
					 (C)--(A)--(A');
				\draw[] (C)--(C') 
					 (D)--(D')
					(B)--(B')--(D') 
					(C)--(C') 
					 (B)--(C)--(D)
					(A')--(B')--(C')--(D')--cycle; 
				\foreach \x/\g in {A/180,B/180,C/0,D/0,A'/180,B'/180,C'/0,D'/0,O/-90} 
				\fill[] (\x)  circle (1pt)
					($(\g:4mm)+(\x)$)  node {$\x$};
			\end{tikzpicture}
		} 
	 \loigiai{
		 Gọi $O=AD\cap BC$. \\
		Ta  có $\heva{&CO\perp BD\\
		&CO\perp BB'} \Rightarrow  CO\perp\left(BDB'D'\right) \Rightarrow \mathrm{d}\left(C;\left(BDB'D'\right)\right)=CO$.\\
		Ta có $CO=\dfrac{1}{2} CA=\dfrac{\sqrt{2}}{2}  a$.
	 }
\end{ex}

% nguồn: Ok_Toan11_TracNghiem_Lan2
\begin{ex}%[1H8H5-3]%[ma:f87e8325]
	Cho hình chóp $S.ABC$ có đáy là tam giác vuông cân tại $C$, $BC=a$, $SA$ vuông  góc với mặt phẳng đáy và $SA=a$. Khoảng cách từ $A$ đến mặt phẳng $(SBC)$ bằng
	 \choice
		{$a\sqrt{2}$} 
		{\True $\dfrac{a\sqrt{2}}{2}$} 
		 {$\dfrac{a\sqrt{2}}{2}$}
		{$\dfrac{\sqrt{3}a}{2}$} 
	\loigiai{ 
		\begin{center}
			\begin{tikzpicture}[scale=1,font=\footnotesize,line join=round,line cap=round,>=stealth]
					\coordinate (A) at (0,0); 
					 \coordinate (B) at (4,0);
					 \coordinate (C) at (1.5,-2);
					 \coordinate (S) at ($(A)+(0,3)$);
					\coordinate (H) at  ($(S)!0.4!(C)$);
					\draw (S)--(C)--(A)--(S) 
						(S)--(B) 
						(B)--(C)--(A)--(H); 
					\draw[dashed,thin]  (A)--(B);
					\pic[draw,thin,angle  radius=6pt] {right angle = S--A--B} pic[draw,thin,angle radius=6pt] {right angle = S--A--C} pic[draw,thin,angle radius=6pt] {right angle = A--C--B} pic[draw,thin,angle radius=6pt] {right angle = A--H--C};
					\foreach \i/\g in {S/90,A/180,B/0,C/-90,H/60}{\draw[fill](\i)  circle (1pt)
						 ($(\i)+(\g:3mm)$) node[scale=1]{$\i$};}
				\end{tikzpicture}
		\end{center}
		Trong mặt phẳng $(SAC)$ kẻ $AH  \perp SC$ theo giao tuyến $SC$.\\
		Khi đó $(SBC) \perp (SAC)$ tại $H$ suy ra  $AH \perp (SBC)$.\\
		Khoảng cách từ $A$ đến  mặt phẳng $(SBC)$ bằng $AH$.\\
		Ta có  $AC=BC=a$, $SA=a$ nên tam giác $SAC$ vuông cân tại $A$.\\
		Theo Py-ta-go $SC  = \sqrt{SA^2+AC^2} = \sqrt{a^2+a^2} = a\sqrt{2}$.\\
		Suy ra $AH  = \dfrac{SA \cdot AC}{SC} = \dfrac{a \cdot a}{a\sqrt{2}} = \dfrac{a\sqrt{2}}{2}$.
	} 
\end{ex}

% nguồn: Ok_Toan11_TracNghiem_Lan2
\begin{ex}%[1H8H5-3]%[ma:c4f7c933]
	Cho hình chóp $S.ABC$ có $SA$, $SB$ và $SC$ đôi một vuông góc với nhau. Biết $SA=SB=SC=3$. Khoảng cách từ điểm $S$ đến mặt  phẳng $(ABC)$ bằng
	\choice 
		 {$\dfrac{\sqrt{3}}{3}$}
		 {$\sqrt{2}$}
		{\True $\sqrt{3}$} 
		{$1$} 
	 \loigiai{
		\begin{center}
			\begin{tikzpicture}[scale=0.85,font=\footnotesize,line join=round,line cap=round,>=stealth]
					 \coordinate (S) at (0,0);
					 \coordinate (A) at (0,3.5);
					\coordinate (B) at (2,-1.5); 
					\coordinate (C)  at (5,0);
					 \draw (A) -- (S) -- (B) -- (C) -- (A);
					 \draw (A) -- (B);
					\draw[dashed] (S)  -- (C);
					\pic[draw,thin,angle radius=6pt] {right angle = A--S--C}  pic[draw,thin,angle radius=6pt] {right angle = A--S--B};
					\foreach \i/\g in {A/90,S/160,B/-90,C/-90}{\draw[fill](\i)  circle (1pt)
						($(\i)+(\g:3mm)$) node[scale=1]{$\i$};} 
				\end{tikzpicture}
		\end{center}
		Gọi $\mathrm{d}(S;(ABC))=h$.\\ 
		Ta có $$\dfrac{1}{h^2} = \dfrac{1}{SA^2} + \dfrac{1}{SB^2}  + \dfrac{1}{SC^2} = \dfrac{1}{3^2} + \dfrac{1}{3^2} + \dfrac{1}{3^2} = \dfrac{1}{9}+\dfrac{1}{9}+\dfrac{1}{9} = \dfrac{3}{9} = \dfrac{1}{3}.$$
		Suy ra  $h^2=3 \Leftrightarrow h=\sqrt{3}$.
	} 
\end{ex}

% nguồn: Ok_Toan11_TracNghiem_Lan2
\begin{ex}%[1H8H5-3]%[ma:065d9ebd]
	Cho hình chóp $S.ABCD$ có đáy  là hình vuông cạnh $2a$, biết $SAB$ là tam giác đều và nằm trong mặt phẳng vuông góc với đáy. Tính khoảng cách từ $A$ tới mặt phẳng $(SCD)$.
	 \choice
		{\True $\dfrac{2a\sqrt{21}}{7}$} 
		 {$\dfrac{a\sqrt{14}}{6}$}
		{$\dfrac{3a\sqrt{14}}{7}$} 
		 {$\dfrac{a\sqrt{21}}{16}$}
	 \loigiai{
		\begin{center}
			\begin{tikzpicture}[scale=1,font=\footnotesize,line join=round,line cap=round,>=stealth]
					\coordinate (A) at (0,0); 
					\coordinate (B) at  (-2.2,-0.8);
					 \coordinate (D) at (3.8,0);
					\coordinate  (C) at (1.6,-0.8);
					\coordinate (K)  at (0,4);
					\coordinate (I)  at ($(A)!0.5!(B)$);
					\coordinate (S)  at ($(I)+(K)-(A)$);
					\coordinate (E) at  ($(C)!0.5!(D)$);
					 \coordinate (H) at ($(S)!0.4!(E)$);
					\draw[dashed] (S) -- (A)  -- (D)
						(A)  -- (B)
						(H)--(I) -- (E) 
						(S)  -- (I);
					\draw (E)--(S)--(B)--(C)--(D)--(S)--(C); 
					\foreach \i/\g in {S/90,A/160,B/-90,C/-90,D/0,I/150,E/-90,H/60}{\draw[fill](\i) circle (1pt) 
						($(\i)+(\g:3mm)$) node[scale=1]{$\i$};} 
				\end{tikzpicture}
		\end{center}
		 Ta có $\mathrm{d}(A,(SCD)) = \mathrm{d}(I,(SCD))$.\\
		Gọi $E$ là trung điểm $CD$.\\ 
		Dựng $IH \perp SE$ thì ta có $$\mathrm{d}(I;(SCD)) = IH = \dfrac{IE \cdot IS}{\sqrt{IE^2+IS^2}} =  \dfrac{2a \cdot a\sqrt{3}}{\sqrt{(2a)^2+(a\sqrt{3})^2}} = \dfrac{2a^2\sqrt{3}}{\sqrt{4a^2+3a^2}} = \dfrac{2a^2\sqrt{3}}{\sqrt{7a^2}} = \dfrac{2a\sqrt{21}}{7}.$$
	 }
\end{ex}

% nguồn: Ok_Toan11_TracNghiem_Lan2
\begin{ex}%[1H8H5-3]%[ma:4d8509c9]
	[MĐ 103 BGD\&ĐT NĂM  2020-2021-ĐỢT 2]
	\immini[thm]{Cho hình lập phương $ABCD'A'B'C'D'$ có cạnh bằng $a$. Khoảng  cách từ $A$ đến mặt phẳng $\left(BDD'B'\right)$ bằng
	\choice 
		{$\sqrt{2} a$} 
		{\True $\dfrac{\sqrt{2}}{2} a$} 
		{$\sqrt{3} a$} 
		{$\dfrac{\sqrt{3}}{2}  a$}}
		 {
		\begin{tikzpicture}[scale=0.9,font=\footnotesize,line join=round,line cap=round,>=stealth]
				 \def\a{3}
				 \def\b{2}
				 \def\h{3}
				 \path (0:0) coordinate (A)
					 ++(0:\a) coordinate (D)
					 ++(-130:\b) coordinate (C)
					 ($(A)+(C)-(D)$) coordinate (B)
					($(A)+(90:\h)$) coordinate  (A')
					($(B)+(90:\h)$) coordinate (B') 
					($(C)+(90:\h)$) coordinate  (C')
					($(D)+(90:\h)$) coordinate  (D')
					 ($(B)!1/2!(D)$) coordinate (O);
				\draw[dashed]  (B)--(A)--(D)--(B)
					 (C)--(A)--(A');
				\draw[] (C)--(C') 
					(D)--(D') 
					 (B)--(B')--(D')
					(C)--(C') 
					 (B)--(C)--(D)
					(A')--(B')--(C')--(D')--cycle; 
				\foreach \x/\g in {A/180,B/180,C/0,D/0,A'/180,B'/180,C'/0,D'/0,O/-90} 
				\fill[]  (\x) circle (1pt)
					($(\g:4mm)+(\x)$)  node {$\x$};
			\end{tikzpicture}
		} 
	 \loigiai{
		Gọi  $AC\cap BD=O$. \\
		Khi đó $AO\perp BD$,  mặt khác $AO\perp BB'$. \\
		Suy ra $AO\perp\left(BDB'D'\right)$ hay $AO$ là khoảng cách từ $A$ đến mặt phẳng $\left(BDD'B'\right)$.\\ 
		Ta có $AC=\sqrt{AB^2+BC^2}=a \sqrt{2}$, $AO=\dfrac{AC}{2}=\dfrac{a \sqrt{2}}{2}$. 
	 }
\end{ex}

% nguồn: Ok_Toan11_TracNghiem_Lan2
\begin{ex}%[1H8H5-3]%[ma:5724285b]
	Cho hình chóp $S.ABCD$ có đáy là  hình chữ nhật, biết $AD=2a$, $SA \perp (ABCD)$, $SA=a$. Khoảng cách từ $A$ đến $(SCD)$ bằng
	 \choice
		 {$\dfrac{3a\sqrt{2}}{2}$}
		{$\dfrac{2a\sqrt{3}}{3}$} 
		{\True $\dfrac{2a}{\sqrt{5}}$} 
		{$\dfrac{3a}{\sqrt{7}}$} 
	\loigiai{ 
		\begin{center}
			\begin{tikzpicture}[scale=0.75,font=\footnotesize,line join=round,line cap=round,>=stealth]
					 \coordinate (A) at (0,0);
					\coordinate  (B) at (-2.2,-0.8);
					\coordinate (D) at (3.8,0); 
					 \coordinate (C) at (1.6,-0.8);
					\coordinate (S) at  (0,3.5);
					\coordinate (H)  at ($(S)!0.6!(D)$);
					\draw[dashed] (S)  -- (A) -- (B)
						 (H) -- (A) -- (D);
					\draw (S)--(B)--(C)--(D)--(S)--(C); 
					\draw (0.25,0) --  (0.25,0.25) -- (0,0.25);
					\foreach \i/\g in {S/90,A/-90,B/-90,C/-90,D/0,H/60}{\draw[fill](\i)  circle (1pt)
						 ($(\i)+(\g:3mm)$) node[scale=1]{$\i$};}
				\end{tikzpicture}
		\end{center}
		Do $SA \perp (ABCD) \Rightarrow SA \perp AD$, $SA \perp CD$.\\ 
		Dựng $AH \perp SD$, $(H \in SD)$.\\ 
		Có $\heva{ &CD  \perp AD \\ &CD \perp SA} \Rightarrow CD \perp (SAD) \Rightarrow CD \perp AH$.\\
		Vậy $AH  \perp (SCD) \Rightarrow \mathrm{d}(A,(SCD))=AH = \dfrac{SA \cdot AD}{\sqrt{SA^2+AD^2}} = \dfrac{a \cdot 2a}{\sqrt{a^2+(2a)^2}} = \dfrac{2a}{\sqrt{5}}$.
	 }
\end{ex}

% nguồn: Ok_Toan11_TracNghiem_Lan2
\begin{ex}%[1H8H5-3]%[ma:c6638e80]
	Cho hình chóp $S.ABCD$  có đáy là hình vuông cạnh $a$, $SA \perp (ABCD)$. Khoảng cách từ $S$ đến mặt phẳng $(ABCD)$ là
	 \choice
		{$a\sqrt{2}$} 
		 {\True $a$}
		 {$\dfrac{a}{2}$}
		{$\dfrac{3a}{4}$} 
	 \loigiai{
		\begin{center}
			\begin{tikzpicture}[scale=0.6,font=\footnotesize,line join=round,line cap=round,>=stealth]
					\coordinate (A) at (0,0); 
					\coordinate  (B) at (-2.2,-0.8);
					\coordinate (D) at  (3.8,0);
					\coordinate (C) at (1.6,-0.8); 
					 \coordinate (S) at (0,3.5);
					\draw[dashed] (S) -- (A) -- (B) 
						 (A) -- (D);
					 \draw (S)--(B)--(C)--(D)--(S)--(C);
					\draw (0.25,0) -- (0.25,0.25)  -- (0,0.25);
					\foreach \i/\g in {S/90,A/-90,B/-90,C/-90,D/0}{\draw[fill](\i)  circle (1pt)
						 ($(\i)+(\g:3mm)$) node[scale=1]{$\i$};}
				\end{tikzpicture}
		\end{center}
		Vì $SA \perp (ABCD)$ nên $\mathrm{d}(S,(ABCD))  = SA=a$.
	} 
\end{ex}

% nguồn: Ok_Toan11_TracNghiem_Lan2
\begin{ex}%[1H8H5-3]%[ma:58cbbc92]
	Cho hình chóp $S.ABC$ có đáy là tam giác đều cạnh $a$ và $SA \perp (ABC)$. Tính khoảng cách  từ $C$ đến $(SAB)$.
	 \choice
		 {$\dfrac{a\sqrt{3}}{4}$}
		 {\True $\dfrac{a\sqrt{3}}{2}$}
		{$\dfrac{a\sqrt{2}}{3}$} 
		 {$a$}
	 \loigiai{
		\begin{center}
			\begin{tikzpicture}[scale=0.85,font=\footnotesize,line join=round,line cap=round,>=stealth]
					\coordinate (A) at  (0,0);
					 \coordinate (S) at (0,3.5);
					\coordinate (B) at  (2,-1.5);
					\coordinate  (C) at (5,0);
					\coordinate  (H) at ($(A)!0.5!(B)$);
					\draw (S) -- (A) -- (B) -- (C) --  (S);
					 \draw (S) -- (B);
					\draw[dashed] (A) -- (C)  -- (H);
					\draw ($(A)!0.3cm!(B)$) -- ($($(A)!0.3cm!(B)$) +  ($(A)!0.2cm!(S)$) - (A)$) -- ($(A)!0.2cm!(S)$);
					 \foreach \i/\g in {S/90,A/160,B/-90,C/-90,H/180}{\draw[fill](\i) circle (1pt)
						 ($(\i)+(\g:3mm)$) node[scale=1]{$\i$};}
				\end{tikzpicture}
		\end{center}
		Gọi $H$ là hình chiếu của $A$ trên  $SB$.\\
		Ta có $BC \perp AB$, $BC \perp SA \Rightarrow BC  \perp (SAB) \Rightarrow BC \perp AH$.\\
		Từ đó $AH \perp (SBC) \Rightarrow \mathrm{d}(A,(SBC)) = AH$.\\ 
		Xét tam giác vuông $SAB$, ta có $AH = \dfrac{SA \cdot AB}{\sqrt{SA^2+AB^2}} = \dfrac{a\sqrt{3}}{2}$.\\ 
		Vậy  $\mathrm{d}(A,(SBC)) = \dfrac{a\sqrt{3}}{2}$.
	 }
\end{ex}

% nguồn: Ok_Toan11_TracNghiem_Lan2
\begin{ex}%[1H8H5-3]%[ma:7ff0fa2c]
	 Cho hình chóp $S.ABCD$ có đáy là hình vuông cạnh $a$, $SA$ vuông góc với mặt phẳng $(ABCD)$ và $SA=a$. Khoảng cách từ điểm $A$ đến mặt phẳng $(SBC)$ bằng bao nhiêu?
	 \choice
		 {$\dfrac{a}{2}$}
		{\True  $\dfrac{a\sqrt{2}}{2}$}
		 {$a\sqrt{2}$}
		{$\dfrac{a\sqrt{3}}{2}$} 
	\loigiai{ 
		\begin{center}
			\begin{tikzpicture}[scale=0.75,font=\footnotesize,line join=round,line cap=round,>=stealth]
					\coordinate (A) at  (0,0);
					\coordinate  (B) at (-2.2,-0.8);
					\coordinate (D) at  (3.8,0);
					\coordinate (C) at  (1.6,-0.8);
					\coordinate (S) at (0,3.5); 
					\coordinate  (H) at ($(S)!0.6!(B)$);
					\draw[dashed] (S) -- (A) -- (B) 
						(H) -- (A)  -- (D);
					\draw  (S)--(B)--(C)--(D)--(S)--(C);
					\pic[draw,thin,angle radius=6pt] {right angle  = S--A--D} pic[draw,thin,angle radius=6pt] {right angle = S--H--A};
					\foreach \i/\g in {S/90,A/-90,B/-90,C/-90,D/0,H/180}{\draw[fill](\i) circle (1pt) 
						($(\i)+(\g:3mm)$) node[scale=1]{$\i$};} 
				\end{tikzpicture}
		\end{center}
		Trong $(SAB)$  vẽ $AH \perp SB$ tại $H$.\\
		 Ta có $BC \perp (SAB) \Rightarrow (SBC) \perp (SAB)$.\\
		Khi đó  $\heva{ &(SAB) \cap (SBC) = SB \\ &AH \perp SB} \Rightarrow AH \perp (SBC)$, hay $\mathrm{d}(A,(SBC))=AH$.\\
		Trong tam  giác vuông $SAB$, ta có $AH = \dfrac{SA \cdot AB}{\sqrt{SA^2+AB^2}} = \dfrac{a \cdot a}{\sqrt{a^2+a^2}} = \dfrac{a^2}{a\sqrt{2}} = \dfrac{a\sqrt{2}}{2}$.
	 }
\end{ex}

% nguồn: Ok_Toan11_TracNghiem_Lan2
\begin{ex}%[1H8H5-3]%[ma:92a79112]
	Cho hình chóp $S.ABC$ có đáy là tam giác vuông tại $A$, $AB=a, AC=a\sqrt{2}$ và $SA \perp (ABC)$, $SA=a$.  Khoảng cách từ $A$ đến mặt phẳng $(SBC)$ bằng
	 \choice
		 {$\dfrac{a\sqrt{3}}{2}$}
		 {$\dfrac{a\sqrt{2}}{2}$}
		{$\dfrac{a\sqrt{5}}{2}$} 
		{\True  $\dfrac{a\sqrt{10}}{5}$}
	\loigiai{ 
		\begin{center}
			\begin{tikzpicture}[scale=1,font=\footnotesize,line join=round,line cap=round,>=stealth]
					\coordinate (A) at  (0,0);
					\coordinate (C) at (4,0); 
					\coordinate (B) at (1.5,-2); 
					\coordinate (S)  at ($(A)+(0,3)$);
					\coordinate (H) at  ($(B)!0.45!(C)$);
					\coordinate (K)  at ($(S)!0.4!(H)$);
					 \draw (H)--(S)--(B)--(A)--(S)
						 (S)--(C)
						 (B)--(C);
					\draw[dashed,thin] (H)--(A)--(C) 
						(A)--(K); 
					\pic[draw,thin,angle radius=6pt] {right angle = S--A--C} pic[draw,thin,angle radius=6pt] {right angle = S--A--B}  pic[draw,thin,angle radius=6pt] {right angle = A--K--H} pic[draw,thin,angle radius=6pt] {right angle = B--A--C} pic[draw,thin,angle radius=6pt] {right angle = A--H--B};
					\foreach \i/\g in {S/90,A/180,C/0,B/-90,K/60,H/0}{\draw[fill](\i) circle (1pt) 
						($(\i)+(\g:3mm)$) node[scale=1]{$\i$};} 
				\end{tikzpicture}
		\end{center}
		 Gọi $H$ là hình chiếu của $A$ trên $BC$, $K$ là hình chiếu của $A$ trên $SH$.\\
		Ta có $\heva{ &SA \perp  BC \\ &AH \perp BC} \Rightarrow BC \perp (SAH)$.\\
		Do đó $AK  \perp (SBC)$.\\
		 Vậy $\mathrm{d}(A,(SBC)) = AK$.\\
		Trong tam giác vuông $ABC$, ta có $$\dfrac{1}{AH^2} = \dfrac{1}{AB^2}+\dfrac{1}{AC^2} = \dfrac{1}{a^2}+\dfrac{1}{(a\sqrt{2})^2} = \dfrac{1}{a^2}+\dfrac{1}{2a^2} = \dfrac{3}{2a^2}.$$ 
		Trong tam giác vuông $SAH$, ta có $$\dfrac{1}{AK^2} = \dfrac{1}{SA^2}+\dfrac{1}{AH^2} = \dfrac{1}{a^2}+\dfrac{3}{2a^2} = \dfrac{5}{2a^2}.$$ 
		$\Rightarrow AK =  \dfrac{a\sqrt{10}}{\sqrt{5}}.$
	 }
\end{ex}

% nguồn: Ok_Toan11_TracNghiem_Lan2
\begin{ex}%[1H8H5-3]%[ma:5e14c00d]
	Cho  hình chóp $S.ABCD$ có đáy là hình thoi tâm $O$, cạnh $a$, góc $\widehat{BAD}=60^\circ$, đường thẳng $SO$ vuông góc với $(ABCD)$ và $SO=a$. Khoảng cách từ điểm $A$ đến mặt phẳng $(SBC)$ bằng
	\choice 
		 {$\dfrac{a\sqrt{21}}{7}$}
		{$\dfrac{a\sqrt{57}}{19}$} 
		{\True $\dfrac{2a\sqrt{57}}{19}$} 
		 {$\dfrac{a\sqrt{21}}{14}$}
	\loigiai{ 
		\begin{center}
			\begin{tikzpicture}[scale=1,font=\footnotesize,line join=round,line cap=round,>=stealth]
					\coordinate (A) at (0,0); 
					\coordinate (D) at (-2,-1); 
					 \coordinate (B) at (4,0);
					\coordinate (C)  at ($(B)+(D)-(A)$);
					\coordinate (O) at  ($(A)!0.5!(C)$);
					\coordinate (S)  at ($(O)+(0,4)$);
					\coordinate (I) at ($(C)!0.4!(B)$); 
					\coordinate  (M) at ($(S)!0.7!(I)$);
					\draw  (S)--(B)
						 (S)--(C)
						 (S)--(D)
						 (B)--(C)--(D)
						(S)--(I); 
					\draw[dashed,thin](S)--(O)--(M) 
						(O)--(I) 
						 (S)--(A)--(C)
						 (A)--(B)
						 (A)--(D)--(B);
					\foreach  \i/\g in {S/90,A/150,D/-90,C/-90,B/0,O/-90,I/-30,M/30}{\draw[fill](\i) circle (1pt)
						($(\i)+(\g:3mm)$) node[scale=1]{$\i$};} 
				\end{tikzpicture}
		\end{center}
		 Gọi $I, M$ lần lượt là hình chiếu của $O$ lên $BC, SI$.\\
		Ta có $AC=2OC \Rightarrow \mathrm{d}(A,(SBC)) = 2\mathrm{d}(O,(SBC))  = 2OM$.\\
		 Do $\widehat{BAD} = 60^\circ$ nên tam giác $ABD$ đều, $OB = \dfrac{a}{2}, OA=OC=\dfrac{a\sqrt{3}}{2}$.\\
		Trong tam giác $OBC$  vuông tại $O$ có $$\dfrac{1}{OI^2} = \dfrac{1}{OB^2} + \dfrac{1}{OC^2} = \dfrac{4}{a^2} + \dfrac{4}{3a^2} = \dfrac{16}{3a^2}.$$
		Trong tam giác $SOI$ vuông tại $O$ có $$\dfrac{1}{OM^2} = \dfrac{1}{OI^2} + \dfrac{1}{OS^2} = \dfrac{16}{3a^2} + \dfrac{1}{a^2} =  \dfrac{19}{3a^2} \Rightarrow OH = \dfrac{a\sqrt{57}}{19}.$$
		 Vậy $\mathrm{d}(A,(SBC)) = \dfrac{2a\sqrt{57}}{19}$.
	} 
\end{ex}

% nguồn: Ok_Toan11_TracNghiem_Lan2
\begin{ex}%[1H8H5-3]%[ma:3523d6c5]
	[MĐ 101  BGD\&ĐT NĂM 2017-2018]
	Cho hình chóp $S. ABC$ có đáy là tam giác vuông đỉnh $B$, $AB=a$, $SA$ vuông  góc với mặt phẳng đáy và $SA=2a$. Khoảng cách từ $A$ đến mặt phẳng $(SBC)$ bằng
	 \choice
		 {\True $\dfrac{2\sqrt{5} a}{5}$}
		{$\dfrac{\sqrt{5} a}{3}$} 
		{$\dfrac{2\sqrt{2} a}{3}$} 
		{$\dfrac{\sqrt{5}  a}{5}$}
	 \loigiai{
	\immini{Ta có  $\heva{&BC\perp AB\\
		 &BC\perp SA} \Rightarrow BC\perp(SAB)$.\\
		Kẻ  $AH\perp SB$. \\
		Khi  đó $AH\perp BC\Rightarrow AH\perp(SBC)$.\\
		Suy ra $AH$ là khoảng cách  từ $A$ đến mặt phẳng $(SBC)$.\\
		Ta có $\dfrac{1}{AH^2}=\dfrac{1}{SA^2}+\dfrac{1}{AB^2}=\dfrac{1}{4a^2}+\dfrac{1}{a^2}=\dfrac{5}{4a^2}$.\\ 
		Suy ra $AH^2=\dfrac{4a^2}{5}\Rightarrow AH=\dfrac{2\sqrt{5}  a}{5}$.}
		{ 
		\begin{tikzpicture}[scale=0.9,font=\footnotesize,line join=round,line cap=round,>=stealth]
				\def\a{4} 
				 \def\h{3}
				\path (0:0)  coordinate (A)
					++(0:\a) coordinate  (C)
					(-40:4*\a/5) coordinate  (B)
					($(A)+(90:\h)$) coordinate  (S)
					($(S)!1/3!(B)$)  coordinate (H);
				\draw[] (A)--(B)--(C) 
					(A)--(H) 
					(A)--(S) 
					 (B)--(S)
					 (C)--(S);
				 \draw[dashed] (A)--(C);
				 \foreach \x /\goc in {A/180,C/0,B/-135,S/90,H/45}
				\fill[] (\x) circle  (1pt)
					($(\x)+(\goc:3mm)$) node  {$\x$};
				\draw pic[draw,angle radius=2mm]{right angle=C--A--S}; 
				\draw pic[draw,angle  radius=2mm]{right angle=C--B--A};
				 \draw pic[draw,angle radius=2mm]{right angle=A--H--B};
			\end{tikzpicture}
		 }
	} 
\end{ex}

% nguồn: Ok_Toan11_TracNghiem_Lan2
\begin{ex}%[1H8H5-3]%[ma:60372fa5]
	Cho hình chóp $S.ABCD$ có đáy là hình chữ nhật, biết $AD=2a$, $SA \perp (ABCD)$, $SA=a$. Khoảng cách  từ $A$ đến $(SCD)$ bằng
	\choice 
		 {$\dfrac{3a}{\sqrt{7}}$}
		{\True  $\dfrac{2a}{3}$}
		 {$\dfrac{2a\sqrt{3}}{3}$}
		{$\dfrac{2a}{\sqrt{5}}$} 
	 \loigiai{
		\begin{center}
			\begin{tikzpicture}[scale=0.75,font=\footnotesize,line join=round,line cap=round,>=stealth]
					\coordinate (A) at  (0,0);
					\coordinate (B) at (-2.2,-0.8); 
					\coordinate (D)  at (3.8,0);
					\coordinate (C) at (1.6,-0.8); 
					\coordinate (S)  at (0,3.5);
					\coordinate (H)  at ($(S)!0.6!(B)$);
					\draw[dashed] (S) -- (A)  -- (B)
						(H) --  (A) -- (D);
					 \draw (S)--(B)--(C)--(D)--(S)--(C);
					\draw (0.25,0) --  (0.25,0.25) -- (0,0.25);
					\foreach \i/\g  in {S/90,A/-90,B/-90,C/-90,D/0,H/120}{\draw[fill](\i) circle (1pt)
						($(\i)+(\g:3mm)$)  node[scale=1]{$\i$};}
				\end{tikzpicture}
		\end{center}
		Gọi $H$ là hình chiếu của $A$ lên cạnh  $SD$.\\
		Ta có 
		$\heva{&CD \perp AD \\ &CD \perp SA }\Rightarrow CD \perp (SAD) \Rightarrow CD \perp  AH$.\\
		Suy ra 
		$\heva{&AH \perp SD  \\ &AH \perp CD} \Rightarrow AH \perp (SCD)$.\\
		Suy ra khoảng cách từ $A$ đến $(SCD)$  bằng $AH$.\\
		Ta có $AH =  \dfrac{SA \cdot AD}{\sqrt{SA^2+AD^2}} = \dfrac{a \cdot 2a}{\sqrt{a^2+(2a)^2}} = \dfrac{2a^2}{\sqrt{5a^2}} = \dfrac{2a}{\sqrt{5}}$.
	 }
\end{ex}

% nguồn: Ok_Toan11_TracNghiem_Lan2
\begin{ex}%[1H8H5-3]%[ma:1622c082]
	Cho hình lăng trụ đứng $ABC.A'B'C'$ có tất  cả các cạnh bằng $a$. Gọi $M$ là trung điểm của $AA'$. Khoảng cách từ $M$ đến mặt phẳng $(AB'C)$ bằng
	 \choice
		{\True  $\dfrac{a\sqrt{21}}{14}$}
		{$\dfrac{a\sqrt{21}}{7}$} 
		{$\dfrac{a\sqrt{2}}{2}$} 
		{$\dfrac{a\sqrt{2}}{4}$} 
	 \loigiai{
		\begin{center}
			\begin{tikzpicture}[scale=1,font=\footnotesize,line join=round,line cap=round,>=stealth]
					\coordinate (A)  at (0,0);
					\coordinate  (B) at (1,-1.5);
					\coordinate (C)  at (4,0);
					\coordinate (A') at ($(A)+(0,3)$); 
					\coordinate (B') at  ($(B)+(0,3)$);
					\coordinate (C')  at ($(C)+(0,3)$);
					\coordinate (K) at ($(A)!0.5!(C)$); 
					 \coordinate (M) at ($(A)!0.5!(A')$);
					 \coordinate (H) at ($(B')!0.45!(K)$);
					 \draw (A)--(B)--(C)--(C')--(B')--(A')--(A)
						(B')--(B) 
						 (A')--(C')
						(A)--(B')--(C); 
					\draw[dashed,thin](A)--(C) 
						 (B)--(K)--(B')
						 (B)--(H);
					 \pic[draw,thin,angle radius=6pt] {right angle = B--K--A} pic[draw,thin,angle radius=2mm] {right angle = B--H--K};
					\foreach  \i/\g in {A'/90,A/-90,B/-90,C/-90,C'/0,B'/180,K/60,M/180,H/30}{\draw[fill](\i) circle (1pt)
						 ($(\i)+(\g:3mm)$) node[scale=1]{$\i$};}
				\end{tikzpicture}
		\end{center}
		Gọi $K$ là trung  điểm của $A'B$, dựng $BH \perp B'K$ tại $H$.\\
		Ta có $\mathrm{d}(M,(A'BC'))  = \dfrac{1}{2}\mathrm{d}(B,(A'BC')) = \dfrac{1}{2}\mathrm{d}(B,(AB'C))$.\\
		$\heva{ &BH \perp B'K \\ &BH \perp AC } \Rightarrow BH \perp (A'BC')$ $\Rightarrow \mathrm{d}(B,(A'BC'))  = BH$.\\
		Xét tam giác vuông  $BB'K$ ta có
		$$\dfrac{1}{BH^2} =  \dfrac{1}{BB'^2} + \dfrac{1}{BK^2} = \dfrac{1}{a^2} + \dfrac{1}{\left(\dfrac{a\sqrt{3}}{2}\right)^2} = \dfrac{1}{a^2}+\dfrac{4}{3a^2} = \dfrac{7}{3a^2}.$$
		 $\Rightarrow BH = \dfrac{a\sqrt{3}}{\sqrt{7}} = \dfrac{a\sqrt{21}}{7}$.\\
		Vậy $\mathrm{d}(M,(A'BC')) = \dfrac{1}{2}BH  = \dfrac{1}{2}\dfrac{a\sqrt{21}}{7} = \dfrac{a\sqrt{21}}{14}$.
	} 
\end{ex}

% nguồn: Ok_Toan11_TracNghiem_Lan2
\begin{ex}%[1H8H5-3]%[ma:577ffbe5]
	Cho hình chóp $S.ABCD$ có đáy là hình thang vuông tại $A$ và  $B$, $AD=2AB=2BC=2a$, cạnh bên $SA$ vuông góc với $(ABCD)$, $SA=a\sqrt{3}$. Khoảng cách từ $A$ đến $(SBC)$ bằng
	\choice 
		{$\dfrac{a\sqrt{5}}{2}$} 
		 {\True $\dfrac{a\sqrt{3}}{2}$}
		{$\dfrac{a\sqrt{21}}{7}$} 
		{$2a$} 
	 \loigiai{
		\begin{center}
			\begin{tikzpicture}[scale=1,font=\footnotesize,line join=round,line cap=round,>=stealth]
					\coordinate (A) at  (0,0);
					\coordinate  (B) at (-1,-0.8);
					\coordinate  (D) at (3.8,0);
					\coordinate (C) at (1.6,-0.8); 
					\coordinate  (S) at (0,3.5);
					\coordinate (H) at ($(S)!0.6!(B)$); 
					\draw[dashed] (S) -- (A)  -- (B)
						(H) --  (A) -- (D);
					\draw (S)--(B)--(C)--(D)--(S)--(C); 
					\pic[draw,thin,angle radius=2mm]  {right angle = D--A--B} pic[draw,thin,angle radius=2mm] {right angle = C--B--A} pic[draw,thin,angle radius=2mm] {right angle = S--H--A};
					\foreach  \i/\g in {S/90,A/30,B/-90,C/-90,D/0,H/180}{\draw[fill](\i) circle (1pt)
						 ($(\i)+(\g:3mm)$) node[scale=1]{$\i$};}
				\end{tikzpicture}
		\end{center}
		Gọi $H$ là hình chiếu của $A$ trên $SB$.\\ 
		Ta có $BC  \perp AB$, $BC \perp SA \Rightarrow BC \perp (SAB) \Rightarrow BC \perp AH$.\\
		Từ đó $\mathrm{d}(A,(SBC)) = AH$.\\ 
		Xét tam giác vuông $SAB$, ta có $$AH = \dfrac{SA \cdot AB}{\sqrt{SA^2+AB^2}}  = \dfrac{a\sqrt{3} \cdot 2a}{\sqrt{(a\sqrt{3})^2+(2a)^2}} = \dfrac{2a^2\sqrt{3}}{\sqrt{3a^2+4a^2}} = \dfrac{2a^2\sqrt{3}}{\sqrt{7a^2}} = \dfrac{2a\sqrt{3}}{\sqrt{7}}.$$
	 }
\end{ex}

% nguồn: Ok_Toan11_TracNghiem_Lan2
\begin{ex}%[1H8H5-3]%[ma:18d9321d]
	Cho  hình chóp đều $S.ABCD$ có cạnh đáy bằng $2a$, cạnh bên bằng $3a$. Khoảng cách từ điểm $A$ đến mặt phẳng $(SCD)$ bằng
	\choice 
		{$\dfrac{a\sqrt{14}}{3}$} 
		 {$\dfrac{a\sqrt{14}}{4}$}
		{$a\sqrt{14}$} 
		 {\True $\dfrac{a\sqrt{14}}{2}$}
	 \loigiai{
		\begin{center}
			\begin{tikzpicture}[scale=1,font=\footnotesize,line join=round,line cap=round,>=stealth]
					\coordinate  (A) at (0,0);
					\coordinate (B) at  (-2,-1);
					 \coordinate (D) at (4,0);
					\coordinate (C) at  ($(D)+(B)-(A)$);
					 \coordinate (O) at ($(A)!0.5!(C)$);
					\coordinate (S)  at ($(O)+(0,4)$);
					\coordinate  (I) at ($(C)!0.5!(D)$);
					\coordinate (M) at ($(S)!0.7!(I)$); 
					 \draw (S)--(D)
						 (S)--(C)
						(S)--(B) 
						 (D)--(C)--(B)
						(S)--(I); 
					\draw[dashed,thin](S)--(O)--(M) 
						 (O)--(I)
						(S)--(A)--(C) 
						(A)--(D) 
						 (A)--(B)--(D);
					\foreach \i/\g in {S/90,A/150,B/-90,C/-90,D/0,O/-90,I/-30,M/30}{\draw[fill](\i)  circle (1pt)
						 ($(\i)+(\g:3mm)$) node[scale=1]{$\i$};}
				\end{tikzpicture}
		\end{center}
	 \loigiai{
		Gọi $O = AC \cap DB$.\\ 
		Vì $S.ABCD$ là hình chóp đều nên $SO \perp (ABCD)$ và  đáy $ABCD$ là hình vuông.\\
		Ta có $\dfrac{\mathrm{d}(A,(SCD))}{\mathrm{d}(O,(SCD))} = \dfrac{AC}{OC} = 2 \Rightarrow \mathrm{d}(A,(SCD))  = 2\mathrm{d}(O,(SCD))$.\\
		Gọi $I$ và $M$ lần lượt là hình chiếu của $O$ trên $CD$ và $SI$.\\ 
		Khi  đó $\mathrm{d}(O;(SCD))=OM$.\\
		Ta  có $OI=\dfrac{1}{2}BC=a$, $SO^2=SC^2-OC^2=9a^2-2a^2=7a^2$.\\
		Lại có $\dfrac{1}{OM^2}=\dfrac{1}{SO^2}+\dfrac{1}{OI^2}=\dfrac{1}{7a^2}+\dfrac{1}{a^2}=\dfrac{8}{7a^2}  \Rightarrow OM=\dfrac{a\sqrt{14}}{4}$.\\
		 Vậy $\mathrm{d}(A,(SCD))=\dfrac{a\sqrt{14}}{2}$.
	 }
	 }
\end{ex}

% nguồn: Ok_TracNghiemLop11_Lan1
\begin{ex}%[1H8H5-3]%[ma:40095466]
	Cho hình chóp tứ giác đều $S.ABCD$ có cạnh đáy bằng $a$, cạnh bên bằng $2a$. Gọi  $O$ là tâm của $ABCD$ và $G$là trọng tâm của tam giác $ACD$. Khoảng cách từ điểm $G$ đến mặt phẳng $(SBC)$ bằng
	\choice 
		{$\dfrac{a\sqrt{210}}{30}$} 
		{\True $\dfrac{2a\sqrt{210}}{45}$} 
		 {$\dfrac{a\sqrt{210}}{45}$}
		 {$\dfrac{a\sqrt{210}}{15}$}
	\loigiai{ 
		\begin{center}
			\begin{tikzpicture}[scale=1,font=\footnotesize,line join=round,line cap=round,>=stealth]
					 \path
						(0,0)  coordinate (A)
						++(-140:3) coordinate (B) 
						 ++(0:4) coordinate (C)
						 ($(A)+(C)-(B)$) coordinate (D)
						 ($(A)!1/2!(C)$) coordinate (O)
						 ($(O)+(0,3.5)$) coordinate (S)
						($(B)!1/2!(C)$)  coordinate (M)
						($(S)!(O)!(M)$) coordinate  (H);
					\foreach  \i in{B,C,D}{\draw (S)--(\i);};
					 \draw (B)--(C)--(D)
						(S)--(M); 
					 \draw[dashed] (S)--(A)--(B)
						(C)--(A)--(D)--(B) 
						(S)--(O)--(M) 
						 (O)--(H);
					\pic[draw,angle eccentricity=1.8,angle  radius=2mm]{right angle=S--O--D};
					\pic[draw,angle  eccentricity=1.8,angle radius=2mm]{right angle=O--H--S};
					\foreach \i/\g in  {A/180,B/180,C/-90,D/0,O/-70,S/90,M/-90,H/180}
					\fill[black] (\i) circle (1pt)+(\g:3mm) node[scale=1]{$\i$}; 
				\end{tikzpicture}
		\end{center}
		Vì $\dfrac{\mathrm{d}(G,(SBC))}{\mathrm{d}(O,(SBC))}=\dfrac{GB}{OB}=\dfrac{4}{3}\Rightarrow  \mathrm{d}(G,(SBC))=\dfrac{4}{3}\mathrm{d}(O,(SBC))$.\\
		Gọi  $M$ là trung điểm của $BC$.\\
		Khi  đó $\heva{
		& BC\perp  OM \\
		& BC\perp SO  \\
		}\Rightarrow BC\perp (SOM)$.\\ 
		Trong  mặt phẳng $(SOM)$, kẻ $OH\perp SM$.\\
		Khi đó $OH\subset (SOM)$ nên $BC\perp OH$, suy ra $OH\perp  (SBC)$ hay $\mathrm{d}(O,(SBC))=OH$.\\
		Ta  có $SO=\sqrt{(2a)^2-{{\left(\dfrac{a\sqrt{2}}{2}\right)}^2}}=\dfrac{a\sqrt{14}}{2}$.\\
		Xét  $\Delta SOM$ vuông tại $O$ có $$\dfrac{1}{OH^2}=\dfrac{1}{OM^2}+\dfrac{1}{SO^2}=\dfrac{1}{{{\left(\dfrac{a}{2}\right)}^2}}+\dfrac{1}{{{\left(\dfrac{a\sqrt{14}}{2}\right)}^2}}=\dfrac{30}{7a^2}\Rightarrow OH=\dfrac{a\sqrt{210}}{30}.$$
	Vậy $\mathrm{d}(G,(SBC))=\dfrac{4}{3}\mathrm{d}(O,(SBC))=\dfrac{2a\sqrt{210}}{45}$.} 
\end{ex}

% nguồn: Ok_TracNghiemLop11_Lan1
\begin{ex}%[1H8H5-3]%[ma:7155f56d]
	Cho hình chóp $S.ABCD$ có đáy là hình vuông cạnh $a$. Biết $SA=2a$ và $SA$ vuông góc với mặt phẳng $(ABCD)$. Khoảng cách từ  điểm $S$ đến mặt phẳng $(ABCD)$ bằng
	 \choice
		 {$a$}
		 {\True $2a$}
		 {$3a$}
		{$a\sqrt{2}$} 
	\loigiai{ 
	Vì  $SA\perp (ABCD)$ nên $\mathrm{d}(S,(ABCD))=SA=2a$}
\end{ex}

% nguồn: Ok_TracNghiemLop11_Lan1
\begin{ex}%[1H8H5-3]%[ma:a50df697]
	Cho hình hộp chữ nhật $ABCD.A'B'C'D'$. Khoảng cách giữa hai mặt  phẳng $(ABCD)$ và $(A'B'C'D')$ bằng
	 \choice
		{$AC'$} 
		 {$AB'$}
		{$81{m}$} 
		 {\True $AA'$}
	\loigiai{ 
	\immini 
		 {
		Khoảng cách giữa hai mặt phẳng $(ABCD)$ và $(A'B'C'D')$ bằng $AA'$. 
		}{ 
		\begin{tikzpicture}[scale=0.6,font=\footnotesize,line join=round,line cap=round,>=stealth]
				 \path
					(0,0) coordinate (A) 
					 ++(-130:3) coordinate (B)
					++(0:4) coordinate  (C)
					 ($(A)+(C)-(B)$) coordinate (D)
					($(A)!1/2!(C)$)  coordinate (O);
				\foreach  \i in {A,B,C,D}{
					\coordinate (\i') at ($(\i)+(0,4)$); 
				 }
				\draw (A')--(B')--(C')--(D')--cycle; 
				\draw  (B)--(B')
					 (C)--(C')
					 (D)--(D')
					 (B)--(C)--(D);
				 \draw[dashed](B)--(A)--(A')
					 (A)--(D);
				 \foreach \i/\g in {A'/90,B'/90,C'/90,D'/90,A/-90,B/-90,C/-90,D/-90}
				\fill[black] (\i) circle (1pt)+(\g:5mm) node[scale=1]{$\i$}; 
			\end{tikzpicture}
		} 
	} 
\end{ex}

% nguồn: Ok_TracNghiemLop11_Lan1
\begin{ex}%[1H8H5-3]%[ma:ae92e834]
	Cho hình  hộp chữ nhật $ABCD.A'B'C'D'$ có $AB=a$, $AD=b$, $AA'=c$. Khẳng định nào sau đây là \textbf{sai}?
	\choice 
		{\True Khoảng cách  giữa đường thẳng $AB$ và $\left(A'B'C'D'\right)$ bằng $a$}
		{Khoảng  cách giữa đường thẳng $AC$ và đường thẳng $B'C'$ bằng $c$}
		{Khoảng cách giữa đường thẳng  $AD$ và $\left(A'B'C'D'\right)$ bằng $c$}
		{Khoảng cách giữa  đường thẳng $A$ và $\left(CDC'\right)$ bằng $b$}
	 \loigiai{
	 \immini
		{Ta  có $AB\parallel A'B'$, suy ra $AB \parallel (A'B'C'D')$.\\
		Do đó $\mathrm{d}\left(AB,\left(A'B'C'C'\right)\right)=\mathrm{d}\left(A,\left(A'B'C'D'\right)\right)=AA'=c$. 
		} 
		{ 
		\begin{tikzpicture}[scale=0.7,font=\footnotesize,line join=round,line cap=round,>=stealth]
				 \def\bc{5}
				 \def\ba{2}
				 \def\h{4}
				\def\gocB{35} 
				 \coordinate[label=below left:$B$] (B) at (0,0);
				\coordinate[label=above left:$A$] (A) at  (\gocB:\ba);
				\coordinate[label=below:$C$] (C)  at (\bc,0);
				 \coordinate[label=right:$D$] (D) at ($(C)-(B)+(A)$);
				\coordinate[label=above left:$A'$]  (A') at ($(A)+(90:\h)$);
				\coordinate[label=left:$B'$] (B') at  ($(B)-(A)+(A')$);
				\coordinate[label=below right:$C'$] (C') at ($(C)-(A)+(A')$); 
				\coordinate[label=right:$D'$] (D') at ($(D)-(A)+(A')$); 
				\draw  (B')--(B)--(C)--(D)--(D')--(A')--(B')--(C')--(D')
					 (C)--(C');
				\draw[dashed]  (A')--(A)--(D)
					 (A)--(B);
				\foreach \diem in {A,B,C,D,A',B',C',D'} \fill (\diem) circle (1pt); 
			\end{tikzpicture}
		} 
	 }
\end{ex}

% nguồn: Ok_TracNghiemLop11_Lan1
\begin{ex}%[1H8H5-3]%[ma:8832e4f1]
	 Cho hình chóp $S.ABC$ có $ABC$ tam giác đều cạnh $2a$, $SA\perp \left( ABC \right)$ và $SA=a\sqrt{6}$. Gọi $M$ là trung điểm của $BC$, khi đó khoảng cách từ $A$ đến đường thẳng $SM$ bằng
	 \choice
		{\True $a\sqrt{2}$} 
		{$a\sqrt{3}$} 
		{$a\sqrt{6}$} 
		 {$a\sqrt{11}$}
	\loigiai{ 
		\begin{center}
			\begin{tikzpicture}[line join=round,line cap=round]
					\coordinate  (A) at (-2,0);
					\coordinate (B) at  (0,-3);
					 \coordinate (C) at (3,0);
					\coordinate (S) at ($(A)+(0,4)$); 
					\coordinate (M) at  ($(B)!0.5!(C)$);
					\coordinate (K)  at ($(S)!0.5!(M)$);
					\draw  (M)--(S)--(A)
						 (S)--(B)
						 (S)--(C)
						(B)--(C) 
						 (A)--(B);
					\draw[dashed](M)--(A)--(C) 
						(A)--(K); 
					\pic[draw,angle radius=2mm] {right angle =  S--A--B}
						pic[draw,angle  radius=2mm] {right angle = A--K--M}
						pic[draw,angle radius=2mm] {right angle =  S--A--C};
					\foreach \i/\g  in {S/90,A/180,B/-90,C/0,M/0,K/0}{\draw[fill=white](\i) circle (1pt)
						($(\i)+(\g:3mm)$)  node[scale=1]{$\i$};}
				\end{tikzpicture}
		\end{center}
		Ta có tam giác $ABC$ đều và $M$ là trung điểm $BC$, suy ra $AM=  a\sqrt{3}$.\\
		 Kẻ $AK \perp SM$, suy ra $\mathrm{d}[A,SM]=AK=\sqrt{\dfrac{SA^2\cdot AM^2}{SA^2+AM^2}}=\sqrt{\dfrac{6a^2\cdot 3a^2}{6a^2 + 3a^2}}=a\sqrt{2}$.
	 }
\end{ex}

% nguồn: Ok_TracNghiemLop11_Lan1
\begin{ex}%[1H8H5-3]%[ma:851084db]
	Cho tứ diện $ABCD$ có các tam giác $ABC$,  $DBC$ vuông cân và nằm trong hai mặt phẳng vuông góc với nhau, $AB = AC = DB = DC = 2a$. Tính khoảng cách từ $B$ đến mặt phẳng $(ACD)$.
	\choice 
		 {\True $\dfrac{2a\sqrt{6}}{3}$}
		{$\dfrac{2a\sqrt{3}}{3}$} 
		{$a\sqrt{6}$} 
		 {$\dfrac{a\sqrt{6}}{2}$}
	 \loigiai{
	 \immini{
		Gọi $E$ là  trung điểm cạnh $BC$. Tam giác $ABC$ vuông cân tại $A$ nên $AE$ vừa là trung tuyến vừa là đường cao. \\
		Suy ra: $AE \perp BC \Rightarrow AE \perp (BCD)$. \\ 
		Gọi $K$ là trung điểm của cạnh  $DC$. Ta có\\ $\heva{EK \parallel BD \\ BD \perp DC} \Rightarrow DC \perp EK$ \hfill(1)\\
		Lại có $AE \perp  DC$ \hfill(2)\\
		Từ (1) và (2), suy ra $DC \perp (AEK) \Rightarrow (AEK) \perp (ACD)$ theo giao tuyến  $AK$.\\
		 }
		 {
		\begin{tikzpicture}[>=stealth,line join=round,line cap=round,font=\footnotesize]
				\path 
					 (0,0) coordinate (B)
					(3,-2) coordinate  (C)
					 (6,0) coordinate (D)
					(1.5,3)  coordinate (A);
				 \coordinate (E) at ($(B)!0.5!(C)$);
				\coordinate (K)  at ($(C)!0.5!(D)$);
				\coordinate (H) at ($(A)!(E)!(K)$); 
				\draw (A)--(B)--(C)--(D)--(A)--(C) 
					(A)--(E) 
					 (A)--(K);
				\draw[dashed] (B)--(D)--(E)--(K)  circle (1pt) node[below]{$K$}
					(E)--(H) circle  (1pt) node[above right]{$H$};
				\fill (A) circle  (1pt) node[above left]{$A$}
					(C) circle (1pt)  node[below right]{$C$}
					(B) circle  (1pt) node[below left]{$B$}
					(D) circle (1pt) node[below right]{$D$} 
					 (E) circle (1pt) node[below left]{$E$}
					 (H) circle (1pt)
					(K)  circle (1pt);
				\draw  pic[draw,angle radius = 5pt] {right angle =A--E--D};
				\draw pic[draw,angle radius = 5pt] {right  angle =A--H--E};
			\end{tikzpicture}
		 }
		 \noindent
		Trong mặt phẳng  $(AEK)$ kẻ $EH \perp AK$ với $H \in AK$. Suy ra $EH \perp (ACD)$.\\
		Do đó  $\mathrm{d}(B, (ACD)) = 2 \mathrm{d}(E, (ACD)) = 2EH$. \\
		Ta có $AE \cdot BC = AB \cdot AC \Rightarrow AE  = \dfrac{AB \cdot AC}{BC} = a \sqrt{2}$. $EK = \dfrac{1}{2}BD = a$. \\
		Tam giác $AEK$ vuông tại $E$  có $EH$ là đường cao. \\
		$EH^2 = \dfrac{AE^2 \cdot EK^2}{AE^2 + EK^2} = \dfrac{2a^2 \cdot a^2}{2a^2 + a^2} = \dfrac{2a^2}{3} \Leftrightarrow EH =  \dfrac{a \sqrt{6}}{3}$. Nên $\mathrm{d}(B, (ACD)) = \dfrac{2a \sqrt{6}}{3}$.
	 }
\end{ex}

% nguồn: Ok_TracNghiemLop11_Lan1
\begin{ex}%[1H8H5-3]%[ma:8eb3a138]
	Cho hình hộp chữ nhật  $ABCD.A'B'C'D'$ có $AB=1$, $BC=2$, $AA'=2$ (tham khảo hình bên).
	\begin{center}
		\begin{tikzpicture}[scale=0.6,font=\footnotesize,line join=round,line cap=round,>=stealth]
				\path 
					(0,0) coordinate (A) 
					++(-130:3)  coordinate (B)
					++(0:4)  coordinate (C)
					($(A)+(C)-(B)$) coordinate  (D)
					($(A)!1/2!(C)$)  coordinate (O);
				 \foreach \i in {A,B,C,D}{
					\coordinate (\i') at  ($(\i)+(0,3)$);
				 }
				\draw  (A')--(B')--(C')--(D')--cycle;
				\draw (B)--(B') 
					(C)--(C') 
					(D)--(D') 
					(B)--(C)--(D); 
				 \draw[dashed](B)--(A)--(A')
					 (A)--(D);
				\foreach \i/\g  in {A'/90,B'/90,C'/90,D'/90,A/-90,B/-90,C/-90,D/-90}
				\fill[black] (\i)  circle (1pt)+(\g:5mm) node[scale=1]{$\i$};
			\end{tikzpicture}
	\end{center}
	Khoảng cách giữa hai đường thẳng  $AD'$ và $DC'$ bằng
	 \choice
		{$\sqrt{2}$} 
		 {$\dfrac{\sqrt{6}}{2}$}
		{$\dfrac{2\sqrt{5}}{5}$.} 
		{\True  $\dfrac{\sqrt{6}}{3}$}
	 \loigiai{
		\begin{center}
			\begin{tikzpicture}[scale=0.6,font=\footnotesize,line join=round,line cap=round,>=stealth]
					 \path
						 (0,0) coordinate (A)
						++(-130:3) coordinate (B) 
						++(0:6) coordinate (C) 
						($(A)+(C)-(B)$)  coordinate (D)
						($(A)!1/2!(C)$)  coordinate (O)
						($(D')!1/2!(C)$) coordinate (M); 
					\foreach \i  in {A,B,C,D}{
						\coordinate (\i')  at ($(\i)+(0,4.5)$);
					 }
					 \coordinate (M) at ($(D')!1/2!(C)$);
					 \draw (A')--(B')--(C')--(D')--cycle;
					\draw (B)--(B') 
						 (C)--(C')
						 (D)--(D')
						(B)--(C)--(D) 
						 (B)--(C')--(D)
						 (C)--(D');
					\draw[dashed](B)--(A)--(A') 
						 (A)--(D)
						 (A)--(D')
						 (M)--(B)--(D);
					\foreach \i/\g  in {A'/90,B'/90,C'/90,D'/90,A/180,B/-90,C/-90,D/-90,M/0}
					\fill[black] (\i) circle (1pt)+(\g:5mm) node[scale=1]{$\i$}; 
				\end{tikzpicture}
		\end{center}
		Vì $AD'{  \parallel }BC'$ nên $AD'{ \parallel }(DBC')$.\\
		Do đó  $\mathrm{d}(AD',DC')=\mathrm{d}(AD',(DBC'))=\mathrm{d}(A,(DBC'))=\mathrm{d}(C,(DBC'))=h$.\\
		Xét tứ  diện $C.BC'D$ có các cạnh $CD,CB,CC'$ đôi một vuông góc nên ta có
		$$\dfrac{1}{h^2}=\dfrac{1}{CB^2}+\dfrac{1}{CD^2}+\dfrac{1}{C{C'}^2}=\dfrac{1}{2^2}+\dfrac{1}{1^2}+\dfrac{1}{2^2}=\dfrac{3}{2}\Rightarrow h=\dfrac{\sqrt{6}}{3}.$$ 
		Vậy  $\mathrm{d}(AD',DC')=\dfrac{\sqrt{6}}{3}$.
	} 
\end{ex}

% nguồn: Ok_TracNghiemLop11_Lan1
\begin{ex}%[1H8H5-3]%[ma:da3401bc]
	Cho hình chóp $S.ABC$ có đáy $ABC$ là tam giác cân tại $B$, cạnh bên  $SA$ vuông góc với đáy, $J$ là hình chiếu của $A$ lên $BC$. Kí hiệu $\mathrm{d}[A,(SBC)]$ là khoảng cách giữa điểm $A$ đến mặt phẳng $(SBC)$. Khẳng định nào sau đây đúng?
	 \choice
		{$\mathrm{d}[A,(SBC)]=AK$ với $K$ là hình chiếu của  $A$ lên $SB$}
		 {\True $\mathrm{d}[A,(SBC)]=AK$ với $K$ là hình chiếu của $A$ lên $SJ$}
		{$\mathrm{d}[A,(SBC)]=AK$ với  $K$ là hình chiếu của $A$ lên $SC$}
		{$\mathrm{d}[A,(SBC)]=AK$ với $K$là hình chiếu  của $A$ lên $SM$}
	\loigiai{ 
	 \immini{
		Ta có $\heva{&SA\perp  BC\\&AJ\perp BC}\Rightarrow BC\perp (SAJ)$.\\
		Gọi $K$ là hình chiếu của $A$ lên  $SJ$ ta có $$\heva{&AK\perp BC\,(BC\perp (SAJ))\\&AJ\perp BC.}$$
		Suy ra  $AK\perp (SBC)$.\\
		 Do đó $\mathrm{d}[A,(SBC)]=AK$.
		 }
		{ 
		\begin{tikzpicture}[line join=round,line cap=round]
				\coordinate (A)  at (-2,0);
				\coordinate  (B) at (0,-3);
				\coordinate (C) at  (3,0);
				\coordinate  (S) at ($(A)+(0,4)$);
				\coordinate (J) at ($(B)!0.4!(C)$); 
				\coordinate  (K) at ($(S)!0.5!(J)$);
				\draw (S)--(A) 
					(S)--(B) 
					 (S)--(C)
					 (B)--(C)
					(A)--(B) 
					 (S)--(J);
				 \draw[dashed](A)--(C)
					 (K)--(A)--(J);
				\pic[draw,angle  radius=2mm] {right angle = S--A--B}
					pic[draw,angle  radius=2mm] {right angle = S--A--C}
					pic[draw,angle radius=2mm] {right  angle = A--J--B}
					pic[draw,angle radius=2mm] {right  angle = A--K--J};
				\foreach \i/\g in {S/90,A/180,B/-90,C/0,J/0,K/0}{\draw[fill=white](\i) circle (1pt) 
					($(\i)+(\g:3mm)$) node[scale=1]{$\i$};} 
			\end{tikzpicture}
		} 
	} 
\end{ex}

% nguồn: Ok_TracNghiemLop11_Lan1
\begin{ex}%[1H8H5-3]%[ma:22c3ecc0]
	Cho hình chóp $S.ABC$ có đáy $ABC$ là tam giác đều cạnh $a$. Cạnh bên $SA=a\sqrt{3}$ và  vuông góc với mặt đáy $(ABC)$. Tính khoảng cách $d$ từ $A$ đến mặt phẳng $(SBC)$.
	\choice 
		 {\True $d=\dfrac{a\sqrt{15}}{5}$}
		{$d=a$} 
		 {$d=\dfrac{a\sqrt{5}}{5}$}
		 {$d=\dfrac{a\sqrt{3}}{2}$}
	\loigiai{ 
	 \immini{
		Gọi $M$ là trung  điểm $BC$, suy ra $AM\perp BC$ và $AM=\dfrac{a\sqrt{3}}{2}$.\\
		Gọi $K$ là hình chiếu của $A$ trên $SM$, suy ra  $AK\perp SM\quad(1)$.\\
		Ta có $\heva{&AM\perp BC\\&BC\perp SA}\Rightarrow BC\perp(SAM)\Rightarrow  BC\perp AK\quad(2)$.\\
		Từ $(1)$  và $(2)$, suy ra $AK\perp(SBC)$ nên $\mathrm{d}\left(A,(SBC)\right)=AK$.\\
		 Trong $\triangle SAM$, có $AK=\dfrac{SA\cdot AM}{\sqrt{SA^2+AM^2}}=\dfrac{3a}{\sqrt{15}}=\dfrac{a\sqrt{15}}{5}$.\\
		Vậy $\mathrm{d}\left(A,(SBC)\right)=AK=\dfrac{a\sqrt{15}}{5}$. 
		}{ 
		\begin{tikzpicture}[scale=0.8,font=\footnotesize,line join=round,line cap=round,>=stealth]
				\path 
					(0,0) coordinate  (A)
					(1.2,-1.5) coordinate  (B)
					(4,0) coordinate (C) 
					($(A)+(0,3.3)$) coordinate (S) 
					($(B)!0.5!(C)$)  coordinate (M)
					($(S)!0.55!(M)$)  coordinate (K);
				\draw (S)--(A)--(B)--(C)--cycle  (M)--(S)--(B);
				 \draw[dashed] (A)--(C)
					 (M)--(A)--(K);
				 \foreach \p/\q in {S/90,A/180,B/-90,C/0,M/-40,K/50}
				 \fill[black] (\p) circle (1.0pt)
					($(\p)+(\q:3mm)$) node{$\p$}; 
				\draw pic[draw,angle  radius=0.16cm] {right angle = B--M--A}
					pic[draw,angle  radius=0.16cm] {right angle = S--K--A};
			\end{tikzpicture}
	 }}
\end{ex}

% nguồn: Ok_TracNghiemLop11_Lan1
\begin{ex}%[1H8H5-3]%[ma:73c5341a]
	Cho  hình chóp $S.ABCD$ có đáy là hình bình hành, cạnh bên $SA$ vuông góc với đáy. Biết khoảng cách từ $A$ đến $(SBD)$ bằng $\dfrac{6a}{7}$. Tính khoảng cách từ $C$ đến mặt phẳng $(SBD)$?
	 \choice
		{$\dfrac{12a}{7}$} 
		 {$\dfrac{3a}{7}$}
		 {$\dfrac{4a}{7}$}
		 {\True $\dfrac{6a}{7}$}
	 \loigiai{
	\immini 
		{Vì $O$ là trung điểm $AC$  nên\\ $\mathrm{d}(C,(SBD))=\mathrm{d}(A,(SBD))=\dfrac{6a}{7}$.
		}{ 
		\begin{tikzpicture}[scale=1,font=\footnotesize,line join=round,line cap=round,>=stealth]
				 \path
					(0,0) coordinate  (A)
					++(-140:2) coordinate  (B)
					++(0:3.5) coordinate (C) 
					($(A)+(C)-(B)$) coordinate  (D)
					($(A)+(0,3)$)  coordinate (S)
					($(A)!.5!(C)$) coordinate (O); 
				\foreach \i  in{B,C,D}{\draw (S)--(\i);};
				\draw  (B)--(C)--(D);
				 \draw[dashed] (S)--(A)--(B)
					 (A)--(D)--(B)
					 (A)--(C)
					 (S)--(O);
				\pic[draw,angle eccentricity=1.8,angle radius=2mm]{right  angle=S--A--D};
				\foreach \i/\g in  {A/-90,B/-90,C/-90,D/-90,S/90,O/-90}
				\fill[black] (\i)  circle (1pt)+(\g:3mm) node[scale=1]{$\i$};
			\end{tikzpicture}
		} 
	 }
\end{ex}

% nguồn: Ok_TracNghiemLop11_Lan1
\begin{ex}%[1H8H5-3]%[ma:8fcbb8b4]
	Cho hình hộp chữ nhật $ABCD.A' B' C' D' $ có $AB=a$, $AD=\sqrt{3}a$, $AA' =2\sqrt{3}a$ (tham khảo hình vẽ). Khoảng cách giữa hai mặt  phẳng $(ABCD)$ và $(A'B'C'D')$ bằng
	\choice 
		{$a$} 
		{\True $2a\sqrt{3}$} 
		{$3a$} 
		 { $a\sqrt{3}$}
	 \loigiai{
	\immini{ 
		 Vì $(ABCD)\parallel (A'B'C'D')$ nên $$\mathrm{d}[(ABCD),(A'B'C'D')]=\mathrm{d}[A,(A'B'C'D')]=AA'=2a\sqrt{3}.$$
		} 
		 {
		\begin{tikzpicture}[font=\footnotesize,line join=round,line cap=round,>=stealth]
				\coordinate (A)  at (0,0);
				\coordinate  (A') at (0,2);
				\coordinate (B) at (2,0); 
				 \coordinate (B') at (2,2);
				 \coordinate (C) at (3,0.5);
				\coordinate (C') at (3,2.5); 
				\coordinate  (D) at (1,0.5);
				\coordinate (D') at (1,2.5); 
				\draw (A) -- (B) -- (C) -- (C') -- (D') -- (A') -- (B') --  (C')
					(A) -- (A') 
					 (B) -- (B');
				\draw[dashed] (A)--(D)  -- (C)
					(D) --(D'); 
				\foreach \i/\g in  {A/180,B/-90,C/0,D/160,A'/180,B'/0,C'/90,D'/90}{\draw[fill=white](\i) circle (1pt)
					 ($(\i)+(\g:3mm)$) node[scale=1]{$\i$};}
			\end{tikzpicture}
		 }
	} 
\end{ex}

% nguồn: Ok_TracNghiemLop11_Lan1
\begin{ex}%[1H8H5-3]%[ma:8a9d7fc2]
	Cho hình lập phương $ABCD.A'B'C'D'$ có cạnh bằng  $a$ (tham khảo hình vẽ).
	\begin{center}
		\begin{tikzpicture}[scale=0.6,font=\footnotesize,line join=round,line cap=round,>=stealth]
				 \path
					(0,0)  coordinate (A)
					++(-130:3) coordinate  (B)
					++(0:4)  coordinate (C)
					($(A)+(C)-(B)$) coordinate  (D)
					($(A)!1/2!(C)$)  coordinate (O);
				\foreach \i  in {A,B,C,D}{
					\coordinate (\i') at ($(\i)+(0,4)$); 
				 }
				\draw (A')--(B')--(C')--(D')--cycle; 
				\draw (B)--(B') 
					 (C)--(C')
					 (D)--(D')
					(B)--(C)--(D); 
				 \draw[dashed](B)--(A)--(A')
					(A)--(D); 
				\foreach \i/\g in {A'/90,B'/90,C'/90,D'/90,A/-90,B/-90,C/-90,D/-90} 
				\fill[black]  (\i) circle (1pt)+(\g:5mm) node[scale=1]{$\i$};
			\end{tikzpicture}
	\end{center}
	Khoảng cách giữa đường thẳng $BD$  và mặt phẳng $(A'B'C'D')$ bằng
	 \choice
		{\True $a$} 
		{$a\sqrt{3}$} 
		{$3a$} 
		{$a\sqrt{2}$} 
	\loigiai{ 
		Vì $BD{ \parallel }B'D'$ nên $BD{ \parallel }(A'B'C'D')$.\\ 
		Do đó  $\mathrm{d}(BD,(A'B'C'D'))=\mathrm{d}(B,(A'B'C'D'))=BB'=a$.\\
	Vì $(ABCD){ \parallel }(A'B'C'D')$ nên $\mathrm{d}((ABCD),(A'B'C'D'))=\mathrm{d}(A,(A'B'C'D'))=AA'=2a\sqrt{3}$} 
\end{ex}

% nguồn: Ok_TracNghiemLop11_Lan1
\begin{ex}%[1H8H5-3]%[ma:c8ef03b1]
	Cho lập phương $ABCD.A'B'C'D'$ có cạnh bằng $a$  (tham khảo hình bên).
	\begin{center}
		\begin{tikzpicture}[scale=0.6,font=\footnotesize,line join=round,line cap=round,>=stealth]
				\path 
					 (0,0) coordinate (A)
					++(-130:3) coordinate  (B)
					++(0:4)  coordinate (C)
					 ($(A)+(C)-(B)$) coordinate (D)
					 ($(A)!1/2!(C)$) coordinate (O);
				 \foreach \i in {A,B,C,D}{
					 \coordinate (\i') at ($(\i)+(0,4)$);
				} 
				\draw  (A')--(B')--(C')--(D')--cycle;
				 \draw (B)--(B')
					 (C)--(C')
					(D)--(D') 
					(B)--(C)--(D) 
					(A')--(C'); 
				\draw[dashed](B)--(A)--(A') 
					 (D)--(B)
					(A)--(D); 
				\foreach \i/\g in  {A'/90,B'/90,C'/90,D'/90,A/-90,B/-90,C/-90,D/-90}
				\fill[black] (\i)  circle (1pt)+(\g:5mm) node[scale=1]{$\i$};
			\end{tikzpicture}
	\end{center}
	 Khoảng cách giữa hai đường thẳng $BD$ và $A'C'$ bằng
	 \choice
		 {\True $a$}
		{$a\sqrt{3}$} 
		{$\dfrac{a\sqrt{3}}{2}$} 
		{$a\sqrt{2}$} 
	 \loigiai{
		Vì $BD{ \parallel  }B'D'$ nên $BD\parallel (A'B'C'D')$.\\
		 Do đó $\mathrm{d}(BD,A'C')=\mathrm{d}(BD,(A'B'C'D'))=\mathrm{d}(B,(A'B'C'D'))=BB'=a$.
	} 
\end{ex}

% nguồn: Ok_TracNghiemLop11_Lan1
\begin{ex}%[1H8H5-3]%[ma:e9c917e3]
	Cho hình chóp $S.ABCD$ có đáy $ABCD$ là hình vuông tâm $O$, $SA\perp (ABCD)$. Gọi $I$ là trung điểm của $SC$. Khoảng cách  từ $I$ đến mặt phẳng $(ABCD)$ bằng độ dài đoạn thẳng nào?
	 \choice
		{$IB$} 
		{$IA$} 
		 {$IC$}
		 {\True $IO$}
	 \loigiai{
	 \immini
		 {
		Vì $OI\parallel SA$ nên khoảng cách từ $I$ đến mặt  phẳng $(ABCD)$ bằng $IO$.
		 }{
		\begin{tikzpicture}[scale=0.7,font=\footnotesize,line join=round,line cap=round,>=stealth]
				 \path
					(0,0) coordinate (A) 
					 ++(-140:2) coordinate (B)
					 ++(0:3.5) coordinate (C)
					 ($(A)+(C)-(B)$) coordinate (D)
					($(A)+(0,3)$)  coordinate (S)
					($(A)!.5!(C)$) coordinate  (O)
					($(S)!.5!(C)$) coordinate (I); 
				\foreach \i in{B,C,D}{\draw  (S)--(\i);};
				\draw  (B)--(C)--(D);
				 \draw[dashed] (S)--(A)--(B)
					 (A)--(D)
					 (A)--(C)
					 (I)--(O)
					 (B)--(D);
				\pic[draw,angle eccentricity=1.8,angle radius=2mm]{right angle=S--A--D}; 
				\pic[draw,angle eccentricity=1.8,angle radius=2mm]{right angle=I--O--C}; 
				\foreach \i/\g in  {A/180,B/-90,C/-90,D/0,S/90,O/-90,I/0}
				\fill[black] (\i) circle (1pt)+(\g:3mm)  node[scale=1]{$\i$};
			\end{tikzpicture}
		} 
	} 
\end{ex}

% nguồn: Ok_Toan11_TracNghiem_Lan2
\begin{ex}%[1H8V5-3]%[ma:ca9b36ec]
	Cho hình chóp $S.ABC$ có đáy $ABC$ là tam giác vuông tại $A$, $AB=2a$, $AC=a$, $\widehat{SBA}=\widehat{SCA}=90^\circ$, góc giữa $SA$ và mặt phẳng $(ABC)$ bằng $45^\circ$.  Tính khoảng cách từ điểm $C$ đến mặt phẳng $(SAB)$ bằng
	\choice 
		{\True  $\dfrac{a\sqrt{3}}{2}$}
		 {$\dfrac{a\sqrt{30}}{6}$}
		{$\dfrac{a\sqrt{30}}{2}$} 
		{$\dfrac{a\sqrt{3}}{6}$} 
	\loigiai{ 
		\begin{center}
			\begin{tikzpicture}[scale=0.75,font=\footnotesize,line join=round,line cap=round,>=stealth]
					\coordinate  (H) at (0,0);
					\coordinate (B)  at (-2.2,-0.8);
					\coordinate (C)  at (3.8,0);
					 \coordinate (A) at (1.6,-0.8);
					\coordinate (S) at  (0,3.5);
					\coordinate  (K) at ($(S)!0.6!(B)$);
					\draw[dashed]  (S) -- (H) -- (B)
						(K) -- (H)  -- (C)
						(A)--(H); 
					\draw (S)--(B)--(A)--(C)--(S)--(A); 
					\pic[draw,thin,angle radius=6pt] {right angle =  S--H--C} pic[draw,thin,angle radius=6pt] {right angle = S--K--H};
					 \foreach \i/\g in {S/90,H/-90,B/-90,A/-90,C/0,K/120}{\draw[fill](\i) circle (1pt)
						($(\i)+(\g:3mm)$) node[scale=1]{$\i$};} 
				\end{tikzpicture}
		\end{center}
		Gọi $H$ là hình chiếu của  $S$ lên $(ABC)$.\\
		Ta có $\heva{ &AB \perp SH \\ &AB \perp SB} \Rightarrow AB \perp (SHB) \Rightarrow  AB \perp HB$.\\
		Mà $AB \perp AC$ nên suy ra $HB //  AC$ \hfill $(1)$.\\
		Mặt khác $\heva{ &AC \perp SH \\ &AC  \perp SC} \Rightarrow AC \perp (SHC) \Rightarrow AC \perp HC$.\\
		Mà $AC \perp AB$ nên suy  ra $HC // AB$ \hfill $(2)$.\\
		 Từ $(1)$, $(2)$ suy ra $ABHC$ là hình bình hành mà $\widehat A=90^\circ$ nên $ABHC$ là hình chữ nhật.\\
		Lại có $(SA,(ABC))=\widehat{SAH} = 45^\circ \Rightarrow SH =  AH = a\sqrt{5}$ và $\mathrm{d}(C,(SAB)) = \mathrm{d}(H,(SAB))$.\\
		Gọi $K$ là hình chiếu của  $H$ lên $SB$.\\
		Mà  $AB \perp (SHB) \Rightarrow AB \perp HK$.\\
		Suy ra  $HK \perp (SAB)$, do đó
		 $\mathrm{d}(C,(SAB)) = HK$.\\
		$\triangle SHB$ vuông tại $H$ có $$\dfrac{1}{HK^2} = \dfrac{1}{SH^2}  + \dfrac{1}{HB^2} = \dfrac{1}{5a^2}+\dfrac{1}{a^2} = \dfrac{6}{5a^2} \Rightarrow HK=\dfrac{a\sqrt{30}}{6}.$$
		Vậy  $\mathrm{d}(C,(SAB)) = \dfrac{a\sqrt{30}}{6}$.
	} 
\end{ex}

% nguồn: Ok_Toan11_TracNghiem_Lan2
\begin{ex}%[1H8V5-3]%[ma:76626050]
	[MĐ 103 BGD\&ĐT NĂM  2018-2019]
	Cho hình chóp $S. ABCD$ có đáy là hình vuông cạnh $a$, mặt bên $SAB$ là tam giác đều và nằm trong mặt phẳng vuông góc với  mặt phẳng đáy. Khoảng cách từ $D$ đến mặt phẳng $(SAC)$ bằng
	\choice 
		 {$\dfrac{a \sqrt{21}}{14}$}
		{$\dfrac{a  \sqrt{21}}{28}$}
		{$\dfrac{a \sqrt{2}}{2}$} 
		{\True  $\dfrac{a \sqrt{21}}{7}$}
	\loigiai{ 
	\immini{Gọi $O=AC\cap BD$ và $G$ là  trọng tâm tam giác $ABD$, $I$ là trung điểm của $AB$ ta có $SI\perp(ABCD)$.\\
		Ta có $\dfrac{\mathrm{d}(D,(SAC))}{\mathrm{d}(I,(SAC))}=\dfrac{DG}{IG}=2$.\\ 
		$\Rightarrow  \mathrm{d}(D,(SAC))=2\mathrm{d}(I,(SAC))$.\\
		 Gọi $K$ là trung điểm của $AO, H$ là hình chiếu của $I$ lên $SK$ ta có $IK\perp AC$, $IH\perp(SAC)$.\\
		Suy ra $\mathrm{d}(D,(SAC))=2\mathrm{d}(I,(SAC))=2IH$.} 
		{ 
		\begin{tikzpicture}[scale=1,font=\footnotesize,line join=round,line cap=round,>=stealth]
				\def\a{4} 
				 \def\h{4}
				\path  (0:0) coordinate (A)
					++(0:\a) coordinate (D) 
					++(-145:\a/1.5)  coordinate (C)
					($(A)+(C)-(D)$)  coordinate (B)
					($(A)!0.5!(B)$) coordinate  (I)
					($(I)+(90:\h)$) coordinate  (S)
					(intersection of A--C and B--D) coordinate  (O)
					 ($(A)!0.5!(O)$) coordinate (K)
					($(A)!2/3!(O)$) coordinate (G) 
					($(S)!0.7!(K)$)  coordinate (H);
				\draw[dashed] (B)--(A)--(D) 
					 (C)--(A)--(S)
					 (H)--(I)--(K)--(S)
					(S)--(I) 
					 (B)--(D);
				 \draw[] (B)--(C)--(D)
					 (B)--(S)
					(C)--(S) 
					 (D)--(S);
				\draw  pic[draw,angle radius=2mm]{right angle=D--I--S};
				\draw pic[draw,angle  radius=2mm]{right angle=S--H--I};
				\foreach \x/\g in {A/45,B/-135,C/-45,D/45,S/90,H/45,O/-90,G/0,K/40,I/180} 
				\fill[] (\x) circle  (1pt)
					($(\g:2mm)+(\x)$) node  {$\x$};
			\end{tikzpicture}
		 }
		Xét tam  giác $SIK$ vuông tại $I$ có $SI=\dfrac{a \sqrt{3}}{2}$, $IK=\dfrac{BO}{2}=\dfrac{a \sqrt{2}}{4}$.
		\[ 
		\dfrac{1}{IH^2}=\dfrac{1}{SI^2}+\dfrac{1}{IK^2}=\dfrac{4}{3a^2}+\dfrac{16}{2a^2}=\dfrac{28}{3a^2}  \Rightarrow IH=\dfrac{a \sqrt{3}}{2\sqrt{7}}.
		\] 
		Vậy  $\mathrm{d}(D,(SAC)=2\mathrm{d}(I,(SAC))=2IH=\dfrac{a \sqrt{21}}{7}$.
	 }
\end{ex}

% nguồn: Ok_Toan11_TracNghiem_Lan2
\begin{ex}%[1H8V5-3]%[ma:c6f7b6ba]
	Cho hình chóp $S.ABCD$ có đáy là hình thoi cạnh $a$. Tam giác $ABC$ là tam giác đều, hình chiếu vuông góc của đỉnh $S$ lên mặt phẳng $(ABCD)$ trùng với trọng tâm tam giác $ABC$. Góc giữa đường thẳng $SD$ và mặt phẳng $(ABCD)$  bằng $30^\circ$. Tính khoảng cách từ điểm $B$ đến mặt phẳng $(SCD)$ theo $a$.
	 \choice
		 {$a$}
		{$\dfrac{2a\sqrt{21}}{7}$} 
		{$a\sqrt{3}$} 
		 {\True $\dfrac{a\sqrt{21}}{7}$}
	\loigiai{ 
		\begin{center}
			\begin{tikzpicture}[scale=1,font=\footnotesize,line join=round,line cap=round,>=stealth]
					\coordinate (B)  at (0,0);
					\coordinate  (C) at (4,0);
					 \coordinate (A) at (-1,-1.5);
					\coordinate (D)  at ($(A)+(C)-(B)$);
					\coordinate  (O) at ($(C)!0.5!(A)$);
					\coordinate (E) at  ($(B)!0.6666!(O)$);
					\coordinate  (S) at ($(E)+(0,4)$);
					\coordinate (H)  at ($(S)!0.7!(C)$);
					 \draw (S)--(A)--(D)--(C)--(S)--(D);
					\draw[dashed,thin]  (S)--(B)--(A)--(C)
						(D)--(B)--(C) 
						 (S)--(E)--(H);
					 \pic[draw,thin,angle radius=6pt] {right angle = S--E--B} pic[draw,thin,angle radius=6pt] {right angle = S--H--E};
					 \foreach \i/\g in {S/90,B/180,A/-90,C/0,H/30,D/-90,O/-90,E/-120}{\draw[fill](\i) circle (1pt)
						($(\i)+(\g:3mm)$) node[scale=1]{$\i$};} 
				\end{tikzpicture}
		\end{center}
		 Gọi $O$ là tâm hình thoi $ABCD$ và $E$ là trọng tâm của tam giác $ABC$.\\
		 Theo giả thiết ta có $SE \perp (ABC)$.\\
		Do  đó góc giữa $SD$ và mặt phẳng $(ABCD)$ là $\widehat{SDE} = 30^\circ$ tại $E$.\\
		Do tam giác $ABC$ đều nên $BD = 2BO = a\sqrt{3}  \Rightarrow DE = \dfrac{2}{3}BD = \dfrac{2a\sqrt{3}}{3}$.\\
		 Khi đó $\tan 30^\circ = \dfrac{SE}{DE} \Rightarrow SE = DE \cdot \tan 30^\circ = \dfrac{2a\sqrt{3}}{3} \cdot \dfrac{1}{\sqrt{3}} = \dfrac{2a}{3}$.\\
		Vì tam giác $ABC$ đều nên $CE \perp AB$ suy  ra $CE \perp CD$.\\
		Mà  $CD \perp SE$ nên $CD \perp (SCE)$.\\
		Kẻ $EH \perp SC$, $\left(H \in SC\right)$ khi đó $EH \perp (SCD)$ hay $\mathrm{d}(E;(SCD))=EH$.\\ 
		$\dfrac{1}{EH^2} = \dfrac{1}{SE^2} + \dfrac{1}{EC^2} = \dfrac{9}{4a^2}+\dfrac{3}{a^2}=\dfrac{21}{4a^2} \Rightarrow EH=\dfrac{2a\sqrt{21}}{21}$.\\ 
		Lại có $\mathrm{d}(B,(SCD)) = \dfrac{3}{2}\mathrm{d}(E,(SCD)) = \dfrac{3}{2}EH =  \dfrac{a\sqrt{21}}{7}$.
	 }
\end{ex}

% nguồn: Ok_Toan11_TracNghiem_Lan2
\begin{ex}%[1H8V5-3]%[ma:36f91fc5]
	[Mã 103 BGD\&ĐT  NĂM 2019-2020]
	Cho hình lăng trụ đứng $ABC.A'B'C'$ có đáy $ABC$ là tam giác đều cạnh $a$ và $A'A=2a$. Gọi $M$ là trung điểm của $A'A$. Khoảng cách từ $M$ đến  mặt phẳng $\left(AB'C\right)$ bằng
	 \choice
		 {\True $\dfrac{\sqrt{57} a}{19}$}
		{$\dfrac{\sqrt{5}  a}{5}$}
		{$\dfrac{2\sqrt{5} a}{5}$} 
		 {$\dfrac{2\sqrt{57} a}{19}$}
	\loigiai{ 
	\immini{Gọi $I=BM\cap AB'$ và $K$  là trung điểm $AC$.\\
		 Ta có $\dfrac{\mathrm{d}\left(M,\left(AB'C\right)\right)}{\mathrm{d}\left(B,\left(AB'C\right)\right)}=\dfrac{MI}{BI}=\dfrac{MA}{BB'}=\dfrac{1}{2}$\\
		$\Rightarrow  \mathrm{d}\left(M,\left(AB'C\right)\right)=\dfrac{1}{2} \mathrm{d}\left(B,\left(AB'C\right)\right)=\dfrac{BH}{2}$.\\
		 Xét tam giác $BB'K$ có $\dfrac{1}{BH^2}=\dfrac{1}{B'B^2}+\dfrac{1}{BK^2}=\dfrac{1}{(2a)^2}+\dfrac{1}{\left(\dfrac{a \sqrt{3}}{2}\right)^2}$.\\
		Suy  ra $BH=\dfrac{2\sqrt{57} a}{19}$.\\
		 Vậy $\mathrm{d}\left(M,\left(AB'C\right)\right)=\dfrac{BH}{2}=\dfrac{\sqrt{57} a}{19}$.}
		{ 
		\begin{tikzpicture}[scale=0.9,font=\footnotesize,line join=round,line cap=round,>=stealth]
				 \def\a{4}
				\def\h{4.5} 
				\path (0:0) coordinate (B) 
					 ++(0:\a) coordinate (C)
					(-30:3*\a/4) coordinate (A) 
					($(A)+(90:\h)$) coordinate  (A')
					 ($(B)+(90:\h)$) coordinate (B')
					($(C)+(90:\h)$)  coordinate (C')
					($(A)!1/2!(C)$) coordinate (K) 
					 ($(A')!1/2!(A)$) coordinate (M)
					($(B')!(B)!(K)$) coordinate (H) 
					(intersection of M--B and A--B') coordinate (I); 
				\draw[dashed] (H)--(B)--(C)--(B')--(K)--(B); 
				\draw[] (B)--(M) 
					 (C)--(C')
					 (A)--(A')
					(B)--(B')--(A) 
					 (B)--(A)--(C)
					 (A')--(B')--(C')--cycle;
				\foreach \x/\g in  {B/180,A/-45,C/0,B'/180,A'/-45,C'/0,M/0,K/-45,H/45,I/90}
				\fill (\x) circle  (1pt)
					($(\g:4mm)+(\x)$) node  {$\x$};
			\end{tikzpicture}
		} 
	 }
\end{ex}

% nguồn: Ok_Toan11_TracNghiem_Lan2
\begin{ex}%[1H8V5-3]%[ma:53f48d0e]
	Cho hình chóp $S.ABC$ có đáy là tam giác vuông tại $A$ có $AB=a$, $AC=2a$, mặt phẳng $(SBC) \perp (ABC)$. Khoảng cách từ $A$  đến $(SBC)$ bằng
	\choice 
		{\True  $\dfrac{2\sqrt{5}a}{5}$}
		 {$3a$}
		 {$a\sqrt{5}$}
		{$a\sqrt{2}$} 
	 \loigiai{
		\begin{center}
			\begin{tikzpicture}[scale=1,font=\footnotesize,line join=round,line cap=round,>=stealth]
					\coordinate (A) at  (0,0);
					\coordinate (C)  at (4,0);
					 \coordinate (B) at (1.5,-2);
					 \coordinate (S) at ($(A)+(2.7,3)$);
					\coordinate (H) at  ($(B)!0.6!(C)$);
					 \draw (S)--(B)--(A)--(S)
						 (S)--(C)
						(B)--(C); 
					 \draw[dashed,thin] (H)--(A)--(C);
					\pic[draw,thin,angle radius=6pt] {right angle =  B--A--C} pic[draw,thin,angle radius=6pt] {right angle = A--H--B};
					\foreach \i/\g in {S/90,A/180,C/0,B/-90,H/0}{\draw[fill](\i) circle (1pt) 
						($(\i)+(\g:3mm)$) node[scale=1]{$\i$};} 
				\end{tikzpicture}
		\end{center}
		 Trong mặt phẳng $(ABC)$ dựng $AH \perp BC$.\\
		Do $(SBC) \perp (ABC) \Rightarrow AH  \perp (SBC)$.\\
		Khoảng  cách từ $A$ đến $(SBC)$ bằng $AH$.\\
		 Trong tam giác vuông $ABC$ ta có $$\dfrac{1}{AH^2} = \dfrac{1}{AB^2}+\dfrac{1}{AC^2} = \dfrac{1}{a^2}+\dfrac{1}{(2a)^2} = \dfrac{1}{a^2}+\dfrac{1}{4a^2} = \dfrac{5}{4a^2} \Rightarrow AH = \sqrt{\dfrac{4a^2}{5}} = \dfrac{2a}{\sqrt{5}} = \dfrac{2a\sqrt{5}}{5}.$$
	} 
\end{ex}

% nguồn: Ok_Toan11_TracNghiem_Lan2
\begin{ex}%[1H8V5-3]%[ma:8eb976f7]
	Cho hình chóp $S . A B C D$ có đáy là hình vuông cạnh $a$, mặt bên $S A B$ là tam giác đều và  nằm trong mặt phẳng vuông góc với mặt phẳng đáy. Khoảng cách từ $A$ đến mặt phẳng $(S B D)$ bằng
	\choice 
		 {$\dfrac{\sqrt{21} a}{14}$}
		{\True  $\dfrac{\sqrt{21} a}{7}$}
		{$\dfrac{\sqrt{2} a}{2}$} 
		{$\dfrac{\sqrt{21} a}{28}$} 
	\loigiai{ 
		\begin{center}
			\begin{tikzpicture}[scale=1,font=\footnotesize,line join=round,line cap=round,>=stealth]
					\coordinate (A) at  (0,0);
					\coordinate (B) at  (-3,-3);
					\coordinate (D) at  (7,0);
					\coordinate (C)  at ($(B)+(D)-(A)$);
					\coordinate (H)  at ($(A)!0.5!(B)$);
					\coordinate (S) at  ($(H)+(0,6)$);
					\coordinate  (O) at ($(A)!0.5!(C)$);
					\coordinate (K) at  ($(O)!0.5!(B)$);
					\coordinate  (I) at ($(K)!0.15!(S)$);
					\draw  (S)--(B)
						 (S)--(C)
						 (S)--(D)
						 (B)--(C)--(D);
					\draw[dashed](S)--(A)--(C) 
						 (A)--(B)
						 (A)--(D)--(B)
						(S)--(H)--(I) 
						 (S)--(K)--(H);
					\foreach \i/\g  in {S/90,A/45,B/-90,C/-90,D/0,H/180,K/-90,I/0,O/90}{\draw[fill=black](\i) circle (1pt)
						($(\i)+(\g:3mm)$)  node[scale=1]{$\i$};}
				\end{tikzpicture}
		\end{center}
		Gọi $H$ là trung điểm của $A  B$. Khi đó $S H \perp(A B C D)$.\\
		Gọi $O$ là giao điểm  của $A C$ và $B D$.\\
		Suy ra $A C \perp B D$. \\  Kẻ $H K \perp B D$ tại $K$ ($K$ là trung điểm $B O$).\\
		 Kẻ $H I \perp S H$ tại I.\\ Khi đó $\mathrm{d}(A,(S B D))=2 \mathrm{d}(H,(S B D))=2 H I$.\\
		Xét tam giác $S H K$ ta có $S H=\dfrac{a \sqrt{3}}{2}$, $H  K=\dfrac{1}{2} A O=\dfrac{a \sqrt{2}}{4}$.\\
		Khi đó $\dfrac{1}{H I^{2}}=\dfrac{1}{S H^{2}}+\dfrac{1}{H K^{2}}=\dfrac{28}{3 a^{2}} \Rightarrow  H I=\dfrac{a \sqrt{21}}{14}$.\\
		Suy ra $\mathrm{d}(A,(S B D))=2 H I=\dfrac{a  \sqrt{21}}{7}$.
	 }
\end{ex}

% nguồn: Ok_Toan11_TracNghiem_Lan2
\begin{ex}%[1H8V5-3]%[ma:69c29857]
	Cho hình chóp tam giác đều $S.ABC$ có cạnh đáy bằng $a$ và $SA=\dfrac{2a\sqrt{3}}{3}$.  Tính khoảng cách từ điểm $A$ đến mặt phẳng $(SBC)$ bằng
	 \choice
		 {$\dfrac{\sqrt{13}}{13}a$}
		{$\dfrac{2\sqrt{13}}{13}a$} 
		{\True  $\dfrac{9\sqrt{13}}{13}a$}
		 {$\dfrac{3\sqrt{13}}{13}a$}
	\loigiai{ 
		\begin{center}
			\begin{tikzpicture}[scale=1,font=\footnotesize,line join=round,line cap=round,>=stealth]
					\coordinate (A)  at (0,0);
					\coordinate (C)  at (4,0);
					\coordinate (B) at (1,-1.5); 
					\coordinate (M) at  ($(C)!0.4!(B)$);
					\coordinate (O) at ($(A)!0.6666!(M)$); 
					 \coordinate (S) at ($(O)+(0,4)$);
					\coordinate (H) at ($(S)!0.7!(M)$); 
					\draw (B)--(S)--(C)--(B)--(A)--(S)--(M); 
					\draw[dashed,thin] (M)--(A)--(C) 
						(S)--(O)--(H); 
					\pic[draw,thin,angle radius=6pt] {right angle = S--O--A} pic[draw,thin,angle radius=6pt] {right  angle = S--M--C} pic[draw,thin,angle radius=6pt] {right angle = A--M--B} pic[draw,thin,angle radius=6pt] {right angle = S--H--O};
					\foreach  \i/\g in {S/90,A/180,B/-90,C/0,H/30,M/-90,O/-90}{\draw[fill](\i) circle (1pt)
						($(\i)+(\g:3mm)$) node[scale=1]{$\i$};} 
				\end{tikzpicture}
		\end{center}
		Gọi $M$ là  trung điểm $BC$, $O$ là trọng tâm tam giác $ABC \Rightarrow SO \perp (ABC)$.\\
		Ta có $\begin{cases} OM \perp BC \\ SO \perp BC
	\end{cases} \Rightarrow BC \perp (SOM)$.\\
	Trong $(SOM)$ kẻ $OH \perp SM$, $\left(H \in SM\right)$ mà $OH \perp  BC$ $\Rightarrow OH \perp (SBC) \Rightarrow \mathrm{d}(O,(SBC))=OH$.\\
	 Ta có $AO = \dfrac{2}{3}AM = \dfrac{2}{3}\cdot\dfrac{a\sqrt{3}}{2} = \dfrac{a\sqrt{3}}{3}$, $OM=\dfrac{1}{3}\cdot AM = \dfrac{1}{3}\cdot \dfrac{a\sqrt{3}}{2} = \dfrac{a\sqrt{3}}{6}$.\\
	Xét tam giác $SAO$ vuông tại  $O$ có $SO = \sqrt{SA^2-AO^2} =a.$
	Xét tam giác $SMO$ vuông tại $O$ có $$\dfrac{1}{OH^2} = \dfrac{1}{SO^2}  + \dfrac{1}{OM^2} = \dfrac{13}{a^2} \Rightarrow OH = \dfrac{a}{\sqrt{13}}.$$
	Lại có $\dfrac{\mathrm{d}(A,(SBC))}{\mathrm{d}(O,(SBC))} = \dfrac{AM}{OM} = 3 \Rightarrow \mathrm{d}(A,(SBC))=3\mathrm{d}(O,(SBC))=3 \cdot \dfrac{a}{2} = 3\dfrac{a}{\sqrt{13}}$. 
 }
\end{ex}

% nguồn: Ok_Toan11_TracNghiem_Lan2
\begin{ex}%[1H8V5-3]%[ma:2a33f520]
	Cho hình chóp $S . A B C$ có đáy $A B C$ là tam giác vuông tại $A$, $A B=a$, $A C=a \sqrt{3}$. Tam giác $S B C$ đều và nằm  trong mặt phẳng vuông với đáy. Tính khoảng cách $d$ từ $B$ đến mặt phẳng $(SAC)$.
	 \choice
		 {$\mathrm{d}=\dfrac{a \sqrt{39}}{13}$}
		 {$\mathrm{d}=a$}
		{$\mathrm{d}=\dfrac{a  \sqrt{3}}{2}$}
		 {\True $\mathrm{d}=\dfrac{2 a \sqrt{39}}{13}$}
	 \loigiai{
		\begin{center}
			\begin{tikzpicture}[scale=1,font=\footnotesize,line join=round,line cap=round,>=stealth]
					\coordinate (S) at  (0,4);
					\coordinate (B) at  (-2,0);
					\coordinate  (C) at (1.5,-3);
					\coordinate (A)  at (3,0);
					\coordinate (H)  at ($(B)!0.5!(C)$);
					\coordinate  (K) at ($(A)!0.7!(C)$);
					 \coordinate (E) at ($(S)!0.63!(K)$);
					\draw (S)--(A) 
						 (S)--(B)
						 (S)--(C)
						 (B)--(C)--(A)
						 (S)--(H)
						(S)--(K); 
					 \draw[dashed](A)--(B)
						(H)--(K) 
						 (H)--(E)
						(A)--(H); 
					 \pic[draw,thin,angle radius=2mm] {right angle = H--E--K};
					\pic[draw,thin,angle radius=2mm] {right  angle = H--K--C};
					\foreach  \i/\g in {S/90,A/0,B/180,C/-90,E/45,K/-45,H/-135}{\draw[fill=black](\i) circle (1pt)
						 ($(\i)+(\g:3mm)$) node[scale=1]{$\i$};}
				\end{tikzpicture}
		\end{center}
		Gọi $H$ là trung điểm của $B C$, suy ra $S H \perp  B C$.\\
		 Mà $(S A B) \perp(A B C)$ theo giao tuyến $B C$.\\
		Do đó $S H \perp(A B C) \Rightarrow S H  \perp A C$.\\
		Gọi $K$  là trung điểm $A C$, suy ra $H K \perp A C$.\\
		Do đó $(S H K) \perp A C \Rightarrow(S H K) \perp(S A C)$ theo giao  tuyến $S K$\\
		Trong $(S H K)$ kẻ $H E \perp S  K$ $(E \in S K)$.\\
		Suy ra $HE \perp(S A C) \Rightarrow \mathrm{d}(H;(S  A C))=H E$.\\
		Ta có $H K=\dfrac{a}{2}$, $BC=\sqrt{A B^{2}+A C^{2}}=2 a \Rightarrow S  H=\dfrac{2a\sqrt{3}}{2}=a \sqrt{3}$.\\
		Xét tam  giác $S H K$ vuông tại $H$ có $S H$ là đường cao nên \[\dfrac{1}{H E^{2}}=\dfrac{1}{S H^{2}}+\dfrac{1}{H K^{2}}=\dfrac{1}{3 a^{2}}+\dfrac{4}{a^{2}}=\dfrac{13}{3 a^{2}}.\]
		Ta được $H  E=\dfrac{a \sqrt{39}}{13}$.\\
		 Khi đó $\mathrm{d}(B,(S A C))=2\mathrm{d}(H,(S A C))=2HE=2\cdot \dfrac{S H \cdot H K}{\sqrt{S H^{2}+H K^{2}}}=\dfrac{2 a \sqrt{39}}{13}$.
	} 
\end{ex}

% nguồn: Ok_Toan11_TracNghiem_Lan2
\begin{ex}%[1H8V5-3]%[ma:316b362c]
	Cho khối hộp chữ nhật $ABCD.A'B'C'D'$ có đáy là hình vuông, $BD=2a$, góc giữa hai mặt phẳng $(A'BD)$ và $(ABCD)$ bằng $30^\circ$. Khoảng cách từ  $A$ đến mặt phẳng $(A'BD)$ bằng
	 \choice
		{$\dfrac{2a\sqrt{13}}{13}$} 
		 {$\dfrac{a}{4}$}
		 {$\dfrac{a\sqrt{14}}{7}$}
		{\True $\dfrac{a}{2}$} 
	 \loigiai{
		\begin{center}
			\begin{tikzpicture}[scale=1,font=\footnotesize,line join=round,line cap=round,>=stealth]
					\coordinate  (A) at (0,0);
					\coordinate (B)  at (-1,-1.5);
					\coordinate (D)  at (4,0);
					\coordinate (C) at ($(B)+(D)-(A)$); 
					\coordinate (A') at ($(A)+(0,3)$); 
					\coordinate (B') at  ($(B)+(0,3)$);
					\coordinate (C')  at ($(C)+(0,3)$);
					\coordinate (D') at ($(D)+(0,3)$); 
					\coordinate (O) at  ($(A)!0.5!(C)$);
					 \coordinate (H) at ($(A')!0.45!(O)$);
					\draw  (B)--(C)--(D)--(D')--(C')--(B')--(A')--(D')
						(B')--(B) 
						 (C')--(C);
					\draw[dashed,thin](A')--(A)--(B)--(A') 
						 (A)--(D)
						(A)--(O)--(A') 
						 (B)--(D)--(A')
						(A)--(H); 
					\pic[draw,thin,angle radius=6pt] {right  angle = A'--H--A} pic[draw,thin,angle radius=6pt] {right angle = A--O--B};
					\foreach \i/\g  in {A'/90,A/-90,B/-90,C/-90,D/0,C'/0,D'/30,B'/180,O/-90,H/0}{\draw[fill](\i) circle (1pt)
						($(\i)+(\g:3mm)$) node[scale=1]{$\i$};} 
				\end{tikzpicture}
		\end{center}
		Gọi $O$ là giao điểm của $AC$ và  $BD$.\\
		Khi đó, góc giữa hai  mặt phẳng $(A'BD)$ và $(ABCD)$ bằng $\angle A'OA=30^\circ$.\\
		Vẽ $AH \perp  A'O$ tại $H$.\\
		Ta  có $BD \perp (AA'O)$, mà $AH \subset (AA'O)$ nên $BD \perp AH$.\\
		Khi đó $\begin{cases} AH \perp A'O \\ AH \perp BD
	\end{cases} \Rightarrow AH \perp (A'BD) \Rightarrow \mathrm{d}(A,(A'BD))=AH$.\\
	 Trong $\triangle AOA'$ vuông tại $A$ có $AH = AO \sin \angle A'OA = a \sin 30^\circ = \dfrac{a}{2}$.\\
	Vậy $\mathrm{d}(A,(A'BD)) =  \dfrac{a}{2}$.
 }
\end{ex}

% nguồn: Ok_Toan11_TracNghiem_Lan2
\begin{ex}%[1H8V5-3]%[ma:4fba82d8]
	[Mã 101  BGD\&ĐT NĂM 2018-2019]
	Cho hình chóp $S. ABCD$ có đáy là hình vuông cạnh $a$, mặt bên $SAB$ là tam  giác đều và nằm trong mặt phẳng vuông góc với mặt phẳng đáy. Khoảng cách từ $A$ đến mặt phẳng $(SBD)$ bằng
	 \choice
		 {$\dfrac{\sqrt{21} a}{14}$}
		{\True $\dfrac{\sqrt{21} a}{7}$} 
		{$\dfrac{\sqrt{2} a}{2}$} 
		 {$\dfrac{\sqrt{21} a}{28}$}
	\loigiai{ 
	\immini{Gọi $H$ là trung điểm $AB$.  Suy ra $SH\perp(ABCD)$.\\
		Ta có $\dfrac{\mathrm{d}(H,(SBD))}{\mathrm{d}(A,(SBD))}=\dfrac{BH}{BA}=\dfrac{1}{2}$.\\Suy ra $\mathrm{d}(A,(SBD))=2\mathrm{d}(H,(SBD))$.\\ 
		Gọi $I$ là trung điểm  $OB$, suy ra $HI\parallel OA$, $HI=\dfrac{OA}{2}=\dfrac{a \sqrt{2}}{4}$. \\
		Lại  có $\heva{&BD\perp HI\\
		&BD\perp SH} \Rightarrow  BD\perp(SHI)$.\\
		Vẽ $HK\perp SI\Rightarrow HK\perp(SBD)$. \\ 
		Ta có $\dfrac{1}{HK^2}=\dfrac{1}{SH^2}+\dfrac{1}{HI^2} \Rightarrow HK=\dfrac{a \sqrt{21}}{14}$.\\ 
		 Suy ra $\mathrm{d}(A,(SBD))=2\mathrm{d}(H,(SBD))=2HK=\dfrac{a \sqrt{21}}{7}$.}
		 {
		\begin{tikzpicture}[scale=1,font=\footnotesize,line join=round,line cap=round,>=stealth]
				\def\a{3.5} 
				 \def\h{4.5}
				\path (0:0) coordinate  (A)
					++(0:\a)  coordinate (D)
					++(-145:\a/1.7)  coordinate (C)
					($(A)+(C)-(D)$) coordinate  (B)
					($(A)!0.5!(B)$)  coordinate (H)
					 ($(H)+(90:\h)$) coordinate (S)
					(intersection  of A--C and B--D) coordinate (O)
					($(B)!0.5!(O)$) coordinate  (I)
					($(S)!0.8!(I)$)  coordinate (K);
				\draw[dashed]  (B)--(A)--(D)
					(C)--(A)--(S) 
					 (K)--(H)--(I)--(S)
					(S)--(H) 
					(B)--(D); 
				 \draw[] (B)--(C)--(D)
					(B)--(S) 
					(C)--(S) 
					 (D)--(S);
				\draw  pic[draw,angle radius=2mm]{right angle=S--K--H};
				\draw pic[draw,angle radius=2mm]{right angle=I--H--S}; 
				\foreach \x/\g in  {A/45,B/-135,C/-45,D/45,S/90,H/135,O/-90,I/-20,K/10}
				 \fill[] (\x) circle (1pt)
					 ($(\g:3mm)+(\x)$) node {$\x$};
			\end{tikzpicture}
		} 
	} 
\end{ex}

% nguồn: Ok_Toan11_TracNghiem_Lan2
\begin{ex}%[1H8V5-3]%[ma:1cf653fc]
	[ĐTK BGD\&ĐT NĂM  2018-2019]
	Cho hình chóp $S.ABCD$ có đáy là hình thoi cạnh $a$, $\widehat{BAD}=60^{\circ}$, $SA=a$ và $SA$ vuông góc với mặt phẳng đáy.  Khoảng cách từ $B$ đến mặt phẳng $(SCD)$ bằng
	 \choice
		{\True $\dfrac{a \sqrt{21}}{7}$} 
		{$\dfrac{a  \sqrt{15}}{7}$}
		{$\dfrac{a \sqrt{21}}{3}$} 
		 {$\dfrac{a \sqrt{15}}{3}$}
	 \loigiai{
		\begin{center}
			\begin{tikzpicture}[scale=1,font=\footnotesize,line join=round,line cap=round,>=stealth]
					\def\a{4} 
					\def\h{4} 
					\path  (0:0) coordinate (B)
						(0:\a) coordinate (C) 
						(-30:\a/1.2)  coordinate (D)
						($(B)+(D)-(C)$) coordinate (A) 
						 ($(A)+(90:\h)$) coordinate (S)
						(intersection of A--C and B--D)  coordinate (O)
						 ($(C)!1/2!(D)$) coordinate (M)
						($(C)!1.4!(D)$) coordinate (K) 
						($(S)!(A)!(K)$) coordinate  (H);
					\draw[dashed] (A)--(B)--(C) 
						(B)--(S) 
						 (A)--(D)--(B)--(M);
					\draw[] (S)-- (A)--(K)--(S)--(C)--(K) 
						(S)--(D) 
						(A)--(H); 
					\foreach \x/\g  in {A/180,B/180,C/45,D/0,S/90,M/0,K/0,H/0}
					 \fill[] (\x) circle (1pt)
						($(\g:3mm)+(\x)$)  node {$\x$};
					\draw  pic[draw,angle radius=6pt]{right angle=D--A--S};
					 \draw pic[draw,angle radius=6pt]{right angle=C--K--A};
					 \draw pic[draw,angle radius=6pt]{right angle=A--H--K};
					\draw pic[draw,angle  radius=6pt]{right angle=C--M--B};
				\end{tikzpicture}
			\quad 
			\begin{tikzpicture}
					\def\a{3} 
					\path (0:0) coordinate  (D)
						 ++(30:\a) coordinate (C)
						++(150:\a) coordinate (B) 
						++(210:\a) coordinate (A) 
						($(C)!1/2!(D)$) coordinate (M) 
						 ($(C)!(A)!(D)$) coordinate (K);
					\draw (A)--(B)--(C)--(D)--cycle (A)--(K)--(D)--(B)--(M) 
						(A)--(C); 
					\draw  pic[draw,angle radius=6pt]{right angle=C--K--A};
					\foreach \x/ \goc in  {A/180,D/-90,C/0,B/90,K/-90,M/-45}
					\fill  (\x) circle (1pt)
						($(\x)+(\goc:3mm)$) node  {$\x$};
					\draw  pic[draw,angle radius=6pt]{right angle=C--M--B};
				\end{tikzpicture}
		\end{center}
		Ta có $AB\parallel CD\Rightarrow AB\parallel(SCD)$. Do  đó $\mathrm{d}(B,(SCD)=\mathrm{d}(A,(SCD))$.\\
		Vì $\widehat{BAD}=60^{\circ}$ nên  $\widehat{BCD}=60^{\circ}$.\\
		Mặt khác tứ giác $ABCD$ là hình thoi cạnh $a$  nên $\triangle BCD$ là tam giác đều cạnh $a$.\\
		Gọi $M$ là trung điểm của $CD$, suy ra $BM\perp  CD$.\\
		Kẻ $AK\parallel BM, K\in CD$, thì $AK\perp  CD$.\\
		Kẻ  $AH\perp SK$ tại $H$.\\
		Ta có  $\heva{&CD\perp AK\\
		&CD\perp SA} \Rightarrow  CD\perp(SAK) \Rightarrow CD\perp AH$, mà $SK\perp AH\Rightarrow AH\perp(SCD)$.\\
		Do  đó $\mathrm{d}(A,(SCD))=AH$.\\
		Ta có tứ giác $ABMK$ là hình chữ nhật nên  $AK=BM=\dfrac{a \sqrt{3}}{2}$.\\
		Khi  đó $AH\cdot SK=SA\cdot AK\Rightarrow AH=\dfrac{SA\cdot AK}{SK}$.\\
		Với $SA=a$,  $AK=\dfrac{a \sqrt{3}}{2}$, $SK=\sqrt{SA^2+AK^2}=\dfrac{a \sqrt{7}}{2}$ suy ra $AH=\dfrac{a \sqrt{21}}{7}$.\\
		Vậy $\mathrm{d}(B,(SCD))=\mathrm{d}(A,(SCD))=AH=\dfrac{a \sqrt{21}}{7}$. 
	} 
\end{ex}

% nguồn: Ok_Toan11_TracNghiem_Lan2
\begin{ex}%[1H8V5-3]%[ma:c2af86b9]
	Cho hình chóp $S.ABCD$ có đáy là hình thang vuông ở $A,B$. $SA \perp (ABCD)$, $SA=a\sqrt{2}$,  $AB=BC=a$, $AD=2a$. Tính khoảng cách từ điểm $B$ đến mặt phẳng $(SCD)$.
	 \choice
		 {$\mathrm{d}(B,(SCD)) = \dfrac{a\sqrt{3}}{3}$}
		{\True $\mathrm{d}(B,(SCD))  = \dfrac{a}{2}$}
		{$\mathrm{d}(B,(SCD)) =  a$}
		{$\mathrm{d}(B,(SCD)) = \dfrac{a\sqrt{6}}{2}$} 
	 \loigiai{
		\begin{center}
			\begin{tikzpicture}[scale=1,font=\footnotesize,line join=round,line cap=round,>=stealth]
					\coordinate (A) at  (0,0);
					 \coordinate (B) at (-1,-0.8);
					\coordinate (D)  at (5.2,0);
					\coordinate (C) at (1.6,-0.8); 
					\coordinate (S)  at (0,3.5);
					\coordinate (H)  at ($(S)!0.6!(C)$);
					\coordinate (J) at ($(A)!2!(B)$); 
					\coordinate  (I) at ($(A)!0.5!(D)$);
					\draw[dashed] (S) --  (A) -- (B)
						(H) -- (A)  -- (D)
						 (C)--(I);
					\draw  (S)--(B)--(C)--(D)--(S)--(C)
						(B)--(J)--(C); 
					\pic[draw,thin,angle radius=2mm] {right angle = D--A--B} pic[draw,thin,angle radius=2mm]  {right angle = C--B--A} pic[draw,thin,angle radius=2mm] {right angle = S--H--A};
					\foreach \i/\g in  {S/90,A/180,B/-90,C/-90,D/0,H/30,J/-90,I/90}{\draw[fill](\i) circle (1pt)
						 ($(\i)+(\g:3mm)$) node[scale=1]{$\i$};}
				\end{tikzpicture}
		\end{center}
		Gọi $J$ là  giao điểm của $AB$ với $CD$, $I$ là trung điểm của $AD$ và $H$ là hình chiếu vuông góc của $A$ trên $SC$.\\
		Ta  có $ABCI$ là hình vuông cạnh $a$.\\
		Ta có $\dfrac{\mathrm{d}(B,(SCD))}{\mathrm{d}(A,(SCD))} = \dfrac{BJ}{AJ} = \dfrac{1}{2} \Rightarrow \mathrm{d}(B,(SCD)) =  \dfrac{1}{2}\mathrm{d}(A,(SCD)) = \dfrac{AH}{2}$.\\
		Mà $\dfrac{1}{AH^2} = \dfrac{1}{SA^2} + \dfrac{1}{AC^2} = \dfrac{1}{(a\sqrt{2})^2} + \dfrac{1}{(a\sqrt{2})^2} =  \dfrac{1}{a^2} \Rightarrow AH = a$.\\
		Vậy $\mathrm{d}(B,(SCD)) = \dfrac{1}{2}AH = \dfrac{a}{2}$. 
	} 
\end{ex}

% nguồn: Ok_Toan11_TracNghiem_Lan2
\begin{ex}%[1H8V5-3]%[ma:d5273f5f]
	Cho lăng trụ tứ giác đều $A B C D.A'B'C'D'$ có cạnh đáy bằng $2 a$, $B'D=3 a$. Khoảng cách từ điểm $C'$ đến  mặt phẳng $(A'B C)$ bằng
	\choice 
		{$\dfrac{2 a  \sqrt{5}}{5}$}
		{\True  $\dfrac{a \sqrt{5}}{5}$}
		{$\dfrac{a  \sqrt{6}}{5}$}
		 {$a \sqrt{5}$}
	\loigiai{ 
		\begin{center}
			\begin{tikzpicture}[scale=1,font=\footnotesize,line join=round,line cap=round,>=stealth]
					 \coordinate (A) at (0,0);
					\coordinate (B) at (-2,-2); 
					 \coordinate (D) at (7,0);
					\coordinate (C) at ($(B)+(D)-(A)$); 
					\coordinate (A') at ($(A)+(0,5)$); 
					\coordinate (B')  at ($(B)+(0,5)$);
					\coordinate  (C') at ($(C)+(0,5)$);
					\coordinate (D')  at ($(D)+(0,5)$);
					\coordinate (H) at  ($(A')!0.6!(B)$);
					\coordinate (O)  at ($(A)!0.5!(C')$);
					\draw (B)--(B')--(A')--(D')--(C')--(B') 
						(C')--(C) 
						 (D')--(D)
						 (B)--(C)--(D);
					 \draw[dashed] (A)--(B)
						(A)--(D) 
						 (A)--(A')
						 (A')--(B)
						(A)--(H) 
						(A)--(C') 
						 (A')--(C);
					\pic[draw,thin,angle radius=6pt]  {right angle = A--H--B};
					\foreach \i/\g in {A'/90,B'/180,C'/90,D'/0,A/-90,B/-90,C/-90,D/0,H/180,O/-90}{\draw[fill=black](\i)  circle (1pt)
						 ($(\i)+(\g:3mm)$) node[scale=1]{$\i$};}
				\end{tikzpicture}
		\end{center}
		Vì  $A B C D.A'B'C'D'$ là lăng trụ tứ giác đều nên là lăng trụ đứng có đáy là hình vuông cạnh $2 a$, suy ra $B D=2 a \sqrt{2}$.\\
		Mà $B'D=3 a$ nên $B'B=\sqrt{B'D^{2}-B D^{2}}=\sqrt{9  a^{2}-8 a^{2}}=a$.\\
		Ta có $A C'  \cap\left(A'B C\right)=O$.\\
		Suy ra  $\mathrm{d}\left(C',\left(A'B C\right)\right)=\mathrm{d}\left(A,(A'BC)\right)$.\\
		Gọi $H$  là hình chiếu của $A$ lên $A'B$, ta chứng minh được $A H \perp\left(A'B C\right)$.\\
		 Suy ra $A H=\mathrm{d}\left(A,\left(A^{\prime} B C\right)\right)$.\\
		Tam giác $A'A B$ vuông tại $A$ và có $A H$ là đường cao  nên
		\[A H=\dfrac{A A'\cdot A B}{\sqrt{A A'^{2}+A B^{2}}}=\dfrac{a \cdot 2 a}{\sqrt{a^{2}+(2  a)^{2}}}=\dfrac{2 a \sqrt{5}}{5}.\]
	 }
\end{ex}

% nguồn: Ok_Toan11_TracNghiem_Lan2
\begin{ex}%[1H8V5-3]%[ma:692782bc]
	 [MĐ 102 BGD\&ĐT NĂM 2019-2020]
	Cho hình lăng trụ đứng $ABC.A'B'C'$ có đáy $ABC$  là tam giác đều cạnh $a$ và $AA'=2a$. Gọi $M$ là trung điểm của $CC'$. Khoảng cách từ $M$ đến mặt phẳng $\left(A'BC\right)$ bằng
	\choice 
		{$\dfrac{\sqrt{5} a}{5}$} 
		{$\dfrac{2\sqrt{5}  a}{5}$}
		{$\dfrac{2\sqrt{57} a}{19}$} 
		 {\True $\dfrac{\sqrt{57} a}{19}$}
	 \loigiai{
	\immini{ 
		Kẻ $AI\perp BC$,  $AH\perp A'I$.\\
		Ta có $\heva{&BC\perp AI\\ & BC\perp  AA'}\Rightarrow BC\perp (AA'I)\Rightarrow BC\perp AH$.\\
		Khi đó $\heva{&AH\perp SI\\ &  AH\perp BC}\Rightarrow AH\perp (A'BC)$.\\
		Ta có \allowdisplaybreaks$\begin{aligned}
		\mathrm{d}\left(M,\left(A'BC\right)\right)&=\dfrac{1}{2} \mathrm{d}\left(C',\left(A'BC\right)\right)=\dfrac{1}{2}  \mathrm{d}\left(A,\left(A'BC\right)\right)\\
		&=\dfrac{1}{2}AH=\dfrac{1}{2} \cdot \dfrac{AA'\cdot AI}{\sqrt{AA'^2+AI^2}}\,\left(*\right) 
	\end{aligned}$\\
	Tam giác $ABC$ đều cạnh $a$ có $AI$ là độ dài  đường trung tuyến nên $AI=\dfrac{a \sqrt{3}}{2}$.\\
	Ta có $(*) \Rightarrow \mathrm{d}\left(M,\left(A'BC\right)\right)=\dfrac{1}{2} \cdot \dfrac{2\cdot  \dfrac{\sqrt{3}}{2}}{\sqrt{4+\dfrac{3}{4}}}.a=\dfrac{\sqrt{3}}{\sqrt{19}} a=\dfrac{a \sqrt{57}}{19}$.}
	 {
	\begin{tikzpicture}[scale=0.9,font=\footnotesize,line join=round,line cap=round,>=stealth]
			 \def\a{4}
			\def\h{4.5} 
			\path (0:0) coordinate (A) 
				 ++(0:\a) coordinate (C)
				(-30:3*\a/4) coordinate  (B)
				($(A)+(90:\h)$) coordinate (A') 
				($(B)+(90:\h)$) coordinate (B') 
				($(C)+(90:\h)$) coordinate  (C')
				($(B)!1/2!(C)$)  coordinate (I)
				($(C')!1/2!(C)$) coordinate  (M)
				($(A')!(A)!(I)$) coordinate (H); 
			 \draw[dashed] (H)--(A)--(C)--(A')--(I)--(A);
			\draw[]  (C)--(C')
				 (B)--(B')
				 (A)--(A')--(B)
				(A)--(B)--(C) 
				 (A')--(B')--(C')--cycle;
			\foreach \x/\g  in {A/180,B/-45,C/0,A'/180,B'/-45,C'/0,M/0,I/-45,H/45}
			\fill (\x) circle (1pt) 
				 ($(\g:4mm)+(\x)$) node {$\x$};
		\end{tikzpicture}
	 }
} 
\end{ex}

% nguồn: Ok_Toan11_TracNghiem_Lan2
\begin{ex}%[1H8V5-3]%[ma:265576d0]
	Cho lăng trụ đều $A B C . A'B'C'$ có cạnh đáy bằng $a$, mặt phẳng $(A'B C)$ tạo với đáy một góc $45^{\circ}$, $M$ là điểm tùy ý thuộc cạnh $B'C'$. Khoảng các  từ điểm $M$ đến mặt phẳng $(A'BC)$ bằng
	\choice 
		 {$\dfrac{a \sqrt{6}}{2}$}
		 {\True $\dfrac{a \sqrt{6}}{4}$}
		 {$\dfrac{a \sqrt{3}}{2}$}
		{$\dfrac{a  \sqrt{3}}{4}$}
	 \loigiai{
		\begin{center}
			\begin{tikzpicture}[scale=1,font=\footnotesize,line join=round,line cap=round,>=stealth]
					\coordinate (A) at (-2,0); 
					\coordinate (B) at  (3.5,-3);
					 \coordinate (C) at (5,0);
					\coordinate (A')  at ($(A)+(0,5)$);
					\coordinate (B') at  ($(B)+(0,5)$);
					\coordinate (C') at ($(C)+(0,5)$); 
					\coordinate  (H) at ($(B)!0.5!(C)$);
					\coordinate (M) at ($(B')!0.5!(C')$); 
					\coordinate (I) at ($(A')!0.3!(H)$); 
					 \coordinate (O) at ($(A')!0.5!(B)$);
					 \draw (A)--(B)--(B')--(A')--cycle (A')--(C')--(B')
						(B)--(C)--(C') 
						 (A')--(B)
						(A)--(B'); 
					 \draw[dashed](A)--(C)
						(A)--(H)--(A') 
						 (A')--(C)
						(A)--(I); 
					\pic[draw,thin,angle radius=2mm] {right angle = A--H--B} pic[draw,thin,angle radius=2mm] {right  angle = A--I--H};
					\foreach \i/\g  in {A/180,B/-90,C/0,A'/90,C'/0,M/-45,B'/-45,I/90,O/-90,H/-45}{\draw[fill=black](\i) circle (1pt)
						 ($(\i)+(\g:3mm)$) node[scale=1]{$\i$};}
				\end{tikzpicture}
		\end{center}
		Vì $A B C.A'B'C'$ là lăng trụ tam giác  đều nên là lăng trụ đứng có đáy $ABC$ là tam giác đều.\\
		Ta có $B'C' \parallel \left(A'B C\right)$ nên  $\mathrm{d}\left(M,\left(A^{\prime} B C\right)\right)=\mathrm{d}\left(B^{\prime},\left(A^{\prime} B C\right)\right)$.\\
		Mà $A B^{\prime} \cap\left(A^{\prime} B C\right)=O$ với $O$ là trung điểm $A B^{\prime}$ nên  $\mathrm{d}\left(B^{\prime},\left(A^{\prime} B C\right)\right)=\mathrm{d}\left(A,\left(A^{\prime} B C\right)\right)$.\\
		Gọi $H$ là hình chiếu của $A$ lên $B C$,  $I$ là hình chiếu của $A$ lên $A^{\prime} H$.\\
		Ta  chứng minh được $A I \perp\left(A^{\prime} B C\right)$, suy ra $\mathrm{d}\left(A,\left(A^{\prime} B C\right)\right)=A I$.\\
		Mà $\left(\widehat{\left(A^{\prime} B C\right),(A B C)}\right)=\widehat{A^{\prime} H A}=45^{\circ}$ nên tam giác $A^{\prime} A H$ vuông cân tại $A$,  do đó
		\[A'H=A H \sqrt{2}=\frac{a \sqrt{3}}{2} \cdot \sqrt{2}=\frac{a \sqrt{6}}{2}.\] 
		Mặt khác $A I$ là đường cao của tam giác $A'A H$ nên $AI=\dfrac{A'H}{2}=\dfrac{\dfrac{a\sqrt{6}}{2}}{2}=\dfrac{a  \sqrt{6}}{4}$.
	 }
\end{ex}

% nguồn: Ok_Toan11_TracNghiem_Lan2
\begin{ex}%[1H8V5-3]%[ma:b2640efc]
	Cho hình chóp $S.ABC$ có đáy $ABC$ là tam giác vuông tại $A$, $AC=a\sqrt{3}$, $\widehat{ABC}=60^\circ$. Gọi $M$ là trung điểm của $BC$.  Biết $SA=SB=SM=\dfrac{2a\sqrt{3}}{3}$. Tính khoảng cách $\mathrm{d}$ từ đỉnh $S$ đến $(ABC)$.
	\choice 
		 {$\mathrm{d} = \dfrac{2a\sqrt{3}}{3}$}
		 {\True $\mathrm{d}=a$}
		 {$\mathrm{d}=2a$}
		 {$\mathrm{d}=a\sqrt{3}$}
	 \loigiai{
		\begin{center}
			\begin{tikzpicture}[scale=1,font=\footnotesize,line join=round,line cap=round,>=stealth]
					 \coordinate (B) at (0,0);
					\coordinate (M)  at (4,0);
					\coordinate  (A) at (1,-1.5);
					\coordinate (C) at (8,0); 
					\coordinate (I)  at ($(A)!0.5!(B)$);
					\coordinate (O)  at ($(M)!0.6666!(I)$);
					\coordinate (S)  at ($(O)+(0,4)$);
					 \draw (A)--(S)--(C)--(A)--(B)--(S);
					\draw[dashed,thin] (A)--(M)--(S) 
						(B)--(C) 
						(I)--(M) 
						 (S)--(O);
					\pic[draw,thin,angle radius=6pt] {right angle  = S--O--I} pic[draw,thin,angle radius=6pt] {right angle = B--A--C};
					\foreach \i/\g in  {S/90,B/180,A/-90,C/0,M/-90,O/-90,I/180}{\draw[fill](\i) circle (1pt)
						 ($(\i)+(\g:3mm)$) node[scale=1]{$\i$};}
				\end{tikzpicture}
		\end{center}
		Từ  giả thiết suy ra $\triangle ABM$ đều.\\
		Lại có $SA=SB=SM$ nên $S.ABM$ là  hình chóp đều.\\
		Xét $\triangle ABC$ vuông tại $A$ có $\sin 60^\circ = \dfrac{AC}{BC} \Rightarrow BC =  \dfrac{AC}{\sin 60^\circ} = \dfrac{a\sqrt{3}}{\dfrac{\sqrt{3}}{2}} = 2a \Rightarrow AM=AB=BM=a$.\\
		Gọi $O$  là trọng tâm của tam giác $ABM$ $\Rightarrow SO \perp (ABM)$, từ đó $\mathrm{d}(S;(ABC))=SO$.\\
		Vì $\triangle ABM$ đều cạnh $a$ nên  $OM=\dfrac{2}{3}MI=\dfrac{2}{3}\cdot\dfrac{a\sqrt{3}}{2}=\dfrac{a\sqrt{3}}{3}$.\\
		Xét $\triangle SOM$  vuông tại $O$ có $SO=\sqrt{SM^2-OM^2}=\sqrt{\dfrac{4a^2}{3}-\dfrac{a^2}{3}}=a$.\\
		Vậy  $\mathrm{d}(S;(ABC))=a$.
	 }
\end{ex}

% nguồn: Ok_Toan11_TracNghiem_Lan2
\begin{ex}%[1H8V5-3]%[ma:778a75b6]
	Cho lăng trụ tam giác đều $ABC.A'B'C'$ có cạnh đáy bằng $4a$.  Góc giữa hai mặt phẳng $(A'BC)$ và $(ABC)$ bằng $30^\circ$. Gọi $M$ là trung điểm của cạnh $AB$. Tìm khoảng cách từ điểm $M$ đến mặt phẳng $(A'BC)$?
	\choice 
		{\True $\dfrac{a\sqrt{3}}{2}$} 
		 {$3a$}
		{$a\sqrt{3}$} 
		 {$\dfrac{3a}{2}$}
	 \loigiai{
		\begin{center}
			\begin{tikzpicture}[scale=1,font=\footnotesize,line join=round,line cap=round,>=stealth]
					\coordinate (A) at  (0,0);
					\coordinate (B) at  (1,-1.5);
					\coordinate (C)  at (4,0);
					\coordinate (A')  at ($(A)+(0,3)$);
					\coordinate (B') at  ($(B)+(0,3)$);
					\coordinate  (C') at ($(C)+(0,3)$);
					\coordinate  (M) at ($(A)!0.5!(B)$);
					\coordinate (N) at ($(C)!0.5!(B)$); 
					 \coordinate (H) at ($(A')!0.6!(N)$);
					\draw  (H)--(A)--(B)--(C)--(C')--(B')--(A')--(A)
						 (B')--(B)
						 (B)--(A')--(C');
					\draw[dashed,thin]  (A')--(N)--(A)--(C)--(A');
					\pic[draw,thin,angle radius=6pt] {right angle =  A'--H--A};
					\foreach \i/\g in {A'/90,A/-90,B/-90,C/-90,C'/0,B'/180,H/0,M/180,N/-90}{\draw[fill](\i) circle (1pt) 
						($(\i)+(\g:3mm)$)  node[scale=1]{$\i$};}
				\end{tikzpicture}
		\end{center}
		Gọi $N$  là trung điểm của $BC$.\\
		Do $ABC.A'B'C'$ là lăng trụ tam giác đều nên $BN \perp (A'BC)$, $AA' \perp AN$ và $AN=2a\sqrt{3}$.\\ Suy ra  $BC \perp (A'AN)$.\\
		Từ đó ta có $((A'BC),(ABC))  = \widehat{A'NA}=30^\circ$.\\
		Gọi $H$ là hình chiếu của $A$ trên $A'N$, do $BC \perp (A'AN)$ nên $AH \perp AN, BC \perp AH \Rightarrow  AH \perp (A'BC)$.\\
		Xét tam giác $AHN$ vuông tại $H$ có $AH = AN\sin\angle A'NA = 2a\sqrt{3} \sin 30^\circ  = 2a\sqrt{3} \dfrac{1}{2} = a\sqrt{3}$.\\
		Mặt khác, $M$ là  trung điểm của cạnh $AB$ nên $\mathrm{d}(M,(A'BC)) = \dfrac{1}{2}\mathrm{d}(A,(A'BC)) = \dfrac{a\sqrt{3}}{2}$.
	 }
\end{ex}

% nguồn: Ok_Toan11_TracNghiem_Lan2
\begin{ex}%[1H8V5-3]%[ma:313fd24e]
	[MĐ 102 BGD\&ĐT  NĂM 2018-2019]
	Cho hình chóp $S. ABCD$ có đáy là hình vuông  cạnh $a$, mặt bên $SAB$ là tam giác đều và nằm trong mặt phẳng vuông góc với mặt phẳng đáy. Khoảng cách từ $C$ đến mặt phẳng $(SBD)$ bằng
	\choice 
		{$\dfrac{\sqrt{21}  a}{28}$}
		{$\dfrac{\sqrt{21}  a}{14}$}
		{$\dfrac{\sqrt{2} a}{2}$} 
		 {\True $\dfrac{\sqrt{21} a}{7}$}
	\loigiai{ 
		\begin{center}
			\begin{tikzpicture}[scale=1,font=\footnotesize,line join=round,line cap=round,>=stealth]
					\def\a{4} 
					\def\h{4} 
					\path (0:0) coordinate  (A)
						++(0:\a)  coordinate (D)
						++(-145:\a/1.7) coordinate (C) 
						($(A)+(C)-(D)$)  coordinate (B)
						($(A)!0.5!(B)$) coordinate (M) 
						($(M)+(90:\h)$) coordinate  (S)
						 (intersection of A--C and B--D) coordinate (O)
						($(B)!0.5!(O)$) coordinate  (K)
						 ($(S)!0.7!(K)$) coordinate (H);
					\draw[dashed]  (B)--(A)--(D)
						(C)--(A)--(S) 
						 (H)--(M)--(K)--(S)
						(S)--(M) 
						 (B)--(D);
					 \draw[] (B)--(C)--(D)
						 (B)--(S)
						(C)--(S) 
						 (D)--(S);
					\draw pic[draw,angle radius=2mm]{right  angle=S--H--M};
					 \draw pic[draw,angle radius=2mm]{right angle=K--M--S};
					\foreach \x/\g in {A/45,B/-135,C/-45,D/45,S/90,H/45,O/-90,K/-20,M/180} 
					 \fill[] (\x) circle (1pt)
						($(\g:4mm)+(\x)$)  node {$\x$};
				\end{tikzpicture}
			\quad 
			\begin{tikzpicture}[scale=0.9,font=\footnotesize,line join=round,line cap=round,>=stealth]
					 \def\a{3}
					\path (0:0) coordinate  (B)
						++(0:\a)  coordinate (C)
						++(90:\a) coordinate (D) 
						++(180:\a) coordinate  (A)
						($(A)!1/2!(B)$)  coordinate (M)
						($(A)!1/2!(C)$)  coordinate (O)
						($(B)!1/2!(O)$) coordinate (K); 
					 \draw pic[draw,angle radius=2mm]{right angle=C--O--D};
					\draw pic[draw,angle  radius=2mm]{right angle=M--K--B};
					\draw (A)--(B)--(C)--(D)--cycle  (A)--(C)
						(B)--(D) 
						(K)--(M)--(C); 
					\foreach \x/ \goc  in {A/135,B/-135,C/-45,D/45,O/90,M/180,K/-90}
					 \fill (\x) circle (1pt)
						($(\x)+(\goc:3mm)$) node  {$\x$};
				\end{tikzpicture}
		\end{center}
		Gọi $M$  là trung điểm của $AB\Rightarrow SM\perp(ABCD)$. Gọi $O=AC\cap BD$.\\
		Ta có $\heva{&AC\cap(SBD)=O\\ &AO=OC} 
		\Rightarrow  \mathrm{d}(C,(SBD))=\mathrm{d}(A,(SBD))$.\\
		Lại có $\heva{&AM\cap  (SBD)=B\\ &AB=2MB} \Rightarrow \mathrm{d}(A,(SBD))=2\mathrm{d}(M,(SBD))$.\\
		 Vậy $\dfrac{\mathrm{d}(C,(SBD))}{\mathrm{d}(M,(SBD))}=2$.\\
		Kẻ $MK\perp BD$  $(K\in BD)$, kẻ $MH\perp SK$ tại $H\Rightarrow MH=\mathrm{d}(M,(SBD))$.\\
		Xét tam giác $SMK$  có $MK=\dfrac{1}{2} AO=\dfrac{1}{2} \cdot \dfrac{a \sqrt{2}}{2}=\dfrac{a \sqrt{2}}{4}$, $SM=\dfrac{a \sqrt{3}}{2}$.\\
		Áp  dụng hệ thức lượng $\dfrac{1}{MH^2}=\dfrac{1}{SM^2}+\dfrac{1}{MK^2}=\dfrac{28}{3a^2} \Rightarrow MH=\dfrac{a \sqrt{21}}{14} \Rightarrow \mathrm{d}(C,(SBD)=\dfrac{a \sqrt{21}}{7}$.
	 }
\end{ex}

% nguồn: Ok_Toan11_TracNghiem_Lan2
\begin{ex}%[1H8V5-3]%[ma:f4420af3]
	 Cho hình chóp $S.ABCD$ có đáy là hình vuông cạnh $a$, tam giác $SAB$ là tam giác đều và mặt phẳng $(SAB)$ vuông góc với mặt phẳng $(ABCD)$. Tính khoảng cách từ điểm $D$ đến mặt phẳng $(SBC)$.
	 \choice
		{\True $\dfrac{a\sqrt{3}}{2}$} 
		 {$\dfrac{a\sqrt{3}}{4}$}
		 {$\dfrac{a}{2}$}
		{$\dfrac{a}{4}$} 
	 \loigiai{
		\begin{center}
			\begin{tikzpicture}[scale=1,font=\footnotesize,line join=round,line cap=round,>=stealth]
					 \coordinate (A) at (0,0);
					\coordinate (B) at  (-2.2,-0.8);
					\coordinate (D) at  (3.8,0);
					\coordinate (C)  at (1.6,-0.8);
					\coordinate (K)  at (0,4);
					\coordinate (H) at ($(A)!0.5!(B)$); 
					\coordinate (S) at  ($(H)+(K)-(A)$);
					\coordinate  (K) at ($(S)!0.4!(B)$);
					\draw[dashed] (S) -- (A) --  (D)
						(K)--(A) --  (B)
						(S) -- (H); 
					 \draw (S)--(B)--(C)--(D)--(S)--(C);
					 \foreach \i/\g in {S/90,A/30,B/-90,C/-90,D/0,H/150,K/180}{\draw[fill](\i) circle (1pt)
						 ($(\i)+(\g:3mm)$) node[scale=1]{$\i$};}
				\end{tikzpicture}
		\end{center}
		Gọi $H$ là  trung điểm $AB$.\\
		Vì tam giác  $SAB$ đều nên $SH \perp AB$.\\
		Ta có $\heva{ &(SAB) \cap (ABCD) = AB \\ &SH \perp (ABCD)} \Rightarrow  SH \perp (ABCD)$.\\
		 Dễ thấy $BC \perp (SAB) \Rightarrow (SBC) \perp (SAB)$.\\
		Kẻ $AK \perp SB \Rightarrow AK  \perp (SBC) \Rightarrow \mathrm{d}(A,(SBC))=AK=\dfrac{a\sqrt{3}}{2}$.\\
		Vì $AD//(SBC)$ nên $\mathrm{d}(D,(SBC)) = \mathrm{d}(A,(SBC))  = \dfrac{a\sqrt{3}}{2}$.
	 }
\end{ex}

% nguồn: Ok_TracNghiemLop11_Lan1
\begin{ex}%[1H8V5-3]%[ma:12be8421]
	 \immini{
	Cho hình chóp $S.ABCD$ có đáy là hình vuông cạnh $a$. Biết $SA$ vuông góc với  đáy và $SA=a$ (tham khảo hình vẽ). Khoảng cách từ điểm $A$ đến mặt phẳng $(SBD)$ bằng
	\choice 
		{$\dfrac{2a}{\sqrt{3}}$} 
		{\True  $\dfrac{a}{\sqrt{3}}$}
		{$\dfrac{a}{2\sqrt{3}}$} 
		{$\dfrac{a\sqrt{2}}{6}$} 
		 }{
		\begin{tikzpicture}[scale=0.8,font=\footnotesize,line join=round,line cap=round,>=stealth]
				\coordinate (A) at (0,0); 
				\coordinate  (B) at (-1.5,-1);
				\coordinate (D) at  (4,0);
				\coordinate (C) at  ($(B)+(D)-(A)$);
				\coordinate  (S) at ($(A)+(0,3)$);
				 \draw (S)--(B)
					 (S)--(C)
					 (S)--(D)
					 (B)--(C)
					(C)--(D); 
				\draw[dashed] (A)--(B) 
					 (A)--(D)
					 (S)--(D)
					(S)--(A); 
				\draw pic[draw,angle  radius=0.2cm] {right angle = S--A--D};
				\foreach \x  in{S,A,B,C,D}{
					\fill  (\x) circle (1pt);
				 }
				\path 
					node at (S)[above]  {$S$}
					node at (A)[above  left] {$A$}
					node at  (B)[below] {$B$}
					node at  (C)[below] {$C$}
					node at (D)[right] {$D$}; 
			\end{tikzpicture}
		 }
	 \loigiai{
	 \immini{
		Gọi  $O$ là tâm hình vuông $ABCD$.\\
		Kẻ  $AH\perp SO$ tại $H$.\\
		Ta có $\heva{&BD\perp AC\\&BD\perp SA}\Rightarrow BD\perp (SAC)\Rightarrow  BD\perp AH$.\\
		Ta có $\heva{&AH\perp  SO\\&AH\perp BD}\Rightarrow AH\perp (SBD)\Rightarrow \mathrm{d}(A,(SBD))=AH$.\\
		Tính $AC=a\sqrt{2}\Rightarrow  AO=\dfrac{AC}{2}=\dfrac{a\sqrt{2}}{2}$;\\
		$\dfrac{1}{AH^2}=\dfrac{1}{SA^2}+\dfrac{1}{AO^2}=\dfrac{1}{a^2}+\dfrac{2}{a^2}\Rightarrow AH^2=\dfrac{a^2}{3}\Rightarrow  AH=\dfrac{a}{\sqrt{3}}$.\\
		Vậy  $\mathrm{d}(A,(SBD))=AH=\dfrac{a}{\sqrt{3}}$.
		 }{
		\begin{tikzpicture}[scale=0.8,font=\footnotesize,line join=round,line cap=round,>=stealth]
				 \coordinate (A) at (0,0);
				\coordinate (B) at (-1.5,-1); 
				\coordinate (D) at  (4,0);
				\coordinate (C) at ($(B)+(D)-(A)$); 
				\coordinate (S) at  ($(A)+(0,3)$);
				\coordinate (O) at  ($(B)!0.5!(D)$);
				\coordinate (H)  at ($(S)!0.6!(O)$);
				\draw (S)--(B) 
					 (S)--(C)
					 (S)--(D)
					 (B)--(C)
					 (C)--(D);
				\draw[dashed]  (A)--(B)
					 (A)--(D)
					(S)--(D) 
					 (S)--(A)
					 (A)--(C)
					(B)--(D) 
					(A)--(H) 
					 (S)--(O);
				\draw pic[draw,angle radius=0.2cm] {right angle =  S--A--D};
				\draw pic[draw,angle radius=0.2cm]  {right angle = A--H--S};
				\foreach \x  in{S,A,B,C,D,O,H}{
					\fill (\x) circle  (1pt);
				} 
				 \path
					node at  (S)[above] {$S$}
					node at (A)[above left]  {$A$}
					 node at (B)[below] {$B$}
					node at (C)[below] {$C$} 
					node  at (D)[right] {$D$}
					 node at (O)[below] {$O$}
					node  at (H)[right] {$H$};
			\end{tikzpicture}
		 }
	} 
\end{ex}

% nguồn: Ok_Toan11_TracNghiem_Lan2
\begin{ex}%[1H8C5-3]%[ma:370d1e04]
	Cho hình chóp $S.ABCD$ có đáy là hình thang vuông tại $A$ và $D$, $AB=AD=a$, $CD=2a$. Hình chiếu của đỉnh $S$ lên mặt $(ABCD)$ trùng với trung điểm $H$ của $BD$. Biết thể tích  tứ diện $SBCD$ bằng $\dfrac{a^3\sqrt{6}}{6}$. Khoảng cách từ đỉnh $A$ đến mặt phẳng $(SBC)$ là?
	 \choice
		{$\dfrac{a\sqrt{3}}{2}$} 
		 {$\dfrac{a\sqrt{2}}{6}$}
		 {$\dfrac{a\sqrt{3}}{6}$}
		 {\True $\dfrac{a\sqrt{6}}{4}$}
	 \loigiai{
	\immini{Gọi $M$ là trung điểm  $CD$.\\
		Ta có $ABMD$ là hình vuông cạnh $a$ do đó  $BC=BD=a\sqrt{2}$.\\
		$\Rightarrow CD^2=4a^2=BC^2+BD^2$ do đó $\triangle BCD$ vuông  cân tại $B$.\\
		Gọi $H$ là trung điểm của  $BD$ thì $SH \perp (ABCD)$.\\
		 Khi đó
		\begin{eqnarray*}
			V_{S.BCD}&=&\dfrac{1}{3}S_{HBD}\cdot BC\\ 
			 &=&\dfrac{1}{3}\cdot \dfrac{1}{2}BD\cdot BC
		\end{eqnarray*}
		Suy  ra $SH=\dfrac{\dfrac{a^3\sqrt{6}}{6}}{\dfrac{1}{2}\cdot a\sqrt{2}\cdot a\sqrt{2}}=\dfrac{a\sqrt{6}}{2}$.
		}{	\begin{tikzpicture}[declare function={r=2; gocD=50;}]
		\path  (0,0) coordinate (A)
			(0:r)  coordinate (B)
			(gocD:0.8*r)  coordinate (D)
			($(B)+(D)-(A)$)  coordinate (M)
			 ($(M)!-1!(D)$) coordinate (C)
			 ($(B)!0.5!(D)$) coordinate (H)
			($(H)+(90:2*r)$)  coordinate (S)
			($(S)!0.8!(B)$) coordinate (I); 
		\foreach \x/\y/\z in {S/H/D,H/I/B}{ 
			 \path pic[draw,angle radius=5pt]{right angle= \x--\y--\z};
		} 
		 \draw (S)--(A)--(B)--(C)--(S)--(B);
		\draw [dashed] (S)--(D)--(B)--(M) 
			 (A)--(D)--(C)
			(S)--(H)--(I); 
		\foreach  \t/\g in {S/90,A/180,B/-60,C/0,D/160,H/-160,M/90,I/20}{
			\fill (\t)  circle (1pt) node[shift={(\g:7pt)},font=\footnotesize]{$ \t $};
		} 
	\end{tikzpicture}
		}\noindent 
		Hạ $HI  \perp SB$.
		Vì $ABMD$ là hình vuông nên $H$  là trung điểm của $AM$.\\
		Ta có $AMCB$ là hình  bình hành do đó $AH \parallel BC$.\\
		 Suy ra $\mathrm{\,d}_{(A;(SBC))}=\mathrm{\,d}_{(H;(SBC))}=HI$.\\
		 Khi đó
		\begin{eqnarray*}
			\dfrac{1}{HI^2}&=&\dfrac{1}{SH^2}+\dfrac{1}{HB^2}\\ 
			&=&\dfrac{1}{\left(\dfrac{a\sqrt{6}}{2}\right)^2}+\dfrac{1}{\left(\dfrac{a\sqrt{2}}{2}\right)^2}\\ 
			 &=&\dfrac{4}{6a^2}+\dfrac{4}{2a^2}\\
			 &=&\dfrac{2}{3a^2}+\dfrac{2}{a^2}\\
			 &=&\dfrac{8}{3a^2}
		\end{eqnarray*}
		 $\Rightarrow HI=\sqrt{\dfrac{3a^2}{8}}=\dfrac{a\sqrt{3}}{2\sqrt{2}}=\dfrac{a\sqrt{6}}{4}$ hay $\mathrm{\,d}_{(A;(SBC))}=\dfrac{a\sqrt{6}}{4}$.
	 }
\end{ex}


%%%%%%%%%%%%%%%% ĐÚNG/SAI %%%%%%%%%%%%%%%%

% nguồn: Ok_dungsai_lop11_Lan1
\begin{ex}%[1H8H5-3]%[ma:0a18870d]
 Cho hình chóp $S.ABCD$  có $SA\perp(ABCD)$, $SA=a\sqrt{3}$, $ABCD$ là hình vuông cạnh bằng $a$. Khi đó:
  \choiceTF
   {$\mathrm{d}(A,(SBC))=\dfrac{\sqrt{3}}{3}a$}
  {\True $AD\parallel(SBC)$} 
  {\True $\mathrm{d}(D,(SBC))=\dfrac{\sqrt{3}}{2}a$} 
  {\True Gọi $M$  là trung điểm của $SA$. Khi đó: $\mathrm{d}(M,(SBC))=\dfrac{\sqrt{3}}{4}a$}
  \loigiai{
  \begin{center}
   \begin{tikzpicture}
      \def\a{4}
     \def\h{4} 
     \path (0:0) coordinate  (A)
      ++(0:\a) coordinate (D) 
      ++(-130:\a/2) coordinate (C) 
      ($(A)+(C)-(D)$)  coordinate (B)
      ($(A)+(90:\h)$) coordinate  (S)
       ($(S)!0.5!(A)$) coordinate (M)
      ($(S)!0.5!(B)$)  coordinate (H)
      (intersection of A--C and B--D) coordinate (O); 
     \draw[dashed,thick]  (B)--(A)--(D)
       (H)--(A)--(S)
       (C)--(A)
       (B)--(D);
     \draw[thick] (B)-- (C)--(D) 
      (B)--(S) 
      (C)--(S) 
       (D)--(S);
      \foreach \x/\g in {A/180,B/-135,C/-45,D/45,S/90,M/0,H/135,O/90}
     \fill[black]  (\x) circle (1pt)
      ($(\g:3mm)+(\x)$)  node {$\x$};
     \draw pic[draw,angle radius=3mm]{right angle=D--A--S}; 
     \draw pic[draw,angle  radius=3mm]{right angle=A--H--S};
    \end{tikzpicture}
  \end{center}
  Kẻ  $SH\perp SB$ tại $H$.\\
  Ta  có $\heva{&BC\perp SA\\&BC\perp AB}\Rightarrow BC\perp (SAB) \Rightarrow BC\perp AH$.\\
  Ta lại có $AH\perp SB  \Rightarrow AH\perp (SBC)\Rightarrow \mathrm{d}(A,(SBC))=AH$.\\
   Ta có $\dfrac{1}{AH^2}=\dfrac{1}{SA^2}+\dfrac{1}{AB^2}=\dfrac{1}{3a^2}+\dfrac{1}{a^2}\Rightarrow AH=\dfrac{a\sqrt{3}}{2}$.\\
  Vậy $\mathrm{d}(A,(SBC))=\dfrac{\sqrt  {3}}{2}a$.\\
  Ta có $AD\parallel(SBC) \Rightarrow  \mathrm{d}(D,(SBC))=\dfrac{\sqrt{3}}{2}a$.\\
  Ta có $MA$ cắt  $(SBC)$ tại $S$. \\
   Suy ra $ \dfrac{\mathrm{d}(M,(SBC))}{\mathrm{d}(A,(SBC))}=\dfrac{MS}{AS}=\dfrac{1}{2}\Rightarrow \mathrm{d}(M,(SBC))=\dfrac{1}{2}\mathrm{d}(A,(SBC))=\dfrac{1}{2} \cdot\dfrac{\sqrt {3}}{2}a=\dfrac{\sqrt{3}}{4}a$.
  \begin{itemchoice}
    \itemch Sai. Vì $\mathrm{d}(A,(SBC))=\dfrac{\sqrt {3}}{2}a$.
   \itemch Đúng. Vì  $AD\parallel(SBC)$.
   \itemch Đúng. Vì $\mathrm{d}(D,(SBC))=\dfrac{\sqrt{3}}{2}a$. 
   \itemch Đúng. Vì $\mathrm{d}(M,(SBC))=\dfrac{\sqrt{3}}{4}a$ 
  \end{itemchoice}
  }
\end{ex}

% nguồn: Ok_dungsai_lop11_Lan2
\begin{ex}%[1H8H5-3]%[ma:ce5ef609]
	Cho hình chóp $S.ABC$ có mặt bên $(SAB)$ vuông góc với mặt đáy và $\triangle SAB$ đều cạnh $2a$. Biết  $\triangle ABC$ vuông tại $C$ và cạnh $AC=a\sqrt{3}$.
	 \choiceTF
		{\True$SH \perp (ABC)$} 
		{\True$\mathrm{d}\left(S,(ABC)\right)=a\sqrt{3}$} 
		{Khoảng cách từ $C$  đến mặt phẳng $(SAB)$ bằng $\dfrac{a\sqrt{3}}{3}$}
		{Thể tích của khối chóp $S.ABC$  bằng $\dfrac{a^3}{6}$}
	\loigiai{\begin{center}
		\begin{tikzpicture}[scale=1,font=\footnotesize,line join=round,line cap=round,>=stealth]
				\coordinate (A)  at (0,0);
				\coordinate (B)  at (1.5,-2);
				\coordinate (C)  at (5,0);
				\coordinate (H) at  ($(A)!0.5!(B)$);
				 \coordinate (S) at ($(H)+(0,5)$);
				 \coordinate (K) at ($(A)!0.7!(B)$);
				\draw[dashed]  (S) -- (A) -- (B)
					(K) --  (A) -- (H)
					(A)--(C)--(K); 
				\draw (S)--(A)--(B)--(C)--(S)--(B) 
					 (S)--(H);
				\pic[draw,thin,angle radius=3mm] {right angle = S--H--A} pic[draw,thin,angle radius=3mm] {right angle = B--C--A} pic[draw,thin,angle radius=3mm] {right angle = C--K--A}; 
				\foreach \i/\g in {S/90,H/180,B/-90,A/-90,C/0,K/-90}{\draw[fill](\i)  circle (1pt)
					 ($(\i)+(\g:3mm)$) node[scale=1]{$\i$};}
			\end{tikzpicture}
	\end{center}
	\begin{itemchoice}
		\itemch Gọi $H$ là trung điểm $AB$, mà  $\triangle SAB$ đều nên $SH \perp AB$.\\
		Ngoài ra $(SAB) \perp (ABC)$ nên  $SH \perp (ABC)$.\\
		\itemch Ta có $\mathrm{d}\left(S,(ABC)\right)=SH=\dfrac{2a\sqrt{3}}{2}=a\sqrt{3}$ (do $\triangle SAB$ đều cạnh $2a$). 
		\itemch Kẻ đường  cao $CK$ của $\triangle ABC$.\\
		 Ta có
		$\heva{&CK \perp AB \\&CK \perp SH} \Rightarrow  CK \perp (SAB) \Rightarrow \mathrm{d}\left(C,(SAB)\right)=CK$.\\
		Xét  $\triangle ABC$ vuông tại $C$, ta có\\
		$BC=\sqrt{AB^2-AC^2}=\sqrt{(2a)^2-\left(a\sqrt{3}\right)^2}=\sqrt{4a^2-3a^2}=\sqrt{a^2}=a$\\ 
		$CK=\dfrac{CA \cdot CB}{AB}=\dfrac{a\sqrt{3}\cdot a}{2a}=\dfrac{a\sqrt{3}}{2}.$\\ 
		Vậy  $\mathrm{d}\left(C,(SAB)\right)=CK=\dfrac{a\sqrt{3}}{2}$.
		\itemch Diện tích đáy hình chóp là $S_{ABC}=\dfrac{1}{2} AC  \cdot BC=\dfrac{1}{2} a\sqrt{3} \cdot a=\dfrac{a^2\sqrt{3}}{2}$.\\
		Thể tích khối chóp là $V_{S.ABC}=\dfrac{1}{3} SH \cdot  S_{ABC}=\dfrac{1}{3} \cdot a\sqrt{3} \cdot \dfrac{a^2\sqrt{3}}{2}=\dfrac{a^3}{2}$.
	\end{itemchoice}
 }
\end{ex}

% nguồn: Ok_dungsai_lop11_Lan2
\begin{ex}%[1H8H5-3]%[ma:14a6076b]
	Cho hình chóp $S.ABC$ có đáy là tam giác $ABC$ vuông tại $B$ có $AB=1$, $\widehat{ACB}=30^\circ$. Biết $SA$  vuông góc với mặt đáy và $SA=2$. Gọi $H$ là hình chiếu của $A$ trên $SB$.
	\choiceTF 
		{\True $\mathrm{d}(A,SB)=AH$} 
		{$\mathrm{d}(B,(SAC))=\dfrac{\sqrt{3}}{3}$} 
		{\True $BC=\sqrt{3}$} 
		{Thể  tích khối chóp $S.ABC$ bằng $\dfrac{\sqrt{3}}{6}$}
	\loigiai{ 
		\begin{center}
			\begin{tikzpicture}[scale=1,font=\footnotesize,line join=round,line cap=round,>=stealth]
					\coordinate (A) at (-2,0); 
					\coordinate (B)  at (0,-3);
					\coordinate (C)  at (3,0);
					 \coordinate (S) at ($(A)+(0,4)$);
					\coordinate (H) at ($(S)!0.7!(B)$); 
					\coordinate (I) at  ($(A)!0.4!(C)$);
					\draw  (S)--(A)
						 (S)--(B)
						 (S)--(C)
						 (B)--(C)
						 (H)--(A)--(B);
					 \draw[dashed](A)--(C)
						(B)--(I); 
					 \pic[draw,thin,angle radius=2mm] {right angle = A--B--C} pic[draw,thin,angle radius=2mm] {right angle = S--A--C} pic[draw,thin,angle radius=2mm] {right angle = B--I--C}
						pic[draw,thin,angle radius=2mm]  {right angle = A--H--B};
					\foreach \i/\g in {S/90,A/180,B/-90,C/0,H/0,I/90}{\draw[fill=black](\i) circle  (1pt)
						 ($(\i)+(\g:3mm)$) node[scale=1]{$\i$};}
				\end{tikzpicture}
		\end{center}
		\begin{itemchoice}
			\itemch Vì $AH\perp SB$  nên $\mathrm{d}(A,SB)=AH$.
			\itemch Tam giác $SAB$ vuông tại $A$  có $AH$ là đường cao nên
			\begin{eqnarray*}
				 \dfrac{1}{AH^2}&=&\dfrac{1}{SA^2}+\dfrac{1}{AB^2}
				\\ 
				 AH&=&\dfrac{SA\cdot AB}{\sqrt{SA^2+AB^2}}
				 \\
				AH&=&\dfrac{2\cdot 1}{\sqrt{2^2+1^2}} 
				 \\
				 AH&=&\dfrac{2\sqrt{5}}{5}.
			\end{eqnarray*}
			Trong mặt phẳng $(ABC)$ kẻ $BI\perp AC$  tại $I$.
			\\ 
			Mặt khác $BI\perp SA$ (do $SA\perp (ABC)$, $BI\subset (ABC)$).\\ 
			Vì vậy  $BI\perp (SAC)$ hay $\mathrm{d}(B,(SAC))=BI$.
			 \\
			 Tam giác $ABI$ vuông tại $I$ có $\sin \widehat{BAC}=\dfrac{BI}{AB}\Rightarrow BI=AB\cdot \sin 60^\circ=\dfrac{\sqrt{3}}{2}$.
			 \\
			Vậy  $\mathrm{d}(B,(SAC))=BI=\dfrac{\sqrt{3}}{2}$.
			\itemch Tam giác $ABC$ vuông tại $B$  có $\tan \widehat{ACB}=\dfrac{AB}{AC}\Rightarrow BC=\dfrac{AB}{\tan 30^\circ}=\sqrt{3}$.
			\itemch Diện tích đáy hình chóp  là $S=S_{\triangle ABC}=\dfrac{1}{2}BA\cdot BC=\dfrac{\sqrt{3}}{2}$.\\
			Chiều  cao hình chóp là $h=SA=2$.
			\\ 
			Thể tích khối chóp là  $V_{S.ABC}=\dfrac{1}{3}S\cdot h=\dfrac{1}{3}\cdot \dfrac{\sqrt{3}}{2}\cdot 2=\dfrac{\sqrt{3}}{3}$ (đơn vị thể tích).
		\end{itemchoice}
	} 
\end{ex}

% nguồn: Ok_dungsai_lop11_Lan2
\begin{ex}%[1H8H5-3]%[ma:2dc4d345]
	Cho hình lăng trụ đứng $ABC.A'B'C'$ có đáy $ABC$ là tam giác đều cạnh $a$ và $A'A=2a$. Gọi $M$ là  trung điểm của $AA'$.
	\choiceTF 
		{\True Khoảng cách từ $A'$ đến mặt phẳng $(ABC)$ bằng  $2a$}
		{Khoảng  cách giữa hai đường thẳng $B'C'$ và $AB$ là $3a$}
		{\True Khoảng  cách từ $M$ đến mặt phẳng $(BCC'B')$ bằng $\dfrac{a\sqrt{3}}{2}$}
		{Khoảng cách từ $M$  đến mặt phẳng $(ACB)$ bằng $\dfrac{2a\sqrt{57}}{19}$}
	\loigiai{\begin{center}
		\begin{tikzpicture}[line join=round,line cap=round,scale=1]
				\coordinate  (A) at (0,0);
				\coordinate  (B) at (2,-2);
				\coordinate (C) at (5,0); 
				 \coordinate (v) at (0,4);
				\coordinate (A') at ($(A)  + (v)$);
				 \coordinate (B') at ($(B) + (v)$);
				\coordinate (C') at ($(C)  + (v)$);
				\coordinate (K) at ($(A)!0.5!(B)$); 
				\coordinate (H)  at ($(C)!0.5!(B)$);
				\coordinate  (N) at ($(C')!0.35!(K)$);
				\coordinate (M)  at ($(A)!0.5!(A')$);
				 \coordinate (I) at (intersection of C--M and A--C');
				 \draw[dashed] (C')--(A)--(C)--(K)--cycle (N)--(C)--(M)
					(A)--(H); 
				 \draw (A')--(A)--(B)--(C)--(C')--(B')--(B)
					(B')--(A')--(C')--(B); 
				 \foreach \p/\pos in {A/above left,B/below,B'/above,C/right,C'/above,
					 I/below,M/left,A'/above,K/below,N/left,H/right}
				\fill (\p) circle (1pt) node[\pos] {$\p$}; 
			\end{tikzpicture}
	\end{center}
	\begin{itemchoice}
		\itemch Vì  $AA' \perp (ABC)$ nên $\mathrm{d}\left(A',(ABC)\right)=AA'=2a$.
		 \itemch Ta có $B'C' \parallel (ABC)$ suy ra $\mathrm{d}\left(B'C', AB\right)=\mathrm{d}\left(B'C',(ABC)\right)=\mathrm{d}\left(B',(ABC)\right)=BB'=2a$.
		 \itemch Vì $M \in AA' \parallel (BCC'B')$ nên $\mathrm{d}\left(M,(BCC'B')\right)=\mathrm{d}\left(A,(BCC'B')\right)$.\\
		Dựng $AH \perp BC$ trong $(ABC)$.\\ Khi  đó
		$\heva{&AH \perp BC\\&AH \perp BB'} \Rightarrow AH  \perp (BCC'B')$\\
		$\Rightarrow \mathrm{d}\left(A,(BCC'B')\right)=AH=\dfrac{a\sqrt{3}}{2}$ (đường  cao trong tam giác đều cạnh $a$).
		\itemch Gọi $I=CM \cap AC'$  trong $(ACC'A')$ và $K$ là trung điểm $AB$.\\
		Dựng $CN \perp  C'K$ trong $(CKC')$.\\
		Ta có 
		$\heva{&AB \perp CK \\&AB \perp CC'} \Rightarrow AB \perp (CKC') \Rightarrow AB \perp CN \subset (CKC')$  mà $CN \perp C'K$ (gt)\\
		$\Rightarrow CN \perp  (AC'B) \Rightarrow \mathrm{d}\left(C,(AC'B)\right)=CN$.\\
		Ta có $\dfrac{\mathrm{d}\left(M,(AC'B)\right)}{\mathrm{d}\left(C,(AC'B)\right)}=\dfrac{MI}{CI}=\dfrac{MA}{CC'}=\dfrac{1}{2}$.  \\
		 $\Rightarrow \mathrm{d}\left(M,(AC'B)\right)=\dfrac{1}{2} \mathrm{d}\left(C,(AC'B)\right)=\dfrac{CN}{2}$.\\
		Xét $\triangle CC'K$ có: \[\dfrac{1}{CN^2}=\dfrac{1}{CC'^2}+\dfrac{1}{CK^2}=\dfrac{1}{(2a)^2}+\dfrac{1}{\left(\frac{a\sqrt{3}}{2}\right)^2}\Rightarrow  CN=\dfrac{2a\sqrt{57}}{19}\].
	\end{itemchoice}
 }
\end{ex}

% nguồn: Ok_dungsai_lop11_Lan2
\begin{ex}%[1H8H5-3]%[ma:c8514c95]
	Cho hình chóp $S.ABCD$ có đáy là vuông cạnh $a$. Biết $SA$ vuông góc với mặt phẳng đáy và $SA=a\sqrt{3}$.  Vẽ đường cao $AH$ của $\triangle SAB$. Vẽ đường cao $AK$ của $\triangle SAD$.
	\choiceTF 
		{\True $BC  \perp AH$}
		 {\True Khoảng cách từ $A$ đến mặt phẳng $(SBC)$ bằng $\dfrac{a\sqrt{3}}{2}$}
		{Khoảng cách từ $A$ đến mặt phẳng $(SBD)$ bằng $\dfrac{a\sqrt{2}}{7}$} 
		 {Khoảng cách từ $C$ đến mặt phẳng $(AHK)$ bằng $\dfrac{a\sqrt{5}}{5}$}
	\loigiai{\begin{center}
		\begin{tikzpicture}[scale=1.2,line join=round,line cap=round]
				\coordinate (A) at (0,0); 
				 \coordinate (B) at (4,0);
				 \coordinate (D) at (-2,-1);
				\coordinate (C)  at ($(B)-(A)+(D)$);
				\coordinate (S) at ($(A)+(0,4)$); 
				\coordinate (O) at  (intersection of A--C and B--D);
				 \coordinate (H) at ($(S)!0.4!(B)$);
				\coordinate (K) at ($(S)!0.4!(D)$); 
				\coordinate (J)  at (intersection of S--O and K--H);
				\coordinate  (F) at ($(S)!0.25!(C)$);
				\coordinate (E) at ($(S)!0.75!(O)$); 
				\draw[dashed] (S) --  (A)--(D)
					(A)--(B)  --(D)
					(C)--(A)--(K) 
					(S)--(O) 
					(H)--(K) 
					 (H)--(A)--(F)
					(A)--(E); 
				 \draw (S) -- (D) --(C) --(B)-- cycle (S)--(C)
					(H)--(F)--(K); 
				\pic [draw,angle radius=2mm,angle eccentricity=1.3] {right angle = S--A--B}; 
				\pic  [draw,angle radius=2mm,angle eccentricity=1.3] {right angle = S--A--D};
				 \pic [draw,angle radius=2mm,angle eccentricity=1.3] {right angle = A--O--B};
				 \node[below left] at (A) {$A$};
				\node[below] at  (B) {$B$};
				\node[below] at (C)  {$C$};
				 \node[above] at (S) {$S$};
				\node[left] at (K)  {$K$};
				\node[left] at (D) {$D$}; 
				\node[below]  at (O) {$O$};
				 \foreach \p/\pos in {H/right,F/above right,E/right}
				\fill  (\p) circle (1pt) node[\pos] {$\p$};
				\foreach  \P in {A,B,C,S,H,D,O,K}{
					 \fill (\P) circle (1pt);
				 }
			\end{tikzpicture}
	\end{center}
	\begin{itemchoice}
		 \itemch Ta có
		$\heva{&BC \perp SA\\ &BC \perp AB} \Rightarrow BC \perp (SAB)  \Rightarrow BC \perp AH$.\\
		\itemch 
		Ta có $\heva{&SB \perp  AH\\ &BC \perp AH} \Rightarrow AH \perp (SBC)$ hay $\mathrm{d}\left(A,(SBC)\right)=AH$.\\
		Tam giác  $SAB$ vuông tại $A$ có đường cao $AH$ nên
		\begin{align*}
			\dfrac{1}{AH^2}&=\dfrac{1}{AB^2}+\dfrac{1}{SA^2}\\ 
			\Rightarrow AH&=\dfrac{AB \cdot SA}{\sqrt{AB^2+SA^2}}=\dfrac{a \cdot a\sqrt{3}}{\sqrt{a^2+(a\sqrt{3})^2}}=\dfrac{a\sqrt{3}}{2}. 
		\end{align*}
		 Vậy $\mathrm{d}\left(A,(SBC)\right)=AH=\dfrac{a\sqrt{3}}{2}$.
		\itemch Gọi $O$  là tâm hình vuông $ABCD$ thì $AO \perp BD$\\
		Ta có $SA \perp BD$ nên $BD \perp  (SAC). \quad(*)$\\
		Kẻ đường cao $AE$ của $\triangle SAO$ thì  $AE \perp BD$ (do $(*)$).\\
		Nên  $AE \perp (SBD)$ hay $\mathrm{d}\left(A,(SBD)\right)=AE$.\\
		Ta có $AC=a\sqrt{2}$ (đường chéo hình  vuông), suy ra $OA=\dfrac{AC}{2}=\dfrac{a\sqrt{2}}{2}$.\\
		Tam giác $SAO$ vuông tại $A$ có $\dfrac{1}{AE^2}=\dfrac{1}{SA^2}+\dfrac{1}{AO^2}\Rightarrow AE=\dfrac{a\sqrt{21}}{7}$.\\ 
		 Vậy $\mathrm{d}\left(A,(SBD)\right)=AE=\dfrac{a\sqrt{21}}{7}$.
		\itemch Ta có $\heva{&AK \perp  SD\\&CD\perp AK}\Rightarrow AK\perp(SCD)$.\\
		Khi đó $\heva{&SC\perp AH\\&SC\perp AK}\Rightarrow  SC\perp (AHK)$.\\
		Trong $(SBD)$ gọi $I=KH\cap SO$.\\ 
		Trong $(SAC)$ gọi $F=AI\cap  SC$.\\
		Nên $\mathrm{d}\left(C,  (AHK)\right)=CF$.\\
		Ta có 
		\[SC=\sqrt{SA^2+AC^2}=\sqrt{\left(a\sqrt{3}\right)^2+\left(a\sqrt{2}\right)^2}=\sqrt{3a^2+2a^2}=\sqrt{5a^2}=a\sqrt{5}.\] 
		Tam giác  $SAC$ vuông tại $A$ có đường cao $AF$ nên
		\[CF=\dfrac{AC^2}{SC}=\dfrac{\left(a\sqrt{2}\right)^2}{a\sqrt{5}}=\dfrac{2a\sqrt{5}}{5}.\] 
		 Vậy $\mathrm{d}\left( C,(AHK)\right)=\dfrac{2a\sqrt{5}}{5}.$
\end{itemchoice}	}
\end{ex}

% nguồn: Ok_dungsai_lop11_Lan2
\begin{ex}%[1H8H5-3]%[ma:b8e04107]
	Cho hình chóp $S.ABCD$ có đáy $ABCD$ là hình chữ nhật với $AB=2a$, $AD=a$. Hình chiếu của $S$ lên mặt phẳng $(ABCD)$ là trung điểm $H$ của  $AB$ và $\widehat{SCH}=45^\circ$.
	\choiceTF 
		{\True$BC \perp (SAB)$} 
		{\True$\mathrm{d}\left(H,(SBC)\right)=\dfrac{a\sqrt{6}}{3}$} 
		{\True Gọi  $K$ là trung điểm $CD$ khi đó: $CD \perp (SHK)$}
		{Khoảng cách từ  $H$ đến mặt phẳng $(SCD)$ bằng $\dfrac{a\sqrt{6}}{2}$}
	\loigiai{\begin{center}
		\begin{tikzpicture}[scale=1,font=\footnotesize,line join=round,line cap=round,>=stealth]
				\coordinate (A) at  (0,0);
				\coordinate  (B) at (-2,-2);
				 \coordinate (D) at (7,0);
				\coordinate (C) at ($(B)+(D)-(A)$); 
				\coordinate  (H) at ($(A)!0.5!(B)$);
				\coordinate (K) at  ($(D)!0.5!(C)$);
				\coordinate (S) at ($(H)+(0,6)$); 
				\coordinate  (E) at ($(S)!0.6!(B)$);
				\coordinate  (I) at ($(S)!0.6!(K)$);
				\draw  (S)--(B)
					(S)--(C) 
					(S)--(D) 
					(B)--(C)--(D); 
				\draw[dashed](S)--(A) 
					 (B)--(A)--(D)
					 (E)--(H)--(I)
					(H)--(K) 
					(S)--(H)--(C) 
					 (D)--(K)--(S);
				\pic[draw,thin,angle radius=3mm] {right angle = H--I--K}; 
				\pic[draw,thin,angle  radius=3mm] {right angle = D--K--H};
				\pic[draw,thin,angle radius=3mm] {right angle =  S--E--H};
				\foreach  \i/\g in {S/90,A/45,B/180,C/-90,D/-90,H/180,I/180,K/0,E/180}{\draw[fill=black](\i) circle (1pt)
					($(\i)+(\g:3mm)$)  node[scale=1]{$\i$};}
			\end{tikzpicture}
	\end{center}
	\begin{itemchoice}
		\itemch Kẻ đường cao $HE$ trong $\triangle  SBH$. $\quad(1)$\\
		Ta có 
		$\heva{&BC \perp AB \\&BC \perp SH} \Rightarrow  BC \perp (SAB).$
		\itemch Ta có $BC \perp  (SAB) \Rightarrow BC \perp HE$. $\quad(2)$\\
		Từ $(1)$ và $(2)$ suy ra $HE \perp (SBC)$ hay  $\mathrm{d}\left(H,(SBC)\right)=HE$.\\
		Xét $\triangle BCH$ vuông tại $B$, ta có\\ $BH=\dfrac{AB}{2}=\dfrac{2a}{2}=a$.\\ 
		 $CH=\sqrt{BC^2+BH^2}=\sqrt{a^2+a^2}=a\sqrt{2}$.\\
		Xét  $\triangle SCH$ vuông tại $H$, ta có\\
		 $\tan \widehat{SCH}=\dfrac{SH}{CH} \Rightarrow SH=CH \tan 45^\circ=a\sqrt{2} \cdot 1=a\sqrt{2}$.\\
		Xét $\triangle SBH$ vuông tại $H$ có  đường cao $HE$, ta có\\
		$\dfrac{1}{HE^2}=\dfrac{1}{SH^2}+\dfrac{1}{BH^2}=\dfrac{1}{(a\sqrt{2})^2}+\dfrac{1}{a^2} 
		\Rightarrow  HE=\dfrac{a\sqrt{6}}{3}$.\\
		 Vậy $\mathrm{d}\left(H,(SBC)\right)=HE=\dfrac{a\sqrt{6}}{3}$.
		\itemch  Gọi $K$ là trung điểm $CD$.\\
		Khi đó $HK$ là đường trung bình của hình chữ nhật $ABCD$ nên  $HK \parallel BC \parallel AD \Rightarrow HK \perp CD$.\\
		Ta có 
		$\heva{&CD \perp HK \\&CD \perp SH} \Rightarrow CD \perp (SHK)$. 
		\itemch Kẻ đường cao $HI$  của $\triangle SHK$.\\
		 Ta có
		$\heva{&HI \perp  SK \\&HI \perp CD} \Rightarrow HI \perp (SCD)$.\\
		Xét $\triangle SHK$ vuông  tại $H$ có đường cao $HI$, ta có
		 \[dfrac{1}{HI^2}=\dfrac{1}{SH^2}+\dfrac{1}{HK^2}
		 \Rightarrow HI=\dfrac{a\sqrt{6}}{3}.\]
\end{itemchoice}}
\end{ex}

% nguồn: Ok_dungsai_lop11_Lan1
\begin{ex}%[1H8V5-3]%[ma:f9dd6eac]
 Cho hình chóp $S. A B C D$  có đáy là vuông cạnh $a$. Biết $S A$ vuông góc với mặt phẳng đáy và $S A=a \sqrt{3}$. Vẽ đường cao $A H$ của tam giác $S A B$.Vẽ đường cao $A K$ của tam giác $S A D$. Khi đó:
  \choiceTF
   {\True $BC\perp AH$}
  {\True Khoảng cách từ  $A$ đến mặt phẳng $(S B C)$ bằng: $\dfrac{a\sqrt{3}}{2}$}
  {Khoảng cách từ $A$  đến mặt phẳng $(S B D)$ bằng: $\dfrac{a\sqrt{2}}{7}$}
  {Khoảng cách từ $C$ đến mặt  phẳng $(A H K)$ bằng: $\dfrac{a\sqrt{5}}{5}$}
  \loigiai{
  \begin{center}
   \begin{tikzpicture}[scale=1,font=\footnotesize,line join=round,line cap=round,>=stealth]
     \path 
      (0,0) coordinate (A) 
      ++(-140:2)  coordinate (B)
      ++(0:4)  coordinate (C)
      ($(A)+(C)-(B)$) coordinate  (D)
      ($(A)+(0,4)$) coordinate  (S)
       ($(A)!.5!(C)$) coordinate (O)
      ($(S)!(A)!(O)$) coordinate (E) 
       ($(S)!(A)!(B)$) coordinate (H)
      ($(S)!(A)!(D)$) coordinate (K) 
      (intersection of S--O and H--K)  coordinate (I)
      ($(A)!2!(I)$) coordinate  (J)
      (intersection of A--J and S--C) coordinate  (F);;
      \foreach \i in{B,C,D}{\draw (S)--(\i);};
      \draw (B)--(C)--(D)
      (H)--(F)--(K); 
     \draw[dashed] (S)--(A)--(B) 
      (A)--(D) 
      (A)--(C) 
      (D)--(B) 
       (K)--(A)--(H)--(K)
      (S)--(O) 
      (E)--(A)--(F); 
     \pic[draw,angle eccentricity=1.8,angle radius=2mm]{right  angle=S--A--D};
     \pic[draw,angle eccentricity=1.8,angle radius=2mm]{right  angle=A--E--O};
     \pic[draw,angle eccentricity=1.8,angle  radius=2mm]{right angle=A--O--B};
     \foreach \i/\g in {A/-90,B/-90,C/-90,D/-90,S/90,H/180,K/50,F/60,E/0} 
     \fill[black]  (\i) circle (1pt)+(\g:3mm) node[scale=1]{$\i$};
    \end{tikzpicture}
  \end{center}
  \begin{itemchoice}
   \itemch Ta có $\heva{&B C \perp S A \text{  (do } S A \perp(A B C D)) \\& B C \perp A B} \Rightarrow B C \perp(S A B) \Rightarrow B C \perp A H$.
   \itemch Ta có  $B C \perp A H$
   mà $S B \perp A H$ nên $A H \perp(S B C)$ hay $\mathrm{d}(A,(S B C))=A  H$.\\
   Tam giác $S A B$ vuông tại $A$ có đường cao  $A H$ nên
   $$\dfrac{1}{A H^2}=\dfrac{1}{A B^2}+\dfrac{1}{S A^2}  \Rightarrow A H=\dfrac{A B \cdot S A}{\sqrt{A B^2+S A^2}}=\dfrac{a \cdot a \sqrt{3}}{\sqrt{a^2+3 a^2}}=\dfrac{a \sqrt{3}}{2}.$$
    Vậy $\mathrm{d}(A,(S B C))=A H=\dfrac{a \sqrt{3}}{2}$.
   \itemch  Gọi $O$ là tâm hình vuông $A B C D$ thì $A O \perp B D$, ta lại có $S A \perp B D$ nên $B D \perp(S A C)$.\tagEX{*}
   \noindent Kẻ đường cao $A E$ của $\Delta S A O$ thì $A E \perp B D$  (do $(*))$.\\
   Vậy $A E  \perp(S B D)$ hay $\mathrm{d}(A,(S B D))=A E$.\\
   Ta có $A C=a \sqrt{2}$ (đường  chéo hình vuông), suy ra $O A=\dfrac{A C}{2}=\dfrac{a \sqrt{2}}{2}$.\\
   Tam giác $S A  O$ vuông tại $A$ có $A E=\dfrac{S A \cdot A O}{\sqrt{S A^2+A O^2}}=\dfrac{a \sqrt{3} \cdot \dfrac{a \sqrt{2}}{2}}{\sqrt{3 a^2+\dfrac{2 a^2}{4}}}=\dfrac{a \sqrt{21}}{7}$.\\
   Vậy $\mathrm{d}(A,(S B D))={AE}=\dfrac{a \sqrt{21}}{7}$. 
   \itemch Ta chứng minh được $A K \perp(S C D)$. Khi đó: $\heva{&S C \perp  A H \\ &S C \perp A K} \Rightarrow S C \perp(A H K)$.\\
   Gọi $F=S C \cap(A  H K)$ thì $S C \perp A F$.
   Khi đó: $\mathrm{d}(C,(A  H K))=C F$.\\
   Ta có $S C=\sqrt{S A^2+A C^2}=\sqrt{3  a^2+2 a^2}=a \sqrt{5}$.\\
   Tam giác $S  A C$ vuông tại $A$ có đường cao $A F$ nên:
   $$CF\cdot  CS=AC^2\Rightarrow CF=\dfrac{AC^2}{CS}=\dfrac{2a^2}{a\sqrt{5}}=\dfrac{2a\sqrt{5}}{5}.$$
   Vậy  $\mathrm{d}(C,(AHK))=CF=\dfrac{2a\sqrt{5}}{5}$.
  \end{itemchoice}
 } 
\end{ex}

% nguồn: Ok_dungsai_lop11_Lan1
\begin{ex}%[1H8V5-3]%[ma:5028b1ab]
 Cho hình chóp $S.ABCD$ có đáy $ABCD$ là hình chữ nhật với $AB=2a$, $AD=a$. Hình chiếu của $S$ lên mặt phẳng $(ABCD)$ là trung điểm $H$ của $AB$, kẻ đường cao $HE$ trong tam giác $SBH$. Gọi $K$ là trung điểm $CD$, biết rằng $\widehat{SCH}=45^{\circ}$. Khi đó,  hãy xét tính đúng sai của các phát biểu sau.
 \choiceTF 
   {$(BC,(SAB))=60^{\circ}$}
  {\True  $CD\perp (SHK)$}
  {\True $\mathrm{d}(H,(SBC))=HE$} 
   {$\mathrm{d}(H,(SCD))=\dfrac{a\sqrt{6}}{2}$}
  \loigiai{
  \begin{center}
   \begin{tikzpicture}[scale=1,font=\footnotesize,line join=round,line cap=round,>=stealth]
     \path 
      (0,0) coordinate  (A)
      (-1.4,-1.6) coordinate  (B)
      (2.5,-1.6) coordinate  (C)
      ($(A)+(C)-(B)$) coordinate (D) 
      ($(A)!0.5!(B)$) coordinate  (H)
      ($(H)+(0,3.5)$) coordinate  (S)
      ($(C)!0.5!(D)$)  coordinate (K)
      ($(S)!0.65!(K)$) coordinate (I) 
      ($(S)!0.75!(B)$) coordinate  (E);
     \draw (S)--(B)--(C)--(D)--cycle (K)--(S)--(C); 
      \draw[dashed] (K)--(H)--(S)--(A)--(D)
      (A)--(B) 
      (I)--(H)--(E) 
      (H)--(C); 
      \foreach \p/\q in {S/90,A/170,B/-90,C/-90,D/0,H/-90,K/-20,I/50,E/150}
     \fill[black] (\p)  circle (1.0pt) node[shift={(\q:3mm)}]{$\p$};
     \draw pic[draw=black,angle  radius=0.2cm] {right angle = S--H--A}
      pic[draw=black,angle radius=0.2cm]  {right angle = H--I--S}
      pic[draw=black,angle radius=0.2cm]  {right angle = H--E--S};
    \end{tikzpicture}
  \end{center}
  \begin{itemchoice}
    \itemch \textbf{Sai}.\\
   Ta có $\heva{&BC\perp  AB\\&BC\perp SH(\text{do } SH\perp (ABCD))}\Rightarrow BC\perp (SAB)$, suy ra $(BC,(SAB))=90^{\circ}$.
    \itemch \textbf{Đúng}.\\
   Ta có $HK\parallel BC\parallel AD\Rightarrow HK\perp CD$.\\ 
   Khi đó $\heva{&CD\perp HK\\&CD\perp SH}\Rightarrow  CD\perp (SHK)$.
   \itemch \textbf{Đúng}.\\ 
    Ta có $BC\perp(SAB)\Rightarrow BC\perp HE$\quad(1).\\
   Do $HE$  la đường cao trong tam giác $SBH$ nên $HE\perp SB$ \quad(2).\\
   Từ(1) và (2) suy ra $HE\perp (SBC)$ hay $\mathrm{d}(H,(SBC))=HE$. 
   \itemch \textbf{Sai}.\\ 
   Kẻ đường cao $HI$ của tam  giác $SHK$.\\
   Ta có $\heva{&HI\perp  SK\\&HI\perp CD(\text{do }CD\perp (SHK),HI\subset (SHK))}\Rightarrow HI\perp (SCD)$.\\
   Tam giác  $BCH$ vuông tại $B$ có
    $CH=\sqrt{BC^2+BH^2}=\sqrt{a^2+a^2}=a\sqrt{2}$.\\
   Tam giác $SCH$ vuông  tại $H$ có
   $\tan\widehat{SCH}=\dfrac{SH}{CH}\Rightarrow  SH=a\sqrt{2}$.\\
   Tam giác $SHK$ vuông tại $H$ có đường  cao $HI$ nên.\\
    $\dfrac{1}{HI^2}=\dfrac{1}{SH^2}+\dfrac{1}{HK^2}\Rightarrow HI=\dfrac{SH\cdot HK}{\sqrt{SH^2+HK^2}}=\dfrac{a\sqrt{2}\cdot a}{\sqrt{2a^2+a^2}}=\dfrac{a\sqrt{6}}{3}$.\\
   Vậy $\mathrm{d}(H,(SCD))=HI=\dfrac{a\sqrt{6}}{3}$. 
 \end{itemchoice}}
\end{ex}

% nguồn: Ok_dungsai_lop11_Lan2
\begin{ex}%[1H8V5-3]%[ma:8e7e52d7]
	Cho hình chóp $S . A B C D$ có đáy $A B C D$ là hình vuông cạnh $a$. Cạnh bên $S A=2 a$ và $S A$ vuông góc với đáy. Gọi $M$ là trung điểm của  $C D$.
	 \choiceTF
		{\True Khoàng cách từ $D$ đến mặt phẳng $(S A  B)$ là độ dài đoạn $D A$}
		{\True Khoảng cách giữa hai đường thẳng $S  B$ và $C D$ là độ dài đoạn $D A$}
		{\True Khoảng cách từ $M$  đến mặt phẳng $(S B D)$ là $\dfrac{a}{3}$}
		{Khoảng cách từ $C$ đến mặt phẳng  $(S B D)$ là $\dfrac{a}{3}$}
	 \loigiai{
		\begin{center}
			\begin{tikzpicture}[scale=0.8,font=\footnotesize,line join=round,line cap=round,>=stealth]
					\coordinate (A)  at (0,0);
					\coordinate  (B) at (-2,-2);
					\coordinate (D) at (7,0); 
					\coordinate  (C) at ($(B)+(D)-(A)$);
					 \coordinate (S) at ($(A)+(0,5)$);
					 \coordinate (O) at ($(A)!0.5!(C)$);
					\coordinate (M) at  ($(D)!0.5!(C)$);
					\coordinate (H) at  ($(S)!0.75!(O)$);
					\draw  (S)--(B)
						 (S)--(C)
						(S)--(D) 
						 (B)--(C)--(D);
					 \draw[dashed](S)--(A)
						(B)--(A)--(D) 
						 (S)--(O)
						 (A)--(C)
						(B)--(D) 
						 (A)--(H);
					\foreach \i/\g in  {S/90,A/180,B/-90,C/-90,D/0,O/-90,H/0,M/-45}{\draw[fill=black](\i) circle (1pt)
						 ($(\i)+(\g:3mm)$) node[scale=1]{$\i$};}
				\end{tikzpicture}
		\end{center}
		\begin{itemchoice}
			\itemch $\heva{&DA\perp AB\\&DA\perp SA}\Rightarrow DA\perp (SAB)\Rightarrow \mathrm{d}(A,(SAB))=DA$. 
			\itemch $C D\parallel A B \Rightarrow C D \parallel(SAB) \Rightarrow \mathrm{d}(S B, C D)=\mathrm{d}(C  D,(S A B))=\mathrm{d}(D,(S A B))=D A$.
			\itemch Vì $M$ là trung  điểm của $C D$.\\
			$\mathrm{d}(M,(S B D))=\dfrac{1}{2} \mathrm{d}(C,(S  B D))=\dfrac{1}{2} \mathrm{d}(A,(SBD))=\dfrac{1}{2} A H=\dfrac{a}{3}$.
			\itemch Gọi  $O$ là tâm đáy.\\
			 Vì
			$A C  \cap(S B D)=O$, $O$ là trung điểm $A C$ nên $\mathrm{d}(C,(S B D))=\mathrm{d}(A,(S B D))$.\\
			Gọi $H$  là hình chiếu vuông góc của $A$ lên $S O$, khi đó $A H \perp S O$. \hfill $(1)$\\
			Mà $A H \subset(S A C)$.\\ 
			$\Rightarrow A H \perp B D$. \hfill $(2)$\\ 
			Mà $S  O$ cắt $B D$ trong mặt phẳng $(S B D)$. \hfill$(3)$\\
			 Từ $(1)$, $(2)$ và $(3)$ suy ra $A H \perp(S B D) \Rightarrow \mathrm{d}(A;(S B D))=A H$.\\
			Áp dụng hệ thức lượng trong tam giác vuông $S O A$ ta có \[\dfrac{1}{A H^{2}}=\dfrac{1}{A S^{2}}+\dfrac{1}{A  O^{2}}=\dfrac{1}{4 a^{2}}+\dfrac{4}{2 a^{2}}=\dfrac{9}{4 a^{2}}.\]
			$\Rightarrow  A H=\dfrac{2 a}{3}$.
		\end{itemchoice}
	 }
\end{ex}


%%%%%%%%%%%%%%%% TRẢ LỜI NGẮN %%%%%%%%%%%%%%%%

% nguồn: Ok_TraLoiNgan_Lop11_Lan1
\begin{ex}%[1H8H5-3]%[ma:0bb2899a]
	 Cho hình hộp chữ nhật $ABCD.A'B'C'D'$ có ba kích thước $AB=7$, $AD=14$, $AA'=21$. Khoảnh cách từ $A$ đến mặt phẳng $(A'BD)$ bằng bao nhiêu?
	\shortans{6} 
	 \loigiai{
		\begin{center}
			\begin{tikzpicture}
					 \def\bc{4}
					 \def\ba{2}
					 \def\gocB{35}
					\coordinate[label=below  left:$B$] (B) at (0,0);
					\coordinate[label=above left:$A$] (A) at (\gocB:\ba); 
					\coordinate[label=below:$C$]  (C) at (\bc,0);
					\coordinate[label=right:$D$] (D) at ($(C)-(B)+(A)$); 
					\coordinate[label=above left:$A'$] (A') at ($(A)+(90:\bc)$); 
					 \coordinate[label=left:$B'$] (B') at ($(B)-(A)+(A')$);
					\coordinate[label=below right:$C'$]  (C') at ($(C)-(A)+(A')$);
					\coordinate[label=right:$D'$] (D')  at ($(D)-(A)+(A')$);
					\coordinate[label=below right:$M$] (M) at  ($(B)!0.5!(D)$);
					\coordinate[label=right:$K$] (K)  at ($(A')!0.7!(M)$);
					 \draw (B')--(B)--(C)--(D)--(D')--(A')--(B')--(C')--(D')
						(C)--(C'); 
					\draw[dashed] (A')--(A)--(D) 
						 (A)--(B)--(D)
						 (K)--(A)--(M)--(A');
					\draw pic[draw,angle radius=2mm]{right angle=A--K--M}; 
					\foreach \diem in  {A,B,C,D,A',B',C',D',M,K} \fill (\diem) circle (1pt);
				\end{tikzpicture}
		\end{center}
		Đặt $AB=a$, $AD=2a$,  $AA'=3a$.\\
		Trong tam giác $ABD$ kẻ $AM\perp BD$  suy ra $\dfrac{1}{AM^2}=\dfrac{1}{AB^2}+\dfrac{1}{AD^2}$.\\
		Trong tam giác $A'AM$ kẻ  $AK\perp A'M$ $\Rightarrow \dfrac{1}{AK^2}=\dfrac{1}{A'A^2}+\dfrac{1}{AM^2}=\dfrac{1}{a^2}+\dfrac{1}{4a^2}+\dfrac{1}{9a^2}$.\\
		 $\Rightarrow AK=\dfrac{6}{7}a$.\\
		Khi đó $AK\perp \left(A'BD\right)$ hay $\mathrm{d}\left( A,\left(A'BD\right)\right)=AK=\dfrac{6}{7}a=6$  với $a=7$.
	} 
\end{ex}

% nguồn: Ok_TraLoiNgan_Lop11_Lan1
\begin{ex}%[1H8H5-3]%[ma:0ae8654b]
	Cho hình chóp $S.ABCD$ có đáy $ABCD$ là hình vuông cạnh $\sqrt{5}$. Biết $(SAB)$ vuông góc với mặt phẳng đáy, $SA=SB$, góc giữa đường thẳng $SC$ và mặt phẳng đáy  bằng $45^\circ$. Tính khoảng cách từ điểm $S$ đến $(ABCD)$.
	 \shortans{2,5}
	\loigiai{ 
	 \immini{
		Gọi $M$ là trung điểm  của $AB$, do $\triangle SAB$ cân nên $SM \perp AB$.\\
		Ta có $\heva{&(SAB)\perp (ABCD)\\&(SAB)\cap (ABCD)=AB\\&SM\perp AB}\Rightarrow SM\perp (ABCD)$.\\ 
		$\triangle BCM$ vuông tại $B$  có $MC=\sqrt{BM^2 + BC^2}=\dfrac{5}{2}$.\\
		$MC$ là hình chiếu vuông góc của $SC$ lên $(ABCD)$ nên  góc giữa $SC$ và $(ABCD)$ là $\widehat{SCM}=45^\circ$.
		 }{
		\begin{tikzpicture}[font=\footnotesize,line join=round,line cap=round,>=stealth,scale=0.8]
				\foreach \x/\y/\pos in {0/0/A,-1.5/-1.7/B,4/0/D} \path ($(\x,\y)$)  coordinate (\pos);
				\path ($(B)+(D)$) coordinate (C) 
					($(B)!1/2!(D)$) coordinate  (O)
					($(A)!.5!(B)$) coordinate  (M)
					($(M)+(0,4)$)  coordinate (S);
				 \draw[dashed] (B)--(A)--(D)--(B)
					 (S)--(A)--(C)--(M)
					(S)--(M); 
				 \draw (C)--(S)--(B)--(C)--(D)--(S);
				\foreach \x/\pos in {S/90,A/180,B/-150,C/-30,D/-30,M/180} \fill (\x) circle (1pt) node[{shift=(\pos:0.25)}]{$\x$}; 
				\draw pic[draw,angle radius=2mm]  {right angle=S--M--A};
			\end{tikzpicture}
		} 
		\noindent 
		Do đó, $\triangle SMC$ vuông cân  tại $M$. Suy ra $MC=SM=\dfrac{5}{2}$.\\
		Vậy $\mathrm{d}(S,(ABCD))=SM=\dfrac{5}{2}=2{,}5$. 
	} 
\end{ex}

% nguồn: Ok_TraLoiNgan_Lop11_Lan1
\begin{ex}%[1H8H5-3]%[ma:715ff822]
	Cho hình chóp $S.ABCD$ có đáy là hình vuông cạnh bằng $2$, cạnh bên $SA$ vuông góc với đáy và $SA=2\sqrt{3}$. Tính khoảng cách từ điểm $B$ đến mặt phẳng $(SCD)$ (làm tròn kết  quả đến $2$ chữ số sau dấu phẩy).
	\shortans{1,73} 
	 \loigiai{
	 \immini{
		Vì $AB$ song song với $(SCD)$ nên khoảng cách từ $B$ đến $(SCD)$ bằng khoảng cách từ $A$ đến $(SCD)$. Vẽ $AH$ vuông góc với $SD$ tại $H$. Ta có $CD\perp AD$ và $CD\perp SA$, suy ra $CD\perp AH$. Mặt khác $AH\perp SD$ nên suy ra $AH\perp (SCD)$. Do đó khoảng cách từ $A$ đến  $(SCD)$ bằng độ dài $AH$. Ta có $$SD=\sqrt{SA^2+AD^2}=\sqrt{12+4}=4.$$ Từ hệ thức lượng $AH\cdot SD=SA\cdot AD$ ta tính được $$AH=\dfrac{SA\cdot AD}{SD}=\sqrt{3}.$$ Vậy khoảng cách cần tính bằng $\sqrt{3}\approx 1{,}73$.
		}{\begin{tikzpicture}[scale=1,font=\footnotesize,line join=round,line cap=round,>=stealth]
		\path (0,0) coordinate (B)++(3,0) coordinate  (C)++(1.3,1) coordinate (D)
			 ($(B)+(D)-(C)$) coordinate (A)
			 (0,3) coordinate (h)
			($(A)+(h)$) coordinate  (S)
			 ($(S)!0.6!(D)$) coordinate (H);
		 \draw (S)--(B)--(C)--(D)--cycle (S)--(C);
		\draw[dashed] (S)--(A)--(B) 
			 (H)--(A)--(D);
		\foreach \p/\g  in {A/150,B/240,C/-60,D/0,S/90,H/45} \fill[black] (\p) circle (1pt)+(\g:0.3) node{$\p$};
		\foreach \x/\y/\z  in {S/A/B,S/A/D,S/H/A,A/B/C,B/C/D,C/D/A,D/A/B}\draw pic[draw,angle radius=0.2cm]{right angle=\x--\y--\z};
	\end{tikzpicture}
		} 
	 }
\end{ex}

% nguồn: Ok_TraLoiNgan_Lop11_Lan1
\begin{ex}%[1H8H5-3]%[ma:0a6e5618]
	Cho khối lăng trụ đều $ABC.A'B'C'$ có cạnh đáy bằng $\sqrt{5}$, $A'A=2\sqrt{5}$. Khoảng cách từ  điểm $A'$ đến mặt phẳng $(AB'C')$ bằng
	 \shortans{2}
	 \loigiai{
		\begin{center}
			\begin{tikzpicture}
					 \def\a{4}
					\def\h{4.5} 
					\path (0:0) coordinate  (A)
						++(0:\a) coordinate  (C)
						++(-150:3*\a/4) coordinate (B) 
						($(A)+(90:\h)$)  coordinate (A')
						($(B)+(90:\h)$) coordinate (B') 
						 ($(C)!0.5!(B)$) coordinate (M)
						($(A')!0.4!(M)$) coordinate (H) 
						($(C)+(90:\h)$) coordinate (C'); 
					 \draw[dashed] (H)--(A)--(C)--(A')--(M)--(A);
					\draw (C)--(C') 
						 (B)--(B')
						 (A)--(A')
						(A)--(B)--(C) 
						 (A')--(B)
						(A')--(B')--(C')--cycle; 
					\draw pic[draw,angle radius=2mm]{right  angle=A--H--M};
					\foreach \x/\g in  {A/180,B/-45,C/0,A'/180,B'/90,C'/0,M/-20}
					\fill[black] (\x)  circle (1pt)
						($(\g:4mm)+(\x)$) node {$\x$}; 
				\end{tikzpicture}
		\end{center}
		Gọi cạnh đáy là $a \Rightarrow AA'=2a$.\\ 
		Gọi $M$ là trung điểm  của $B'C'$.\\
		Ta có $\heva{&AA'\perp B'C'\\&A'M\perp B'C'}\Rightarrow B'C'\perp  \left(AA'M\right)\Rightarrow\left(AB'C'\right) \perp \left(AA'M\right)$
		theo giao  tuyến $AM$.\\
		Kẻ $A'H\perp AM$ trong mặt phẳng $\left(AA'M\right) \Rightarrow A'H\perp  \left(AB'C'\right)$.\\
		Vậy khoảng cách  từ $A'$ đến mặt phẳng $\left(AA'M\right)$ là $A'H$.\\
		Ta  có $\dfrac{1}{A'H^2}=\dfrac{1}{A'A^2}+\dfrac{1}{A'M^2}=\dfrac{1}{4a^2}+\dfrac{1}{a^2}\Rightarrow A'H=\dfrac{2a\sqrt{5}}{5}=2$ với $a=\sqrt{5}$.
	} 
\end{ex}

% nguồn: Ok_TraLoiNgan_Lop11_Lan2
\begin{ex}%[1H8H5-3]%[ma:f69d2058]
	Cho hình chóp $S.ABC$ có đáy là tam giác vuông tại $A$, $AB=1$, $AC=\sqrt{3}$, $SA$ vuông góc  với mặt phẳng đáy và $SA=2$. Khoảng cách từ điểm $A$ đến mặt phẳng $(SBC)$ bằng bao nhiêu? (kết quả làm tròn đến hàng phần trăm)
	 \shortans{0,79}
	 \loigiai{
		\begin{center}
			\begin{tikzpicture}[scale=0.8,font=\footnotesize,line join=round,line cap=round,>=stealth]
					 \coordinate (A) at (0,0);
					\coordinate (S)  at (0,3.5);
					\coordinate (B) at (2,-1); 
					\coordinate (C) at  (5,0.2);
					 \coordinate (I) at ($(B)!0.4!(C)$);
					 \coordinate (H) at ($(S)!0.4!(I)$);
					\draw (S) -- (A) -- (B) --  (C) -- (S)
						(I)--(S)  -- (B);
					 \draw[dashed] (I)--(A) -- (C)
						(A)--(H); 
					\pic[draw,thin,angle radius=2mm] {right  angle = S--A--C} pic[draw,thin,angle radius=2mm] {right angle = B--A--C} pic[draw,thin,angle radius=2mm] {right angle = A--H--I} pic[draw,thin,angle radius=2mm] {right angle = A--I--B};
					\foreach \i/\g in {S/90,A/-90,B/-90,C/-90,H/60,I/-60}{\draw[fill](\i) circle (1pt) 
						($(\i)+(\g:3mm)$)  node[scale=1]{$\i$};}
				\end{tikzpicture}
		\end{center}
		Gọi $I$ và $H$ lần lượt  là hình chiếu của $A$ trên $BC$ và $SI$.\\
		Suy  ra $\mathrm{d}(A,(SBC))=AH$.\\
		Ta có $\dfrac{1}{AH^2}=\dfrac{1}{SA^2}+\dfrac{1}{AI^2}=\dfrac{1}{SA^2}+\dfrac{1}{AB^2}+\dfrac{1}{AC^2}=\dfrac{1}{4}+1+\dfrac{1}{3}=\dfrac{19}{12}$.\\ 
		Suy ra  $AH=\dfrac{2\sqrt{57}}{19}$.\\
		Vậy $\mathrm{d}(A;(SBC)) \approx  0{,}79$.
	} 
\end{ex}

% nguồn: Ok_TraLoiNgan_Lop11_Lan2
\begin{ex}%[1H8H5-3]%[ma:c33192e1]
	Cho hình chóp $S.ABCD$ có $SA \perp (ABCD)$, đáy $ABCD$ là hình chữ nhật. Biết $AD=2$, $SA=1$. Khoảng cách từ $A$ đến $(SCD)$  bằng (kết quả làm tròn đến hàng phần trăm)
	\shortans{0,89} 
	 \loigiai{
		\begin{center}
			\begin{tikzpicture}[scale=1,font=\footnotesize,line join=round,line cap=round,>=stealth]
					\coordinate  (A) at (0,0);
					\coordinate (B) at  (-2,-2);
					\coordinate (D) at  (4,0);
					\coordinate (C)  at ($(B)+(D)-(A)$);
					\coordinate (S) at (0,3.5); 
					\coordinate (H) at  ($(S)!0.6!(D)$);
					\draw[dashed] (S)  -- (A) -- (B)
						(H)  -- (A) -- (D);
					\draw  (S)--(B)--(C)--(D)--(S)--(C);
					\pic[draw,thin,angle radius=2mm] {right angle =  D--A--B} pic[draw,thin,angle radius=2mm] {right angle = C--B--A} pic[draw,thin,angle radius=2mm] {right angle = S--H--A};
					\foreach \i/\g in {S/90,A/180,B/-90,C/-90,D/0,H/30}{\draw[fill=black](\i)  circle (1pt)
						($(\i)+(\g:3mm)$) node[scale=1]{$\i$};} 
				\end{tikzpicture}
		\end{center}
		Gọi $H$ là  hình chiếu của $A$ lên $SD$.\\
		Ta  chứng minh được $AH \perp (SCD)$.\\
		 Xét $\triangle SAD$ vuông tại $A$ có $\dfrac{1}{AH^2} = \dfrac{1}{SA^2}+\dfrac{1}{AD^2} \Rightarrow AH=\dfrac{2}{\sqrt{5}} \approx 0{,}89$.\\
		Vậy $\mathrm{d}(A,(SCD)) = 0{,}89$. 
	} 
\end{ex}

% nguồn: Ok_TraLoiNgan_Lop11_Lan2
\begin{ex}%[1H8H5-3]%[ma:73626c8c]
	Cho hình chóp tứ giác đều $S.ABCD$ có độ dài cạnh đáy bằng $2$ và độ dài cạnh bên bằng $3$ (tham khảo hình bên). Khoảng cách từ $S$ đến mặt phẳng $\left(ABCD\right)$ bằng (kết quả làm tròn  đến hàng phần trăm)
	\begin{center}
		\begin{tikzpicture}[scale=0.75,font=\footnotesize,line join=round,line cap=round,>=stealth]
				\coordinate  (A) at (0,0);
				 \coordinate (B) at (-2,-1);
				\coordinate (D) at  (4,0);
				 \coordinate (C) at ($(B)+(D)-(A)$);
				\coordinate (O) at ($(A)!0.5!(C)$); 
				\coordinate (S)  at ($(O)+(0,4)$);
				\draw (S)--(B) 
					 (S)--(C)
					 (S)--(D)
					(B)--(C)--(D); 
				 \draw[dashed](S)--(O)
					(S)--(A)--(C) 
					 (A)--(B)
					(A)--(D)--(B); 
				\foreach \i/\g in {S/90,A/150,B/-90,C/-90,D/0,O/-90}{\draw[fill](\i)  circle (1pt)
					($(\i)+(\g:3mm)$) node[scale=1]{$\i$};} 
			\end{tikzpicture}
	\end{center}
	\shortans{2,56} 
	\loigiai{ 
		Gọi $O$ là giao điểm của $AC$ và $BD \Rightarrow SO \perp  (ABCD)$. \\
		$\Rightarrow \mathrm{d}(S,(ABCD)) =  SO$. \\
		Ta  có $AC = 2\sqrt{2} \Rightarrow OC=\sqrt{2}$. \\
		 Xét tam giác $SOC$ vuông tại $O$ có $SO = \sqrt{SC^2-OC^2} = \sqrt{3^2-2} = \sqrt{7}$. \\
		 Vậy $\mathrm{d}(S,(ABCD)) = \sqrt{7} \approx 2{,}65$.
	 }
\end{ex}

% nguồn: Ok_TraLoiNgan_Lop11_Lan2
\begin{ex}%[1H8H5-3]%[ma:6638e747]
	Cho hình  chóp $S.ABCD$ có đáy $ABCD$ là hình vuông cạnh $1$, $SA\perp \left(ABCD\right)$. Tính khoảng cách từ điểm $B$ đến $\mathrm{mp}\left(SAC\right)$. (kết quả làm tròn đến hàng phần trăm).
	\shortans{0,71} 
	\loigiai{ 
	 \immini{
		 Gọi $O=AC\cap BD$.\\
		Vì $ABCD$  là hình vuông nên $BO\perp AC$.\hfill (1)\\
		Ta lại  có $SA\perp \left(ABCD\right)\Rightarrow SA\perp BO$.\hfill (2)\\
		Từ $\left(1\right)$ và $\left(2\right)\Rightarrow BO\perp  (SAC)$.\\
		Mà $BO=\dfrac{1}{2}\cdot BD=\dfrac{\sqrt{2}}{2}$.\\ 
		 Vậy $\mathrm{d}_{\left(B,\left(SAC\right)\right)}=BO=\dfrac{\sqrt{2}}{2}\approx 0{,}71$.
		}{		\begin{tikzpicture}[declare function={r=2.5; gocD=-140;},font=\footnotesize]
		\path (0,0)  coordinate (A)
			(90:r) coordinate (S) 
			(0:1.4*r)  coordinate (B)
			 (gocD:0.6*r) coordinate (D)
			($(B)+(D)-(A)$) coordinate  (C)
			($(B)!0.5!(D)$) coordinate  (O);
		 \foreach \x/\y/\z in {S/A/B,S/A/D,B/A/D,B/O/C}{
			 \path pic[draw=red,angle radius=5pt]{right angle= \x--\y--\z};
		} 
		\draw  (S)--(B)--(C)--(D)--(S)--(C);
		 \draw[dashed] (S)--(A)--(C)
			(B)--(A)--(D)--(B); 
		\fill[ball color = red] (A) circle (1pt) node[shift={(160:9pt)}]{$A$}; 
		\foreach  \t/\g in {B/0,C/-60,S/90,D/-140,O/-90}{
			\fill[ball color = red] (\t)  circle (1pt) node[shift={(\g:7pt)}]{$ \t $};
		} 
	\end{tikzpicture}
		} 
	} 
\end{ex}

% nguồn: Ok_TraLoiNgan_Lop11_Lan2
\begin{ex}%[1H8H5-3]%[ma:8bf341de]
	Cho hình chóp $S.ABC$ có $SA\perp \left(ABC\right)$, $SA=\sqrt{3}$ và $\triangle ABC$ vuông tại $B$ có cạnh $BC=1$, $AC=\sqrt{5}$. Tính khoảng cách  từ $A$ đến $\left(SBC\right)$. (kết quả làm tròn đến hàng phần trăm).
	 \shortans{0,76}
	 \loigiai{
	 \immini{
		Gọi $H$ là hình  chiếu của $A$ lên $SB$.\\
		Vì  $SA\perp \left(ABC\right)\Rightarrow SA\perp BC$.\\
		Ta  có $\heva{&SA\perp BC\\&AB\perp BC}\Rightarrow BC\perp \left(SAB\right)\Rightarrow BC\perp AH$.\\
		Lại có $\heva{&AH\perp BC\\&AH\perp SB}\Rightarrow AH\perp  \left(SBC\right)\Rightarrow {\mathrm{d}_{(A,(SBC))}}=AH$.\\
		Mặt khác  $AB=\sqrt{AC^2-BC^2}=\sqrt{5-1}=2$.\\
		Xét $\triangle SAB$ vuông tại  $A$ có $AH$ là đường cao nên ta có
		 $$AH=\dfrac{SA\cdot AB}{\sqrt{SA^2+AB^2}}=\dfrac{\sqrt{3}\cdot 2}{\sqrt{3+4}}=\dfrac{2\sqrt{21}}{7}\approx 0{,}76.$$
		Vậy khoảng  cách từ $A$ đến $\left(SBC\right)$ là $0{,}76$.
		}{		\begin{tikzpicture}[declare function={r=3; gocB=-40;},font=\footnotesize]
		 \path (0,0) coordinate (A)
			 (90:r) coordinate (S)
			(0:1.4*r)  coordinate (C)
			 (gocB:0.56*r) coordinate (B)
			 ($(S)!0.63!(B)$) coordinate (H);
		\foreach \x/\y/\z  in {A/B/C,A/H/B,S/A/C,S/A/B}{
			 \path pic[draw=red,angle radius=5pt]{right angle= \x--\y--\z};
		} 
		\draw  (S)--(A)--(B)--(C)--(S)--(B)
			 (A)--(H);
		\draw[dashed] (A)--(C); 
		 \foreach \t/\g in {A/180,B/-90,C/0,S/90,H/30}{
			\fill[ball color = red] (\t) circle  (1pt) node[shift={(\g:7pt)}]{$ \t $};
		 }
	\end{tikzpicture}
		} 
	} 
\end{ex}

% nguồn: Ok_TraLoiNgan_Lop11_Lan1
\begin{ex}%[1H8V5-3]%[ma:0945c4e9]
	Cho hình chóp $S.ABC$ có đáy  là tam giác đều cạnh $a$, $SA\perp (ABC)$ và $SB=2a$. Gọi $G$ là trọng tâm tam giác $ABC$. Với $a=\sqrt{15}$, tính khoảng cách từ $G$ đến mặt phẳng $(SBC)$.
	 \shortans{1}
	\loigiai{ 
		\begin{center}
			\begin{tikzpicture}
					 \def\a{4}
					\def\h{4} 
					\path  (0:0) coordinate (A)
						++(0:\a)  coordinate (B)
						++(-150:4*\a/5) coordinate  (C)
						($(A)+(90:\h)$) coordinate (S) 
						($(A)!0.5!(B)$)  coordinate (M)
						($(C)!0.5!(B)$)  coordinate (I)
						 ($(S)!0.3!(I)$) coordinate (H)
						($(A)!2/3!(I)$)  coordinate (G);
					\draw (A)--(C)--(B) 
						(A)--(S) 
						(B)--(S)--(I) 
						(C)--(S); 
					\draw[dashed] (I)--(A)--(B) 
						(H)--(A); 
					\foreach \x  /\goc in {A/180,B/0,C/-135,S/90,I/-20,G/-90,H/0}
					\fill[black] (\x) circle  (1pt)
						($(\x)+(\goc:3mm)$) node  {$\x$};
					 \draw pic[draw,angle radius=2mm]{right angle=A--I--C};
					\draw pic[draw,angle radius=2mm]{right angle=A--H--I}; 
				\end{tikzpicture}
		\end{center}
		Kẻ  $AI\perp BC$, kẻ $AH\perp SI$ tại $H$.\\
		Ta có $\heva{&BC\perp SA\\&BC\perp AI}\Rightarrow BC\perp  (SAI)\Rightarrow BC\perp AH$.\\
		Ta lại có $AH\perp SI \Rightarrow AH\perp (SBC) \Rightarrow \mathrm{d}(A,(SBC))=AH$.\\ 
		Ta  có $SA=\sqrt{SB^2-BA^a2}=\sqrt{(2a)^2-a^2}=\sqrt{3}a$.\\
		Ta có $\dfrac{1}{AH^2}=\dfrac{1}{SA^2}+\dfrac{1}{AI^2}=\dfrac{1}{3a^2}+\dfrac{4}{3a^2}\Rightarrow AH=\dfrac{\sqrt{15}}{5}a$.\\ 
		Vậy  $\mathrm{d}(A,(SBC))=\dfrac{\sqrt{15}}{5}a$.\\
		 Ta có $GA$ cắt $BC$ tại $I$.\\
		$\Rightarrow\dfrac{\mathrm{d}(G,(SBC))}{{\mathrm{d}(A,(SBC))}} =\dfrac{GI}{AI} =\dfrac{1}{3}\Rightarrow \mathrm{d}(G,(SBC))=\dfrac{1}{3}\mathrm{d}(A,(SBC))=\dfrac{\sqrt{15}}{15}a=1$  với $a=\sqrt{15}$.
	} 
\end{ex}

% nguồn: Ok_TraLoiNgan_Lop11_Lan1
\begin{ex}%[1H8V5-3]%[ma:0eab15ca]
	Cho hình chóp ${S. A B C}$ có đáy là tam giác vuông cân tại ${A, S H \perp(A B C)}$ với ${H}$ là trung điểm ${B C}$. Biết ${A B=S C=\sqrt{6}}$. Tính khoảng  cách từ ${B}$ đến mặt phẳng ${(S A C)}$.
	\shortans{2} 
	\loigiai{ 
		\begin{center}
			\begin{tikzpicture}[scale=1,font=\footnotesize,line join=round,line cap=round,>=stealth]
					 \path
						(0,0) coordinate (B) 
						++(-45:2) coordinate  (A)
						 (5,0) coordinate (C)
						 ($(B)!1/2!(C)$) coordinate (H)
						($(A)!1/2!(C)$)  coordinate (I)
						($(H)+(0,4)$)  coordinate (S)
						($(S)!1/2!(I)$)  coordinate (K);
					\foreach  \i in {A,B,C}{\draw (S)--(\i);}
					 \draw (B)--(A)--(C)
						(S)--(I); 
					\draw[dashed,thin](B)--(C) 
						(S)--(H)--(I) 
						 (H)--(K);
					\pic[draw,thin,angle radius=2mm] {right  angle = S--H--C};
					 \pic[draw,thin,angle radius=2mm] {right angle = B--A--C};
					 \path (A)--(B) node[midway,left]{$a$};
					\path (S)--(C)  node[midway,right]{$a$};
					\path (B)--(A)  node[midway,sloped,rotate=90]{$=$};
					\path (I)--(A)  node[midway,sloped,rotate=90]{$=$};
					\path (B)--(H) node[midway,sloped]{$|$}; 
					\path (C)--(H)  node[midway,sloped]{$|$};
					\foreach \i/\g in  {S/90,A/-90,B/180,C/0,H/-90,I/-50,K/0}
					\fill[black] (\i) circle (1pt)+(\g:4mm)  node[scale=1]{$\i$};
				\end{tikzpicture}
		\end{center}
		Kẻ ${H I \perp A C}$, kẻ ${H  K \perp S I}$ tại ${K}$.\\
		Ta có ${\left\{\begin{array}{l}A C \perp S H \\ A C \perp H I
	\end{array} \Rightarrow A C \perp(S H I) \Rightarrow A C \perp H K\right.}$\\
	Ta lại có ${H K \perp S I \Rightarrow H K \perp(S A C) \Rightarrow \mathrm{d}(H,(S A C))=H  K}.$\\
	Ta có ${{HI}=\dfrac{1}{2}  a}$;
	${S  H=\sqrt{S C^2-H C^2}=\sqrt{a^2-\left(\dfrac{a \sqrt{2}}{2}\right)^2}=\dfrac{\sqrt{2}}{2} a}$.\\
	Ta có ${H K=\dfrac{1}{\sqrt{\dfrac{1}{S H^2}+\dfrac{1}{H  I^2}}}=\dfrac{1}{\sqrt{\dfrac{1}{\sqrt{2})^2}+\dfrac{1}{(1)^2}}}=\dfrac{\sqrt{6}}{6} a}$.\\
	Vậy ${\mathrm{d}(H,(S A C))=\dfrac{\sqrt{6}}{6}  a}$.\\
	Ta có ${B H}$ cắt ${(S  A C)}$ tại ${C}$
	 $\Rightarrow \dfrac{\mathrm{d}(B,(S A C))}{\mathrm{d}(H,(S A C))}=\dfrac{B C}{H C}=2$.\\
	$ \Rightarrow \mathrm{d}(B,(S A C))=2 \mathrm{d}(H,(S A C))=\dfrac{\sqrt{6}}{3} a=2$ với $a=\sqrt{6}$. 
	} 
\end{ex}

% nguồn: Ok_TraLoiNgan_Lop11_Lan1
\begin{ex}%[1H8V5-3]%[ma:d95c89ec]
	\immini{Cho hình chóp $S.ABCD$ có đáy $ABCD$ là hình chữ nhật, biết $AB=a=\sqrt{10}$,  $AD=a\sqrt{2}$, $SA$ vuông góc với mặt phẳng đáy và $SA=a$ (tham khảo hình vẽ bên). Tính khoảng cách từ $A$ đến mặt phẳng $\left(SBD\right)$.
	} 
	{\begin{tikzpicture}[scale=0.7,font=\footnotesize,line join=round,line cap=round,>=stealth]
		\coordinate[label=below left:$B$] (B) at  (0,0);
		\coordinate[label=above left:$A$] (A) at  (45:2);
		 \coordinate[label=below:$C$] (C) at (4,0);
		\coordinate[label=right:$D$] (D) at  ($(C)-(B)+(A)$);
		\coordinate[label=above:$S$] (S)  at ($(A)+(90:4)$);
		 \draw (B)--(C)--(D)--(S)--cycle (S)--(C);
		 \draw[dashed] (A)--(D)
			 (S)--(A)--(B)
			(A)--(C) 
			 (B)--(D);
		\foreach \diem in {A,B,C,D,S} \fill (\diem) circle (1pt); 
		\path  pic[draw,angle radius=5]{right angle=S--A--D};
	\end{tikzpicture}
	 }
	\shortans{2} 
	\loigiai{\immini{Trong hình chữ nhật $ABCD$  kẻ $AK\perp BD$ ($K$ thuộc $BD$).\\
		 Trong tam giác $SAK$ kẻ $AH\perp SK$.\\
		Ta có $AK=\dfrac{AB\cdot AD}{\sqrt{AB^2+AD^2}}=\dfrac{a\cdot a\sqrt{2}}{\sqrt{a^2+2a^2}}=\dfrac{a\sqrt{6}}{3}$.\\ 
		Do  $\heva{&SA\perp BD\\&AK\perp BD}\Rightarrow\left(SAK\right)\perp\left(SBD\right)\Rightarrow AH\perp\left(SBD\right)$.\\
		Vậy $\mathrm{d}\left(A,(SBD)\right)=AH=\dfrac{AS\cdot  AK}{\sqrt{AS^2+AK^2}}=\dfrac{\dfrac{a\sqrt{6}}{3}\cdot a}{\sqrt{\dfrac{2a^2}{3}+a^2}}=\dfrac{a\sqrt{10}}{5}=2$.
		} 
		{\begin{tikzpicture}[scale=0.8,font=\footnotesize,line join=round,line cap=round,>=stealth]
		\coordinate[label=below left:$B$]  (B) at (0,0);
		 \coordinate[label=above left:$A$] (A) at (45:2);
		\coordinate[label=below:$C$] (C)  at (4,0);
		\coordinate[label=right:$D$] (D) at ($(C)-(B)+(A)$); 
		\coordinate[label=above:$S$] (S) at ($(A)+(90:4)$); 
		\coordinate[label=below:$K$] (K)  at ($(B)!0.35!(D)$);
		\coordinate[label=right:$H$] (H)  at ($(K)!0.35!(S)$);
		\draw (B)--(C)--(D)--(S)--cycle  (S)--(C);
		 \draw[dashed] (A)--(D)
			(K)--(S)--(A)--(B) 
			 (A)--(C)
			(B)--(D) 
			(A)--(H) 
			(A)--(K); 
		\foreach \diem in {A,B,C,D,S,K,H} \fill  (\diem) circle (1pt);
	\end{tikzpicture}
		 }
	 }
\end{ex}

% nguồn: Ok_TraLoiNgan_Lop11_Lan1
\begin{ex}%[1H8V5-3]%[ma:8458b87c]
	Cho hình chóp $S.ABC$ có đáy $ABC$ là tam giác vuông cân tại $A$, $AB=1$, $SA=SB=SC$. Góc giữa đường thẳng $SA$ và mặt phẳng $(ABC)$ bằng $45^{\circ}$. Tính khoảng cách từ điểm $S$ đến  mặt phẳng $(ABC)$ (kết quả làm tròn đến hàng phần trăm).
	\shortans{0,71} 
	 \loigiai{
		Gọi $H$ là hình chiếu  vuông góc của $S$ trên $(ABC)$.
	\immini{$ \Rightarrow SH\perp (ABC) \Rightarrow  SH=\mathrm{d}(S;(ABC))$. \\
		Ta có $\triangle SHA=\triangle SHB=\triangle SHC$  \\
		$ \Rightarrow HA=HB=HC \Rightarrow H$ là tâm  đường tròn ngoại tiếp $\triangle ABC$. \\
		Mặt khác $\triangle ABC$  vuông tại $A$ nên $H$ là trung điểm của $BC$. \\
		 Do $SH\perp (ABC) \Rightarrow (SA;(ABC))=(SA;AH)=\widehat{SAH}=45^{\circ}$. \\
		 $ \Rightarrow SH=AH=\dfrac{1}{2}BC=\dfrac{\sqrt{2}}{2}\approx 0{,}71$.}
		{\begin{tikzpicture}[line width=0.6pt,line cap=round,line join=round,>=stealth,scale=1]
		\coordinate[label=left:{$A$}] (A) at (0,0); 
		\coordinate[label=below:{$B$}] (B) at (1.5,-1); 
		\coordinate[label=right:{$C$}] (C) at  (4,0);
		 \coordinate[label=below:{$H$}] (H) at ($(B)!0.5!(C)$);
		\coordinate[label={$S$}]  (S) at ($(H)+(0,3)$);
		\draw  (B)--(A)--(S)--(B)--(C)--(S)--(H);
		\draw[dashed] (H)--(A)--(C); 
		\foreach \x/\dinh/\y in {B/A/C,S/H/A} \draw[fill = gray!50] ($(\dinh)!3pt!(\x)$)--($(\dinh)!3pt!(\x)+(\dinh)!3pt!(\y)-(\dinh)$)--($(\dinh)!3pt!(\y)$)--(\dinh)--cycle; 
		\foreach \p in {A,B,C,S,H} 
		\fill  (\p) circle (1pt);
	\end{tikzpicture}
		} 
	 }
\end{ex}

% nguồn: Ok_TraLoiNgan_Lop11_Lan1
\begin{ex}%[1H8V5-3]%[ma:f00b0166]
	Cho hình chóp đều $S.ABC$ có đáy cạnh $a$ và  cạnh bên $2a$. Với $a=10$, tính khoảng cách từ $A$ đến mặt phẳng $(SBC)$ (kết quả làm tròn hai chữ số thập phân).
	\shortans{8,56} 
	 \loigiai{
		\begin{center}
			\begin{tikzpicture}
					\def\a{4} 
					\path  (0:0) coordinate (B)
						++(0:\a)  coordinate (A)
						 ++(-120:\a/2) coordinate (C)
						($(B)+(70:\a)$)  coordinate (S)
						 ($(B)!0.5!(C)$) coordinate (I)
						($(S)!0.8!(I)$) coordinate (H) 
						($(A)!0.7!(I)$) coordinate (O); 
					 \draw[dashed] (B)--(A)--(I)
						(O)--(H); 
					 \draw (B)--(S)--(A)
						(I)--(S)--(C) 
						(B)--(C)--(A); 
					\foreach  \x/\g in {S/90,B/180,C/-45,A/0,I/-90,H/135,O/-45}
					\fill[black] (\x) circle (1pt) 
						($(\g:3mm)+(\x)$) node  {$\x$};
					\draw pic[draw,angle radius=2mm]{right angle=O--H--S}; 
					\draw pic[draw,angle  radius=2mm]{right angle=A--I--C};
				\end{tikzpicture}
		\end{center}
		Kẻ $OH\perp SI$ tại $H$. Ta có $\heva{&BC\perp OI\\&BC\perp SO}\Rightarrow BC\perp (SOI)\Rightarrow  BC\perp OH$.\\
		Ta lại  có $OH\perp SI \Rightarrow OH\perp (SBC)\Rightarrow \mathrm{d}(O,(SBC))=OH$.\\
		Ta có $OI=\dfrac{1}{3}\cdot\dfrac{a\sqrt{3}}{2}=\dfrac{\sqrt{3}}{6}a$;  $SO=\sqrt{SB^2-OB^2}=\sqrt{(2a)^2-\left(\dfrac{2}{3}\cdot\dfrac{a\sqrt{3}}{2}\right)^2}=\dfrac{\sqrt{33}}{3}a$.\\
		$OH=\dfrac{1}{\sqrt{\dfrac{1}{SO^2}+\dfrac{1}{OI^2}}}=\dfrac{1}{\sqrt{\dfrac{1}{\left(\dfrac{\sqrt{33}}{3}a\right)^2}+\dfrac{1}{\left(\dfrac{\sqrt{3}}{6}a\right)^2}}}=\dfrac{\sqrt{165}}{45}a$.\\ 
		 Vậy $\mathrm{d}(O,(SBC))=\dfrac{\sqrt{165}}{45}a$.\\
		Ta có $AO$ cắt $(SBC)$  tại $I$.\\
		$\Rightarrow \dfrac{\mathrm{d}(A,(SBC))}{\mathrm{d}(O,(SBC)}=\dfrac{AI}{OI}=3\Rightarrow \mathrm{d}(A,(SBC))=3\mathrm{d}(O,(SBC))=\dfrac{\sqrt {165}}{15}a \approx 8{,}56$ với $a=10$. 
	 }
\end{ex}

% nguồn: Ok_TraLoiNgan_Lop11_Lan1
\begin{ex}%[1H8V5-3]%[ma:a10f39dc]
	Cho hình chóp $S.ABCD$ có đáy là hình vuông cạnh $a$, tam giác $SAB$ đều và nằm trong mặt phẳng vuông góc với đáy. Với $a=\sqrt{21}$, khoảng cách từ $A$  đến $(SDC)$ là
	\shortans{3} 
	\loigiai{ 
		\begin{center}
			\begin{tikzpicture}
					 \def\a{4}
					\def\h{4.5} 
					\path (0:0) coordinate (A) 
						++(0:\a) coordinate  (D)
						++(-130:\a/2)  coordinate (C)
						($(A)+(C)-(D)$) coordinate  (B)
						 ($(A)!0.5!(B)$) coordinate (H)
						($(C)!0.5!(D)$) coordinate  (K)
						($(H)+(90:\h)$) coordinate  (S)
						($(S)!0.6!(K)$)  coordinate (I)
						 (intersection of A--C and B--D) coordinate (O);
					\draw[dashed]  (B)--(A)--(D)
						 (A)--(S)
						 (H)--(I)
						 (S)--(H)--(K);
					 \draw (B)--(C)--(D)
						 (B)--(S)--(K)
						 (C)--(S)
						(D)--(S); 
					\draw pic[draw,angle radius=2mm]{right  angle=H--I--S};
					\foreach \x/\g in {A/135,B/-135,C/-45,D/45,S/90,H/135,K/-90,I/90} 
					\fill[black]  (\x) circle (1pt)
						($(\g:4mm)+(\x)$) node  {$\x$};
				\end{tikzpicture}
		\end{center}
		Gọi $H$ là trung  điểm $AB$, suy ra $SH\perp AB$ (do tam giác $SAB$ đều).\\
		 Mặt khác $(SAB)\perp (ABCD)$ nên $SH\perp (ABCD)$.\\
		Ta có $AH\parallel CD$ nên $AH\parallel  (SCD)$.\\
		Kẻ $HK\perp CD$ và  $HI\perp SK$, suy ra $\mathrm{d}(H,(SCD))=HI$.\\
		Nên  $\mathrm{d}(A,(SCD))=\mathrm{d}(H,(SCD))=HI$.\\
		Ta  có $HK=a$, đường cao hình chóp là $SH=\dfrac{a\sqrt{3}}{2}$.\\
		Xét tam giác $SHK$ vuông tại $H$ có $HI$  là đường cao, nên ta có\\ $\dfrac{1}{HI^2}=\dfrac{1}{HK^2}+\dfrac{1}{HS^2}\Rightarrow \dfrac{1}{HI^2}=\dfrac{1}{a^2}+\dfrac{4}{3a^2}\Rightarrow HI=\dfrac{a\sqrt{21}}{7}=3$ với $a=\sqrt{21}$.
	} 
\end{ex}

% nguồn: Ok_TraLoiNgan_Lop11_Lan1
\begin{ex}%[1H8V5-3]%[ma:f08288a8]
	Cho hình chóp $S.ABC$ có $SA \perp (ABC)$, $SA = a\sqrt{2}$ và tam  giác $ABC$ đều cạnh $a$. Khoảng cách từ $A$ đến mặt phẳng $(SBC)$ bằng $\dfrac{a\sqrt{x}}{y}$, tính $x+y$.
	 \shortans{77}
	\loigiai{ 
	\immini 
	 {
	Gọi $M$ là trung điểm  của $BC$.\\
	Ta có 
	\begin{align*}
		\heva{& BC \perp AM \\& BC \perp SA}  \Rightarrow BC \perp (SAM).
	\end{align*}
	Mà  $BC \subset (SBC)$ nên $(SBC) \perp (SAM)$.
	Kẻ $AH \perp SM$ tại $H$. Suy ra $AH  \perp (SBC)$.\\
	Vậy khoảng cách  từ $A$ đến mặt phẳng $(SBC)$ là $AH$.\\
	} 
	 {
	\begin{tikzpicture}[scale=0.8,font=\footnotesize,line join=round,line cap=round,>=stealth]
			\def\a{4} \def\b{2}  \def\h{4}
			 \path
				(0:0) coordinate  (A)
				(0:\a)  coordinate (B)
				(-40:\b) coordinate  (C)
				(90:\h) coordinate  (S)
				($(B)!0.5!(C)$) coordinate (M) 
				($(S)!0.6!(M)$) coordinate (H); 
			 \draw
				(S)--(A)--(C)--(B)--(S) 
				(S)--(M) 
				 (S)--(C);
			\draw[dashed] (B)--(A)--(M) 
				 (A)--(H);
			\pic [draw,,angle radius=6] {right angle  = A--H--S};
			\pic [draw,,angle radius=6] {right angle  = A--M--C};
			\foreach  \i/\g in {A/180,B/0,C/-90,M/-70,S/90,H/45}{\filldraw[black](\i) circle (1pt)
				($(\i)+(\g:3mm)$)  node[scale=1]{$\i$};}
		\end{tikzpicture}
	 }
	\noindent 
	Ta lại có $AM  = \dfrac{a \sqrt{3}}{2}$ (đường cao trong tam giác đều).\\
	Khi  đó
	\begin{align*}
		& \dfrac{1}{AH^2} = \dfrac{1}{SA^2} + \dfrac{1}{AM^2} = \dfrac{1}{2a^2} + \dfrac{4}{3a^2}  = \dfrac{11}{6a^2}\\
		\Rightarrow &  AH = \dfrac{a \sqrt{66}}{11}.
	\end{align*}
	Do đó  $d(A, (SBC)) = \dfrac{a \sqrt{66}}{11}$. Suy ra $x=66$, $y=11$ nên $x+y=66+11=77$.
	} 
\end{ex}

% nguồn: Ok_TraLoiNgan_Lop11_Lan1
\begin{ex}%[1H8V5-3]%[ma:32888ea8]
	Cho hình chóp ${S. A B C D}$ có đáy là hình vuông cạnh $6$, hai mặt phẳng ${(S A B)}$ và ${(S  B C)}$ cùng vuông góc với mặt phẳng ${(A B C D)}$ và ${S C=6\sqrt{5}}$. Tính khoảng cách từ ${D}$ đến mặt phẳng ${(S A C)}$.
	 \shortans{4}
	 \loigiai{
		\begin{center}
			\begin{tikzpicture}[scale=1,font=\footnotesize,line join=round,line cap=round,>=stealth]
					\path 
						(0,0) coordinate (B) 
						++(-140:3)  coordinate (A)
						 ++(0:5) coordinate (D)
						($(B)+(D)-(A)$) coordinate  (C)
						 ($(B)+(0,4)$) coordinate (S)
						($(A)!.5!(C)$) coordinate (O) 
						($(S)!(B)!(O)$) coordinate (H); 
					\foreach \i  in{A,C,D}{\draw (S)--(\i);};
					\draw  (A)--(D)--(C);
					 \draw[dashed] (S)--(B)--(A)
						 (A)--(C)--(B)--(D)
						 (S)--(O)
						 (B)--(H);
					\pic[draw,angle eccentricity=1.8,angle  radius=2mm]{right angle=B--H--O};
					\path (S)--(C)  node[right,midway]{$a\sqrt{5}$};
					\path  (D)--(A) node[below,midway]{$a$};
					\foreach \i/\g  in {A/180,B/180,C/-90,D/-90,S/90,O/-90,H/0}
					\fill[black] (\i)  circle (1pt)+(\g:3mm) node[scale=1]{$\i$};
				\end{tikzpicture}
		\end{center}
		Đặt  $AB=a, SC=a\sqrt{5}$.\\
		Kẻ $BH\perp  SO$ tại ${H}$.\\
		Ta có $\heva{&A C \perp S B  \\ &A C \perp B D} \Rightarrow AC \perp(SBD) \Rightarrow AC \perp BH$.\\
		Ta lại có ${B H \perp S O \Rightarrow  B H \perp(S A C) \Rightarrow \mathrm{d}(B,(S A C))=B H}$.\\
		Ta có ${S B=\sqrt{S C^2-B C^2}=\sqrt{(a  \sqrt{5})^2-a^2}=2 a}$.\\
		Ta có ${B H=\dfrac{1}{\sqrt{\dfrac{1}{S B^2}+\dfrac{1}{O B^2}}}=\dfrac{1}{\sqrt{\dfrac{1}{4a^2}+\dfrac{1}{\frac{a^2}{2}}}}}=\dfrac{2a}{3}$. \\ 
		Vậy  ${\mathrm{d}(B,(S A C))=\dfrac{2}{3} a}$.\\
		 Ta có ${D B}$ cắt ${(S A C)}$ tại ${O}$
		$\Rightarrow \dfrac{\mathrm{d}(D,(S A C))}{\mathrm{d}(B,(S  A C))}=\dfrac{D O}{B O}=1$\\
		$ \Rightarrow \mathrm{d}(D,(S A C))=\mathrm{d}(B,(S  A C))=\dfrac{2}{3} a=4$ với $a=6$.
	} 
\end{ex}

% nguồn: Ok_TraLoiNgan_Lop11_Lan1
\begin{ex}%[1H8V5-3]%[ma:5798a062]
	Cho hình lăng trụ $ABC.A'B'C'$ có tất cả các cạnh đều bằng $1$. Góc tạo bởi cạnh  bên và mặt phẳng đáy bằng $30^{\circ}$. Hình chiếu $H$ của $A'$ lên mặt phẳng $(ABC)$ thuộc đường thẳng $BC$. Tính khoảng cách từ $B$ đến mặt phẳng $(ACC'A')$ (làm tròn đến hàng phần trăm).
	 \shortans{0,65}
	 \loigiai{
		Ta có $\widehat{A'AH}=30^{\circ}$ do đó $AH=\cos  30^{\circ}=\dfrac{\sqrt{3}}{2}$ và $A'H=\dfrac{\sqrt{3}}{2}\cdot \dfrac{1}{\sqrt{3}}=\dfrac{1}{2}$. \\
	\immini{Gọi $K$ là hình chiếu của $H$ lên $AC$ và $I$ là hình chiếu của $H$ lên $A'K$. \\ 
		 Do tam giác $ABC$ đều cạnh bằng $a$ nên $H$ thuộc đoạn $BC$ và $H$ là trung điểm $BC$. \\
		$\mathrm{d}(B,(ACC'A'))=2\mathrm{d}(H,(ACC'A'))=2HI$.  \\
		 $HK=\dfrac{1}{2}\cdot\dfrac{a\sqrt{3}}{2}=\dfrac{\sqrt{3}}{4}$ nên\\
		 $2HI=2\dfrac{HA'\cdot HK}{\sqrt{H{A'}^2+HK^2}}=2\dfrac{\dfrac{1}{2}\cdot \dfrac{\sqrt{3}}{4}}{\sqrt{\dfrac{1}{4}+\dfrac{3}{16}}}=\dfrac{\dfrac{\sqrt{3}}{4}}{\sqrt{\dfrac{7}{16}}}=\dfrac{\sqrt{21}}{7}\approx 0{,}65$.}
		 {
		\begin{tikzpicture}[scale=1,font=\footnotesize,line join=round,line cap=round,>=stealth]
				\coordinate[label=left:{$A$}] (A) at (0,0); 
				\coordinate[label=below:{$C$}]  (B) at (0.8,-1);
				\coordinate[label=right:{$B$}] (C) at (3,0); 
				 \coordinate[label=below:{$H$}] (H) at ($(B)!0.5!(C)$);
				\coordinate[label={$A'$}] (A') at ($(H)+(0,3.5)$); 
				\coordinate[label={$C'$}] (B') at ($(A')+(B)-(A)$); 
				 \coordinate[label={$B'$}] (C') at ($(A')+(C)-(A)$);
				\coordinate[label=left:{$K$}]  (K) at ($(A)!3/4!(B)$);
				\coordinate[label=left:{$I$}] (I)  at ($(A')!3/4!(K)$);
				\draw (B)--(B')--(A')--(A)--(B)--(C)--(C')--(B') 
					 (C')--(A')--(K);
				\draw[dashed]  (A')--(H)--(A)--(C)
					(K)--(H)--(I); 
				\foreach \x/\dinh/\y in {A'/H/K,H/K/A,H/I/A'}  \draw[fill = gray!50] ($(\dinh)!3pt!(\x)$)--($(\dinh)!3pt!(\x)+(\dinh)!3pt!(\y)-(\dinh)$)--($(\dinh)!3pt!(\y)$)--(\dinh)--cycle;
				 \foreach \p in {A,B,C,H,I,K,A',B',C'}
				\fill  (\p) circle (1pt);
			\end{tikzpicture}
	 }}
\end{ex}

% nguồn: Ok_TraLoiNgan_Lop11_Lan2
\begin{ex}%[1H8V5-3]%[ma:a33df975]
	Cho hình chóp $S.ABC$ có đáy là tam giác vuông đỉnh $B$, $AB=1$, $SA$  vuông góc với mặt phẳng đáy và $SA=2$. Khoảng cách từ $A$ đến mặt phẳng $(SBC)$ bằng (kết quả làm tròn đến hàng phần trăm).
	\begin{center}
		\begin{tikzpicture}[scale=0.65,font=\footnotesize,line join=round,line cap=round,>=stealth]
				\coordinate (A) at  (0,0);
				 \coordinate (S) at (0,3.5);
				\coordinate (B) at  (2,-1);
				\coordinate  (C) at (5,0.2);
				\draw (S) --  (A) -- (B) -- (C) -- (S);
				\draw (S)  -- (B);
				\draw[dashed]  (A) -- (C);
				\pic[draw,thin,angle radius=2mm] {right angle = S--A--C} pic[draw,thin,angle radius=3mm] {right angle  = A--B--C};
				\foreach \i/\g in {S/90,A/-90,B/-90,C/-90}{\draw[fill](\i)  circle (1pt)
					($(\i)+(\g:3mm)$) node[scale=1]{$\i$};} 
			\end{tikzpicture}
	\end{center}
	\shortans{0,89} 
	 \loigiai{
		\begin{center}
			\begin{tikzpicture}[scale=0.65,font=\footnotesize,line join=round,line cap=round,>=stealth]
					\coordinate (A) at (0,0); 
					\coordinate  (S) at (0,3.5);
					\coordinate (B) at (2,-1); 
					 \coordinate (C) at (5,0.2);
					\coordinate (H)  at ($(S)!0.4!(B)$);
					\draw (S) -- (A) -- (B) -- (C)  -- (S)
						 (S) -- (B)
						 (A)--(H);
					\draw[dashed]  (A) -- (C);
					\pic[draw,thin,angle radius=2mm] {right  angle = S--A--C} pic[draw,thin,angle radius=2mm] {right angle = A--B--C} pic[draw,thin,angle radius=2mm] {right angle = A--H--B};
					 \foreach \i/\g in {S/90,A/-90,B/-90,C/-90,H/60}{\draw[fill](\i) circle (1pt)
						($(\i)+(\g:3mm)$)  node[scale=1]{$\i$};}
				\end{tikzpicture}
		\end{center}
		 Ta có $\heva{ &BC \perp AB \\ &BC \perp SA} \Rightarrow BC \perp (SAB)$.\\
		Kẻ $AH \perp  SB$.\\
		Khi đó  $AH \perp BC \Rightarrow AH \perp (SBC)$.\\
		$\Rightarrow AH$ là khoảng  cách từ $A$ đến mặt phẳng $(SBC)$.\\
		Ta có $\dfrac{1}{AH^2}=\dfrac{1}{SA^2}+\dfrac{1}{AB^2}=\dfrac{1}{2^2}+\dfrac{1}{1^2}=\dfrac{1}{4}+1=\dfrac{5}{4} \Rightarrow AH^2=\dfrac{4}{5}  \Rightarrow AH=\dfrac{2\sqrt{5}}{5}\approx 0,89$.
	} 
\end{ex}

% nguồn: Ok_TraLoiNgan_Lop11_Lan2
\begin{ex}%[1H8V5-3]%[ma:f6d068f8]
	Cho khối chóp $S.ABCD$ có đáy $A B C D$ là hình vuông cạnh $1$, $S A \perp(A B C D)$ và $S A=\sqrt{2}$. Gọi $M$ là trung điểm cạnh $SC$. Khoảng cách từ điểm $M$ đến  mặt phẳng $(SBD)$ bằng (kết quả làm tròn đến hàng phần trăm)
	 \shortans{0,32}
	 \loigiai{
		\begin{center}
			\begin{tikzpicture}[scale=1,font=\footnotesize,line join=round,line cap=round,>=stealth]
					\coordinate (A) at  (0,0);
					\coordinate (D)  at (-2,-2);
					\coordinate (B) at (5,0); 
					\coordinate (C)  at ($(B)+(D)-(A)$);
					\coordinate  (S) at ($(A)+(0,5)$);
					\coordinate (M) at ($(S)!0.5!(C)$); 
					\coordinate (O)  at ($(A)!0.5!(C)$);
					\coordinate  (H) at ($(S)!0.6!(O)$);
					\draw  (S)--(B)
						 (S)--(C)
						(S)--(D) 
						 (B)--(C)--(D);
					\draw[dashed](S)--(A) 
						(B)--(A)--(D) 
						(S)--(A) 
						 (A)--(H)
						 (A)--(C)
						 (B)--(D)
						(S)--(O); 
					\foreach \i/\g  in {S/90,A/-90,B/-90,C/-90,D/-90,H/0,M/0,O/-90}{\draw[fill=black](\i) circle (1pt)
						 ($(\i)+(\g:3mm)$) node[scale=1]{$\i$};}
				\end{tikzpicture}
		\end{center}
		Do $M$ là trung điểm $S C$ nên $\mathrm{d}(M,(S B D))=\dfrac{1}{2} \mathrm{d}(C,(S B  D))=\dfrac{1}{2}\mathrm{d}(A;(S B D))$.\\
		Gọi $H$  là hình chiếu của $A$ lên $(S B D)$.\\
		$\Rightarrow \mathrm{d}(A,(S B D))=A H$.\\ 
		Lại có $A S$, $A B$,  $A D$ đôi một vuông góc nên
		\[\dfrac{1}{A H^{2}}=\dfrac{1}{A S^{2}}+\dfrac{1}{A B^{2}}+\dfrac{1}{A  D^{2}}=\dfrac{1}{1}+\dfrac{1}{1}+\dfrac{1}{2}=\dfrac{5}{2}.\]\\
		$\Rightarrow  A H=\dfrac{\sqrt{10}}{5}$.
		Vậy $\mathrm{d}(M;(S B D))=\dfrac{\sqrt{10}}{10} \approx  0{,}32$.
	} 
\end{ex}

% nguồn: Ok_TraLoiNgan_Lop11_Lan2
\begin{ex}%[1H8V5-3]%[ma:e58815bf]
	Cho hình lăng trụ đứng $ABC.A'B'C'$ có đáy $ABC$ là tam giác đều cạnh 1 và $AA'=2$. Gọi $M$ là trung điểm của $A'A$ (tham khảo  hình bên). Khoảng cách từ $M$ đến mặt phẳng $(A'BC)$ bằng (kết quả làm tròn đến hàng phần mười)
	\begin{center}
		\begin{tikzpicture}[scale=1,font=\footnotesize,line join=round,line cap=round,>=stealth]
				\coordinate (A) at  (0,0);
				 \coordinate (C) at (1,-1.5);
				 \coordinate (B) at (4,0);
				 \coordinate (A') at ($(A)+(0,3)$);
				\coordinate (C') at ($(C)+(0,3)$); 
				\coordinate (B')  at ($(B)+(0,3)$);
				 \coordinate (M) at ($(A')!0.5!(A)$);
				 \draw (A)--(C)--(C')--(B')--(A')--(A)
					(B')--(B)--(C) 
					 (A')--(C')
					(B')--(C); 
				\draw[dashed,thin](B)--(A)--(B'); 
				\foreach  \i/\g in {A'/90,A/-90,B/-90,C/-90,C'/60,B'/90,M/0}{\draw[fill=white](\i) circle (1.5pt)
					($(\i)+(\g:3mm)$) node[scale=1]{$\i$};} 
			\end{tikzpicture}
	\end{center}
	 \shortans{0,4}
	\loigiai{ 
		\begin{center}
			\begin{tikzpicture}[scale=1,font=\footnotesize,line join=round,line cap=round,>=stealth]
					\coordinate  (A) at (0,0);
					\coordinate  (C) at (1,-1.5);
					 \coordinate (B) at (4,0);
					\coordinate (A')  at ($(A)+(0,3)$);
					\coordinate  (C') at ($(C)+(0,3)$);
					\coordinate (B')  at ($(B)+(0,3)$);
					\coordinate (M) at  ($(A')!0.5!(A)$);
					\coordinate (K) at  ($(A)!0.5!(C)$);
					\coordinate (H) at ($(K)!0.6!(B')$); 
					\coordinate (I)  at (intersection of B--M and A--B');
					 \draw (A)--(C)--(C')--(B')--(A')--(A)
						(B')--(B)--(C) 
						 (A')--(C')
						(B')--(C); 
					\draw[dashed,thin](K)--(B)--(A)--(B') 
						(H)--(B)--(M) 
						(K)--(B'); 
					\pic[draw,thin,angle radius=3mm] {right angle = B--H--K}; 
					 \foreach \i/\g in {A'/90,A/-90,B/-90,C/-90,C'/60,B'/90,M/180,H/90,K/-120,I/-90}{\draw[fill=white](\i) circle (1.5pt)
						($(\i)+(\g:3mm)$)  node[scale=1]{$\i$};}
				\end{tikzpicture}
		\end{center}
		Gọi $I=BM \cap AB'$ và $K$ là trung điểm $AC$.  \\
		Ta có $\dfrac{\mathrm{d}(M,(A'BC))}{\mathrm{d}(B,(A'BC))}=\dfrac{MI}{BI}=\dfrac{MA}{Bb'}=\dfrac{1}{2}\Rightarrow \mathrm{d}(M,(A'BC))=\dfrac{1}{2}\mathrm{d}(B,(A'BC))=\dfrac{BH}{2}$.\\ 
		Xét tam  giác $BB'K$ có
		 \[\dfrac{1}{BH^2}=\dfrac{1}{B'B^2}+\dfrac{1}{BK^2}=\dfrac{1}{2^2}+\dfrac{1}{\left(\dfrac{\sqrt{3}}{2}\right)^2}=\dfrac{1}{4}+\dfrac{4}{3}=\dfrac{19}{12}\Rightarrow BH=\dfrac{2\sqrt{57}}{19}.\]
		Vậy $\mathrm{d}(M,(A'BC))  = \dfrac{BH}{2} = \dfrac{\dfrac{2\sqrt{57}}{19}}{2} = \dfrac{\sqrt{57}}{19} \approx 0{,}4$.
	 }
\end{ex}

% nguồn: Ok_TraLoiNgan_Lop11_Lan2
\begin{ex}%[1H8V5-3]%[ma:90c52596]
	Cho  hình chóp $S.ABCD$ có đáy là hình vuông cạnh $1$, mặt bên $SAB$ là tam giác đều và nằm trong mặt phẳng vuông góc với mặt phẳng đáy (minh họa như hình vẽ bên). Khoảng cách từ $A$ đến mặt phẳng $(SBD)$ bằng (kết quả làm tròn đến hàng phần trăm)
	 \shortans{0,65}
	\begin{center}
		\begin{tikzpicture}[scale=1,font=\footnotesize,line join=round,line cap=round,>=stealth]
				\coordinate (A) at (0,0); 
				\coordinate (B) at (-2,-2); 
				 \coordinate (D) at (4,0);
				\coordinate (I)  at ($(B)!0.5!(A)$);
				\coordinate  (C) at ($(B)+(D)-(A)$);
				\coordinate (S) at ($(I)+(0,4)$); 
				\draw (S)--(B) 
					 (S)--(C)
					 (S)--(D)
					 (B)--(C)--(D);
				\draw[dashed,thin](S)--(A)--(C) 
					 (A)--(B)
					(A)--(D)--(B); 
				 \foreach \i/\g in {S/90,A/150,B/-90,C/-90,D/0}{\draw[fill=black](\i) circle (1pt)
					 ($(\i)+(\g:3mm)$) node[scale=1]{$\i$};}
			\end{tikzpicture}
	\end{center}
	 \loigiai{
		\begin{center}
			\begin{tikzpicture}[scale=1,font=\footnotesize,line join=round,line cap=round,>=stealth]
					\coordinate  (A) at (0,0);
					\coordinate (B) at  (-2,-2);
					\coordinate (D) at (6,0); 
					\coordinate (H) at ($(B)!0.5!(A)$); 
					\coordinate (C)  at ($(B)+(D)-(A)$);
					\coordinate (O) at  ($(A)!0.5!(C)$);
					\coordinate (S)  at ($(H)+(0,4)$);
					\coordinate  (K) at ($(B)!0.4!(O)$);
					\coordinate  (I) at ($(S)!0.6!(K)$);
					\draw (S)--(B) 
						 (S)--(C)
						 (S)--(D)
						(B)--(C)--(D); 
					\draw[dashed,thin] (S)--(H)--(K)--(S)--(A)--(C) 
						 (A)--(B)
						(A)--(D)--(B) 
						(I)--(H); 
					\foreach \i/\g in  {S/90,A/150,B/-90,C/-90,D/0,O/-90,I/150,K/-60,H/180}{\draw[fill=black](\i) circle (1pt)
						($(\i)+(\g:3mm)$)  node[scale=1]{$\i$};}
				\end{tikzpicture}
		\end{center}
		Gọi $H$ là trung điểm  $AB$. Khi đó $SH \perp (ABCD)$.\\
		 Gọi $O$ là giao điểm của $AC$ và $BD$ suy ra $AC \perp BD$.\\
		Kẻ $HK \perp BD$  tại $K$ (K là trung điểm $BO$).\\
		Kẻ $HI \perp SH$ tại  $I$.\\
		Khi đó $\mathrm{d}(A,(SBD)) = 2\mathrm{d}(H,(SBD))=2HI$.\\ 
		Xét tam giác $SHK$  có $SH=\dfrac{\sqrt{3}}{2}$, $HK=\dfrac{1}{2}AO = \dfrac{\sqrt{2}}{4}$.\\
		Khi đó: $\dfrac{1}{HI^2}=\dfrac{1}{SH^2}+\dfrac{1}{HK^2}=\dfrac{4}{3}+8  = \dfrac{28}{3} \Rightarrow HI = \dfrac{\sqrt{21}}{14}$.\\
		Suy ra $\mathrm{d}(A,(SBD)) =  2HI = \dfrac{\sqrt{21}}{7} \approx 0{,}65$.
	} 
\end{ex}

% nguồn: Ok_TraLoiNgan_Lop11_Lan2
\begin{ex}%[1H8V5-3]%[ma:7f15c57a]
	Cho lăng trụ đứng $ABC.A'B'C'$ có đáy  $ABC$ là tam giác đều cạnh $1$ và $AA'=2$. Gọi $M$ là trung điểm của $CC'$ (tham khảo hình bên). Khoảng cách từ $M$ đến mặt phẳng $(A'BC)$ bằng (kết quả làm tròn đến hàng phần trăm) \shortans{0,79}
	\begin{center}
		\begin{tikzpicture}[scale=1,font=\footnotesize,line join=round,line cap=round,>=stealth]
				\coordinate (A) at (0,0); 
				\coordinate  (B) at (1,-1.5);
				\coordinate  (C) at (4,0);
				 \coordinate (A') at ($(A)+(0,3)$);
				\coordinate (B')  at ($(B)+(0,3)$);
				 \coordinate (C') at ($(C)+(0,3)$);
				\coordinate (M) at ($(C')!0.5!(C)$); 
				\draw  (A)--(B)--(C)--(C')--(B')--(A')--(A)
					(B')--(B) 
					(A')--(C') 
					(A')--(B); 
				 \draw[dashed,thin](A)--(C)--(A');
				 \foreach \i/\g in {A'/90,A/-90,B/-90,C/-90,C'/0,B'/180,M/0}{\draw[fill=black](\i) circle (1pt)
					 ($(\i)+(\g:3mm)$) node[scale=1]{$\i$};}
			\end{tikzpicture}
	\end{center}
	 \loigiai{
		\begin{center}
			\begin{tikzpicture}[scale=1,font=\footnotesize,line join=round,line cap=round,>=stealth]
					\coordinate (A)  at (0,0);
					 \coordinate (B) at (1,-1.5);
					 \coordinate (C) at (4,0);
					\coordinate (A') at  ($(A)+(0,3)$);
					\coordinate (B') at  ($(B)+(0,3)$);
					 \coordinate (C') at ($(C)+(0,3)$);
					\coordinate (M) at  ($(C')!0.5!(C)$);
					 \coordinate (H) at ($(B)!0.5!(C)$);
					\coordinate  (K) at ($(A')!0.6!(H)$);
					\draw (A)--(B)--(C)--(C')--(B')--(A')--(A) 
						 (B')--(B)
						(A')--(C') 
						 (A')--(B);
					 \draw[dashed,thin](H)--(A)--(C)--(A')--(H)
						(A)--(K); 
					\pic[draw,thin,angle  radius=2mm] {right angle = B--H--A} pic[draw,thin,angle radius=2mm] {right angle = A--K--H};
					\foreach \i/\g in {A'/90,A/-90,B/-90,C/-90,C'/0,B'/180,M/0,H/-90,K/0}{\draw[fill=black](\i)  circle (1pt)
						($(\i)+(\g:3mm)$)  node[scale=1]{$\i$};}
				\end{tikzpicture}
		\end{center}
		Gọi $H$ là hình chiếu của $A$ lên $BC$, $K$ là hình chiếu  của $A$ lên $A'H$. \\
		Ta có $\mathrm{d}(M,(A'BC)) = \dfrac{1}{2}\mathrm{d}(C,(A'BC)) =  \dfrac{1}{2}\mathrm{d}(A,(A'BC)) = \dfrac{1}{2}AK$. \\
		Mà $AH =  \dfrac{\sqrt{3}}{2}$; $AA'=2$ nên $$AK = \dfrac{AH \cdot AA'}{\sqrt{AH^2+AA'^2}} = \dfrac{\dfrac{\sqrt{3}}{2} \cdot 2}{\sqrt{\left(\dfrac{\sqrt{3}}{2}\right)^2+2^2}} = \dfrac{\sqrt{3}}{\sqrt{\dfrac{3}{4}+4}} = \dfrac{\sqrt{3}}{\sqrt{\dfrac{19}{4}}} = \dfrac{2\sqrt{3}}{\sqrt{19}} = \dfrac{2\sqrt{57}}{19} \approx 0{,}79.$$
		Vậy  $\mathrm{d}(M,(A'BC))=0{,}79$.
	} 
\end{ex}

% nguồn: Ok_TraLoiNgan_Lop11_Lan2
\begin{ex}%[1H8V5-3]%[ma:2df942d7]
	Cho hình chóp $S.ABCD$ có đáy là hình thoi  cạnh 1, $\widehat{BAD}=60^\circ$, $SA=1$ và $SA$ vuông góc với mặt phẳng đáy. Khoảng cách từ $B$ đến $(SCD)$ bằng (kết quả làm tròn đến hàng phần trăm)
	\shortans{0,65} 
	\loigiai{ 
		\begin{center}
			\begin{tikzpicture}[scale=1,font=\footnotesize,line join=round,line cap=round,>=stealth]
					 \coordinate (A) at (0,0);
					\coordinate (B)  at (-2,-2);
					\coordinate (D) at (4,0); 
					\coordinate  (C) at ($(B)+(D)-(A)$);
					\coordinate (S) at  ($(A)+(0,3)$);
					\coordinate (H) at ($(S)!0.4!(B)$); 
					\coordinate  (M) at ($(D)!1.4!(C)$);
					 \coordinate (H) at ($(S)!0.4!(M)$);
					 \coordinate (G) at (intersection of A--M and B--C);
					 \draw (S)--(B)
						(S)--(C)--(M)--(S)--(D) 
						 (B)--(G)
						(C)--(D) 
						 (G)--(M);
					\draw[dashed,thin](S)--(A) 
						(A)--(B) 
						 (A)--(D)
						(G)--(A)--(H) 
						(G)--(C); 
					\pic[draw,thin,angle radius=2mm] {right angle = A--H--M} pic[draw,thin,angle  radius=2mm] {right angle = S--A--B};
					\foreach \i/\g in {S/90,A/-90,B/-90,C/-90,D/0,H/180,M/-90}{\draw[fill=white](\i) circle  (1.5pt)
						($(\i)+(\g:3mm)$) node[scale=1]{$\i$};} 
				\end{tikzpicture}
		\end{center}
		Ta có $AB\parallel CD \Rightarrow \mathrm{d}(B,(SCD))  = \mathrm{d}(A,(SCD))$.\\
		Kẻ $AM  \perp CD$ $(M \in CD)$.\\
		Kẻ  $AH \perp SM \Rightarrow AH \perp (SCD) \Rightarrow \mathrm{d}(A,(SCD)) = SH$.\\
		Ta có $AM = \dfrac{2S_{ACD}}{CD} =  \dfrac{S_{ABCD}}{CD}=\dfrac{\sqrt{3}}{2}$.\\
		Xét $\triangle SAM$ vuông tại $A$ có $\dfrac{1}{AH^2}  = \dfrac{1}{SA^2}+\dfrac{1}{AM^2}=\dfrac{7}{3} \Rightarrow AH = \dfrac{\sqrt{21}}{7} \approx 0{,}65$.\\
		Vậy $\mathrm{d}(B,(SCD)) = 0{,}65$. 
	 }
\end{ex}

% nguồn: Ok_TraLoiNgan_Lop11_Lan2
\begin{ex}%[1H8V5-3]%[ma:e38c54dc]
	Cho hình chóp tứ  giác đều $S . A B C D$ có cạnh đáy bằng 1 và chiều cao bằng $\sqrt{2}$. Tính khoảng cách $d$ từ tâm $O$ của đáy $(A B C D)$ đến một mặt bên (kết quả làm tròn đến hàng phần trăm).
	\shortans{0,47} 
	 \loigiai{
		\begin{center}
			\begin{tikzpicture}[scale=.7,font=\footnotesize,line join=round,line cap=round,>=stealth]
					 \coordinate (A) at (0,0);
					\coordinate (B) at (-2,-2); 
					\coordinate (D) at  (7,0);
					\coordinate (C) at  ($(B)+(D)-(A)$);
					\coordinate (O) at ($(A)!0.5!(C)$); 
					\coordinate (S) at ($(O)+(0,7)$); 
					 \coordinate (H) at ($(D)!0.5!(C)$);
					\coordinate  (K) at ($(S)!0.7!(H)$);
					\draw (S)--(B) 
						 (S)--(C)
						 (B)--(C)
						(S)--(D) 
						 (S)--(H)
						(C)--(D); 
					 \draw[dashed](S)--(A)--(C)
						(A)--(D) 
						 (C)--(D)
						 (S)--(O)
						(B)--(D) 
						(O)--(K) 
						 (S)--(O)--(H)
						 (A)--(B)
						 (O)--(K);
					\pic[draw,thin,angle radius=2mm] {right angle = H--K--O} pic[draw,thin,angle radius=2mm]  {right angle = C--H--A};
					\foreach \i/\g in {S/90,A/180,B/-90,C/-90,D/0,O/-90,H/0,K/0}{\draw[fill=black](\i) circle (1pt) 
						($(\i)+(\g:3mm)$) node[scale=1]{$\i$};} 
				\end{tikzpicture}
		\end{center}
		Vì $S.ABCD$ là hình chóp tứ giác đều nên $ABCD$  là hình vuông và $SO \perp(ABCD)$.\\
		Vẽ $OH$ vuông góc với $C D$ tại $H$ thì $H$ là trung điểm  $C D$ và $O H=\dfrac{1}{2}$.\\
		Ta có $\heva{&CD\perp SO\\&CD\perp OH}\Rightarrow C D \perp(S O H) \Rightarrow(S C D) \perp(S O H)$.\\ Kẻ $O K$ vuông góc với  $S H$ tại $K$. \\
		Ta có $\heva{&OK\perp CD\\&OK\perp SH}\Rightarrow O K \perp(S C  D)$.\\
		$\Rightarrow \mathrm{d}(O,(S  C D))=O K$.\\
		Tam giác vuông $S O H$  có $O K$ là đường cao nên
		\[O K=\dfrac{O S \cdot O H}{\sqrt{O  S^{2}+O H^{2}}}=\dfrac{\sqrt{2} \cdot \dfrac{1}{2}}{\sqrt{2+\dfrac{1}{4}}}=\dfrac{\sqrt{2}}{3}.\]
		Vậy $\mathrm{d}(O,(S C  D))=\dfrac{\sqrt{2}}{3} \approx 0,47$.
	} 
\end{ex}

% nguồn: Ok_TraLoiNgan_Lop11_Lan2
\begin{ex}%[1H8V5-3]%[ma:75455716]
	Một hình lăng trụ đứng $ABC.A'B'C'$ có đáy $ABC$ là tam giác vuông tại $B$, $AB=1$, $AA'=2$.  Khoảng cách từ điểm $A$ đến mặt phẳng $(A'BC) $ bằng bao nhiêu? (kết quả làm tròn đến hàng phần trăm)
	 \shortans{0,89}
	 \loigiai{
		\begin{center}
			\begin{tikzpicture}[scale=1,font=\footnotesize,line join=round,line cap=round,>=stealth]
					 \def\h{3} \def\d{4}
					\path 
						(0,0)  coordinate (A)
						(2.5,-1.25) coordinate  (B)
						(A)++(\d,0)  coordinate (C)
						(A)++(0,\h) coordinate (A') 
						(C)++(0,\h) coordinate  (C')
						(B)++(0,\h)  coordinate (B')
						($(A')!.6!(B)$)  coordinate (H);
					\draw[dashed] (A)--(C) 
						(A')--(C); 
					\draw  (A)--(A')
						 (C)--(B)--(B')--(C')--(A')
						(C)--(C') 
						 (A)--(B)--(B')--(A')--(B)
						 (A)--(H);
					\foreach  \p/\r in {A/180,A'/180,B/-45,B'/-45,C/0,C'/0,H/45}
					\fill (\p) circle  (1pt) node[shift={(\r:3mm)}]{$\p$};
					\foreach \x/\y/\z in {A/B/C,A/H/B} 
						{ 
					\path ($(\y)!2mm!(\x)$)  coordinate (\y\x)
						 ($(\y)!2mm!(\z)$) coordinate (\y\z);
					\draw (\y\x)--($(\y\x)+(\y\z)-(\y)$)--(\y\z); 
					 }
				\end{tikzpicture}
		\end{center}
		Trong mặt phẳng  $(A'AB)$ kẻ $AH\perp A'B\quad (1)$.\\
		Ta  có $\heva{& AB\perp BC \\ & AA'\perp BC}\Rightarrow BC\perp (A'AB)\Rightarrow BC\perp AH\quad (2)$.\\
		 Từ (1) và (2) suy ra $AH\perp (A'AB) \Rightarrow \mathrm{d}\left(A,(A'BC)\right)=AH$.\\
		 Trong $\triangle A'AB$ vuông tại $A$ có đường cao $AH$ ta có
		$$\dfrac{1}{AH^2}=\dfrac{1}{AB^2}+\dfrac{1}{AA'^2} \Rightarrow AH=\dfrac{AB\cdot AA'}{\sqrt{AB^2+AA'^2}}=\dfrac{1\cdot 2}{\sqrt{1^2+2^2}}=\dfrac{2\sqrt{5}}{5} \approx 0{,}89.$$ 
	 }
\end{ex}

% nguồn: Ok_TraLoiNgan_Lop11_Lan2
\begin{ex}%[1H8V5-3]%[ma:831116d1]
	Cho hình chóp $S.A B C D$ có đáy là hình thang vuông tại $A$ và $D$, $S  D$ vuông góc với mặt đáy $(A B C D)$, $A D=2$, $SD=\sqrt{2}$. Tính khoảng cách giữa đường thẳng $C D$ và mặt phẳng $(SAB)$ (kết quả làm tròn đến hàng phần trăm).
	 \shortans{1,15}
	\loigiai{ 
		\begin{center}
			\begin{tikzpicture}[scale=1,font=\footnotesize,line join=round,line cap=round,>=stealth]
					\coordinate (D)  at (0,0);
					\coordinate  (A) at (-2,-2);
					\coordinate (B)  at (3,-2);
					\coordinate (C) at (7,0); 
					\coordinate (S)  at ($(D)+(0,5)$);
					 \coordinate (H) at ($(S)!0.6!(A)$);
					\draw  (S)--(B)
						 (S)--(C)
						(S)--(A) 
						 (A)--(B)--(C);
					\draw[dashed](S)--(D) 
						 (A)--(D)--(C)
						(D)--(H); 
					\pic[draw,thin,angle radius=2mm] {right  angle = S--H--D};
					\foreach \i/\g in {S/90,A/-90,B/-90,C/0,D/45,H/135}{\draw[fill=black](\i)  circle (1pt)
						($(\i)+(\g:3mm)$)  node[scale=1]{$\i$};}
				\end{tikzpicture}
		\end{center}
		Ta có $\heva{&A B \perp A D\\&A B \perp S D}\Rightarrow  A B \perp(S A D)$.\\
		Kẻ $D  H \perp S A$ tại $H$. Do $D H \subset(S A D)$ nên $A B \perp D H$.\\
		Ta có $\heva{&D H \perp S A\\&D  H \perp A B}\Rightarrow D H \perp(S A B)$.\\
		Do $D C\parallel A B$ nên $D  C\parallel (S A B)$.\\
		Vậy khoảng cách giữa đường thẳng $C D$ và mặt  phẳng $(S A B)$ là $D H$.\\
		Xét $\triangle  S A D$ vuông tại $D$ có $\dfrac{1}{D H^{2}}=\dfrac{1}{S D^{2}}+\dfrac{1}{A D^{2}}=\dfrac{1}{(\sqrt{2})^{2}}+\dfrac{1}{2^{2}}=\dfrac{3}{4}$.\\
		$\Rightarrow D H=\dfrac{2 \sqrt{3}}{3}$. \\ 
		Vậy khoảng cách giữa đường thẳng $C D$ và mặt  phẳng $(S A B)$ là $\dfrac{2\sqrt{3}}{3} \approx 1{,}15$.
	} 
\end{ex}

% nguồn: Ok_TraLoiNgan_Lop11_Lan2
\begin{ex}%[1H8V5-3]%[ma:5acc2041]
	Cho hình chóp $S.ABCD$ có đáy $ABCD$ là hình thoi tâm $O$, cạnh $1$, góc $\widehat{BAD}=60^\circ$, cạnh $SO$ vuông góc với $(ABCD)$ và $SO=1$. Khoảng cách từ $O$ đến $(SBC)$ bằng bao nhiêu? (kết quả  làm tròn đến hàng phần mười)
	\shortans{0,4} 
	\loigiai{ 
		\begin{center}
			\begin{tikzpicture}[scale=0.6,line join = round,line cap = round,font=\footnotesize]
					 \foreach \x/\y/\d in{0/0/D,7/0/C,3/3/A} \coordinate (\d) at (\x,\y);
					\coordinate (B) at ($(A)+(C)-(D)$); 
					\coordinate  (O) at ($(A)!0.5!(C)$);
					 \coordinate (S) at ($(O)+(0,7)$);
					\coordinate  (M) at ($(B)!0.35!(C)$);
					\coordinate (H) at ($(S)!0.65!(M)$); 
					\draw  (S)--(D)--(C)--(B)--(S)--(C)
						 (S)--(M);
					\draw[dashed]  (D)--(A)--(S)--(O)
						(A)--(B)--(D) 
						 (A)--(C)
						(O)--(M) 
						(O)--(H); 
					\foreach \d in {S,A,B,C,D,O,M,H}  \fill (\d) circle (1pt);
					\foreach \d/\l in {A/left,D/left,B/right,C/right,O/below,S/above,M/below right,H/above  right} \draw (\d) node[\l]{$\d$};
				\end{tikzpicture}
		\end{center}
		Vẽ $OM\perp BC$ tại $M$ thì $(SMO)\perp  BC\Rightarrow(SMO)\perp(SBC)$, vẽ $OH\perp SM$ tại $H$\\
		$\Rightarrow  OH\perp(SBC) \Rightarrow\mathrm{d}(O,(SBC))=OH$.\\
		 Ta có $AC=\sqrt{3}$, $OC=\dfrac{\sqrt{3}}{2}$, $OB=\dfrac{1}{2}$, $OM \cdot BC=OB\cdot OC\Rightarrow OM=\dfrac{OB\cdot OC}{BC}=\dfrac{\sqrt{3}}{4}$.\\
		Ta có $OH=\dfrac{SO\cdot MO}{\sqrt{SO^2+MO^2}}=\dfrac{1\cdot  \dfrac{\sqrt{3}}{4}}{\sqrt{1^2+\dfrac{3\cdot 1^2}{16}}}=\dfrac{\sqrt{57}}{19} \approx 0{,}4$.
	} 
\end{ex}

% nguồn: Ok_TraLoiNgan_Lop11_Lan1
\begin{ex}%[1H8C5-3]%[ma:78717686]
	Cho hình chóp $S.ABCD$ có đáy  $ABCD$ là hình vuông cạnh bằng $7$, $SA\perp (ABCD)$ và $SA=7\sqrt{6}$. Gọi $M$ là trung điểm của $SB$. Tính khoảng cách từ điểm $B$ đến mặt phẳng $(AMC)$ (làm tròn kết quả đến hàng phần trăm).
	\shortans{4,76} 
	 \loigiai{
	\immini{ 
		Gọi  $E$, $F$ lần lượt là trung điểm của $AB$ và $AO$.\\
		Ta có $ME\parallel SA$ nên  $ME\perp (ABCD)$.\\
		Ta cũng có $EF\parallel BO$  nên $EF\perp AC$.\\
		Gọi $H$  là hình chiếu vuông góc của $E$ trên $MF$.\\
		Khi đó, $\mathrm{d}\left(E, (AMC)\right)=EH=\sqrt{\dfrac{EF^2\cdot EM^2}{EF^2+  EM^2}}$.\\
		 Ta tính được $EM=\dfrac{SA}{2}=\dfrac{7\sqrt{6}}{2}$ và $EF=\dfrac{BO}{2}=\dfrac{7\sqrt{2}}{4}$.\\
		 Do đó, $EH=\dfrac{7\sqrt{78}}{26}$.\\
		Vì $BE$ cắt $(AMC)$ tại $A$  và $E$ là trung điểm của $BA$ nên $$\mathrm{d}\left(B, (AMC)\right)=2\mathrm{d}\left(E, (AMC)\right)=\dfrac{7\sqrt{78}}{13}.$$
		Vậy $\mathrm{d}\left(B, (AMC)\right)=\dfrac{7\sqrt{78}}{13}\approx 4{,}76$.} 
		{\begin{tikzpicture}[scale=1,font=\footnotesize,line join=round,line cap=round,>=stealth]
		\coordinate (A) at (0,0); 
		\coordinate (B)  at (-2.9,-2);
		\coordinate  (D) at (4,0);
		\coordinate (C)  at ($(B)-(A)+(D)$);
		\coordinate (S) at ($(A)+(90:3.5)$); 
		\coordinate  (M) at ($(S)!.5!(B)$);
		\coordinate  (E) at ($(A)!.5!(B)$);
		\coordinate (O) at  ($(A)!.5!(C)$);
		 \coordinate (F) at ($(A)!.5!(O)$);
		\coordinate (H) at  ($(M)!.35!(F)$);
		\draw (S)--(B)--(C)--(D)--(S)--(C)--(M); 
		\draw  [dashed] (S)--(A)--(D)
			(B)--(A)--(M) 
			(A)--(C) 
			 (B)--(D)
			(M)--(E)--(F)--(M) 
			 (E)--(H);
		\foreach \x/\g in {A/50,B/-90,C/-90,S/90,D/0,M/150,E/170,F/30,O/150,H/50} \fill (\x) circle  (1pt) +(\g:3mm) node{$\x$};
	\end{tikzpicture}
		 }
	} 
\end{ex}

% nguồn: Ok_TraLoiNgan_Lop11_Lan1
\begin{ex}%[1H8C5-3]%[ma:be1980ea]
	Cho hình chóp tứ giác đều $ S.ABCD $ có tất cả các cạnh đều bằng $ a $. Gọi $ O $ là giao điểm của $ AC $ và $ BD $. Khoảng cách từ  điểm $ A $ tới mặt phẳng $ (SCD) $ có dạng $\dfrac{a\sqrt{m}}{n}$. Tính $m+n$.
	\shortans{9} 
	\loigiai{ 
	\immini 
		 {
		Theo đề bài thì $SO\perp (ABCD)$. Dựng  $OK\perp CD$ tại $K$.\\
		Ta có $ \heva{& CD\perp OK \text{ (cmt)}\\ &CD\perp SO \text{ (vì $  SO\perp (ABCD)$)}.} $\\
		Suy ra $CD\perp (SKO)\Rightarrow (SCD)  \perp (SKO)$.\\
		Hai mặt phẳng này  vuông góc với nhau theo giao tuyến $SK$.\\
		Dựng  $OH \perp SK$ tại $H$.\\ Suy ra $OH\perp (SAB)$. Vậy $\mathrm{d}\left( O,(SCD)\right)= OH$.\\
		 $OK$ là đường trung bình của $\triangle ACD$.\\ $\Rightarrow OK=\dfrac{1}{2}AD=\dfrac{a}{2}$.\\
		Ta có $ SK=\dfrac{a\sqrt{3}}{2} $  ( do $\triangle SCD $ đều).
		} 
		 {
		\begin{tikzpicture} [scale=.8,font=\footnotesize,line join=round,line cap=round,>=stealth]
				 \coordinate (z) at (0,5);
				 \coordinate (A) at (0,0);
				 \coordinate (B) at (-2,-1.5);
				\coordinate (D)  at (5,0);
				\coordinate (C) at  ($(B)+(D)-(A)$);
				\coordinate  (O) at (intersection of A--C and B--D);
				\coordinate  (S) at ($(O)+(z)$);
				\coordinate (I)  at ($(A)!.5!(B)$);
				 \coordinate (K) at ($(C)!.5!(D)$);
				\coordinate (H) at ($(S)!.7!(K)$); 
				 \draw (S)--(B)--(C)--(D)--(S)--(C)
					 (S)--(K);
				\draw[dashed] (B)--(A)--(D)--(B) 
					 (S)--(A)--(C)
					(O)--(K) 
					 (S)--(O)--(H);
				\foreach \p/\pos in  {S/above,A/left,B/below,C/below,D/right,O/below,K/right,H/right}
				\fill  (\p) circle (1pt) node[\pos] {$\p$};
			\end{tikzpicture}
		 }
		\noindent Trong $ \triangle SOK $ có $ SO=\sqrt{SK^2  - OK^2}=\dfrac{a\sqrt{2}}{2} $.\\
		Do đó $  OH=\dfrac{OK\cdot SO}{\sqrt{OK^2+SO^2}}=\dfrac{a\sqrt{6}}{6} $.\\
		Hai  điểm $A$ và $O$ nằm trên đường thẳng có giao điểm với mặt phẳng $(SCD)$ tại $C$, nên có $\dfrac{\mathrm{d}\left( A, (SCD)\right)}{\mathrm{d}\left( O, (SCD)\right)}= \dfrac{AC}{OC}= 2$.\\
		Vậy $\mathrm{d}\left( A, (SCD)\right)=2\cdot  \mathrm{d}\left( O, (SCD)\right)= \dfrac{a\sqrt{6}}{3}$.\\
		Suy ra $m=6,n=3$.  Khi đó $m+n=9$.
	} 
\end{ex}

% nguồn: Ok_TraLoiNgan_Lop11_Lan1
\begin{ex}%[1H8C5-3]%[ma:e5358cb5]
	Cho hình chóp $S.ABCD$ có  đáy là hình thoi cạnh $a$, cạnh $SA$ vuông góc với mặt phẳng $(ABCD)$. Gọi $M$ là trung điểm $BC$. Biết rằng $\widehat{BAD}=120^\circ$, $\widehat{SMA}=45^\circ$. Khoảng cách từ $D$ đến mặt phẳng $(SBC)$ có dạng $\dfrac{a\sqrt{m}}{n}$ (phân số đã được rút gọn), khi đó $m^2-n^2$ bằng?
	\shortans{20} 
	\loigiai{ 
	\immini{- Dễ thấy các tam giác $ABC$ và $ADC$ là các tam giác đều, cạnh bằng $a$. Suy ra  $AM=\dfrac{a\sqrt{3}}{2}$.\\
		- Từ  $\widehat{SMA}=45^\circ$ suy ra tam giác $SAM$ vuông cân.\\
		- Dễ thấy $AM\perp BC$, mà $BC\perp SA$, suy ra $BC\perp (SAM)$. Kẻ $AH\perp SM$ thì $BC\perp AH$, do đó $AH\perp  (SBC)$.\\
		- Vì $AD\parallel  (SBC)$ nên $\mathrm{d}(D,(SBC))=\mathrm{d}(A,(SBC))=AH$.\\
		-  $\triangle SAM$ vuông cân nên $AH=AM\cdot \sin 45^\circ=\dfrac{a\sqrt{6}}{4}$.\\
		Vậy $m=6$, $n=4$.  Do đó $m^2-n^2=20$.
		 }{
		\begin{tikzpicture}[scale=0.5]
				\coordinate (a) at (0,0); 
				 \coordinate (b) at (8,0);
				 \coordinate (d) at (-2,-3);
				\coordinate  (c) at (6,-3);
				 \coordinate (s) at ($(a)+(0,6)$);
				\coordinate (m) at  ($(b)!1/2!(c)$);
				\coordinate (h) at ($(s)!1/2!(m)$); 
				 \draw [dashed] (a) node[left]{$A$}--(b) node[right]{$B$}
					(a)--(d) node[below]{$D$} 
					(s)  node[above]{$S$}--(a);
				\draw (b)--(c)  node[below]{$C$}
					(m)--(s)--(b) 
					(s)--(c)--(d)--cycle; 
				\draw [dashed] (a)--(h) node[right]{$H$} 
					(a)--(m) node[right]{\ $M$}; 
				\draw  pic[draw=black,angle radius=0.2cm] {right angle = a--m--c};
				\draw  pic[draw=black,angle radius=0.2cm] {right angle = a--h--m};
				\draw pic[draw=black,angle  radius=0.2cm] {right angle = s--a--m};
				\foreach \x in{a,b,c,d,s,m,h}{ 
					 \fill (\x) circle (1pt);
				 }
			\end{tikzpicture}
		 }
	} 
\end{ex}

% nguồn: Ok_TraLoiNgan_Lop11_Lan2
\begin{ex}%[1H8C5-3]%[ma:df428809]
	Cho hình chóp $S.ABCD$ có đáy là hình vuông cạnh $1$, mặt bên $SAB$ là tam giác đều và nằm trong mặt phẳng vuông góc với mặt phẳng đáy (minh họa như hình vẽ bên). Khoảng cách từ $D$ đến mặt phẳng $(SAC)$  bằng (kết quả làm tròn đến hàng phần trăm)
	\begin{center}
		\begin{tikzpicture}[scale=1,font=\footnotesize,line join=round,line cap=round,>=stealth]
				\coordinate (A) at  (0,0);
				 \coordinate (B) at (-2,-1);
				\coordinate (D) at  (4,0);
				\coordinate (I) at  ($(B)!0.5!(A)$);
				\coordinate (C)  at ($(B)+(D)-(A)$);
				 \coordinate (S) at ($(I)+(0,4)$);
				\draw (S)--(B) 
					(S)--(C) 
					(S)--(D) 
					 (B)--(C)--(D);
				 \draw[dashed,thin](S)--(A)--(C)
					 (A)--(B)
					 (A)--(D)--(B);
				 \foreach \i/\g in {S/90,A/150,B/-90,C/-90,D/0}{\draw[fill=black](\i) circle (1pt)
					($(\i)+(\g:3mm)$) node[scale=1]{$\i$};} 
			\end{tikzpicture}
	\end{center}
	 \shortans{0,65}
	\loigiai{ 
		\begin{center}
			\begin{tikzpicture}[scale=1,font=\footnotesize,line join=round,line cap=round,>=stealth]
					 \coordinate (A) at (0,0);
					 \coordinate (B) at (-2,-2);
					\coordinate  (D) at (6,0);
					\coordinate (I)  at ($(B)!0.5!(A)$);
					\coordinate (C)  at ($(B)+(D)-(A)$);
					\coordinate (O) at  ($(A)!0.5!(C)$);
					\coordinate  (S) at ($(I)+(0,4)$);
					\coordinate (G)  at (intersection of A--C and I--D);
					\coordinate (K) at  ($(A)!0.15!(C)$);
					 \coordinate (H) at ($(S)!0.6!(K)$);
					 \draw (S)--(B)
						 (S)--(C)
						(S)--(D) 
						(B)--(C)--(D); 
					 \draw[dashed,thin](D)--(I)--(S)--(O)
						 (K)--(S)--(A)--(C)
						(A)--(B) 
						(A)--(D)--(B) 
						 (I)--(H);
					 \foreach \i/\g in {S/90,A/150,B/-90,C/-90,D/0,O/-90,I/180,G/-90,K/30,H/60}{\draw[fill=black](\i) circle (1pt)
						($(\i)+(\g:3mm)$) node[scale=1]{$\i$};} 
				\end{tikzpicture}
		\end{center}
		Gọi $O=AC \cap BD$ và $G$ là trọng tâm tam  giác $ABD$, $I$ là trung điểm của $AB$.\\
		 Ta có
		$SI \perp (ABCD)$ và $\dfrac{\mathrm{d}(D,(SAC))}{\mathrm{d}(I,(SAC))} = \dfrac{DG}{IG} = 2 \Rightarrow \mathrm{d}(D,(SAC)) =  2\mathrm{d}(I,(SAC))$.\\
		Gọi $K$ là trung  điểm của $AO$, $H$ là hình chiếu của $I$ lên $SK$.\\
		 Ta có $IK \perp AC$, $IH \perp (SAC)$.\\
		$\Rightarrow \mathrm{d}(D,(SAC))=2\mathrm{d}(I,(SAC))=2IH$.\\ 
		Xét tam giác $SIK$ vuông tại $I$ có $SI=\dfrac{\sqrt{3}}{2}$, $IK=\dfrac{BO}{2}=\dfrac{\sqrt{2}}{4}$.\\ 
		Suy  ra $\dfrac{1}{IH^2} = \dfrac{1}{SI^2}+\dfrac{1}{IK^2} = \dfrac{4}{3}+\dfrac{16}{2} = \dfrac{4}{3}+8=\dfrac{28}{3}$.\\
		$\Rightarrow  IH=\dfrac{\sqrt{3}}{\sqrt{28}} = \dfrac{\sqrt{21}}{14}$.\\
		Vậy $\mathrm{d}(D,(SAC))=2IH  = \dfrac{\sqrt{21}}{7} \approx 0{,}65$.
	} 
\end{ex}

% nguồn: Ok_TraLoiNgan_Lop11_Lan1
\begin{ex}%[1H8B5-3]%[ma:82eebffb]
	Cho hình chóp $S.ABCD$ có $SA\perp(ABCD)$, đáy $ABCD$ là hình thang vuông có chiều cao $AB=3$. Gọi $I$ là $J$ lần lượt là trung điểm của $AB$ và $CD$. Khoảng cách giữa $IJ$ và $(SAD)$ bằng bao  nhiêu (kết quả làm tròn đến hàng phần mười)?
	 \shortans{1,5}
	\loigiai{\immini{Ta có $BC\parallel  AD$ nên $\mathrm{d}(IJ,(SAD))=\mathrm{d}(I,(SAD))=AI=\dfrac{3}{2}.$}{
		\begin{tikzpicture}[scale=0.7,font=\footnotesize,line join=round,line cap=round,>=stealth]
				 \path
					(0,0)  coordinate (A)
					(-1.3,-1.6)  coordinate (B)
					 (1.5,-1.6) coordinate (C)
					(4,0)  coordinate (D)
					($(A)+(0,3)$) coordinate  (S)
					($(A)!0.5!(B)$) coordinate (I) 
					($(C)!0.5!(D)$)  coordinate (J);
				\draw (S)--(B)--(C)--(D)--cycle  (S)--(C);
				 \draw[dashed] (S)--(A)--(D)
					(A)--(B) 
					 (I)--(J);
				\foreach \p/\q  in {S/90,A/180,B/-90,C/-90,D/0,I/-80,J/0}
				\fill[black]  (\p) circle (1.0pt) node[shift={(\q:2.5mm)}]{$\p$};
				\draw  pic[draw=black,angle radius=0.2cm] {right angle = S--A--D};
			\end{tikzpicture}
	 }}
\end{ex}


%%%%%%%%%%%%%%%% TỰ LUẬN %%%%%%%%%%%%%%%%

% nguồn: Ok_TuanLuan_Toan11_Lan2
\begin{ex}%[1H8N5-3]%[ma:87ab512a]
	Cho hình chóp $S.ABC$ có $SA$ vuông góc với mặt phẳng $(ABC)$, $\triangle ABC$ là tam giác vuông tại $B$, $BC=2a$. Khoảng cách  từ $C$ đến mặt phẳng $(SAB)$ bằng
	\loigiai{ 
	\immini{ 
		Do $SA \perp (ABC)$ nên $(SAB)  \perp (ABC)$.\\
		Mặt khác do $BC \perp  AB \Rightarrow BC \perp (SAB)$.\\
		Suy ra $\mathrm{d}\left(C,(SAB)\right)=CB=2a$.}{\begin{tikzpicture}[scale=1,line join=round,line cap=round]
		\coordinate (A)  at (0,0);
		\coordinate (B)  at (1,-1);
		 \coordinate (C) at (3,0);
		\coordinate (S) at  (0,3);
		\draw (S) -- (A)--(B)--(C)--cycle -- (B); 
		\draw[dashed]  (A) -- (C);
		\pic [draw,angle  radius=2mm,angle eccentricity=1.3] {right angle = S--A--B};
		\pic [draw,angle radius=2mm,angle eccentricity=1.3] {right angle  = S--A--C};
		\pic [draw,angle radius=2mm,angle eccentricity=1.3] {right angle = C--B--A}; 
		\node[left] at  (A) {$A$};
		\node[below] at  (B) {$B$};
		\node[right] at (C)  {$C$};
		\node[above]  at (S) {$S$};
		\foreach  \P in {A,B,C,S}{
			\fill  (\P) circle (1pt);
		} 
	\end{tikzpicture}
	}} 
\end{ex}

% nguồn: Ok_TuanLuan_Toan11_Lan2
\begin{ex}%[1H8H5-3]%[ma:3b34a979]
	Cho hình chóp $S.ABCD$ có $SA$ vuông góc với mặt  phẳng $(ABCD)$, $SA=a$, $ABCD$ là hình thoi cạnh $a$, $\widehat{ABC}=60^\circ$. Gọi $G$ là trọng tâm tam giác $SBC$. Khoảng cách từ $G$ đến mặt phẳng $(SCD)$ bằng
	 \loigiai{\immini{ $\triangle ABC$ cân tại $B$ và $\widehat{ABC}=60^\circ \Rightarrow \triangle ABC$, $\triangle ACD$ đều.\\
		 Gọi $M$ là trung điểm $CD \Rightarrow CD \perp AM$.\\
		Mà $CD \perp SA$ nên $CD \perp (SAM)$.\\ 
		 Dựng $AH \perp SM \Rightarrow AH \perp (SCD) \Rightarrow \mathrm{d}(A,(SCD))=AH$.\\
		 Xét $\triangle SAM$ vuông tại $A$:\\
		$\dfrac{1}{AH^2}=\dfrac{1}{AS^2}+\dfrac{1}{AM^2} \Rightarrow  AH=\dfrac{a\sqrt{21}}{7}$.\\
		Do $AB\parallel(SCD) \Rightarrow  \mathrm{d}(B,(SCD))=\mathrm{d}(A,(SCD))=\dfrac{a\sqrt{21}}{7}$.\\
		 Gọi $N$ là trung điểm $BC$ nên $\dfrac{GS}{NS}=\dfrac{2}{3}$.\\
		Suy ra 
		\begin{align*}
			\mathrm{d}(G,(SCD))&=\dfrac{2}{3}  \mathrm{d}(N,(SCD))\\
			&=\dfrac{2}{3} \cdot\dfrac{1}{2}  \mathrm{d}(B,(SCD))\\
			&=\dfrac{1}{3} \mathrm{d}(A,(SCD))\\ 
			 &=\dfrac{a\sqrt{21}}{21}.
		\end{align*}}
		{\begin{tikzpicture}[scale=1,line join=round,line cap=round]
		\coordinate (A) at  (0,0);
		\coordinate (D)  at (4,0);
		\coordinate (B) at  (-1.5,-1);
		\coordinate  (C) at ($(B)-(A)+(D)$);
		\coordinate (M) at ($(C)!0.5!(D)$); 
		\coordinate  (S) at ($(A)+(0,3)$);
		 \coordinate (H) at ($(S)!0.6!(M)$);
		\coordinate (N) at ($(C)!0.5!(B)$); 
		\coordinate  (G) at ($(S)!0.6!(N)$);
		\draw[dashed] (S) --  (A)--(D)
			(A)--(B) --(D) 
			(C)--(A)--(H) 
			(A)--(M); 
		\draw (S) -- (D)  --(C) --(B)-- cycle (S)--(C)
			 (S)--(M)
			 (S)--(N);
		\pic [draw,angle radius=2mm,angle  eccentricity=1.3] {right angle = S--A--B};
		\pic [draw,angle  radius=2mm,angle eccentricity=1.3] {right angle = S--A--D};
		\pic [draw,angle radius=2mm,angle  eccentricity=1.3] {right angle = A--M--C};
		\node[above  left] at (A) {$A$};
		\node[below] at (B) {$B$}; 
		 \node[below] at (C) {$C$};
		 \node[above] at (S) {$S$};
		\node[right] at (H)  {$H$};
		 \node[right] at (D) {$D$};
		\node[below right] at (M)  {$M$};
		\node[below] at (N) {$N$}; 
		\node[right] at  (G) {$G$};
		\foreach \P in {A,B,C,S,H,D,M,G,N}{ 
			\fill (\P)  circle (1pt);
		} 
	\end{tikzpicture}
		} 
	 }
\end{ex}

% nguồn: Ok_TuanLuan_Toan11_Lan2
\begin{ex}%[1H8H5-3]%[ma:4fbbf0cb]
	Cho hình chóp $S.ABCD$ có $(SAB)$ và $(SAD)$ cùng vuông góc với mặt phẳng $(ABCD)$, $ABCD$ là hình thoi  cạnh $a$, $\widehat{BAD}=120^\circ$, biết $SC$ hợp với đáy một góc $45^\circ$. Khoảng cách từ $B$ đến mặt phẳng $(SCD)$ bằng
	 \loigiai{\immini{ Ta có $\heva{
		&(SAB)  \perp (ABCD) \\
		 &(SAD) \perp (ABCD) \\
		 &(SAB)\cap(SAD)=SA}
		\Rightarrow SA \perp  (ABCD)$.\\
		Tam giác $ABC$ cân  tại $B$ và $\widehat{BAC}=60^\circ$.\\
		Suy ra $\triangle ABC$, $\triangle  ACD$ đều.\\
		 Vậy $(SC,(ABCD))=\widehat{SCA}=45^\circ \Rightarrow SA=AC=a$.\\
		Gọi $M$ là  trung điểm của $CD$\\
		 $\Rightarrow \heva{
		&CD  \perp AM \\
		&CD \perp SA}\Rightarrow CD \perp (SAM)$.\\ 
		 Dựng $AH \perp SM$\\
		$\Rightarrow AH \perp (SCD)  \Rightarrow \mathrm{d}(A,(SCD))=AH$.\\
		Xét $\triangle SAM$ vuông tại  $A$:\\
		$\dfrac{1}{AH^2}=\dfrac{1}{AS^2}+\dfrac{1}{AM^2} \Rightarrow  AH=\dfrac{a\sqrt{21}}{7}$.\\
		 Do $AB\parallel(SCD)$ nên\\ $\mathrm{d}(B,(SCD))=\mathrm{d}(A,(SCD))=\dfrac{a\sqrt{21}}{7}$.}
		{\begin{tikzpicture}[scale=1,line join=round,line cap=round]
		\coordinate (A)  at (0,0);
		\coordinate  (D) at (4,0);
		\coordinate (B) at (-1.5,-1); 
		\coordinate (C)  at ($(B)-(A)+(D)$);
		\coordinate (M) at ($(C)!0.5!(D)$); 
		\coordinate  (S) at ($(A)+(0,3)$);
		 \coordinate (H) at ($(S)!0.6!(M)$);
		 \draw[dashed] (S) -- (A)--(D)
			(A)--(B) --(D) 
			(C)--(A)--(H) 
			 (A)--(M);
		 \draw (S) -- (D) --(C) --(B)-- cycle (S)--(C)
			 (S)--(M);
		\pic [draw,angle radius=2mm,angle eccentricity=1.3] {right  angle = S--A--B};
		\pic [draw,angle  radius=2mm,angle eccentricity=1.3] {right angle = S--A--D};
		\pic [draw,angle radius=2mm,angle eccentricity=1.3] {right angle = A--H--M}; 
		\pic[draw,angle radius=5mm,angle eccentricity=1.4,"$45^\circ$"]  {angle = S--C--A};
		\node[above left] at (A)  {$A$};
		\node[below] at (B)  {$B$};
		\node[below] at (C) {$C$}; 
		\node[above] at  (S) {$S$};
		\node[left]  at (H) {$H$};
		\node[right] at (D) {$D$}; 
		\node[below] at (M) {$M$}; 
		\foreach \P in {A,B,C,S,H,D,M}{ 
			\fill (\P) circle  (1pt);
		} 
	\end{tikzpicture}
		 }
	 }
\end{ex}

% nguồn: Ok_TuanLuan_Toan11_Lan2
\begin{ex}%[1H8H5-3]%[ma:77076d2f]
	Cho hình chóp $S.ABC$ có $SA$ vuông góc với mặt phẳng $(ABC)$, $\triangle ABC$ là tam giác  vuông tại $B$, $AB=a$, $BC=2a$, $SA=a$. Gọi $G$ là trọng tâm tam giác $SAB$. Khoảng cách từ $G$ đến mặt phẳng $(SAC)$ bằng
	 \loigiai{\immini{
		 Do $SA \perp (ABC)$ nên $(SAC) \perp (ABC)$.\\
		 Trong mặt phẳng $(ABC)$, dựng $BH \perp AC$.\\
		Ta có  $BH \perp (SAC)$. Suy ra $\mathrm{d}(B,(SAC))=BH$.\\
		Xét $\triangle ABC$ vuông tại $B$  nên
		 $\dfrac{1}{BH^2}=\dfrac{1}{BA^2}+\dfrac{1}{BC^2} \Rightarrow BH=\dfrac{2a\sqrt{5}}{5}$.\\
		Do $G$ là trọng tâm $\triangle SAB$  nên $\dfrac{NG}{NB}=\dfrac{1}{3}$.\\
		Suy ra $\mathrm{d}(G, (SBC))=\dfrac{1}{3} \mathrm{d}(A,(SBC))=\dfrac{2a\sqrt{5}}{15}$.} 
		{\begin{tikzpicture}[scale=1,line join=round,line cap=round]
		\coordinate  (A) at (0,0);
		\coordinate (B) at  (2,-1);
		\coordinate  (C) at (4,0);
		\coordinate (S) at  ($(A)+(0,4)$);
		\coordinate (H) at ($(A)!0.6!(C)$); 
		 \coordinate (N) at ($(A)!0.5!(S)$);
		\coordinate (G)  at ($(B)!2/3!(N)$);
		\draw  (S) -- (A)--(B)--(C)--cycle -- (B) -- (N);
		\draw[dashed] (A) -- (C) 
			 (B)--(H);
		\pic [draw,angle radius=2mm,angle eccentricity=1.3] {right angle =  S--A--B};
		\pic [draw,angle radius=2mm,angle eccentricity=1.3]  {right angle = S--A--C};
		\pic [draw,angle radius=2mm,angle  eccentricity=1.3] {right angle = A--H--B};
		\pic [draw,angle  radius=2mm,angle eccentricity=1.3] {right angle = A--B--C};
		 \pic [draw,angle radius=2mm,angle eccentricity=1.3] {right angle = C--B--A};
		\node[left]  at (A) {$A$};
		 \node[below] at (B) {$B$};
		\node[right] at (C)  {$C$};
		\node[above]  at (S) {$S$};
		\node[above] at (H)  {$H$};
		\node[left]  at (N) {$N$};
		 \node[below left] at (G) {$G$};
		\foreach \P in  {A,B,C,S,H,G,N}{
			 \fill (\P) circle (1pt);
		 }
	\end{tikzpicture}
	}} 
\end{ex}

% nguồn: Ok_TuanLuan_Toan11_Lan2
\begin{ex}%[1H8H5-3]%[ma:271ab3ec]
	Cho hình chóp $S.ABC$ có $SA$ vuông góc với mặt phẳng $(ABC)$, $\triangle ABC$ là tam giác  đều cạnh bằng $a$, $SA=2a$. Gọi $M$ là trung điểm $BC$. Khoảng cách từ $M$ đến mặt phẳng $(SAB)$ bằng
	\loigiai{ 
	 \immini{
		Dựng $CN \perp AB \Rightarrow CN \perp (SAB)  \Rightarrow \mathrm{d}(C, (SAB))=CN$.\\
		Do $\triangle ABC$ đều cạnh bằng $a$  nên $CN=\dfrac{a\sqrt{3}}{2}$.\\
		Do $M$ là trung điểm $BC$ nên 
		$\mathrm{d}\left( M,  (SAB)\right) =\dfrac{1}{2} \mathrm{d}\left( C, (SAB)\right) =\dfrac{a\sqrt{3}}{4}$.}
		{\begin{tikzpicture}[scale=1,line join=round,line cap=round]
		\coordinate (A) at  (0,0);
		\coordinate (B) at  (1,-1);
		\coordinate  (C) at (3,0);
		\coordinate (S)  at (0,3);
		\coordinate (N) at ($(A)!0.5!(B)$); 
		\draw (S)  -- (A)--(B)--(C)--cycle -- (B);
		\draw[dashed] (A)  -- (C)--(N);
		\pic [draw,angle radius=2mm,angle eccentricity=1.3]  {right angle = S--A--B};
		\pic [draw,angle  radius=2mm,angle eccentricity=1.3] {right angle = S--A--C};
		\pic [draw,angle radius=2mm,angle eccentricity=1.3] {right angle =  C--N--B};
		\node[left] at (A) {$A$}; 
		\node[below]  at (B) {$B$};
		\node[right] at  (C) {$C$};
		\node[above] at (S)  {$S$};
		\node[below  left] at (N) {$N$};
		\foreach \P in {A,B,C,S,N}{ 
			\fill  (\P) circle (1pt);
		 }
	\end{tikzpicture}
	 }}
\end{ex}

% nguồn: Ok_TuanLuan_Toan11_Lan2
\begin{ex}%[1H8H5-3]%[ma:3fdbff58]
	Cho hình chóp $S.ABCD$ có đáy là hình bình hành với $BC=a\sqrt{2}$, $\widehat{ABC}=60^\circ$. Tam giác $SAB$ nằm trong mặt phẳng vuông góc  với mặt phẳng đáy. Khoảng cách từ điểm $D$ đến mặt phẳng $(SAB)$ bằng
	 \loigiai{
	\immini{Dựng  $SH \perp AB$.\\
		Do  $(SAB) \perp (ABCD)$ nên $SH \perp (ABCD)$.\\
		Dựng $CK \perp AB$.\\ 
		Vì $CK \parallel SH$ nên $CK \perp  (SAB)$.\\
		 Do $CD\parallel AB$ nên
		\begin{align*}
			 \mathrm{d}(D,(SAB))=\mathrm{d}(C,(SAB))=CK
			=BC  \sin 60^\circ
			=\dfrac{a\sqrt{6}}{2}. 
		\end{align*}}{
		\begin{tikzpicture}[scale=1.1,>=stealth]
				 \coordinate (B) at (0,0);
				\coordinate  (C) at (3,0);
				\coordinate  (A) at (.75,2);
				\coordinate (D) at ($(A) + (C) -  (B)$);
				 \coordinate (H) at ($(A)!0.5!(B)$);
				\coordinate (S) at ($(H)+(0,3)$); 
				\coordinate (K) at  ($(H)!0.5!(B)$);
				\draw[dashed] (D) --  (A) --(B)
					 (C) -- (K)
					 (H) -- (S) -- (A)
					(C) --(K); 
				\draw(S) --  (D) -- (C) -- (B) --cycle (S) --(C);
				\foreach  \P in {A,B,C,D,S,H,K}{
					\fill (\P) circle (1pt); 
				} 
				\node[above] at (S)  {$S$};
				\node[above right] at (A)  {$A$};
				\node[left] at  (B) {$B$};
				\node[right]  at (C) {$C$};
				\node[above  right] at (D) {$D$};
				 \node[above right] at (H) {$H$};
				\node[above right] at (K)  {$K$};
			\draw($(B)+(0.25,0)$) arc[start angle=0,end angle=60,radius=0.25];
			\node at ($(B)+(0.6,0.16)$)  {$60^\circ$};
		\end{tikzpicture}
		} 
	} 
\end{ex}

% nguồn: Ok_TuanLuan_Toan11_Lan2
\begin{ex}%[1H8H5-3]%[ma:7a52a529]
	Cho tứ diện đều $ABCD$, biết khoảng cách $A$ đến mặt  phẳng $(BCD)$ bằng $a\sqrt{6}$. Diện tích tam giác $ABC$ bằng
	 \loigiai{
	\immini{Gọi $M$ là trung điểm $CD$ và $H$ là hình chiếu vuông góc của $A$ trên  $BM$.\\
		Ta  có $\heva{&AM\perp CD\\&BM\perp CD\\&CD\not\subset (ABM)\\&AM, CD\subset (ABM)}\Rightarrow CD\perp(ABM)\Rightarrow CD\perp AH$\\
		Mà $CD\perp BM$ nên  $AH\perp (BCD)$\\
		Suy ra $\mathrm{d}\left(A,(BCD)\right)=AH$ và $H$  là trọng tâm $ABCD$.\\
		Xét $\triangle ABH$ vuông tại $H$:\\ 
		\begin{align*}
			&AH^2=AB^2 -  BH^2\\
			&\Leftrightarrow AH^2=AB^2-\left(\dfrac{2}{3}\cdot  \dfrac{\sqrt{3}}{2} AB\right)^2\\
			&\Leftrightarrow 6a^2=\dfrac{2}{3}  AB^2\\ &\Rightarrow AB=3a.
		\end{align*}
		Vậy $S_{\triangle ABC}=\dfrac{\sqrt{3}\cdot (3a)^2}{4}=\dfrac{9\sqrt{3}a^2}{4}$.} 
		{\begin{tikzpicture}[scale=1,line join=round,line cap=round]
		\coordinate (B)  at (0,0);
		\coordinate (C) at (3,-1); 
		\coordinate (D)  at (4,0);
		\coordinate (M)  at ($(C)!0.5!(D)$);
		 \coordinate (H) at ($(B)!2/3!(M)$);
		\coordinate  (A) at ($(H)+(0,4)$);
		\draw  (A) --(B)--(C)--(D)--cycle (A)--(C)
			 (A)--(M);
		 \draw[dashed] (B) -- (M)
			 (B)--(D)
			(A) --  (H);
		\pic [draw,angle radius=2mm,angle eccentricity=1.3] {right  angle = A--M--D};
		\pic  [draw,angle radius=2mm,angle eccentricity=1.3] {right angle = B--M--C};
		\pic [draw,angle radius=2mm,angle eccentricity=1.3]  {right angle = A--H--B};
		 \node[above] at (A) {$A$};
		\node[below]  at (B) {$B$};
		\node[right]  at (C) {$C$};
		\node[right] at (D)  {$D$};
		\node[below] at (H) {$H$}; 
		\node[right]  at (M) {$M$};
		\foreach  \P in {A,B,C,D,H,M}{
			\fill  (\P) circle (1pt);
		} 
	\end{tikzpicture}
	 }}
\end{ex}

% nguồn: Ok_TuanLuan_Toan11_Lan2
\begin{ex}%[1H8H5-3]%[ma:7e78938f]
	Cho hình chóp $S.ABCD$ có  đáy $ABCD$ là hình thoi tâm $O$, cạnh bằng $a$, $\widehat{BAD}=60^\circ$, $SO \perp (ABCD)$, $SO=a$. Khoảng cách từ đường thẳng $AD$ đến mặt phẳng $(SBC)$ bằng
	\loigiai{\immini{Kẻ $OK \perp BC(K \in BC)$, $OH \perp SK  (H\in SK)$.\\
		 Ta có $AD\parallel BC \Rightarrow AD\parallel (SBC)$.\\
		Khi đó $\mathrm{d}(AD, (SBC))=\mathrm{d}(M,(SBC))$ (với $M$ là giao  điểm của $AD$ và $OK$).\\
		Kẻ $MN\parallel OH (N \in  SK)$.\\
		Ta có $(SOK) \perp (SBC)$ theo giao  tuyến $SK$ nên $OH \perp (SBC)$.\\
		Suy ra $MN \perp (SBC)$.\\ 
		Suy ra  $\mathrm{d}(AD,(SBC))=\mathrm{d}(M,(SBC))=MN=2OH=\dfrac{a\sqrt{3}}{2}$.}
		{\begin{tikzpicture}[scale=1.2,line join=round,line cap=round]
		\coordinate (A) at (0,0); 
		 \coordinate (B) at (4,0);
		\coordinate  (D) at (-2,-1);
		 \coordinate (C) at ($(B)-(A)+(D)$);
		\coordinate  (O) at (intersection of A--C and B--D);
		\coordinate (M) at ($(A)!0.5!(D)$); 
		\coordinate (S)  at ($(O)+(0,4)$);
		\coordinate  (K) at ($(B)!0.5!(C)$);
		\coordinate (N)  at ($(S)!0.4!(K)$);
		\coordinate  (H) at ($(N)!0.5!(K)$);
		\draw[dashed]  (S) -- (A)--(D)
			(A)--(B)  --(D)
			(N)--(M)--(K) 
			 (S)--(O)--(H);
		\draw (S) -- (D) --(C) --(B)--  cycle (S)--(C)
			(S)--(K); 
		\pic [draw,angle radius=2mm,angle eccentricity=1.3] {right  angle = S--O--D};
		\pic [draw,angle radius=2mm,angle eccentricity=1.3] {right angle  = S--O--A};
		\pic [draw,angle radius=2mm,angle eccentricity=1.3] {right angle = S--N--M}; 
		 \pic [draw,angle radius=2mm,angle eccentricity=1.3] {right angle = S--H--O};
		\pic [draw,angle radius=2mm,angle eccentricity=1.3] {right angle  = M--K--C};
		\node[above right] at (A) {$A$}; 
		\node[right] at (B) {$B$}; 
		\node[below] at  (C) {$C$};
		\node[above]  at (S) {$S$};
		\node[right]  at (H) {$H$};
		\node[left]  at (D) {$D$};
		\node[above left] at (M)  {$M$};
		\node[right] at (N)  {$N$};
		\node[below] at (O) {$O$}; 
		\node[right]  at (K) {$K$};
		\foreach \P in {A,B,C,S,H,D,M,O,N,K}{ 
			 \fill (\P) circle (1pt);
		 }
	\end{tikzpicture}
		} 
	 }
\end{ex}

% nguồn: Ok_TuanLuan_Toan11_Lan2
\begin{ex}%[1H8H5-3]%[ma:c8c7b751]
	Cho hình  chóp $S.ABCD$ có $SA$ vuông góc với mặt phẳng $(ABCD)$, $ABCD$ là hình vuông cạnh $a$. Biết góc giữa đường thẳng $SB$ và mặt phẳng $(ABCD)$ bằng $60^\circ$. Khoảng cách từ $B$ đến mặt phẳng $(SCD)$ bằng
	\loigiai{ 
	\immini{Do $SA \perp (ABCD)$ nên $AB$ là hình chiếu vuông góc của $SB$ trên  mặt phẳng $(ABCD)$\\
		 Nên $(SB,(ABCD))=\widehat{SBA}$.\\
		Ta có $\heva{ 
		&BC \perp SA  \\
		&BC \perp AB} \Rightarrow BC \perp (SAB) \Rightarrow (SAB) \perp  (SBC)$.\\
		Xét $\triangle SAB$ vuông tại $A$:  $SA=AB \tan \widehat{SBA}=a\sqrt{3}$.\\
		Dựng $AH \perp SD \Rightarrow AH \perp  (SCD) \Rightarrow \mathrm{d}(A,(SCD))=AH$.\\
		Xét $\triangle SAD$  vuông tại $A$:\\ $\dfrac{1}{AH^2}=\dfrac{1}{SA^2}+\dfrac{1}{AD^2} \Rightarrow AH=\dfrac{a\sqrt{3}}{2}$.\\
		Do $AB\parallel CD$ nên $\mathrm{d}(B,(SCD))=\mathrm{d}(A,(SCD))=\dfrac{a\sqrt{3}}{2}$.} 
		{\begin{tikzpicture}[scale=1,line join=round,line cap=round]
		\coordinate (A) at (0,0); 
		\coordinate (B)  at (4,0);
		\coordinate (D) at (-2,-1); 
		\coordinate  (C) at ($(B)-(A)+(D)$);
		\coordinate  (S) at ($(A)+(0,4)$);
		\coordinate (H) at  ($(S)!0.4!(D)$);
		\draw[dashed] (S) --  (A)--(D)
			(A)--(B) 
			(A)--(H); 
		\draw (S) -- (D)  --(C) --(B)-- cycle (S)--(C);
		\pic [draw,angle radius=2mm,angle eccentricity=1.3] {right angle =  S--A--B};
		 \pic [draw,angle radius=2mm,angle eccentricity=1.3] {right angle = S--A--D};
		\pic [draw,angle radius=2mm,angle eccentricity=1.3] {right  angle = A--H--D};
		 \node[below] at (A) {$A$};
		\node[below]  at (B) {$B$};
		\node[below] at (C) {$C$}; 
		\node[above] at (S) {$S$}; 
		\node[left] at (H)  {$H$};
		\node[left] at (D) {$D$}; 
		\foreach \P in {A,B,C,S,H,D}{ 
			\fill (\P)  circle (1pt);
		} 
	\end{tikzpicture}
		} 
	 }
\end{ex}

% nguồn: Ok_TuanLuan_Toan11_Lan2
\begin{ex}%[1H8H5-3]%[ma:d44cd3d1]
	 Cho hình chóp $S.ABC$ có $SA$ vuông góc với mặt phẳng $(ABC)$, $\triangle ABC$ là tam giác đều cạnh bằng $a$, $SA=2a$. Gọi $G$ là trọng tâm $\triangle ABC$. Khoảng cách từ $G$ đến mặt phẳng $(SBC)$ bằng
	\loigiai{ 
	\immini{Gọi $M$  là trung điểm $BC$, $H$ là hình chiếu vuông góc của $A$ trên $SM$.\\
		 Ta có
		$\heva{&BC\perp AM\\  &BC\perp SA} \Rightarrow BC \perp (SAM) \Rightarrow (SBC) \perp (SAM)$.\\
		 $\Rightarrow AH \perp (SBC) \Rightarrow \mathrm{d}(A,(SBC))=AH$.\\
		 Xét $\triangle SAM$ vuông tại $A$ có
		 $\dfrac{1}{AH^2}=\dfrac{1}{SA^2}+\dfrac{1}{AM^2} \Rightarrow AH=\dfrac{a\sqrt{57}}{6}$.\\
		Do $G$ là  trọng tâm $\triangle ABC$ nên $\dfrac{GM}{MA}=\dfrac{1}{3}$.\\
		Suy  ra $\mathrm{d}(G,(SBC))=\dfrac{1}{3} \mathrm{d}(A,(SBC))=\dfrac{a\sqrt{57}}{18}$.}
		{\begin{tikzpicture}[scale=1,line join=round,line cap=round]
		\coordinate (A)  at (0,0);
		\coordinate (B)  at (1,-1);
		\coordinate (C) at (3.5,0); 
		 \coordinate (S) at ($(A)+(0,3)$);
		 \coordinate (M) at ($(B)!0.5!(C)$);
		\coordinate (G) at  ($(A)!0.666!(M)$);
		\coordinate (H)  at ($(S)!0.4!(M)$);
		 \draw (S) -- (A)--(B)--(C)--cycle -- (B)
			 (S)--(M);
		\draw[dashed] (H)--(A)  --(C)
			 (A) -- (M);
		\pic [draw,angle  radius=2mm,angle eccentricity=1.3] {right angle = S--A--B};
		\pic [draw,angle radius=2mm,angle eccentricity=1.3] {right angle =  S--A--C};
		 \pic [draw,angle radius=2mm,angle eccentricity=1.3] {right angle = S--M--B};
		\pic [draw,angle radius=2mm,angle eccentricity=1.3]  {right angle = A--H--M};
		\node[left] at  (A) {$A$};
		\node[below] at (B)  {$B$};
		\node[right] at (C) {$C$}; 
		\node[above] at (S) {$S$}; 
		\node[below right] at (M)  {$M$};
		\node[below  left] at (G) {$G$};
		 \node[left] at (H) {$H$};
		\foreach \P in  {A,B,C,G,S,H,M}{
			\fill  (\P) circle (1pt);
		} 
	\end{tikzpicture}
	}} 
\end{ex}

% nguồn: Ok_TuanLuan_Toan11_Lan2
\begin{ex}%[1H8H5-3]%[ma:b6c57026]
	Cho hình chóp $S.ABCD$ có $SA$ vuông góc với mặt phẳng $(ABCD)$, $ABCD$ là hình vuông tâm $O$ có cạnh $a$. Biết góc giữa hai mặt phẳng $(SBC)$ và $(ABCD)$ bằng $60^\circ$. Khoảng cách từ $O$ đến  mặt phẳng $(SBC)$ bằng
	\loigiai{\immini{ Ta có $\heva{ 
		&BC \perp  AB \\
		&BC \perp  SA
		 }
		\Rightarrow BC  \perp (SAB) \Rightarrow BC \perp SB$.\\
		Suy ra $((SBC),(ABCD))=\widehat{SBA}$.\\ 
		 Xét $\triangle SAB$ vuông tại $A$:\\ $SA=AB\cdot\tan \widehat{SBA}=a\sqrt{3}$.\\
		Vì $BC \perp (SAB)$ nên $(SAB) \perp  (SBC)$.\\
		Dựng  $AH \perp SB$\\
		$\Rightarrow  AH \perp (SBC) \Rightarrow \mathrm{d}(A,(SBC))=AH$.\\
		 Xét $\triangle SAB$ vuông tại $A$ nên\\
		$\dfrac{1}{AH^2}=\dfrac{1}{SA^2}+\dfrac{1}{AB^2} \Rightarrow AH=\dfrac{a\sqrt{3}}{2}$.\\ 
		Do $C \in (SBC)$ và $O$  là trung điểm $AC$ nên
		$\mathrm{d}(O,(SBC))=\dfrac{1}{2} \mathrm{d}(A,(SBC))=\dfrac{a\sqrt{3}}{4}$.} 
		{\begin{tikzpicture}[scale=1.2,line join=round,line cap=round]
		\coordinate (A) at  (0,0);
		\coordinate (B) at (4,0); 
		 \coordinate (D) at (-2,-1);
		\coordinate  (C) at ($(B)-(A)+(D)$);
		\coordinate (S)  at ($(A)+(0,4)$);
		\coordinate (O) at  (intersection of A--C and B--D);
		\coordinate  (H) at ($(S)!0.4!(B)$);
		\draw[dashed]  (S) -- (A)--(D)
			 (A)--(B) --(D)
			 (C)--(A)--(H)
			(S)--(O); 
		\draw (S) -- (D) --(C) --(B)-- cycle  (S)--(C);
		\pic [draw,angle  radius=2mm,angle eccentricity=1.3] {right angle = S--A--B};
		\pic [draw,angle radius=2mm,angle  eccentricity=1.3] {right angle = S--A--D};
		\pic [draw,angle radius=2mm,angle eccentricity=1.3] {right angle = A--H--B}; 
		\node[below  left] at (A) {$A$};
		\node[below] at  (B) {$B$};
		 \node[below] at (C) {$C$};
		\node[above] at (S)  {$S$};
		 \node[left] at (H) {$H$};
		\node[left] at (D) {$D$}; 
		\node[below]  at (O) {$O$};
		\foreach \P in  {A,B,C,S,H,D,O}{
			\fill (\P) circle  (1pt);
		 }
	\end{tikzpicture}
		} 
	} 
\end{ex}

% nguồn: Ok_TuanLuan_Toan11_Lan2
\begin{ex}%[1H8H5-3]%[ma:bafdebfd]
	Cho hình hộp chữ nhật $ABCD.A'B'C'D'$ có  $AB=a$, $BC=2a$, $BB'=a\sqrt{3}$. Khoảng cách từ điểm $B$ đến mặt phẳng $(ACC'A')$ bằng
	\loigiai{\begin{center}
		\begin{tikzpicture}[scale=1,line join=round,line cap=round]
				\coordinate (A) at (0,0); 
				\coordinate (D)  at (3,0);
				\coordinate (B)  at (-1.5,-1);
				\coordinate (C)  at ($(D)-(A)+(B)$);
				\coordinate (A') at ($(A)+(0,3)$); 
				\coordinate  (B') at ($(A')-(A)+(B)$);
				 \coordinate (C') at ($(A')-(A)+(C)$);
				\coordinate  (D') at ($(A')-(A)+(D)$);
				 \coordinate (H) at ($(A)!1/3!(C)$);
				\draw (A')--(B')--(C')--(D') --(D) --(C)--(B) 
					 (C)--(C')--(A')
					 (B)--(B')
					 (A')--(D');
				\draw[dashed] (A')--(A) -- (B)-- (H) 
					(C)--(A)--(D); 
				\pic  [draw,angle radius=2mm,angle eccentricity=1.3] {right angle = A--H--B};
				 \node[above left] at (A) {$A$};
				 \node[below] at (B) {$B$};
				\node[below] at  (C) {$C$};
				 \node[above right] at (D) {$D$};
				\node[left] at (B')  {$B'$};
				\node[above] at (A') {$A'$}; 
				 \node[below right] at (C') {$C'$};
				\node[below] at  (H) {$H$};
				\node[above] at (D')  {$D'$};
				\foreach \P in  {A,B,C,A',B',C',H,D,D'}{
					\fill (\P) circle  (1pt);
				 }
			\end{tikzpicture}
	\end{center}
	Kẻ $BH \perp AC(H \in AC)$.\\ 
	Lại có $BH \perp AA'$  (do $AA' \perp (ABCD)$).\\
	Suy ra $BH  \perp (ACC'A') \Rightarrow \mathrm{d}(B,(ACC'A'))=BH$.\\
	Xét  $\triangle ABC$ vuông tại $B$ có\\
	 $\dfrac{1}{BH^2}=\dfrac{1}{AB^2}+\dfrac{1}{BC^2}\Rightarrow BH=\dfrac{2a\sqrt{5}}{5}$.\\
Vậy $\mathrm{d}(B,(ACC'A'))=\dfrac{2a\sqrt{5}}{5}$.} 
\end{ex}

% nguồn: Ok_TuanLuan_Toan11_Lan2
\begin{ex}%[1H8H5-3]%[ma:b355b217]
	Cho lăng trụ $ABC.A'B'C'$ có đáy là tam giác vuông cân tại $A$ với $AB=AC=3a$. Hình chiếu vuông góc của $B'$ lên mặt đáy là điểm $H$ thuộc $BC$ sao cho $HC=2HB$. Biết cạnh bên của lăng trụ bằng $2a$.  Khoảng cách từ $B$ đến mặt phẳng $(B'AC)$ bằng
	 \loigiai{\immini{ Ta có $BC=3a\sqrt{2} \Rightarrow HB=a\sqrt{2}$.\\
		Lại có $B'H=\sqrt{BB'^2  - HB^2}=a\sqrt{2}$.\\
		Dựng $HE \perp AC$; $HF \perp B'E$.\\ 
		Suy ra $HF \perp (B'AC)  \Rightarrow \mathrm{d}(H,(B'AC))=HF$.\\
		Ta có $\dfrac{HE}{AB}=\dfrac{CH}{BC}=\dfrac{2}{3} \Rightarrow HE=2a$.\\ 
		 Suy ra $\dfrac{1}{HF^2}=\dfrac{1}{HE^2}+\dfrac{1}{B'H^2}$\\ $\Rightarrow HF=\dfrac{2a\sqrt{3}}{3}$.\\
		Mặt khác $\dfrac{\mathrm{d}(B,(B'AC))}{\mathrm{d}(H,(B'AC))}=\dfrac{BC}{HC}=\dfrac{3}{2}$.\\ 
		 Do đó $\mathrm{d}(B,(B'AC))=\dfrac{3}{2}HF=a\sqrt{3}$.}
		{\begin{tikzpicture}[scale=1,line join=round,line cap=round]
		\coordinate (B) at  (0,0);
		\coordinate (A) at  (2.5,-1.5);
		\coordinate (C) at (4,0); 
		\coordinate (B')  at (1,4);
		\coordinate (A') at ($(B')-(B)+(A)$); 
		\coordinate (C') at  ($(B')-(B)+(C)$);
		\coordinate  (H) at ($(B)!1/3!(C)$);
		\coordinate (E) at ($(A)!1/3!(C)$); 
		\coordinate (F)  at ($(B')!2/3!(E)$);
		 \draw (B)--(A)--(C)--(C') --(A') --(B')--cycle (B')--(A)
			 (B')--(C')
			(A)--(A'); 
		\draw[dashed] (H)--(B')  -- (E)
			(H)--(F) 
			 (B')--(C)--(B)
			(H)--(E); 
		\node[below] at (A) {$A$}; 
		\node[left] at  (B) {$B$};
		 \node[right] at (C) {$C$};
		\node[above]  at (B') {$B'$};
		\node[right] at (A')  {$A'$};
		 \node[above] at (C') {$C'$};
		\node[below]  at (H) {$H$};
		\node[below] at  (E) {$E$};
		\node[right]  at (F) {$F$};
		\foreach \P  in {A,B,C,A',B',C',H,E,F}{
			\fill (\P)  circle (1pt);
		} 
	\end{tikzpicture}
		} 
	} 
\end{ex}

% nguồn: Ok_TuanLuan_Toan11_Lan2
\begin{ex}%[1H8H5-3]%[ma:44c45af4]
	Cho hình chóp $S.ABCD$ có  đáy $ABCD$ là hình chữ nhật, $AB=a$, $BC=a\sqrt{3}$, $SA \perp (ABCD)$. Góc giữa $SC$ và mặt đáy bằng $45^\circ$. Khoảng cách từ điểm $A$ đến mặt phẳng $(SCD)$ bằng
	 \loigiai{\immini{ Vì $SA \perp (ABCD)$ nên $(SC,(ABCD))=\widehat{SCA}=45^\circ$.\\
		Kẻ  $AH \perp SD (H \in SD) \quad(1)$.\\
		Ta có $CD \perp AD$ và $CD  \perp SA$.\\
		Suy ra $CD  \perp (SAD) \Rightarrow CD \perp AH\quad (2)$.\\
		Từ $(1)$ và $(2)$ suy ra $AH \perp  (SCD)$.\\
		Do  đó $\mathrm{d}(A,(SCD))=AH$.\\
		Xét $\triangle ABC$ vuông tại $B$ có\\ 
		$AC=\sqrt{AB^2+BC^2}=2a$.\\ 
		Xét $\triangle  SAC$ vuông tại $A$ có\\
		$SA=AC \tan  45^\circ=2a$.\\
		 Xét $\triangle SAD$ vuông tại $A$ có\\
		$\dfrac{1}{AH^2}=\dfrac{1}{SA^2}+\dfrac{1}{AD^2}=\dfrac{7}{12a^2} \Rightarrow AH=\dfrac{2a\sqrt{21}}{7}$.\\ 
		 Vậy $\mathrm{d}(B,(SCD))=AH=\dfrac{2a\sqrt{21}}{7}$.}
		{\begin{tikzpicture}[scale=0.9,line join=round,line cap=round]
		\coordinate (A) at (0,0); 
		 \coordinate (B) at (4,0);
		 \coordinate (D) at (-2,-1);
		 \coordinate (C) at ($(B)-(A)+(D)$);
		\coordinate (S)  at ($(A)+(0,4)$);
		\coordinate (H) at ($(S)!0.4!(B)$); 
		\draw[dashed]  (S) -- (A)--(D)
			(A)--(B) 
			(C)--(A)--(H); 
		\draw (S) -- (D) --(C)  --(B)-- cycle (S)--(C);
		\pic [draw,angle  radius=2mm,angle eccentricity=1.3] {right angle = S--A--B};
		\pic  [draw,angle radius=2mm,angle eccentricity=1.3] {right angle = S--A--D};
		\pic [draw,angle radius=2mm,angle eccentricity=1.3] {right  angle = A--H--B};
		\node[above left] at (A)  {$A$};
		 \node[below] at (B) {$B$};
		\node[below] at (C)  {$C$};
		 \node[above] at (S) {$S$};
		\node[right] at (H)  {$H$};
		\node[left]  at (D) {$D$};
		 \foreach \P in {A,B,C,S,H,D}{
			 \fill (\P) circle (1pt);
		 }
	\end{tikzpicture}
		} 
	 }
\end{ex}

% nguồn: Ok_TuanLuan_Toan11_Lan2
\begin{ex}%[1H8H5-3]%[ma:6784b172]
	Cho hình chóp $S.ABC$ có đáy $ABCD$ là tam giác đều cạnh $a$. Cạnh bên $SB$ vuông góc mặt phẳng $(ABC)$ và $SB=2a$. Gọi $M$ là trung điểm của cạnh $BC$. Khoảng cách từ điểm $B$  đến mặt phẳng $(SAM)$ bằng
	\loigiai{\immini{ Ta có $AM \perp BC$ ($\triangle ABC$ đều); $AM  \perp SB$ (do $SB \perp (ABC)$)\\
		Do đó $AM  \perp (SBC)$.\\
		Trong mặt phẳng $(SBM)$, kẻ  $BH \perp SM$.\\
		Mà $BH \perp AM$ nên $BH  \perp (SAM)$.\\
		 Suy ra $\mathrm{d}(B,(SAM))=BH$.\\
		Xét $\triangle SBM$ vuông tại $B$  có
		 $\dfrac{1}{BH^2}=\dfrac{1}{SB^2}+\dfrac{1}{BM^2}$\\
		 Do đó $BH=\dfrac{2a\sqrt{17}}{17}$.}
		{\begin{tikzpicture}[scale=1,line join=round,line cap=round]
		\coordinate (B) at (0,0); 
		\coordinate (A) at  (1,-1.5);
		\coordinate (C) at  (4,0);
		\coordinate (S) at ($(B)+(0,3)$); 
		\coordinate  (M) at ($(C)!0.5!(B)$);
		\coordinate (H) at  ($(S)!0.6!(M)$);
		\draw (S)  -- (B)--(A)--(C)--cycle --(A);
		\draw[dashed] (H)--(B) --  (C)
			 (S)--(M)--(A);
		\pic [draw,angle radius=2mm,angle eccentricity=1.3]  {right angle = B--H--M};
		 \pic [draw,angle radius=2mm,angle eccentricity=1.3] {right angle = A--M--B};
		\node[below] at  (A) {$A$};
		\node[left]  at (B) {$B$};
		\node[right] at (C)  {$C$};
		\node[above] at  (S) {$S$};
		\node[below right]  at (M) {$M$};
		\node[above]  at (H) {$H$};
		\foreach \P  in {A,B,C,S,H,M}{
			 \fill (\P) circle (1pt);
		} 
	\end{tikzpicture}
		} 
	} 
\end{ex}

% nguồn: Ok_TuanLuan_Toan11_Lan2
\begin{ex}%[1H8H5-3]%[ma:28437460]
	Cho hình chóp $S.ABC$ có $SA$ vuông góc với mặt phẳng $(ABC)$, $\triangle ABC$ là tam giác vuông tại $B$, $AB=a$, $BC=2a$. Biết góc giữa đường thẳng $SB$ và mặt phẳng $(ABC)$ bằng $45^\circ$. Khoảng cách từ $A$ đến mặt phẳng $(SBC)$  bằng
	 \loigiai{\immini{
		Do $SA \perp (ABC)$ nên $AB$ là  hình chiếu vuông góc của $SB$ trên $(ABC)$\\
		 $\Rightarrow \left(SB,(ABC)\right)=\widehat{SBA}=45^\circ$.\\
		Vậy $\triangle SAB$ vuông cân tại  $A$\\
		$\Rightarrow SA=AB=a$.\\ 
		Dựng $AH \perp SB$, ta có 
		$(SAB) \perp  (SBC) \Rightarrow AH \perp (SBC)$\\ $\Rightarrow \mathrm{d}\left(A,(SBC)\right)=AH$.\\
		Xét $\triangle  SAB$ vuông tại $A$ nên
		$\dfrac{1}{AH^2}=\dfrac{1}{SA^2}+\dfrac{1}{AB^2}$ $\Rightarrow AH=\dfrac{a\sqrt{2}}{2}$.}{\begin{tikzpicture}[scale=1,line join=round,line cap=round]
		\coordinate (A) at  (0,0);
		\coordinate (B) at (2,-1); 
		\coordinate (C) at  (4,0);
		\coordinate  (S) at (0,3);
		 \coordinate (H) at ($(S)!0.4!(B)$);
		\draw (S) -- (A)--(B)--(C)--cycle --  (B)
			 (A)--(H);
		\draw[dashed] (A)  -- (C);
		\pic [draw,angle radius=2mm,angle eccentricity=1.3]  {right angle = S--A--B};
		\pic [draw,angle radius=2mm,angle eccentricity=1.3] {right angle =  S--A--C};
		\pic [draw,angle radius=2mm,angle eccentricity=1.3] {right angle =  A--H--B};
		\pic [draw,angle radius=2mm,angle  eccentricity=1.3] {right angle = C--B--A};
		\node[left] at (A)  {$A$};
		\node[below] at (B)  {$B$};
		\node[right] at (C) {$C$}; 
		 \node[above] at (S) {$S$};
		 \node[right] at (H) {$H$};
		\foreach \P in {A,B,C,S,H}{ 
			\fill (\P) circle (1pt); 
		 }
	\end{tikzpicture}
	}} 
\end{ex}

% nguồn: Ok_TuanLuan_Toan11_Lan2
\begin{ex}%[1H8H5-3]%[ma:f8bc11c1]
	Cho hình chóp $S.ABCD$ có đáy $ABCD$  là hình thang, $\widehat{ABC}=\widehat{BAD}=90^\circ$, $BA=BC=a$; $AD=2a$. Cạnh bên $SA$ vuông góc với đáy. Góc tạo bởi giữa $SC$ và $(SAD)$ bằng $30^\circ$. Khoảng cách từ $A$ đến $(SCD)$ bằng
	 \loigiai{
	\immini{Gọi $E$ là trung điểm $AD$.\\ 
		Khi đó  $ABCE$ là hình vuông cạnh $a$.\\
		 Suy ra $CE \perp AD$.\\
		 Lại có $CE \perp SA$.\\
		Do đó $CE \perp  (SAD) \Rightarrow \widehat{CSE}=(SC,(SAD))=30^\circ$.\\
		 Lại có $SC\cdot\sin 30^\circ=CE=a \Rightarrow SC=2a$.\\
		$\triangle ABC$ vuông cân  tại $B$ nên $AC=a\sqrt{2}$.\\
		Ta có $SA=\sqrt{SC^2-AC^2}=a\sqrt{2}$.\\ 
		Do $CE=\dfrac{1}{2} AD$  nên $\triangle ACD$ vuông tại $C \Rightarrow AC \perp CD$.\\
		Dựng  $AF \perp SC$.\\
		 Ta có $\mathrm {d}(A,(SCD))=AF=\dfrac{SA\cdot SC}{SC}=\dfrac{a\sqrt{2}\cdot 2a}{2a}=a$.}
		{  }
	} 
\end{ex}

% nguồn: Ok_TuanLuan_Toan11_Lan2
\begin{ex}%[1H8H5-3]%[ma:22d4779d]
	Cho hình chóp $S.ABCD$ có $SA$ vuông góc với mặt phẳng $(ABCD)$, $ABCD$ là hình vuông cạnh  $a$, $SA=a\sqrt{3}$. Khoảng cách từ $A$ đến mặt phẳng $(SBD)$ bằng
	\loigiai{\immini{Gọi $O$ là tâm hình  vuông $ABCD$\\ $\Rightarrow \heva{
		 &BD \perp OA \\
		&BD  \perp SA}
		\Rightarrow BD \perp (SAO) \Rightarrow  (SBD) \perp (SAO)$.\\
		Dựng $AK  \perp SO \Rightarrow AK \perp (SBD)$.\\
		Suy ra  $\mathrm{d}(A,(SBD))=AK$.\\
		Xét $\triangle SAO$  vuông tại $A$:\\ $\dfrac{1}{AK^2}=\dfrac{1}{SA^2}+\dfrac{1}{AO^2}$\\ $\Rightarrow AK=\dfrac{a\sqrt{21}}{7}$.}{
		\begin{tikzpicture}[scale=1,line join=round,line cap=round]
				\coordinate (A) at  (0,0);
				\coordinate  (B) at (4,0);
				\coordinate (D)  at (-2,-1);
				\coordinate  (C) at ($(B)-(A)+(D)$);
				\coordinate (S) at  ($(A)+(0,4)$);
				\coordinate (O) at (intersection of A--C and  B--D);
				\coordinate (K) at ($(S)!0.6!(O)$); 
				 \draw[dashed] (S) -- (A)--(D)
					(A)--(B) --(D) 
					 (C)--(A)--(K)
					(S)--(O); 
				\draw (S) -- (D) --(C) --(B)-- cycle  (S)--(C);
				\pic [draw,angle radius=2mm,angle eccentricity=1.3] {right  angle = S--A--B};
				\pic [draw,angle radius=2mm,angle  eccentricity=1.3] {right angle = S--A--D};
				\pic [draw,angle radius=2mm,angle eccentricity=1.3] {right  angle = A--K--O};
				\node[below left] at (A)  {$A$};
				\node[below]  at (B) {$B$};
				\node[below] at (C) {$C$}; 
				\node[above]  at (S) {$S$};
				\node[left] at (K)  {$K$};
				\node[left] at (D) {$D$}; 
				\node[below] at  (O) {$O$};
				\foreach \P  in {A,B,C,S,K,D,O}{
					\fill (\P) circle (1pt); 
				} 
			\end{tikzpicture}
		 }
	} 
\end{ex}

% nguồn: Ok_TuanLuan_Toan11_Lan2
\begin{ex}%[1H8H5-3]%[ma:e8ae2392]
	Cho hình lăng trụ $ABCD.A'B'C'D'$ có đáy là hình chữ nhật với $AD=a\sqrt{3}$. Tam giác $A'AC$ vuông cân tại $A$  và thuộc mặt phẳng vuông góc với đáy. Biết rằng $A'A=a\sqrt{2}$. Tính khoảng cách từ $D'$ đến mặt phẳng $(A'ACC')$
	 \loigiai{
	\immini{Trong $(A'AC)$, kẻ  $A'I \perp AC$.\\
		Vì $(A'AC) \perp (ABCD)$ và $(A'AC)\cap(ABCD)=AC$ nên $A'I  \perp (ABCD)$.\\
		Vì $DD'\parallel AA'$ nên  $DD'\parallel (A'ACC')$\\ $\Rightarrow \mathrm{d}(D',(A'AC))=\mathrm{d}(D,(A'AC))$\\
		Kẻ $DH  \perp AC$.\\
		 Ta có $AC=A'A\sqrt{2}=2a \Rightarrow CD=a$.\\
		 Suy ra $\mathrm{d}(D,(A'AC))=DH=\dfrac{a\sqrt{3}}{2}$.}
		{\begin{tikzpicture}[scale=1,line join=round,line cap=round]
		\coordinate  (A) at (0,0);
		\coordinate (B)  at (-1.5,-2);
		\coordinate (C)  at (2,-2);
		\coordinate  (D) at ($(A)-(B)+(C)$);
		\coordinate (A') at  ($(A)+(0,3)$);
		\coordinate (B')  at ($(B)+(0,3)$);
		\coordinate (C')  at ($(C)+(0,3)$);
		 \coordinate (D') at ($(D)+(0,3)$);
		\coordinate (I) at (intersection  of A--C and B--D);
		\coordinate  (H) at ($(I)!0.6!(C)$);
		\draw[dashed]  (B)--(A)--(D)
			(A) -- (A'); 
		\draw (D) -- (C)--  (B);
		 \draw (A')--(B')--(C')--(D')--cycle;
		\draw  (B)--(B')
			 (C)--(C')
			 (D)--(D');
		\draw[dashed] (A)-- (A')  --(C)-- cycle;
		 \draw[dashed] (B)--(D)--(H)
			 (A')--(I);
		\draw[dashed] (I)--(C); 
		\node at (I) [below]{$I$}; 
		\node at (H) [left]{$H$}; 
		\node[above left] at (A) {$A$}; 
		\node[below left] at (B)  {$B$};
		\node[below right] at  (C) {$C$};
		\node[above right] at (D) {$D$}; 
		\node[above  left] at (A') {$A'$};
		\node[above  left] at (B') {$B'$};
		\node[below  right] at (C') {$C'$};
		 \node[above right] at (D') {$D'$};
		\foreach \P in  {A,B,C,D,A',B',C',D',I,H}{
			 \fill (\P) circle (1pt);
		 }
	\end{tikzpicture}
		 }
	 }
\end{ex}

% nguồn: Ok_TuanLuan_Toan11_Lan1
\begin{bt}%[1H8C5-3]%[ma:89e3e097]
	Cho hình chóp $S.ABC$ có cạnh bên $SA\perp (ABC)$, tam giác $ABC$ vuông tại $B$. Biết  $SA=AB=BC$, $AC=4a$.
	\begin{enumerate}
		\item Chứng minh $CB\perp (SAB)$. 
		\item Kẻ $BM$  vuông góc với $AC$ tại $M$. Chứng minh $BM\perp (SAC)$.
		\item Gọi $N$ là trung điểm của $SB$. Chứng minh $AN\perp  SC$.
		\item Trên đoạn thẳng $BC$ lấy điểm $H$  sao cho $HC=2HB$. Gọi $K$ là giao điểm của $AH$ và $BM$. Tính khoảng các từ điểm $K$ đến mặt phẳng $(SBC)$.
	\end{enumerate}
	\loigiai{ 
		\begin{center}
			\begin{tikzpicture}[scale=1,line cap=round,line join=round,font=\footnotesize,>=stealth]
					\path 
						 (0,0) coordinate (A)
						($(A)+(-60:2)$) coordinate (B) 
						($(A)+(0:4)$) coordinate (C) 
						($(A)+(90:3.5)$) coordinate  (S)
						($(A)!0.5!(C)$) coordinate (M) 
						 ($(S)!0.5!(B)$) coordinate (N)
						($(B)!1/3!(C)$) coordinate  (H)
						(intersection of A--H and B--M) coordinate (K); 
					\draw  (S)--(A)--(B)--(S)--(C)--(B)
						 (A)--(N);
					\draw[dashed]  (A)--(C)
						(A)--(H) 
						(B)--(M); 
					\foreach \p/\r in  {A/-90,B/-90,C/-90,S/90,M/60,N/60,H/-10,K/0}
					\fill (\p) circle (1.25pt) node[shift={(\r:3mm)}]{$\p$}; 
					\draw 
						pic [draw,angle radius=2mm]  {right angle = C--B--A}
						 pic [draw,angle radius=2mm] {right angle = S--N--A}
						pic [draw,angle radius=2mm]  {right angle = A--M--B};
				\end{tikzpicture}
		\end{center}
		\begin{enumerate}
			\item  Ta có $\heva{&BC\perp AB\\&BC\perp SA}\Rightarrow BC\perp (SAB)$.
			 \item Ta có $\heva{&BM\perp AC\\&BM\perp SA}\Rightarrow BM\perp (SAC)$.
			\item Theo câu a, ta có $BC\perp (SAB)$  nên $BC\perp AN$, mà $AN\perp SB$ (vì tam giác $SAB$ cân tại $A$). Do đó $AN\perp (SBC)\Rightarrow AN\perp SC$.
			\item Áp dụng định lý Menelaus vào tam giác $AHC$, ta  có $$ \dfrac{HK}{KA}\cdot \dfrac{AM}{AC}\cdot \dfrac{CB}{BH}=1\Leftrightarrow \dfrac{HK}{KA}\cdot \dfrac{1}{2}\cdot 3=1\Leftrightarrow \dfrac{HK}{KA}=\dfrac{2}{3}\Rightarrow \dfrac{KH}{AH}=\dfrac{2}{5}.$$
			Suy  ra
			$\mathrm{d}(K,(SBC))=\dfrac{2}{5}\mathrm{d}(A,(SBC))=\dfrac{2}{5}\cdot AN $.\\ 
			Ta có $\triangle ABC$ vuông cân tại  $B$, nên $AB=BC=\dfrac{AC}{\sqrt{2}} =2a\sqrt{2}$.\\
			Tam giác $SAB$ vuông  cân nên $AN=\dfrac{1}{2}SB=2a$.\\
			Do  đó $\mathrm{d}(K,(SBC))=\dfrac{4a}{5}$.
		\end{enumerate}
	} 
\end{bt}

% nguồn: Ok_TuanLuan_Toan11_Lan2
\begin{ex}%[1H8C5-3]%[ma:cff9f489]
	Cho  hình chóp $S.ABC$ có cạnh bên $SA\perp (ABC)$, tam giác $ABC$ vuông tại $B$. Biết $SA=AB=BC$, $AC=4a$.
	\begin{enumerate}
		\item Chứng minh $CB\perp (SAB)$. 
		\item Kẻ $BM$ vuông góc với $AC$ tại $M$. Chứng minh  $BM\perp (SAC)$.
		 \item Gọi $N$ là trung điểm của $SB$. Chứng minh $AN\perp SC$.
		\item Trên đoạn thẳng $BC$ lấy điểm $H$ sao cho $HC=2HB$. Gọi $K$ là giao điểm của $AH$ và $BM$. Tính khoảng các từ điểm $K$ đến  mặt phẳng $(SBC)$.
	\end{enumerate}
	 \loigiai{
		\begin{center}
			\begin{tikzpicture}[scale=1,line cap=round,line join=round,font=\footnotesize,>=stealth]
					 \path
						(0,0) coordinate  (A)
						($(A)+(-60:2)$) coordinate  (B)
						 ($(A)+(0:4)$) coordinate (C)
						($(A)+(90:3.5)$) coordinate (S) 
						($(A)!0.5!(C)$)  coordinate (M)
						($(S)!0.5!(B)$) coordinate  (N)
						($(B)!1/3!(C)$)  coordinate (H)
						(intersection of A--H and B--M)  coordinate (K);
					 \draw (S)--(A)--(B)--(S)--(C)--(B)
						 (A)--(N);
					 \draw[dashed] (A)--(C)
						 (A)--(H)
						(B)--(M); 
					\foreach \p/\r  in {A/-90,B/-90,C/-90,S/90,M/60,N/60,H/-10,K/0}
					 \fill (\p) circle (1.25pt) node[shift={(\r:3mm)}]{$\p$};
					 \draw
						pic  [draw,angle radius=2mm] {right angle = C--B--A}
						pic [draw,angle radius=2mm]  {right angle = S--N--A}
						pic [draw,angle  radius=2mm] {right angle = A--M--B};
				\end{tikzpicture}
		\end{center}
		\begin{enumerate}
			\item Ta có $\heva{&BC\perp AB\\&BC\perp SA}\Rightarrow BC\perp (SAB)$. 
			\item Ta có $\heva{&BM\perp AC\\&BM\perp SA}\Rightarrow BM\perp (SAC)$. 
			\item Theo câu  a, ta có $BC\perp (SAB)$ nên $BC\perp AN$, mà $AN\perp SB$ (vì tam giác $SAB$ cân tại $A$). Do đó $AN\perp (SBC)\Rightarrow AN\perp SC$.
			\item Áp dụng định lý Menelaus vào tam giác $AHC$, ta có $$ \dfrac{HK}{KA}\cdot \dfrac{AM}{AC}\cdot \dfrac{CB}{BH}=1\Leftrightarrow \dfrac{HK}{KA}\cdot \dfrac{1}{2}\cdot 3=1\Leftrightarrow  \dfrac{HK}{KA}=\dfrac{2}{3}\Rightarrow \dfrac{KH}{AH}=\dfrac{2}{5}.$$
			 Suy ra
			 $\mathrm{d}(K,(SBC))=\dfrac{2}{5}\mathrm{d}(A,(SBC))=\dfrac{2}{5}\cdot AN $.\\
			Ta có $\triangle ABC$ vuông cân tại $B$, nên $AB=BC=\dfrac{AC}{\sqrt{2}} =2a\sqrt{2}$.\\ 
			Tam giác $SAB$ vuông cân  nên $AN=\dfrac{1}{2}SB=2a$.\\
			Do  đó $\mathrm{d}(K,(SBC))=\dfrac{4a}{5}$.
		\end{enumerate}
	 }
\end{ex}
